% !TEX program = xelatex
% 编译方式：xelatex opf_manual.tex （需两遍以生成目录）
% 注意：MiKTeX 在含中文的路径下编译会崩溃，请复制到纯英文路径（如 C:\latex_build\）后编译。
\documentclass[UTF8,fontset=fandol,a4paper,11pt]{ctexart}

\usepackage[margin=2.2cm]{geometry}
\usepackage{amsmath,amssymb}
\usepackage{booktabs,longtable,array}
\usepackage{graphicx}
\usepackage{xcolor}
\usepackage{float}
\usepackage{tikz}
\usetikzlibrary{arrows.meta,positioning,shapes.geometric}
\usepackage{hyperref}
\hypersetup{colorlinks=true,linkcolor=black,urlcolor=blue}

% ── 统一参数表环境 ────────────────────────────────────────────────
% 六列：字段名 | 符号 | 单位 | 典型范围 | 默认值 | 说明
\emergencystretch=3em
\newcolumntype{L}[1]{>{\raggedright\arraybackslash}p{#1}}
\newenvironment{paramtable}{%
  \begin{footnotesize}
  \begin{longtable}{@{}L{3.2cm}L{1.3cm}L{1.6cm}L{2.3cm}L{1.5cm}L{4.3cm}@{}}
  \toprule
  \textbf{字段名} & \textbf{符号} & \textbf{单位} & \textbf{典型范围} & \textbf{默认值} & \textbf{说明}\\
  \midrule
  \endfirsthead
  \multicolumn{6}{@{}l}{\footnotesize（续表）}\\
  \toprule
  \textbf{字段名} & \textbf{符号} & \textbf{单位} & \textbf{典型范围} & \textbf{默认值} & \textbf{说明}\\
  \midrule
  \endhead
  \midrule
  \multicolumn{6}{r@{}}{\footnotesize（接下页）}\\
  \endfoot
  \bottomrule
  \endlastfoot
}{%
  \end{longtable}
  \end{footnotesize}
}

% 字段名排版（下划线需转义后传入；\_ 处允许换行以防超长字段名溢出版心）
\makeatletter
\newcommand{\fld}[1]{\texttt{\fld@scan#1\fld@stop}}
\def\fld@scan#1{%
  \ifx#1\fld@stop
  \else
    \ifx#1\_\_\allowbreak\else#1\fi
    \expandafter\fld@scan
  \fi}
\makeatother
% 元件元信息行：#1=结构体/类名  #2=头文件路径
\newcommand{\compmeta}[2]{\noindent\textbf{代码结构体：}\texttt{#1}\hfill
  \textbf{定义位置：}\texttt{#2}\par\smallskip}

% 可断行的源码路径/函数名（/ 与 \_ 处允许换行）
\makeatletter
\newcommand{\srcpath}[1]{\texttt{\src@scan#1\src@stop}}
\def\src@scan#1{%
  \ifx#1\src@stop
  \else
    \ifx#1\_\_\allowbreak\else
    \ifx#1/#1\allowbreak\else
    #1\fi\fi
    \expandafter\src@scan
  \fi}
\makeatother

\title{\textbf{交直流混合高弹性能源电力系统仿真平台}\\[0.5em]
       {\Large 最优潮流（OPF）技术手册}\\[0.3em]
       {\large 数学模型、求解算法、接口参数与验证}}
\author{HySim-XJTU-HRPES（西安交通大学 HRPES 团队）}
\date{\today}

\begin{document}

\maketitle
\begin{abstract}
\noindent 本手册依据 HySim-XJTU-HRPES 平台 C++ 源代码中的最优潮流模块
（\texttt{src/\allowbreak optimal\_\allowbreak power\_\allowbreak flow/} 与
\texttt{include/\allowbreak hacdcpf/\allowbreak optimal\_\allowbreak power\_\allowbreak flow/}）
整理，按三个主类组织：
\textbf{AC OPF}（Native IPM 紧致路径）、\textbf{DC OPF}（无损线性化 LP/QP）、
\textbf{交直流混合 OPF}（Parity 全空间建模层 + Parity IPM / 嵌入式 Ipopt 求解核，
承载含 DC/VSC/LCC 的算例）；另含两个专项——含 OLTC 离散档位的无功优化（RPO）与
三相混合 OPF（相域）。每个求解器给出功能定位、
完整数学模型（目标函数与约束逐条注明源码出处）、算法流程、全部选项与结果字段的
单位制和默认值、验证与校核依据及已知局限。文档风格与《元件模型技术手册》一致。
\end{abstract}

\tableofcontents
\newpage

\section{阅读指南与约定}

\subsection{数据来源}
本手册全部模型、算法与参数取自以下源代码文件，参数表覆盖每个公开结构体的
\textbf{全部对外字段}：
\begin{itemize}
  \item \texttt{src/optimal\_power\_flow/} 与
        \texttt{include/hacdcpf/optimal\_power\_flow/} —— OPF 求解器本体
        （AC/DC/RPO/三相混合）；
  \item \texttt{tests/} 下 OPF 相关测试与 \texttt{docs/} 下 OPF 相关设计/验证文档
        —— 验证与校核一章的对照依据。
\end{itemize}

\subsection{单位制约定}
与《元件模型技术手册》一致：平台统一采用 MW·kV·MVA 单位制配合标幺值，系统基准功率
$S_{\mathrm{base}}=100$~MVA（\texttt{Defaults::kBaseMva}），基准频率 50~Hz。
OPF 内部方程一律以标幺值组装；目标函数中的成本项按有名值（MW、MVAr）计价后求和，
量纲在相应章节注明。布尔量、枚举量、计数型整数及 ID 引用等无量纲参数，单位列填
“—”；成本/价格类参数的币种在代码中未固定，单位列填“—”。

\subsection{求解器家族与选择逻辑}
平台 OPF 按建模口径分为三个主类与三个专项：
\begin{enumerate}
  \item \textbf{AC OPF}（纯交流，第\ref{sec:acopfnative}章）：Native IPM 紧致变量
        路径（极坐标 $V/\theta$，障碍--罚混合原生内点循环），入口
        \srcpath{opf::solve\_ac\_opf}；
  \item \textbf{DC OPF}（第\ref{sec:dcopf}章）：无损线性化 LP/QP（纯交流口径，
        不建直流电网，遇 LCC 直接拒绝），入口 \srcpath{opf::solve\_dc\_opf}；
  \item \textbf{交直流混合 OPF}（第\ref{sec:parityform}--\ref{sec:parityipm}章）：
        Parity 全空间建模层（显式直流电压块、VSC 三端口耦合、DC/DC、LCC 准稳态、
        能量路由器）+ Parity IPM 或嵌入式 Ipopt 求解核；统一入口
        \srcpath{opf::solve\_ac\_opf} 在 Auto 后端下检测到任何 DC/VSC/LCC 组件时
        一律分派到此路径，显式选择不兼容后端则拒绝而非静默丢弃直流网络；
  \item \textbf{专项}：含 OLTC 离散档位的无功优化 RPO（第\ref{sec:rpo}章）与
        相域三相混合 OPF（第\ref{sec:threephase}章，活跃研发）。基于
        MIPSolvers AML 的求解器无关建模层属于独立的 \srcpath{power\_models}
        模块，其源码等价公式仅在
        \srcpath{docs/modules/power\_models/power\_models\_manual.tex} 维护。
\end{enumerate}
运行时可用后端由 \srcpath{include/hacdcpf/api/solver\_capabilities.hpp} 查询；
嵌入式 Ipopt 仅在编译期开启 CMake 选项 \srcpath{HACDCPF\_ENABLE\_IPOPT}
（仅 macOS 默认 ON，映射为编译宏 \srcpath{HACDCPF\_HAVE\_IPOPT}）时可用。

\subsection{诚实结果口径}
平台铁律：近似、fallback、time-limit、模型覆盖不足必须写进结果的
\fld{model\_scope}、\fld{ValidityFlags} 或 \fld{model\_\allowbreak limitations} 字段。
各章“边界与局限”小节如实列出当前实现的已知近似与未覆盖项，不以后续路线图代替
现状描述。

\subsection{验证通道总览}
OPF 模块的验证通道（细节见第\ref{sec:verification}章）：
\begin{itemize}
  \item \textbf{手算解析解算例}：\srcpath{external\_data/opf\_hand\_cases/}
        （该目录当前不在仓库检出中）小型
        AC/DC OPF 手算基准（当前作为 GUI 人工校核材料保存，未接入 ctest，
        见第\ref{sec:verification}章）；
  \item \textbf{MATPOWER 算例}：标准算例潮流/OPF 基准对照；
  \item \textbf{跨求解器交叉验证}：Native IPM、Parity IPM、AML/Ipopt 三路结果互相比对
        （\srcpath{tests/test\_acopf\_dcopf\_crossval.cpp}、
        \srcpath{tests/test\_opf\_solver\_backends.cpp}）；
  \item \textbf{BPA/DSP 算例}：LCC 直流算例的 OPF 对照报告
        （\srcpath{docs/lcc\_dat\_opf\_report/}）；
  \item \textbf{专项验证文档}：无功优化验证（\srcpath{docs/reactive\_\allowbreak
        power\_\allowbreak optimization\_\allowbreak validation.md}）、大规模混合 IPM
        诊断（\srcpath{docs/large\_\allowbreak hybrid\_\allowbreak ipm\_\allowbreak
        diagnostics.md}）。
\end{itemize}

\newpage
% 最优潮流总览与公共接口层
% 本章文件由 \input 引入，不含导言区。
\section{最优潮流总览与公共接口层}\label{sec:overview}

本章给出平台最优潮流（OPF）模块的全景与公共接口层：求解器家族的组成与
数据流、Native AC 与交直流混合引擎的对比、公共选项 \fld{ACOPFOptions}/\fld{DCOPFOptions}
的全部字段、结果类型 \fld{ACOPFResult}/\fld{DCOPFResult} 的全部字段及其诚实
口径（\fld{model\_scope}、\fld{ValidityFlags}、\fld{model\_limitations}）、
问题描述 \fld{OPFProblem} 与求解器抽象 \fld{IOPFSolver}、嵌入式 Ipopt 后端的
编译门控与运行时探测、OPF 在平台中的主要调用方，以及求解器能力的运行时查询。
各引擎的内部数学（建模层、IPM 求解核、RPO 离散搜索、三相混合 OPF）
不在本章展开，分别见第\ref{sec:acopfnative}、\ref{sec:parityform}、
\ref{sec:parityipm}、\ref{sec:dcopf}、\ref{sec:rpo}、\ref{sec:threephase}、
节；验证与校核通道见第\ref{sec:verification}节。AML 建模层属于独立的
\srcpath{power\_models} 模块，其公式见
\srcpath{docs/modules/power\_models/power\_models\_manual.tex}。

除特别注明外，本章源码范围是
\srcpath{include/hacdcpf/optimal\_power\_flow/} 目录下的公共头文件、
对外门面 \srcpath{include/hacdcpf/api/hacdcpf.hpp} 与
\srcpath{src/api/hacdcpf.cpp}、以及三个求解器实现文件中的路径分派代码。
单位制遵循全平台约定：MW$\cdot$kV$\cdot$MVA，$S_{\mathrm{base}}=100$~MVA，
内部方程一律标幺值；成本类量的币种在代码中未固定，参数表单位列填“—”。

\subsection{求解器家族全景与数据流}\label{subsec:overview-family}

\subsubsection*{家族成员与入口}

平台 OPF 家族由六个求解/建模组件构成，全部挂在命名空间
\srcpath{hacdcpf::opf}（三相混合 OPF 为 \srcpath{hacdcpf::opf::phase\_hybrid}）
之下：

\begin{itemize}
  \item \textbf{AC OPF（纯交流）}：非线性交流 OPF，统一入口
        \srcpath{opf::solve\_ac\_opf}
        （\srcpath{include/hacdcpf/optimal\_power\_flow/ac\_opf\_solver.hpp:solve\_ac\_opf}；
        实现 \srcpath{src/optimal\_power\_flow/ac\_opf.cpp:solve\_ac\_opf}）。
        纯交流算例默认走 Native 紧致路径（第\ref{sec:acopfnative}节）；
        后端选择规则见第\ref{subsec:overview-engines}节。
  \item \textbf{交直流混合 OPF}：含 DC/VSC/LCC 组件的算例由同一入口
        \srcpath{opf::solve\_ac\_opf} 在 Auto 后端下一律分派到 Parity 全空间
        建模（第\ref{sec:parityform}节）+ Parity IPM 或嵌入式 Ipopt 求解核
        （第\ref{sec:parityipm}节）；显式选择不兼容后端时拒绝求解而非静默
        丢弃直流网络（见第\ref{sec:verification}节"混合系统禁 AC-only 回退"）。
  \item \textbf{DC OPF}：无损线性化 LP/QP（纯交流口径，不建直流电网），入口
        \srcpath{opf::solve\_dc\_opf}
        （\srcpath{include/hacdcpf/optimal\_power\_flow/dc\_opf\_solver.hpp:solve\_dc\_opf}；
        实现 \srcpath{src/optimal\_power\_flow/dc\_opf.cpp:solve\_dc\_opf}），另提供
        独立于求解的可行性复核 \srcpath{check\_dc\_opf\_feasibility}
        （同头文件:22--25）。详见第\ref{sec:dcopf}节。
  \item \textbf{RPO}：无功优化 MINLP（连续内层 NLP + 离散分接头/并联补偿
        搜索），入口 \srcpath{opf::solve\_rpo}
        （\srcpath{src/optimal\_power\_flow/reactive\_power\_opt.cpp:solve\_rpo}；
        选项 \fld{RPOOptions} 定义于
        \srcpath{include/hacdcpf/optimal\_power\_flow/reactive\_power\_opt.hpp:RPOOptions}）。
        详见第\ref{sec:rpo}节。
  \item \textbf{三相混合 OPF}：相域（不平衡）AC/DC 单体式非线性 OPF，入口
        \srcpath{solve\_three\_phase\_hybrid\_opf}
        （\srcpath{include/hacdcpf/optimal\_power\_flow/three\_phase\_hybrid\_opf.hpp:solve\_three\_phase\_hybrid\_opf}），
        其输入 \fld{ThreePhaseHybridOPFCase}（同文件:ThreePhaseHybridOPFCase）是已装配的相域
        算例而非 Rich 系统；从 Rich 系统的转换由适配器
        \srcpath{build\_three\_phase\_hybrid\_model} 完成
        （\srcpath{include/hacdcpf/optimal\_power\_flow/three\_phase\_hybrid\_adapter.hpp:build\_three\_phase\_hybrid\_model}）。
        详见第\ref{sec:threephase}节。
  \item \textbf{AML 建模层}：基于 MIPSolvers AML 的求解器无关建模层
        （\srcpath{include/hacdcpf/power\_models/} 下
        \srcpath{ac\_pf\_model\_builder.hpp}、
        \srcpath{hybrid\_opf\_model\_builder.hpp}、
        \srcpath{dc\_opf\_model\_builder.hpp}、
        \srcpath{lindistflow\_builder.hpp}、
        \srcpath{scuc\_builder.hpp}；实现于 \srcpath{src/power\_models/}），
        供外部求解器与对照实验使用。该实现不属于本模块；其唯一数学规范见
        \srcpath{docs/modules/power\_models/power\_models\_manual.tex}。
  \item \textbf{公共接口层}：选项、问题描述、结果类型与求解器抽象（本章
        第\ref{subsec:overview-options}--\ref{subsec:overview-interface}节）。
\end{itemize}

对外门面 \srcpath{include/hacdcpf/api/hacdcpf.hpp:solve\_ac\_opf} 等四处声明
（:solve\_ac\_opf、:solve\_dc\_opf、:solve\_rpo、:verify\_opf\_result）暴露四个 OPF
相关入口：\allowbreak\srcpath{hacdcpf::solve\_ac\_opf}、\allowbreak\srcpath{hacdcpf::solve\_dc\_opf}、
\allowbreak\srcpath{hacdcpf::solve\_rpo} 与求解后独立校核 \srcpath{verify\_opf\_result}。
门面 \srcpath{solve\_ac\_opf} 在调用内核返回后，额外在 OPF 调度点上重算
开关/断路器端潮流并写入 \fld{ac\_switch\_flows} 与
\fld{ac\_circuit\_breaker\_flows}（\srcpath{src/api/hacdcpf.cpp:solve\_ac\_opf}，
其中支路潮流经 \fld{bus\_merge\_map} 反投影回原始元件，见:1872--1884）；
\srcpath{solve\_dc\_opf} 与 \srcpath{solve\_rpo} 门面是直接转发
（同文件:solve\_dc\_opf）。

\subsubsection*{数据流架构}

\begin{figure}[H]
\centering
\resizebox{\textwidth}{!}{%
\begin{tikzpicture}[
  font=\small,
  box/.style={draw, rectangle, minimum height=0.95cm, align=center, inner sep=4pt},
  arr/.style={-{Stealth[length=2.6mm]}, thick},
  node distance=0.9cm]
  \node[box, minimum width=2.9cm] (rich) {Rich\\ \fld{HybridPowerSystem}};
  \node[box, right=1.0cm of rich, minimum width=2.9cm] (proj) {校验与投影\\ canonical projection};
  \node[box, right=1.0cm of proj, minimum width=2.4cm] (data) {canonical\\ \fld{SolverData}};
  \node[box, right=1.4cm of data, yshift=2.05cm, minimum width=5.1cm] (acopf)
    {AC OPF / 交直流混合 OPF（NLP）\\ Native IPM（纯交流）/ Parity IPM / 嵌入式 Ipopt};
  \node[box, below=0.32cm of acopf, minimum width=5.1cm] (dcopf)
    {DC OPF（线性化 LP/QP）};
  \node[box, below=0.32cm of dcopf, minimum width=5.1cm] (rpo)
    {RPO 无功优化（MINLP）};
  \node[box, below=0.32cm of rpo, minimum width=5.1cm] (tp)
    {三相混合 OPF（相域）};
  \node[box, below=0.32cm of tp, minimum width=5.1cm] (aml)
    {AML 建模层（\srcpath{src/power\_models/}）};
  \node[box, right=1.4cm of dcopf, minimum width=3.6cm] (res)
    {结果归因（反投影）\\ \fld{ACOPFResult} / \fld{DCOPFResult}};
  \node[box, below=0.55cm of res, minimum width=3.6cm] (audit)
    {独立校核\\ \srcpath{verify\_opf\_result}};
  \draw[arr] (rich) -- (proj);
  \draw[arr] (proj) -- (data);
  \draw[arr] (data.east) -- (acopf.west);
  \draw[arr] (data.east) -- (dcopf.west);
  \draw[arr] (data.east) -- (rpo.west);
  \draw[arr] (rich.south) |- (tp.west);
  \draw[arr] (rich.south) |- (aml.west);
  \draw[arr] (acopf.east) -- (res.north west);
  \draw[arr] (dcopf.east) -- (res.west);
  \draw[arr] (rpo.east) -- (res.south west);
  \draw[arr] (res) -- (audit);
\end{tikzpicture}}
\caption{OPF 求解器家族数据流。Rich 模型经校验与 canonical 投影装配为
\fld{SolverData} 后送入各求解引擎；结果经映射归因回原始组件稳定 ID，并可由
\srcpath{verify\_opf\_result} 做独立于求解器的可行性校核。三相混合 OPF 与
AML 建模层不经 \fld{SolverData} 主通道，前者经适配器直接装配相域算例，后者
自建 AML 模型。}
\label{fig:overview-family}
\end{figure}

数据流遵循全平台统一主线：Rich \fld{HybridPowerSystem}（工程语义层）
$\rightarrow$ 校验与 canonical 投影（开关收缩、零阻抗合并，产生
\fld{BusMergeMap}/\fld{BranchExpandMap}）$\rightarrow$ 装配 \fld{SolverData}
$\rightarrow$ 求解 $\rightarrow$ 结果反投影归因回原始组件稳定 ID。DC OPF
显式复用同一投影算子
（\srcpath{src/optimal\_power\_flow/dc\_opf.cpp:solve\_dc\_opf}，调用
\srcpath{projection::RichToCanonicalOperator::apply}），其母线/支路结果经
\srcpath{unproject\_bus\_vector} 与 \fld{branch\_orig\_to\_proj} 映射回原
始 ID 空间（同文件:separate\_external\_grid\_dispatch（lambda））；AC OPF 的结果归因集中在
\srcpath{separate\_external\_grid\_dispatch}
（\srcpath{src/optimal\_power\_flow/ac\_opf.cpp:separate\_external\_grid\_dispatch（lambda）}），外部电网注入被
从扩展发电机数组中剥离回填 \fld{external\_grid\_p\_mw}/
\fld{external\_grid\_q\_mvar}（:separate\_external\_grid\_dispatch（lambda））。

需要如实说明的三点：

\begin{itemize}
  \item \textbf{DC OPF 拒绝 LCC 算例}：只要系统中存在在运 LCC 换流站，
        \srcpath{solve\_dc\_opf} 直接返回拒绝状态而不求解，\fld{status} 为
        ``DC OPF rejected: this AC-only linear formulation does not model
        LCC AC/DC coupling\ldots''，\fld{model\_scope} 置
        \fld{dc-opf:rejected-lcc}
        （\srcpath{src/optimal\_power\_flow/dc\_opf.cpp:solve\_dc\_opf}）。
  \item \textbf{混合算例被强制走向全空间路径}：只要系统含 DC/VSC 子系统
        或 LCC，AC OPF 在路径分派处强制打开 \fld{enable\_primal\_dual} 与
        \fld{use\_parity\_ipm}（除非显式选择了 \fld{EconomicDispatch} 后端），
        即传统 Native AC 路径实际上只服务纯交流算例
        （\srcpath{src/optimal\_power\_flow/ac\_opf.cpp:solve\_ac\_opf\_impl}；混合标志
        计算于:3261--3265）。显式选择 \fld{EconomicDispatch} 的混合算例
        不会被静默降级，而是直接报错并给出提示
        （同文件:solve\_ac\_opf\_impl）。
  \item \textbf{AML 建模层与 Parity 建模层是两套独立装配}：Parity IPM 与
        嵌入式 Ipopt 共用 \srcpath{parity::build\_problem} 装配的全空间
        模型（\srcpath{src/optimal\_power\_flow/ac\_opf.cpp:solve\_with\_parity\_ipm}），并不
        经过 \srcpath{src/power\_models/} 的 AML builder。
        （\srcpath{src/optimal\_power\_flow/ac\_opf.cpp:solve\_with\_parity\_ipm} 的注释
        称“Parity-IPM uses the AML hybrid OPF builder”，与:2540 的实际
        调用不符，本章以代码为准。）
\end{itemize}

\subsection{非线性引擎对比：Native AC 与交直流混合两条路线}\label{subsec:overview-engines}

\fld{ACOPFOptions::ac\_solver\_backend} 在四个取值中选择非线性引擎
（枚举定义于
\srcpath{include/hacdcpf/optimal\_power\_flow/opf\_options.hpp:ACOPFSolverBackend}）：
\fld{Auto}（默认，Parity 全空间 IPM 优先、不收敛且编译了 Ipopt 时回退
Ipopt）、\fld{ParityIPM}（自研全空间原始对偶 IPM）、\fld{Ipopt}（嵌入式
Ipopt，滤波线搜索 NLP 求解器，作用于同一 Parity 建模）、
\fld{EconomicDispatch}（merit-order 经济调度加一次 AC 潮流，快速但次优、
仅限纯交流系统）。分派逻辑把该枚举翻译成 legacy 标志
\fld{enable\_primal\_dual}/\fld{use\_parity\_ipm} 与内部枚举
\srcpath{ParityInnerSolver}
（\srcpath{src/optimal\_power\_flow/ac\_opf.cpp:solve\_ac\_opf\_impl}；内部枚举定义于
:1870--1873）。其中 Parity IPM 与嵌入式 Ipopt 作用于同一份 Parity 全空间
建模，是\textbf{交直流混合 OPF 的两条承载路线}；Native AC 紧致路径只承担
纯交流算例。三条非线性引擎的对比见表~\ref{tab:overview-engines}。

\begin{table}[H]
\centering
\footnotesize
\begin{tabular}{@{}p{2.1cm}p{3.6cm}p{3.9cm}p{3.9cm}@{}}
\toprule
\textbf{维度} & \textbf{Native AC（传统路径）} & \textbf{Parity IPM（自研全空间）} & \textbf{嵌入式 Ipopt}\\
\midrule
建模口径 &
降阶传统建模：\srcpath{core::SolverData} + \fld{ACOPFIndex}，潮流雅可比上下文
（\srcpath{src/optimal\_power\_flow/ac\_opf.cpp:solve\_ac\_opf\_impl}） &
Parity 全空间建模：全部电压/功率/换流器变量为显式原始变量
（\srcpath{include/hacdcpf/optimal\_power\_flow/formulation.hpp:Problem}） &
与 Parity IPM 完全同一模型，1:1 映射为 \fld{engine::NLPModel}
（\srcpath{src/optimal\_power\_flow/ac\_opf.cpp:build\_parity\_nlp\_input}）\\
\midrule
坐标与变量布局 &
极坐标 $(\theta,V)$ 加发电机功率；非线性工程限值经外部栅栏/罚外环处理
（结果 \fld{model\_limitations} 自述，同文件:separate\_external\_grid\_dispatch（lambda）） &
极坐标 AC + 直流电压 + 换流器/储能/可控资源全变量块，布局见 \fld{VarIndex}
（\srcpath{include/hacdcpf/optimal\_power\_flow/formulation.hpp:VarIndex}） &
同 Parity IPM（共享 \fld{Problem::vidx}/\fld{cidx}）\\
\midrule
求解核 &
内环 KKT 牛顿 + 外环罚迭代；\fld{solver\_path=NativeAC}
（\srcpath{src/optimal\_power\_flow/ac\_opf.cpp:solve\_ac\_opf\_impl}） &
自研原始对偶 IPM \srcpath{solve\_primal\_dual\_ipm}
（\srcpath{include/hacdcpf/optimal\_power\_flow/native\_ipm\_solver.hpp:solve\_primal\_dual\_ipm}）；
\fld{solver\_path=ParityIPM}（src/optimal\_power\_flow/ac\_opf.cpp:solve\_with\_parity\_ipm） &
MIPSolvers 内嵌 Ipopt（滤波线搜索），经 \fld{engine::IpoptAdapter}
（\srcpath{src/optimal\_power\_flow/ac\_opf.cpp:solve\_parity\_with\_ipopt}）\\
\midrule
换流器/LCC 建模 &
换流器为自由 $P_{ac}/Q_{ac}/P_{dc}$ 盒式变量 + 容量圆，不钉 VDC\_Q 直流
电压（\fld{model\_scope}=\allowbreak\texttt{ac-opf-\allowbreak native:\allowbreak vsc-free-\allowbreak pq+\allowbreak capacity-\allowbreak circle}，
同文件:solve\_ac\_opf\_impl）；不支持 LCC（被强制转走，:solve\_ac\_opf\_impl） &
VDC\_Q 直流电压等式恒强制，容量圆/电流/调制度约束族可开关
（\fld{model\_scope} 标签按实际强制的约束族拼接，同文件:solve\_with\_parity\_ipm）；
支持 LCC 准稳态（:solve\_with\_parity\_ipm） &
与 Parity IPM 相同（共享建模）\\
\midrule
节点电价 LMP &
不提取节点电价：\fld{lmp\_p}/\fld{lmp\_q} 留空、\fld{lmp\_valid=false}，
由结果归因层统一补写原因（同文件:separate\_external\_grid\_dispatch（lambda）） &
从等式对偶 $\lambda$ 提取 \fld{lmp\_p}/\fld{lmp\_q} 并置 \fld{lmp\_valid}
（同文件:solve\_with\_parity\_ipm） &
适配器不回传约束乘子，\fld{lmp\_valid=false} 并注明原因
（同文件:solve\_with\_parity\_ipm）\\
\midrule
编译门控与可用性 &
始终可用 &
始终可用（自研代码） &
受 \srcpath{HACDCPF\_HAVE\_IPOPT} 门控；未编译时 \fld{available=false}、
\texttt{status=``Ipopt not compiled in''}
（同文件:solve\_parity\_with\_ipopt），显式请求时回退 Parity IPM 并在后端标签注明
（:solve\_with\_parity\_ipm）。详见第\ref{subsec:overview-ipopt}节\\
\midrule
适用场景 &
纯交流中小算例、回归对照；混合算例不可用 &
生产默认：交直流混合/LCC 算例的统一路径 &
大系统/困难算例的 Auto 回退与独立对照\\
\bottomrule
\end{tabular}
\caption{三条 AC 非线性引擎对比。第四条路径 \fld{EconomicDispatch} 为
merit-order 经济调度加一次牛顿 AC 潮流（实现
\srcpath{src/optimal\_power\_flow/ac\_opf.cpp:try\_dispatch\_pf\_fallback}，入口
:3325--3329、3520--3525），不是非线性优化，故不列入。}
\label{tab:overview-engines}
\end{table}

Auto 模式的完整分派顺序如下（
\srcpath{src/optimal\_power\_flow/ac\_opf.cpp}）：
\begin{enumerate}
  \item \fld{ac\_solver\_backend} 枚举翻译成 legacy 标志与
        \srcpath{ParityInnerSolver}（:solve\_ac\_opf\_impl）；
  \item 含 DC/VSC 或 LCC 的算例强制启用 Parity 全空间路径（:solve\_ac\_opf\_impl）；
  \item \fld{objective\_homotopy} 为真且非 Ipopt 后端时走两阶段目标同伦
        （$t:0\to1$ 的成本同伦 + Davidenko 切向预测，:solve\_ac\_opf\_impl）；显式
        Ipopt 与同伦不组合，退化为独立单解（:solve\_ac\_opf\_impl）；
  \item Parity 路径内部再按 \srcpath{ParityInnerSolver} 选内层牛顿引擎：
        \fld{Auto} 且启用 \fld{ac\_pf\_warm\_start} 时先用独立 Ipopt 再以
        原生 IPM 精化（:solve\_with\_parity\_ipm）；显式 \fld{Ipopt} 但不可用时回退原生
        IPM（:solve\_with\_parity\_ipm）；其余 \fld{Auto} 情形先原生 Parity IPM，不收敛
        且 \fld{allow\_fallback} 时回退 Ipopt（:solve\_with\_parity\_ipm）；
  \item 未启用原始对偶路径时：混合算例报错（:solve\_ac\_opf\_impl），纯交流算例走
        经济调度快速路径（:solve\_ac\_opf\_impl）或传统 Native AC 路径
        （:solve\_ac\_opf\_impl）。
\end{enumerate}

\subsection{公共选项：ACOPFOptions 与 DCOPFOptions}\label{subsec:overview-options}

\compmeta{ACOPFOptions}{\srcpath{include/hacdcpf/optimal\_power\_flow/opf\_options.hpp:ACOPFOptions}}

\fld{ACOPFOptions} 是 AC OPF（三条引擎共用）的全部用户选项。注意两点实现
口径：其一，\srcpath{solve\_ac\_opf} 入口处会把非正（越界）的数值字段重置
为默认值（\srcpath{src/optimal\_power\_flow/ac\_opf.cpp:solve\_ac\_opf\_impl}），因此
传 0 等价于取默认；其二，\fld{ac\_solver\_backend} 与 legacy 标志
\fld{enable\_primal\_dual}/\fld{use\_parity\_ipm} 并存，\fld{Auto} 时由
legacy 标志继续控制路径选择（字段注释，
\srcpath{include/hacdcpf/optimal\_power\_flow/opf\_options.hpp:ac\_solver\_backend}）。

\begin{paramtable}
\fld{ac\_solver\_backend} & — & — & Auto/ParityIPM/\allowbreak Ipopt/\allowbreak EconomicDispatch & Auto & 非线性引擎选择（枚举 \fld{ACOPFSolverBackend}，include/hacdcpf/optimal\_power\_flow/opf\_options.hpp:ACOPFSolverBackend）\\
\fld{objective} & — & — & Economic/\allowbreak VoltageDeviation/\allowbreak ActiveLoss/\allowbreak VoltageDeviation\allowbreak AndLoss & Economic & 连续目标（枚举 \fld{ACOPFObjective}，:ACOPFObjective）；后三者供 RPO 内层 NLP 使用\\
\fld{voltage\_target\_pu} & $V_{\mathrm{ref}}$ & pu & 0.95--1.05（经验值） & 1.0 & 电压偏差目标值\\
\fld{voltage\_deviation\_weight} & $w_V$ & — & $\ge 0$ & 1.0 & 电压偏差目标项权重\\
\fld{active\_loss\_weight} & $w_{\mathrm{loss}}$ & — & $\ge 0$ & 1.0 & 有功网损目标项权重\\
\fld{max\_inner\_iterations} & — & — & 1--$10^{5}$（HTTP 层 clamp 至 1--100000，\srcpath{runtime\_api\_v1.cpp:ac\_opf\_options}） & 80 & IPM 内层最大迭代数\\
\fld{max\_outer\_iterations} & — & — & 1--$10^{4}$（HTTP 层 clamp，\srcpath{runtime\_api\_v1.cpp:ac\_opf\_options}） & 8 & 外环（罚/栅栏）最大迭代数\\
\fld{max\_line\_search\_steps} & — & — & $\ge 1$ & 20 & 单次线搜索最大步数\\
\fld{feasibility\_tol} & $\varepsilon_{\mathrm{feas}}$ & pu & $10^{-12}$--1（HTTP 层 clamp，\srcpath{runtime\_api\_v1.cpp:ac\_opf\_options}） & $10^{-6}$ & 原始可行性容差\\
\fld{stationarity\_tol} & $\varepsilon_{\mathrm{stat}}$ & pu & 同上 & $10^{-6}$ & 驻点/对偶容差\\
\fld{barrier\_mu0} & $\mu_0$ & — & $>0$（典型 $10^{-3}$--$10^{-1}$） & $10^{-2}$ & 初始栅栏参数\\
\fld{barrier\_mu\_reduction} & — & — & $(0,1)$ & 0.2 & 栅栏参数缩减因子\\
\fld{barrier\_mu\_min} & $\mu_{\min}$ & — & $>0$ & $10^{-8}$ & 栅栏参数下限\\
\fld{merit\_penalty} & — & — & $>0$ & 10.0 & merit 函数罚系数（传统路径）\\
\fld{regularization} & — & — & $>0$ & $10^{-6}$ & KKT 系统正则化\\
\fld{step\_backoff} & — & — & $(0,1)$ & 0.5 & 线搜索步长回退因子\\
\fld{interior\_fraction} & — & — & $(0,1)$ & 0.995 & fraction-to-boundary 内点保持系数\\
\fld{ac\_eval\_threads} & — & — & $\ge 1$ & 1 & AC 方程/雅可比求值线程数\\
\fld{enable\_primal\_dual} & — & — & true/false & false & legacy 标志：启用原始对偶路径\\
\fld{use\_parity\_ipm} & — & — & true/false & false & legacy 标志：选用 Parity 全空间建模\\
\fld{allow\_fallback} & — & — & true/false & true & 允许 Auto 模式在不收敛时回退（Ipopt 或同伦末段全目标重试）\\
\fld{verbose} & — & — & true/false & false & 迭代日志\\
\fld{enable\_phase\_one} & — & — & true/false & true & Parity Native IPM 的有界 primal/dual warm-start 构造\\
\fld{phase\_one\_time\_limit\_ms} & — & ms & $>0$ 协作截止；$\le0$ 不限 & 5000 & Phase I 墙钟预算\\
\fld{phase\_one\_max\_iterations} & — & — & $\ge0$ & 12 & Phase I Newton 迭代硬上限\\
\fld{phase\_one\_max\_factorizations} & — & — & $\ge0$ & 14 & Phase I primal/dual 共用分解上限\\
\fld{phase\_one\_max\_backtracks} & — & — & $\ge0$ & 12 & 每次 Phase I 步的回溯上限\\
\fld{phase\_one\_barrier\_mu} / \fld{phase\_one\_admission\_mu\_factor} & $\mu_0,\eta_a$ & — & $(0,0.1]$ / $\ge0$ & 0.1 / 1.0 & central warm-start 的障碍尺度与原坐标准入比例\\
\fld{phase\_one\_primal\_mu\_factor} / \fld{phase\_one\_centrality\_tolerance} & $\eta_p,\eta_c$ & — & $\ge0$ & 0.1 / 0.5 & primal handoff corridor 与 $\max_i|s_i z_i/\mu_0-1|$ 门限\\
\fld{warm\_start} & — & — & 指针/nullptr & nullptr & 指向前次 \fld{ACOPFResult} 的非拥有暖启动指针；维度不兼容的变量族各自回退物理初值（字段注释，include/hacdcpf/optimal\_power\_flow/opf\_options.hpp:ACOPFOptions）\\
\fld{ac\_pf\_warm\_start} & — & — & true/false & false & 先解一次混合潮流，以 AC 可行点初始化 Parity IPM（字段注释：PEGASE 受压算例由此解锁，:ac\_pf\_warm\_start；实现 src/optimal\_power\_flow/ac\_opf.cpp:solve\_with\_parity\_ipm）\\
\fld{ac\_pf\_dc\_phase\_one} & — & — & true/false & true & 已启用 AC-PF warm start 时，对至少 5000 母线的纯 AC 模型按 baseline-dual 预筛选运行有限步 compact NativeLCQP dispatch seed\\
\fld{ac\_pf\_dc\_phase\_one\_time\_limit\_ms} & $T_{\mathrm{DC},I}$ & ms & $>0$ 协作截止；$\le0$ 不限 & 2000 & DC Phase I 端到端预算，含 projection/formulation/structural/symbolic/numeric；允许一个在途稀疏分解超时\\
\fld{ac\_pf\_dc\_phase\_one\_max\_iterations} / \fld{ac\_pf\_dc\_phase\_one\_tolerance} & — & — / pu & $\ge1$ / $\ge0$ & 5 / $10^{-2}$ & DC warm-start-only 迭代预算与可用残差门限\\
\fld{ac\_pf\_dc\_phase\_one\_min\_dual\_improvement} / \fld{ac\_pf\_dc\_phase\_one\_baseline\_dual\_threshold} & — & — & $[0,1]$ / $\ge0$ & 0.2 / 100 & 候选 dual 最小相对改善与是否支付 DC Phase I 的 baseline 门限\\
\fld{compute\_homotopy\_tangent} & — & — & true/false & false & 在返回点计算 Davidenko 目标同伦切向量 $\mathrm{d}w/\mathrm{d}t$（一次额外惯性受控 KKT 解，:compute\_homotopy\_tangent）\\
\fld{homotopy\_t} & $t$ & — & $(0,1]$ & 1.0 & 所解问题的成本尺度 $t$（切向量右端为 $\nabla f/t$）\\
\fld{objective\_homotopy} & — & — & true/false & false & 两阶段目标同伦求解（$t:0\to1$，:objective\_homotopy；实现 src/optimal\_power\_flow/ac\_opf.cpp:solve\_ac\_opf\_impl）\\
\fld{homotopy\_dt0} & $\Delta t_0$ & — & $(0,1]$ & 0.10 & 初始延拓步长；非正值回退 0.10（src/optimal\_power\_flow/ac\_opf.cpp:solve\_ac\_opf\_impl）\\
\fld{enforce\_branch\_limits} & — & — & true/false & true & 约束族开关：支路容量\\
\fld{enforce\_converter\_capacity} & — & — & true/false & true & 约束族开关：换流器容量圆\\
\fld{enforce\_converter\_current\_limits} & — & — & true/false & true & 约束族开关：换流器电流限值\\
\fld{enforce\_converter\_modulation\_limits} & — & — & true/false & true & 约束族开关：调制度窗口与 DC/DC 占空比\\
\end{paramtable}

\compmeta{DCOPFOptions}{\srcpath{include/hacdcpf/optimal\_power\_flow/opf\_options.hpp:DCOPFOptions}}

\fld{DCOPFOptions} 是 DC OPF 的全部用户选项。后端枚举
\fld{DCOPFSolverBackend}（:DCOPFSolverBackend）取
\fld{Auto}/\fld{Native}/\fld{NativeQP}/\fld{HiGHS}/\fld{Gurobi}；实际回退
链见第\ref{sec:dcopf}节，Auto 链为 Gurobi-QP $\to$ NativeLCQP $\to$ HiGHS
$\to$ NativeDualSimplex
（\srcpath{src/optimal\_power\_flow/dc\_opf.cpp:solve\_dc\_opf}），每步尝试记录
进 \fld{solver\_chain}（:solve\_dc\_opf）。

\begin{paramtable}
\fld{feasibility\_tol} & $\varepsilon_{\mathrm{feas}}$ & pu & $10^{-12}$--1（HTTP 层 clamp，\srcpath{runtime\_api\_v1.cpp:dc\_opf\_options}） & $10^{-6}$ & LP/QP 可行容差\\
\fld{optimality\_tol} & $\varepsilon_{\mathrm{opt}}$ & pu & $>0$ & $10^{-6}$ & 最优性（对偶）容差\\
\fld{max\_iterations} & — & — & 1--$10^{6}$（HTTP 层 clamp，\srcpath{runtime\_api\_v1.cpp:dc\_opf\_options}） & 10000 & 单纯形/IPM 最大迭代数\\
\fld{verbose} & — & — & true/false & false & 日志\\
\fld{solver} & — & — & Auto/Native/\allowbreak NativeQP/\allowbreak HiGHS/\allowbreak Gurobi & Auto & LP/QP 后端选择（枚举 \fld{DCOPFSolverBackend}，include/hacdcpf/optimal\_power\_flow/opf\_options.hpp:DCOPFSolverBackend）\\
\fld{pwl\_segments} & — & — & 1--1000（HTTP 层 clamp，\srcpath{runtime\_api\_v1.cpp:dc\_opf\_options}） & 4 & 二次成本分段线性化段数（LP 路径）\\
\fld{include\_branch\_limits} & — & — & true/false & true & 是否计入支路潮流限值\\
\fld{branch\_limit\_margin} & — & pu & $>0$ & 1.0 & 支路限值缩放系数\\
\fld{load\_shedding} & — & — & true/false & true & 是否启用切负荷松弛变量\\
\fld{voll} & VOLL & — & $\ge 0$ & 0.0 & 失负荷价值（成本币种未定）；$\le 0$ 时按 $\max(100,\,1.1\times$ 最大边际成本$)$ 自动折算（\srcpath{dc\_opf.cpp:build\_dc\_opf\_lp}）\\
\fld{compute\_lmp} & — & — & true/false & true & 是否在主解后补解支撑 LP 以恢复约束对偶（LMP 与阻塞价格）；关闭可省一次 LP（字段注释，include/hacdcpf/optimal\_power\_flow/opf\_options.hpp:DCOPFOptions；实现 src/optimal\_power\_flow/dc\_opf.cpp:populate\_missing\_duals\_from\_supporting\_lp）\\
\fld{structural\_warm\_start} & — & — & true/false & true & NativeLCQP 的逐连通分量配平与约化 Laplacian 初值；其他后端忽略\\
\end{paramtable}

\subsection{结果类型与诚实口径}\label{subsec:overview-result}

\subsubsection*{ACOPFResult 字段}

\compmeta{ACOPFResult}{\srcpath{include/hacdcpf/optimal\_power\_flow/opf\_result.hpp:ACOPFResult}}

结果向量的索引空间一律为\textbf{原始（authored）组件顺序}：母线向量按
\fld{sys.ac.buses}/\fld{sys.dc.buses} 顺序，设备向量按对应容器顺序；归因
映射经 \texttt{ComponentRef\{original\_index, source\_type\}}（:ComponentRef）
回指组件稳定 \fld{index}。量纲除注明外为有名值（MW/MVAr/kV 制下的 pu）。

\begin{paramtable}
\fld{vm} & $V_m$ & pu & 约 0.9--1.1 & 空 & AC 母线电压幅值（authored bus 顺序）\\
\fld{va} & $\theta$ & rad & $-\pi$--$\pi$ & 空 & AC 母线电压相角（HTTP 序列化名 \fld{va\_rad}，\srcpath{runtime\_api\_v1.cpp:serialize\_ac\_opf}）\\
\fld{vdc} & $V_{dc}$ & pu & 约 0.9--1.1 & 空 & DC 母线电压（标幺）\\
\fld{pg\_mw} & $P_g$ & MW & $P_{\min}$--$P_{\max}$ & 空 & 发电机有功（原始发电机顺序，外部电网已剥离）\\
\fld{qg\_mvar} & $Q_g$ & MVAr & $Q_{\min}$--$Q_{\max}$ & 空 & 发电机无功\\
\fld{external\_grid\_p\_mw} & $P_{\mathrm{eg}}$ & MW & — & 空 & 外部电网有功（由扩展发电机行剥离回填，src/optimal\_power\_flow/ac\_opf.cpp:separate\_external\_grid\_dispatch（lambda））\\
\fld{external\_grid\_q\_mvar} & $Q_{\mathrm{eg}}$ & MVAr & — & 空 & 外部电网无功\\
\fld{dpd\_mw} & $\Delta P^d$ & MW & $\ge 0$ & 空 & 各母线有功切负荷量\\
\fld{dqd\_mvar} & $\Delta Q^d$ & MVAr & $\ge 0$ & 空 & 各母线无功切负荷量\\
\fld{pac\_mw} & $P^{ac}_c$ & MW & $\pm S_N$ & 空 & VSC 交流侧有功\\
\fld{qac\_mvar} & $Q^{ac}_c$ & MVAr & $\pm S_N$ & 空 & VSC 交流侧无功\\
\fld{lcc\_transfers} & — & — & — & 空 & LCC 准稳态运行点（按 \fld{LCCConverter::index} 键控，字段注释 include/hacdcpf/optimal\_power\_flow/opf\_result.hpp:ACOPFResult）\\
\fld{pren\_mw} / \fld{qren\_mvar} & $P^{ren},Q^{ren}$ & MW/MVAr & $\ge 0$ & 空 & 可再生能源出力\\
\fld{pstor\_mw} / \fld{qstor\_mvar} & $P^{stor},Q^{stor}$ & MW/MVAr & $\pm S_N$ & 空 & 储能充放功率\\
\fld{pdcdc\_mw} & $P^{dcdc}$ & MW & $\pm P_N$ & 空 & DC/DC 变换器传输功率\\
\fld{pflex\_mw} & $P^{flex}$ & MW & — & 空 & 柔性负荷调节量\\
\fld{er\_port\_p\_mw} / \fld{er\_port\_q\_mvar} & $P^{er},Q^{er}$ & MW/MVAr & — & 空 & 能量路由器端口功率\\
\fld{lmp\_p} & $\lambda^P$ & — & — & 空 & AC 有功节点电价（币种未定；由等式对偶换算，src/optimal\_power\_flow/ac\_opf.cpp:solve\_with\_parity\_ipm）\\
\fld{lmp\_q} & $\lambda^Q$ & — & — & 空 & AC 无功节点电价\\
\fld{lmp\_valid} & — & — & true/false & false & 仅当价格来自本解所出建模的认证等式乘子时为真（字段注释，include/hacdcpf/optimal\_power\_flow/opf\_result.hpp:ACOPFResult）\\
\fld{lmp\_validity\_reason} & — & — & — & 空 & \fld{lmp\_valid} 判定的人读原因\\
\fld{ipm\_primal\_state} & $x$ & — & — & 空 & 原生 IPM 延拓状态（原始变量），供 \fld{warm\_start} 复用（字段注释 :109--112）\\
\fld{ipm\_equality\_dual\_state} & $\lambda$ & — & — & 空 & 延拓状态：等式对偶\\
\fld{ipm\_inequality\_dual\_state} & $\mu$ & — & — & 空 & 延拓状态：不等式对偶\\
\fld{ipm\_slack\_state} & $s$ & — & — & 空 & 延拓状态：松弛变量\\
\fld{ipm\_tangent\_primal} & $\mathrm{d}x/\mathrm{d}t$ & — & — & 空 & Davidenko 切向量（原始变量块），仅 \fld{compute\_homotopy\_tangent} 时填充（:ACOPFResult）\\
\fld{ipm\_tangent\_slack} & $\mathrm{d}s/\mathrm{d}t$ & — & — & 空 & 切向量：松弛块\\
\fld{ipm\_tangent\_equality\_dual} & $\mathrm{d}\lambda/\mathrm{d}t$ & — & — & 空 & 切向量：等式对偶块\\
\fld{ipm\_tangent\_inequality\_dual} & $\mathrm{d}\mu/\mathrm{d}t$ & — & — & 空 & 切向量：不等式对偶块\\
\fld{gen\_map} & — & — & — & 空 & 发电机归因映射（\fld{original\_index}+\fld{source\_type}）\\
\fld{ren\_map} & — & — & — & 空 & 可再生能源归因映射\\
\fld{stor\_map} & — & — & — & 空 & 储能归因映射\\
\fld{dcdc\_map} & — & — & — & 空 & DC/DC 归因映射\\
\fld{flex\_map} & — & — & — & 空 & 柔性负荷归因映射\\
\fld{er\_port\_map} & — & — & — & 空 & 能量路由器端口归因映射\\
\fld{converged} & — & — & true/false & false & 收敛标记（Ipopt 路径含“可接受残差”判据，src/optimal\_power\_flow/ac\_opf.cpp:solve\_parity\_with\_ipopt）\\
\fld{iterations} & — & — & $\ge 0$ & 0 & 总迭代数\\
\fld{outer\_iterations} & — & — & $\ge 0$ & 0 & 外环迭代数（传统路径）\\
\fld{objective} & $f$ & — & $\ge 0$ & 0 & 目标值（成本币种未定）\\
\fld{max\_constraint\_violation} & — & pu & $\ge 0$ & 0 & 返回点最大约束违反\\
\fld{max\_stationarity} & — & pu & $\ge 0$ & 0 & 返回点驻点残差\\
\fld{status} & — & — & — & 空 & 人读状态串（引擎标签原样保留，如 ``Ipopt: \ldots''）\\
\fld{solver\_path} & — & — & Unknown/\allowbreak NativeAC/\allowbreak ParityIPM & Unknown & 求解路径枚举（include/hacdcpf/optimal\_power\_flow/opf\_result.hpp:OPFSolverPath）。注意 Ipopt 解也记 \fld{ParityIPM}，真实内层引擎见 \fld{profiling.linear\_solver\_backend}（src/optimal\_power\_flow/ac\_opf.cpp:solve\_with\_parity\_ipm）；经济调度快速路径不置该字段，保持 \fld{Unknown}\\
\fld{profiling} & — & — & — & 默认 & 求解剖析；含 KKT 计数/残差及 Phase I 前后 violation、dual-fit、预算、终止、线性后端与 Phase II 接受状态\\
\fld{infeasibility\_hints} & — & — & — & 空 & 不可行/不收敛的可能原因（\fld{converged=false} 时由求解器和/或 \srcpath{verify\_opf\_result} 填充，include/hacdcpf/optimal\_power\_flow/opf\_result.hpp:ACOPFResult）\\
\fld{audit} & — & — & — & 默认 & 独立求解后校核（\fld{OpfAudit}，见下表）\\
\fld{converter\_model\_scope} & — & — & — & 默认 & 换流器模型保真度声明（见下“诚实口径”）\\
\fld{model\_limitations} & — & — & — & 空 & 本求解路径的能力边界声明（见下“诚实口径”）\\
\fld{ac\_switch\_flows} & — & — & — & 空 & OPF 调度点上的开关端潮流（门面层填充，src/api/hacdcpf.cpp:solve\_ac\_opf）\\
\fld{ac\_circuit\_breaker\_flows} & — & — & — & 空 & OPF 调度点上的断路器端潮流\\
\end{paramtable}

\subsubsection*{OpfAudit 字段}

\compmeta{OpfAudit}{\srcpath{include/hacdcpf/optimal\_power\_flow/opf\_result.hpp:OpfAudit}}

\fld{OpfAudit} 是独立于求解器的求解后可行性审计，由门面函数
\srcpath{verify\_opf\_result} 填充（声明
\srcpath{include/hacdcpf/api/hacdcpf.hpp:verify\_opf\_result}；实现
\srcpath{src/api/hacdcpf.cpp:verify\_opf\_result}）。HTTP v1 接口在序列化 AC OPF
结果前自动执行该审计（\srcpath{src/server/runtime\_api\_v1.cpp:serialize\_ac\_opf}）。

\begin{paramtable}
\fld{audited} & — & — & true/false & false & 审计已执行标记（src/api/hacdcpf.cpp:verify\_opf\_result）\\
\fld{max\_power\_balance\_violation\_mw} & — & MW & $\ge 0$ & 0 & 母线功率平衡最大违反。\textbf{当前实现未填充}：\srcpath{verify\_opf\_result} 只检查电压限值、发电机限值与目标一致性（src/api/hacdcpf.cpp:verify\_opf\_result），该字段保持默认 0\\
\fld{max\_voltage\_limit\_violation\_pu} & — & pu & $\ge 0$ & 0 & 电压越限最大值（src/api/hacdcpf.cpp:verify\_opf\_result）\\
\fld{max\_branch\_limit\_violation\_pu} & — & pu & $\ge 0$ & 0 & 支路负载越限最大值。\textbf{当前实现未填充}（同上）\\
\fld{max\_gen\_limit\_violation\_mw} & — & MW & $\ge 0$ & 0 & 发电机出力越限最大值（src/api/hacdcpf.cpp:verify\_opf\_result）\\
\fld{objective\_recomputed} & — & — & $\ge 0$ & 0 & 由 \fld{pg\_mw}/\fld{qg\_mvar} 重算的目标（src/api/hacdcpf.cpp:verify\_opf\_result）\\
\fld{objective\_reported} & — & — & $\ge 0$ & 0 & 求解器上报的目标（src/api/hacdcpf.cpp:verify\_opf\_result）\\
\fld{objective\_discrepancy\_pct} & — & \% & $\ge 0$ & 0 & 重算与上报目标的相对偏差（src/api/hacdcpf.cpp:verify\_opf\_result）\\
\fld{violations} & — & — & — & 空 & 人读违约描述列表\\
\fld{feasible()} & — & — & true/false & — & 成员函数：\fld{violations} 为空即真（include/hacdcpf/optimal\_power\_flow/opf\_result.hpp:OpfAudit）\\
\end{paramtable}

需要如实指出：头文件注释（include/hacdcpf/api/hacdcpf.hpp:verify\_opf\_result）称该审计“Checks power
balance, voltage limits, branch limits, generator limits, and objective
consistency”，但当前实现仅覆盖电压限值、发电机限值与目标一致性三项
（另在 \fld{converged=false} 时追加无源负荷、限值冲突等
\fld{infeasibility\_hints}，src/api/hacdcpf.cpp:verify\_opf\_result）；功率平衡与支路限值
两个审计字段当前恒为 0，消费方不得据此断言对应可行性。

\subsubsection*{DCOPFResult 字段}

\compmeta{DCOPFResult}{\srcpath{include/hacdcpf/optimal\_power\_flow/opf\_result.hpp:DCOPFResult}}

\begin{paramtable}
\fld{va} & $\theta$ & rad & $-\pi$--$\pi$ & 空 & AC 母线相角（authored bus 顺序）\\
\fld{pg\_mw} & $P_g$ & MW & $P_{\min}$--$P_{\max}$ & 空 & 发电机有功\\
\fld{external\_grid\_p\_mw} & $P_{\mathrm{eg}}$ & MW & — & 空 & 外部电网有功（提升为 slack 发电机后剥离回填，src/optimal\_power\_flow/dc\_opf.cpp:separate\_external\_grid\_dispatch（lambda））\\
\fld{pf\_mw} & $P_f$ & MW & $\pm S_{\max}$ & 空 & 支路有功潮流（映射回原始支路顺序，src/optimal\_power\_flow/dc\_opf.cpp:separate\_external\_grid\_dispatch（lambda））\\
\fld{converged} & — & — & true/false & false & 收敛标记\\
\fld{iterations} & — & — & $\ge 0$ & 0 & 迭代数\\
\fld{objective} & $f$ & — & $\ge 0$ & 0 & 目标值（语义见 \fld{objective\_model}）\\
\fld{status} & — & — & — & 空 & 人读状态串\\
\fld{solver\_name} & — & — & — & 空 & 实际求解后端名（如 \fld{NativeLCQP}、\fld{HiGHS-LP(QP-fallback)}，src/optimal\_power\_flow/dc\_opf.cpp:extract\_dc\_opf\_result、1105）\\
\fld{runtime\_sec} & — & s & $\ge 0$ & 0 & 墙钟耗时\\
\fld{structural\_warm\_start\_*} / \fld{solver\_initial\_primal\_residual} & — & pu/计数/布尔 & — & inf/0/false & NativeLCQP 结构初值的请求、构建、消费、分量/分解、原坐标 residual 和第 0 次迭代 residual\\
\fld{lmp} & $\lambda$ & — & — & 空 & 节点电价（等式行对偶恢复）\\
\fld{branch\_mu\_lower} / \fld{branch\_mu\_upper} & $\mu^-,\mu^+$ & — & — & 空 & 支路限值对偶（当前不可靠，见 \fld{branch\_mu\_valid}）\\
\fld{branch\_mu\_valid} & — & — & true/false & false & 恒为 false：原生支撑 LP 不能稳健恢复盒式变量对偶（字段注释 include/hacdcpf/optimal\_power\_flow/opf\_result.hpp:DCOPFResult；实现 src/optimal\_power\_flow/dc\_opf.cpp:extract\_dc\_opf\_result，设计说明 :621--634）\\
\fld{branch\_mu\_validity\_reason} & — & — & — & 空 & 上述判定的人读原因\\
\fld{lmp\_valid} & — & — & true/false & false & \fld{lmp} 是否来自报告原始解的等式行对偶（不蕴含阻塞对偶有效）\\
\fld{lmp\_validity\_reason} & — & — & — & 空 & 判定原因；\fld{compute\_lmp=false} 时注明被选项关闭（src/optimal\_power\_flow/dc\_opf.cpp:solve\_dc\_opf）\\
\fld{load\_shedding\_mw} & $\Delta P^d$ & MW & $\ge 0$ & 空 & 各母线切负荷量（src/optimal\_power\_flow/dc\_opf.cpp:extract\_dc\_opf\_result）\\
\fld{total\_load\_shedding\_mw} & — & MW & $\ge 0$ & 0 & 总切负荷量\\
\fld{solver\_chain} & — & — & — & 空 & 依次尝试的后端链，形如 ``\textless backend\textgreater:\textless status\textgreater''（字段注释 include/hacdcpf/optimal\_power\_flow/opf\_result.hpp:DCOPFResult；记录点 src/optimal\_power\_flow/dc\_opf.cpp:solve\_dc\_opf）\\
\fld{objective\_model} & — & — & QP/LP/LP-PWL & 空 & 上报目标实际使用的成本模型（实现 src/optimal\_power\_flow/dc\_opf.cpp:solve\_dc\_opf；头文件注释 :222--226 只列 ``QP''/``LP''，代码另发出 ``LP-PWL''，以代码为准）\\
\fld{pwl\_segments\_effective} & — & — & $\ge 0$ & 0 & LP 近似实际使用的线性段数；0 表示真 QP 成本\\
\fld{converter\_model\_scope} & — & — & — & 默认 & 恒为 ``dc-opf:linear-no-converter-model''，全部换流器标志为 false（字段注释 include/hacdcpf/optimal\_power\_flow/opf\_result.hpp:DCOPFResult；实现 src/optimal\_power\_flow/dc\_opf.cpp:extract\_dc\_opf\_result）\\
\fld{model\_limitations} & — & — & — & 空 & 能力边界声明（见下）\\
\end{paramtable}

\subsubsection*{诚实口径：model\_scope、ValidityFlags 与 model\_limitations}

平台约定每个 OPF 结果必须机器可读地声明自己\textbf{实际}建模了哪些换流器
物理、省略了哪些物理或证书，而不是让下游假定全保真。载体是
\fld{ConverterModelScope}
（\srcpath{include/hacdcpf/model/converter\_model\_scope.hpp:ConverterModelScope}）：
短标签 \fld{model\_scope} 加 12 个按特性划分的 \fld{ValidityFlags} 布尔位
（:ValidityFlags，覆盖 VSC 损耗、容量圆、电流限值、调制度、VDC 控制、多源协调、
DC/DC 损耗与占空比、LCC 准稳态与分接头、方程封闭性检查）。

AC OPF Parity 路径的标签\textbf{按本次运行实际强制的约束族拼接}——被
调用方关闭的约束族不出现在标签中
（\srcpath{src/optimal\_power\_flow/ac\_opf.cpp:solve\_with\_parity\_ipm}），例如
\fld{enforce\_converter\_capacity=false} 时标签不含
\fld{+capacity-circle}，对应 \fld{validity} 位同步置 false（:solve\_with\_parity\_ipm）。
传统 Native 路径标签固定为
\fld{ac-opf-native:vsc-free-pq+capacity-circle} 且
\fld{vsc\_vdc\_control\_modelled=false}（:solve\_ac\_opf\_impl）。DC OPF 标签固定为
\fld{dc-opf:linear-no-converter-model}（src/optimal\_power\_flow/dc\_opf.cpp:extract\_dc\_opf\_result），LCC 拒绝情形为
\fld{dc-opf:rejected-lcc}（:solve\_dc\_opf）。

\fld{model\_limitations} 是人读字符串列表，专门记录近似、省略的物理与不可
用证书（字段注释 include/hacdcpf/optimal\_power\_flow/opf\_result.hpp:ACOPFResult）。当前代码会写入的条目包括：

\begin{itemize}
  \item 非凸性声明：``AC OPF is a non-convex nonlinear program; convergence
        certifies a local KKT point, not a global optimum.''——所有 AC OPF
        收敛结果恒带（src/optimal\_power\_flow/ac\_opf.cpp:separate\_external\_grid\_dispatch（lambda））；
  \item 储能为单时段功率注入、不建模跨时段 SOC 动态（AC：:3207--3210；
        DC：src/optimal\_power\_flow/dc\_opf.cpp:extract\_dc\_opf\_result）；
  \item 能量路由器端口平衡为无损模型（src/optimal\_power\_flow/ac\_opf.cpp:separate\_external\_grid\_dispatch（lambda））；
  \item 传统 Native AC 路径经外部栅栏/罚环处理非线性工程限值，建议生产
        混合建模使用 ParityIPM（:separate\_external\_grid\_dispatch（lambda））；
  \item LCC 相关四条：功率/电流/$\alpha$/$\gamma$ 指令与换流变分接头为固定
        输入而非经济调度变量；额定电流饱和建模但 $\alpha$/$\gamma$ 限值为
        求解后审计而非硬约束；换相重叠迭代、分接头控制与主动损耗未建模；
        LCC 等式二阶导用解析雅可比的局部有限差分（:solve\_with\_parity\_ipm）；整流/
        逆变站求解后触发角/关断角窗口违例时逐站追加（:solve\_with\_parity\_ipm）；
  \item DC OPF 无损线性化、不建模电压幅值/无功/换流器物理
        （src/optimal\_power\_flow/dc\_opf.cpp:extract\_dc\_opf\_result）；LCC 拒绝情形（:solve\_dc\_opf）；
  \item LMP/阻塞对偶不可用时的原因说明（src/optimal\_power\_flow/ac\_opf.cpp:solve\_with\_parity\_ipm、
        3215--3218；src/optimal\_power\_flow/dc\_opf.cpp:extract\_dc\_opf\_result）。
\end{itemize}

\subsection{问题描述与求解器接口}\label{subsec:overview-interface}

\subsubsection*{OPFProblem：自包含任务描述}

\compmeta{OPFProblem}{\srcpath{include/hacdcpf/optimal\_power\_flow/opf\_problem.hpp:OPFProblem}}

\fld{OPFProblem} 把“优化什么”与“如何求解”解耦：一个非拥有系统指针
\fld{sys}、一份 \fld{ACOPFOptions}、一份 \fld{DCOPFOptions}，加
\fld{is\_ac}/\fld{is\_dc} 二选一标志。它由投影/调用层产生，可被任何
\fld{IOPFSolver} 实现消费。

\subsubsection*{IOPFSolver 与 DefaultOPFSolver}

\compmeta{IOPFSolver}{\srcpath{include/hacdcpf/optimal\_power\_flow/opf\_solver\_interface.hpp:IOPFSolver}}

求解器抽象只有一个纯虚方法
\texttt{solve(const OPFProblem\&) const}，返回类型
\fld{OPFSolveResult} 是 \texttt{std::variant<ACOPFResult, DCOPFResult>}
（:OPFSolveResult）。默认实现 \fld{DefaultOPFSolver}（:DefaultOPFSolver）不做第二套建模，直接
适配生产入口：\fld{sys} 为空或 \fld{is\_ac==is\_dc} 时抛
\fld{std::invalid\_argument}（:solve），否则按标志转发
\srcpath{hacdcpf::solve\_ac\_opf} 或 \srcpath{hacdcpf::solve\_dc\_opf}
（:solve）。该设计保证接口层与生产路径永远同源（头文件注释 :3--5）。

\subsubsection*{Parity 公共抽象层}

Parity 全空间建模的公共抽象（变量/约束布局与全部 KKT 求值函数）集中
在 \srcpath{include/hacdcpf/optimal\_power\_flow/formulation.hpp}，供
Parity IPM 与嵌入式 Ipopt 共用：

\begin{itemize}
  \item \fld{ParityOptions}（:ParityOptions）：建模族选项，含 \fld{load\_shedding}、
        \fld{voll}（0 表示按最大发电成本自动折算）、\fld{eps\_iac}
        （$|I_{ac}|$ 平滑）、目标与四个约束族开关；
  \item \fld{VarIndex}（:VarIndex）：全空间变量向量 $x$ 的连续分块布局
        （$\theta$、$V$、$P^g$、$Q^g$、$V^{dc}$、$P^{ac}_c$、$Q^{ac}_c$、
        $P^{dc}_c$、$\Delta P^d$、$\Delta Q^d$、$P^{ren}$、$Q^{ren}$、
        $P^{stor}$、$Q^{stor}$、$P^{stor,dc}$、$P^{dcdc}$、$P^{flex}$
        与能量路由器端口块），空块长度为零但保留有效偏移；
  \item \fld{ConstraintIndex}（:ConstraintIndex）：等式行（AC 有功/无功平衡、DC
        平衡、VSC/DC-DC/能量路由器耦合、AC 相角参考、DC 电压参考）与非线性
        不等式行（支路两端容量、VSC 容量圆、DC 支路、VSC 电流、调制度上下
        限、DC/DC 占空比）的连续行块布局；
  \item \fld{Problem}（:Problem）：装配完成的模型 + 全部模型到向量的映射
        （含 ZIP 系数、固定注入、$G/B$ 对角与 DC 电导矩阵）；
  \item 装配与求值函数：\srcpath{build\_problem}（:build\_problem）、
        \srcpath{build\_variable\_bounds}（:build\_variable\_bounds）、
        \srcpath{build\_initial\_point}（严格内点，:build\_initial\_point）、
        \srcpath{objective}（:objective）、
        \srcpath{objective\_gradient\_hessian\_diag}（:objective\_gradient\_hessian\_diag）、
        \srcpath{equality\_constraints}（:equality\_constraints）、
        \srcpath{equality\_jacobian}（:equality\_jacobian）、
        \srcpath{nonlinear\_inequality\_constraints}（:nonlinear\_inequality\_constraints）、
        \srcpath{nonlinear\_inequality\_jacobian}（:nonlinear\_inequality\_jacobian）、
        \srcpath{lagrangian\_hessian\_dense}/\srcpath{lagrangian\_hessian}
        （:lagrangian\_hessian）。数学口径详见第\ref{sec:parityform}节。
\end{itemize}

IPM 求解核的公共契约是 \fld{IPMOptions}/\fld{IPMResult} 与
\srcpath{solve\_primal\_dual\_ipm}
（\srcpath{include/hacdcpf/optimal\_power\_flow/native\_ipm\_solver.hpp:IPMOptions}、
\srcpath{include/hacdcpf/optimal\_power\_flow/native\_ipm\_solver.hpp:IPMResult}、
\srcpath{include/hacdcpf/optimal\_power\_flow/native\_ipm\_solver.hpp:solve\_primal\_dual\_ipm}），
另有 Davidenko 切向量接口 \srcpath{homotopy\_tangent}（:homotopy\_tangent）；详见
第\ref{sec:parityipm}节。高层驱动 \srcpath{solve\_ac\_opf\_parity}
（\srcpath{include/hacdcpf/optimal\_power\_flow/parity\_ipm\_solver.hpp:solve\_ac\_opf\_parity}）
串联装配、求解与结果回填，成功时置
\fld{solver\_path=ParityIPM}。雅可比对照诊断入口
\srcpath{compute\_jacobian\_diagnostics}
（\srcpath{include/hacdcpf/optimal\_power\_flow/ac\_opf\_solver.hpp:compute\_jacobian\_diagnostics}）
用于解析导数对有限差分的回归检查，见第\ref{sec:verification}节。

\subsection{嵌入式 Ipopt 后端：编译门控与运行时探测}\label{subsec:overview-ipopt}

嵌入式 Ipopt 是\textbf{可选}后端，其可用性由“编译期宏 + 运行时环境变量 +
能力查询”三层机制共同决定：

\begin{enumerate}
  \item \textbf{编译期门控}。CMake 选项 \srcpath{HACDCPF\_ENABLE\_IPOPT}
        仅 macOS 默认 ON、其余平台默认 OFF
        （\srcpath{CMakeLists.txt}:HACDCPF\_ENABLE\_IPOPT，帮助串为
        ``when supported by that checkout''）。非 Apple 平台显式开启
        \textbf{不再}触发 \fld{FATAL\_ERROR}：选项仅经
        \srcpath{MIPSOLVERS\_BUILD\_LOCAL\_IPOPT} 转发给 MIPSolvers
        子项目（同文件:MIPSOLVERS\_BUILD\_LOCAL\_IPOPT），能否真正构建
        取决于该 MIPSolvers 检出对内嵌 Ipopt 的支持；当前
        \srcpath{CMakeLists.txt} 中的 \fld{FATAL\_ERROR} 只剩 IPO 与
        测试联用、OpenDSS 库缺失等少数守卫。
        构建成功时 MIPSolvers 子项目向全部依赖目标传播宏
        \srcpath{HACDCPF\_HAVE\_IPOPT=1}
        （\srcpath{MIPSolvers/CMakeLists.txt:HACDCPF\_HAVE\_IPOPT}）。OPF 门面代码
        以 \texttt{\#ifdef HACDCPF\_HAVE\_IPOPT} 包裹引擎头文件包含
        （\srcpath{src/optimal\_power\_flow/ac\_opf.cpp:HACDCPF\_HAVE\_IPOPT}）与整个
        NLP 适配结构 \fld{ParityNLPInput}/\srcpath{build\_parity\_nlp\_input}
        （:ParityNLPInput、:build\_parity\_nlp\_input）。
  \item \textbf{运行时探测与降级}。统一入口
        \srcpath{solve\_parity\_with\_ipopt}
        （\srcpath{src/optimal\_power\_flow/ac\_opf.cpp:solve\_parity\_with\_ipopt}）以
        出参 \fld{available} 报告可用性：未编译时直接返回
        \fld{available=false}、\texttt{status=``Ipopt not compiled in''}
        （:solve\_parity\_with\_ipopt）；已编译时还可用环境变量
        \srcpath{HACDCPF\_DISABLE\_IPOPT\_OPF} 强制禁用（用于后端隔离与
        诊断，:solve\_parity\_with\_ipopt）。调用方按 \fld{available} 降级：显式
        \fld{Ipopt} 请求但不可用时回退原生 Parity IPM 并在后端标签注明
        ``(ipopt unavailable)''（:solve\_with\_parity\_ipm）；\fld{Auto} 模式先跑原生
        Parity IPM，不收敛且 \fld{allow\_fallback=true} 时尝试 Ipopt，
        成功则标记 ``ipopt\_filter\_linesearch (auto-fallback)''
        （:solve\_with\_parity\_ipm）。
  \item \textbf{能力查询}。\fld{SolverCapabilities::has\_ipopt} 在同一宏下
        置真（\srcpath{src/api/hacdcpf.cpp:get\_solver\_capabilities}），供 GUI/脚本在
        提交任务前探测（见第\ref{subsec:overview-capabilities}节）。
\end{enumerate}

Ipopt 路径还有两条必须知晓的诚实口径：其一，Ipopt 在大型混合模型上可能
以 \fld{Maximum\_Iterations\_Exceeded} 停机但已到达高度可用点，门面在原始
可行性严格满足、对偶/互补残差进入“可接受带”（对偶残差
$\le 1000\,\varepsilon_{\mathrm{dual}}$、互补残差
$\le \max(10^{-2},\,1000\,\varepsilon_{\mathrm{compl}})$，src/optimal\_power\_flow/ac\_opf.cpp:solve\_parity\_with\_ipopt）
时接受该点并在 \fld{status} 注明
``acceptable residuals after \ldots''（:solve\_parity\_with\_ipopt）；其二，
适配器不把约束乘子回传到 LMP 提取链路，Ipopt 解的
\fld{lmp\_valid=false}，原因为 ``The embedded Ipopt adapter does not
return constraint multipliers''（:solve\_with\_parity\_ipm）。Ipopt 终止时若对偶向量
维度齐全，门面会把约束/盒式对偶映射回 Parity 行序并做正性修复
（:solve\_parity\_with\_ipopt），供互补性度量使用，但这不改变 LMP 不可用的结论。

\subsection{OPF 在平台中的调用方速览}\label{subsec:overview-callers}

\begin{itemize}
  \item \textbf{时序生产模拟（UC $\to$ OPF $\to$ PF）}：
        \srcpath{src/time\_series/time\_series\_pf.cpp} 的三段流水线在
        UC 出力计划之后对每个时段调用 \srcpath{opf::solve\_ac\_opf}
        （:solve\_time\_series\_pf），上一时段收敛结果经 \fld{warm\_start} 传入下一
        时段（:solve\_time\_series\_pf），OPF 调度回写系统后再做牛顿 PF 校验；OPF 不
        收敛的时段回退为“UC 计划直接 PF”并如实记录（:solve\_time\_series\_pf）。
        按日并行分解时，各日分片强制使用线程安全的 ParityIPM 后端
        （:solve\_time\_series\_pf）。
  \item \textbf{电力市场（SCUC/SCED）}：
        \srcpath{src/market/market\_simulation.cpp} \textbf{不调用}本章
        的 \srcpath{solve\_dc\_opf}/\srcpath{solve\_ac\_opf}，而是在
        engine 层自建 SCUC/SCED 模型并用原生 B\&C/HiGHS/SCIP/Gurobi
        适配器求解（适配器选择 :803--839，SCED 装配与校验
        :2440--2469）。二者同属 DC 线性近似家族，但实现独立，价格/结算
        口径以市场模块自身结果为准。
  \item \textbf{网络重构}：
        \srcpath{src/network\_reconfiguration/topology\_reconfiguration.cpp}
        同样不经过 OPF 模块：其 MILP 自带虚拟潮流变量布局（文件头注释
        :9--17），经 engine 层 B\&C/MILP 适配器求解（:try\_native\_bc（lambda））。
  \item \textbf{其他内部调用方}：可靠性评估以 DC OPF 做削负荷最小化
        （\srcpath{src/reliability/reliability\_assessment.cpp:evaluate\_state}）；
        EV-交通耦合的联合社会福利与 CTM 模型内嵌 DC OPF
        （\srcpath{src/ev\_power\_traffic/joint\_social\_welfare.cpp:solve\_joint\_social\_welfare}、
        \srcpath{src/ev\_power\_traffic/ctm\_joint\_welfare.cpp:simulate\_ev\_power\_traffic\_ctm\_joint}）；
        反事实规划调用 AC OPF
        （\srcpath{src/analysis/counterfactual\_planning.cpp}\allowbreak :144）；
        SPPT 变形测试以 DC OPF 做投影前后对照
        （\srcpath{src/sppt/metamorphic.cpp}\allowbreak :223--225）。
  \item \textbf{HTTP 运行时 API（v1）}：GUI/脚本经
        \texttt{POST /api/v1/sessions/\{id\}/jobs} 提交
        \fld{analysis="optimal\_power\_flow"} 任务
        （路由 \srcpath{src/server/runtime\_api\_v1.cpp:register\_routes}；
        分派 :864--881）。请求体 \fld{solver} 字段映射为
        \fld{ac\_solver\_backend}（\fld{parity}/\fld{ipopt}/\fld{auto}/
        \fld{dispatch}，:ac\_opf\_options）或 DC OPF（\fld{solver="dc"}，
        :874--876）；\fld{network\_model} 当前仅支持
        \fld{balanced\_aggregate}（:execute\_analysis）。AC 结果序列化前自动执行
        \srcpath{verify\_opf\_result} 审计（:serialize\_ac\_opf），响应携带
        \fld{model\_scope}、\fld{validity\_flags}、
        \fld{model\_limitations}、\fld{infeasibility\_hints} 与
        \fld{audit}（:serialize\_ac\_opf）；DC 结果另携带 \fld{solver\_chain}、
        \fld{objective\_model} 与 \fld{branch\_mu\_valid}（:serialize\_dc\_opf）。
\end{itemize}

\subsection{求解器能力运行时查询}\label{subsec:overview-capabilities}

\compmeta{SolverCapabilities}{\srcpath{include/hacdcpf/api/solver\_capabilities.hpp:SolverCapabilities}}

\srcpath{hacdcpf::get\_solver\_capabilities()}（同头文件:44；实现
\srcpath{src/api/hacdcpf.cpp:get\_solver\_capabilities}）返回按编译结果填充的能力结构，
开销极小、可反复调用。恒真字段（内置 Eigen SparseLU、自研 IPM、AC/DC OPF、
三相潮流、二次目标）在:1983--1988 直接置真；可选后端按
\srcpath{HACDCPF\_HAVE\_*} 系列宏逐个置位（:get\_solver\_capabilities），这些宏由
MIPSolvers 子项目在找到对应库时传播（如
\srcpath{MIPSolvers/CMakeLists.txt:HACDCPF\_HAVE\_IPOPT} 对 Ipopt 的传播）。

\begin{paramtable}
\fld{has\_sparse\_lu} & — & — & true/false & true & Eigen SparseLU（恒有）\\
\fld{has\_klu} & — & — & true/false & false & SuiteSparse KLU（\texttt{\#ifdef HACDCPF\_HAVE\_KLU}，src/api/hacdcpf.cpp:get\_solver\_capabilities）\\
\fld{has\_umfpack} & — & — & true/false & false & SuiteSparse UMFPACK（:get\_solver\_capabilities）\\
\fld{has\_accelerate} & — & — & true/false & false & Apple Accelerate 框架（仅 macOS，:get\_solver\_capabilities）\\
\fld{has\_highs} & — & — & true/false & false & HiGHS LP/MIP；同时置 \fld{supports\_integer\_variables}（:get\_solver\_capabilities）\\
\fld{has\_gurobi} & — & — & true/false & false & Gurobi LP/QP/MIP（:get\_solver\_capabilities）\\
\fld{has\_ipopt} & — & — & true/false & false & 嵌入式 Ipopt NLP（:get\_solver\_capabilities；对应 CMake 门控见第\ref{subsec:overview-ipopt}节）\\
\fld{has\_scip} & — & — & true/false & false & SCIP MIP（:get\_solver\_capabilities）\\
\fld{has\_papilo} & — & — & true/false & false & PaPILO 预处理库（:get\_solver\_capabilities）\\
\fld{has\_native\_ipm} & — & — & true/false & true & 内置原始对偶 IPM（恒有）\\
\fld{has\_excel} & — & — & true/false & false & ETAP Excel I/O（\texttt{\#ifdef HACDCPF\_ENABLE\_ETAP}，:get\_solver\_capabilities）\\
\fld{has\_opendss} & — & — & true/false & false & OpenDSS 桥（\texttt{\#ifndef HACDCPF\_NO\_OPENDSS}，:get\_solver\_capabilities）\\
\fld{supports\_quadratic\_objective} & — & — & true/false & true & 二次目标（原生 + HiGHS 路径）\\
\fld{supports\_integer\_variables} & — & — & true/false & false & 整数变量（需 HiGHS/Gurobi/SCIP 任一）\\
\fld{supports\_ac\_opf} & — & — & true/false & true & 非线性 AC OPF（恒有）\\
\fld{supports\_dc\_opf} & — & — & true/false & true & 线性 DC OPF（恒有）\\
\fld{supports\_three\_phase} & — & — & true/false & true & 三相不平衡潮流（恒有）\\
\end{paramtable}

注意该查询只反映\textbf{编译进本构建}的后端，不代表运行时一定可用：
Gurobi 还需许可证（DC OPF 的 \srcpath{try\_gurobi\_qp} 每次调用都做
\fld{available()} 探测，src/optimal\_power\_flow/dc\_opf.cpp:try\_gurobi\_qp（lambda））；Ipopt 还可被环境变量
\srcpath{HACDCPF\_DISABLE\_IPOPT\_OPF} 临时禁用
（src/optimal\_power\_flow/ac\_opf.cpp:solve\_parity\_with\_ipopt）。因此生产代码应把能力查询作为提交前提示，
并以结果中的 \fld{solver\_name}/\fld{solver\_chain}/
\fld{profiling.linear\_solver\_backend} 作为实际所用后端的最终依据。

\subsection{本章小结与阅读路线}\label{subsec:overview-reading}

OPF 模块的公共契约可概括为：统一 Rich 输入、投影装配、按后端分派求解、
结果归因回稳定 ID，并以 \fld{model\_scope}/\fld{ValidityFlags}/
\fld{model\_limitations} 三元组机器可读地声明每次求解实际兑现的物理与
证书边界。后续章节按 AC OPF、DC OPF、交直流混合 OPF 三个主类与各专项
逐一展开：AC OPF（Native 紧致路径）见第\ref{sec:acopfnative}节，
DC OPF 见第\ref{sec:dcopf}节，交直流混合 OPF 的 Parity 建模层与 IPM 求解核
见第\ref{sec:parityform}、\ref{sec:parityipm}节，
RPO 见第\ref{sec:rpo}节，三相混合 OPF 见
第\ref{sec:threephase}节；AML 建模层见独立文档
\srcpath{docs/modules/power\_models/power\_models\_manual.tex}；跨引擎
验证与校核通道见第\ref{sec:verification}节。

% AC 最优潮流：Native IPM 求解器
% 本章文件由 \input 引入，不含导言区。
\section{AC 最优潮流：Native IPM 求解器}\label{sec:acopfnative}

本章给出平台自研 AC 最优潮流（AC OPF）引擎的技术口径。求解入口为
\allowbreak\srcpath{hacdcpf::opf::solve\_ac\_opf}\allowbreak （声明于
\srcpath{include/hacdcpf/optimal\_power\_flow/ac\_opf\_solver.hpp}:42--43，
实现于 \srcpath{src/optimal\_power\_flow/ac\_opf.cpp}:3038--4209）。该实现包含两条
自研数值路径：一条是内嵌于 \srcpath{ac\_opf.cpp} 的紧致变量原生内点循环
（结果标记为 \fld{OPFSolverPath::\allowbreak NativeAC}），另一条是作用于 Parity 全空间
建模的原始对偶内点核 \srcpath{parity::solve\_primal\_dual\_ipm}\allowbreak （接口声明于
\srcpath{include/hacdcpf/optimal\_power\_flow/native\_ipm\_solver.hpp}:90，
结果标记为 \fld{OPFSolverPath::\allowbreak ParityIPM}）；两者共享同一组选项与结果类型，
并可在编译入 Ipopt 时自动回退到嵌入式 Ipopt。需要特别说明：尽管接口文档以
``polar AC balance'' 描述核心等式（\srcpath{ac\_opf\_solver.hpp}:18--31），
两条路径的电压变量均为\textbf{极坐标}（幅值 $V$ 与相角 $\theta$），而非直角坐标
实虚部；本章全部公式以代码实际采用的极坐标形式给出。

\subsection{功能定位与求解路径}\label{subsec:acopfnative:role}

\subsubsection*{功能定位}
Native IPM 是平台交直流混合 OPF 的自研引擎层，承担以下职责
（\srcpath{src/optimal\_power\_flow/ac\_opf.cpp}:3038--4209）：
\begin{itemize}
  \item 对 rich 模型 \fld{HybridPowerSystem} 做合法性门禁（空 AC 母线、LCC 就绪性、
        参考母线资格、图拓扑孤岛），随后经 \srcpath{core::make\_solver\_data}
        装配 canonical \fld{SolverData}（含 Ybus 与直流电导矩阵 \fld{gdc}）；
  \item 在标幺体系（$S_{\mathrm{base}}$ 取算例 \fld{base\_mva}，典型 100~MVA）下组装
        非线性规划：极坐标节点功率平衡、直流节点平衡、换流器耦合等式，
        支路视在功率与换流器容量圆不等式，以及全体决策变量的盒式界；
  \item 以原始对偶内点框架求解，输出电压、发电调度、换流器功率、LMP 及
        诚实口径的 \fld{model\_limitations} / \fld{converter\_model\_scope} 声明。
\end{itemize}
与 Parity/Ipopt 的关系：Parity 全空间建模（见第~\ref{sec:parityform}~章）是
\textbf{生产推荐的建模层}，覆盖 VSC 控制模式、直流电压等式、DER 全家族与 LCC
准稳态；其内层牛顿引擎默认就是本章的 Native IPM 核
\srcpath{parity::solve\_primal\_dual\_ipm}（见第~\ref{sec:parityipm}~章），
Ipopt 仅作为可选回退。代码自身把内嵌于 \srcpath{ac\_opf.cpp} 的紧致路径标记为
legacy：结果中会写入建议改用 ParityIPM 的 \fld{model\_limitations} 条目
（\srcpath{src/optimal\_power\_flow/ac\_opf.cpp}:3219--3222）。Ipopt 的可用性由
编译宏 \srcpath{HACDCPF\_HAVE\_IPOPT} 门控（同文件:33--35、2267--2373），
并可用环境变量 \srcpath{HACDCPF\_DISABLE\_IPOPT\_OPF} 在运行时强制关闭
（同文件:2391--2396）。

\subsubsection*{求解路径选择}
\fld{ACOPFOptions::ac\_solver\_backend}（枚举 \fld{ACOPFSolverBackend}，
\srcpath{include/hacdcpf/optimal\_power\_flow/opf\_options.hpp}:28--33）给出四种
后端：\fld{Auto}、\fld{ParityIPM}、\fld{Ipopt}、\fld{EconomicDispatch}。入口处的
路径翻译规则为（\srcpath{src/optimal\_power\_flow/ac\_opf.cpp}:3309--3350）：
\begin{itemize}
  \item \fld{ParityIPM}：置 \fld{enable\_primal\_dual} 与 \fld{use\_parity\_ipm}
        两个标志；内层固定为自研核 \fld{ParityInnerSolver::\allowbreak NativeIPM}，
        不经 Ipopt 翻译；
  \item \fld{Ipopt}：同样进入 Parity 建模，内层固定为 \fld{ParityInnerSolver::Ipopt}，
        未编译 Ipopt 时回退自研核；
  \item \fld{EconomicDispatch}：清除 primal-dual 标志，强制纯电交流
        ``经济调度 + 单次 AC 潮流'' 快速路径；
  \item \fld{Auto}：保持调用方标志；只要算例含任何 DC/VSC/LCC 组件
        （判定见 \srcpath{contains\_hybrid\_acdc\_components}，同文件:3022--3034）
        或存在在运 LCC，即自动升级为 Parity 路径（同文件:3336--3350）。
        纯 AC 且未显式开启 primal-dual 时走快速路径（同文件:3501--3526）；
        纯 AC 且 \fld{enable\_primal\_dual=true}、\fld{use\_parity\_ipm=false}
        时才进入本章重点描述的 legacy 原生内点循环（同文件:3528--4098）。
\end{itemize}
整体决策流程如图~\ref{fig:acopfnative:flow}~所示。

\begin{figure}[H]
\centering
\begin{tikzpicture}[
  font=\footnotesize,
  node distance=0.55cm and 0.9cm,
  box/.style={draw, rectangle, align=center, minimum height=0.85cm, inner sep=3pt},
  dec/.style={draw, diamond, aspect=2.2, align=center, inner sep=1.5pt},
  arr/.style={-{Stealth[length=2.5mm]}, thick},
  lbl/.style={font=\scriptsize, midway, above}
]
\node[box] (entry) {入口校验\\空 AC / LCC 就绪 / 参考母线 / 拓扑};
\node[dec, below=of entry] (hybrid) {含 DC/VSC/LCC？\\或显式 Parity/Ipopt};
\node[box, right=2.6cm of hybrid] (legacy) {legacy 原生障碍--罚\\内点循环（NativeAC）};
\node[box, below=of hybrid] (parity) {Parity 全空间建模\\内层：Native IPM / Ipopt};
\node[box, left=2.6cm of hybrid] (ed) {经济调度 + AC PF\\（EconomicDispatch）};
\node[box, below=of parity] (homo) {目标同伦延续\\（\fld{objective\_homotopy}，可选）};
\node[box, below=of legacy] (fb) {ED+PF 回退\\（仅纯 AC）};
\node[box, below=of homo] (outp) {结果反投影 / 外部电网拆分 /\\ \fld{model\_limitations} 汇总};
\draw[arr] (entry) -- (hybrid);
\draw[arr] (hybrid) -- node[lbl]{否，且未开 primal-dual} (ed);
\draw[arr] (hybrid) -- node[lbl]{否，开 primal-dual 非 parity} (legacy);
\draw[arr] (hybrid) -- node[lbl, right=1pt]{是} (parity);
\draw[arr] (parity) -- (homo);
\draw[arr] (legacy) -- node[lbl, right=1pt]{未收敛} (fb);
\draw[arr] (homo) -- (outp);
\draw[arr] (ed.south) |- (outp.west);
\draw[arr] (fb.south) |- (outp.west);
\end{tikzpicture}
\caption{\srcpath{solve\_ac\_opf} 的求解路径决策流程（依据
\srcpath{src/optimal\_power\_flow/ac\_opf.cpp}:3309--3526 整理）}
\label{fig:acopfnative:flow}
\end{figure}

\subsubsection*{入口门禁}
在进入任何数值路径之前，求解器执行四组静态检查并\textbf{直接拒绝}不合格算例：
\begin{enumerate}
  \item 空 AC 母线集：返回 \fld{status}=``AC OPF failed: empty AC bus set.''
        （\srcpath{src/optimal\_power\_flow/ac\_opf.cpp}:3089--3092）；
  \item LCC 就绪性：对在运 LCC 逐一校验控制模式配对、\fld{n\_bridges} 与
        \fld{vn\_ac\_kv} 正值性、恒 $\alpha$/恒 $\gamma$ 控制的 \fld{x\_comm\_ohm}
        正值性、交流/直流端母线可解析、直流基准电压可得；任一失败即拒绝并写入
        \fld{infeasibility\_hints}（同文件:3094--3140）。注释明确说明动机：潮流的
        LCC 注入助手对不可用站点返回零以便诊断，而 OPF 静默丢弃固定传输量会改变
        可行性与最优解，故必须前置拒绝（同文件:3094--3098）；
  \item 参考母线资格：\srcpath{validation::validate\_reference\_bus\_eligibility}
        存在错误或``裸理想 DC\_V 边界''警告时拒绝，后者附加说明理想 DC\_V 边界
        没有调度成本模型，优化需要显式可控的直流平衡设备（同文件:3142--3161）；
  \item 图拓扑预检：\srcpath{graph::build\_power\_system\_graph} +
        \srcpath{analyze\_topology}，无任何含松弛母线的有效孤岛时判定不可行；
        对无松弛孤岛写入不可服务提示（同文件:3280--3307）。
\end{enumerate}

\subsubsection*{外部电网提升}
每个在运 \fld{ExternalGrid} 在求解前被提升为一台合成 \fld{Generator} 追加到
工作副本：\fld{vg\_pu} 取外部电网 \fld{vm\_pu}，P/Q 上下界取
$\pm\max(1.0,\;$\fld{s\_sc\_max\_mva}$,\;$\fld{base\_mva}$)$~MW/MVAr，成本系数
\fld{cost\_c0/c1/c2} 原样拷贝（\srcpath{src/optimal\_power\_flow/ac\_opf.cpp}:
3163--3196）。代码注释如实声明：短路容量并非进口限值，只是一个``有用的有限
尺度''，以防止过大的合成界使小算例的 KKT 系统病态（同文件:3173--3177）。
求解完成后由 \srcpath{separate\_external\_grid\_dispatch} 把合成机结果拆回
\fld{external\_grid\_p\_mw}/\fld{external\_grid\_q\_mvar}，并把
\fld{pg\_mw}/\fld{qg\_mvar}/\fld{gen\_map} 还原为仅含用户自建发电机的顺序
（同文件:3203--3259）。

\subsection{数学模型}\label{subsec:acopfnative:model}

以下数学模型为 legacy 紧致路径（\fld{NativeAC}）实际求解的形式；Parity 全空间
模型的扩展变量族（DER、储能、柔性负荷、能量路由器、LCC 等）见
第~\ref{sec:parityform}~章。除注明外，功率与电压均为标幺值
（$S_{\mathrm{base}}=$~\fld{base\_mva}）。

\subsubsection*{决策变量与索引布局}
变量向量 $x\in\mathbb{R}^{n}$ 由 \srcpath{build\_index} 按固定块序布局
（\srcpath{src/optimal\_power\_flow/ac\_opf.cpp}:225--344，结构体
\srcpath{ACOPFIndex}:93--142）：
\begin{align*}
x = \left[\,
\underbrace{\theta}_{n_p},\;
\underbrace{V}_{n_q},\;
\underbrace{P_g}_{n_g},\;
\underbrace{Q_g}_{n_g},\;
\underbrace{V_{dc}}_{n_{dc}},\;
\underbrace{P_{ac}}_{n_c},\;
\underbrace{Q_{ac}}_{n_c},\;
\underbrace{P_{dc}}_{n_c}
\,\right]
\end{align*}
各块偏移由 \fld{i\_theta}=0、\fld{i\_vm}=$n_p$ 起顺序累加（同文件:273--274、
326--332），总维数
\begin{align*}
n = n_p + n_q + 2n_g + n_{dc} + 3n_c
\end{align*}
（同文件:332）。其中 $n_p$ 为非松弛交流母线数、$n_q$ 为 PQ 母线数、$n_g$ 为
在运发电机数、$n_{dc}$ 为直流母线数、$n_c$ 为在运 VSC 换流器数。母线归属规则：
松弛母线不设 $\theta$ 变量；仅 PQ 母线设 $V$ 变量（同文件:261--269）；未找到
显式 SLACK 时退化取第 0 号母线为松弛（\srcpath{find\_slack\_bus}，
同文件:178--185）。发电机与换流器仅在 \fld{in\_service} 且端母线可解析时进入
变量集（同文件:287--318；组件 \fld{bus} 字段为 1-based，代码内统一减 1）。
PV/松弛母线的电压幅值、松弛母线的相角不作为变量，而是由参考向量固定：
\fld{vm\_reference\_by\_bus} 先取母线 \fld{vm\_pu}，再被该机端母线发电机的
\fld{vg\_pu} 覆盖（同文件:243--256），解包时先铺参考值再覆盖变量块
（\srcpath{unpack\_state}，同文件:466--538）。

等式右端向量 $g(x)\in\mathbb{R}^{m}$ 的块布局为
（同文件:334--338）：
\begin{align*}
g = \left[\,
\underbrace{g^{P}}_{n_p},\;
\underbrace{g^{Q}}_{n_q},\;
\underbrace{g^{dc}}_{n_{dc}},\;
\underbrace{g^{conv}}_{n_c}
\,\right],\qquad
m = n_p + n_q + n_{dc} + n_c
\end{align*}
故优化自由度恒为 $n-m = 2n_g + 2n_c$：发电机 P/Q 与换流器 P/Q 是仅有的
真正自由度，$V_{dc}$ 与直流平衡一一对应。

\subsubsection*{变量界与初值}
初值与盒式界由 \srcpath{initialize\_bounds\_and\_state} 给出
（\srcpath{src/optimal\_power\_flow/ac\_opf.cpp}:366--464），所有界先经
\srcpath{sanitize\_bounds} 清洗：非有限值回退默认、上下颠倒交换、宽度不足
\fld{kMinInteriorWidth}=$10^{-6}$ 时以中点撑开（同文件:187--203，常数定义于
:41--44）。分块口径：
\begin{itemize}
  \item 相角 $\theta\in[-\pi,\pi]$（硬编码），初值取母线 \fld{va\_deg} 化弧度
        （同文件:375--382）；
  \item 电压幅值 $V\in[$\fld{vmin\_pu}$,$\fld{vmax\_pu}$]$，缺省回退 $[0.5,1.5]$，
        初值取 \fld{vm\_pu}（同文件:384--392）；
  \item 发电机 $P_g\in[$\fld{pmin\_mw}$/S_{\mathrm{base}},\,$
        \fld{pmax\_mw}$/S_{\mathrm{base}}]$，$Q_g$ 同理取
        \fld{qmin\_mvar}/\fld{qmax\_mvar}，缺省回退 $\pm$\fld{kHugeBound}
        （$10^4$），初值取 \fld{pg\_mw}/\fld{qg\_mvar}（同文件:394--413）；
  \item 直流电压普通母线 $V_{dc}\in[0.8,1.2]$（硬编码），DC\_V 型母线收缩为
        设定值 $\pm0.05$，初值取 \fld{vm\_pu}（同文件:415--428）；
  \item 换流器 $P_{ac}$、$Q_{ac}$ 取自身 \fld{pmin\_mw}/\fld{pmax\_mw}、
        \fld{qmin\_mvar}/\fld{qmax\_mvar}；$P_{dc}$ 为双向推断对称界
        $\pm\max(10^{-3},\;|\,$\fld{pmax\_mw}$|,|\,$\fld{pmin\_mw}$|,|\,$
        \fld{qmax\_mvar}$|)/S_{\mathrm{base}}$（同文件:430--459）；三者初值均为
        0 向界内的截断（同文件:456--458）。
\end{itemize}
最后全体变量经 \srcpath{clamp\_in\_interior} 压入严格内点：边距
$\varepsilon=\min(\max(10^{-8},\,10^{-3}w),\;0.49w)$，$w$ 为该变量界宽
（同文件:205--223、461--463）；无法构造严格内点时直接失败返回
（同文件:3558--3562）。

\subsubsection*{目标函数}
目标为发电成本二次多项式之和，系数取自 canonical 发电机记录的
\fld{cost\_c2}/\fld{cost\_c1}/\fld{cost\_c0}（外部电网提升时同源拷贝），
\textbf{以 MW 有名值计价}
（\srcpath{src/optimal\_power\_flow/ac\_opf.cpp}:540--550，函数
\srcpath{evaluate\_objective}）：
\begin{align*}
f = \sum_{k=1}^{n_g}\left(c_{2,k}\,P_{g,k}^{2} + c_{1,k}\,P_{g,k} + c_{0,k}\right),
\qquad P_{g,k}\ \text{以 MW 计}
\end{align*}
内部变量 $P_g$ 为标幺值，故对 $x$ 的梯度与 Hessian 对角为
（同文件:552--572，函数 \srcpath{accumulate\_objective\_grad\_hdiag}）：
\begin{align*}
\frac{\partial f}{\partial P_{g,k}} &=
  2c_{2,k}S_{\mathrm{base}}^{2}\,P_{g,k} + c_{1,k}S_{\mathrm{base}} &
\frac{\partial^{2} f}{\partial P_{g,k}^{2}} &= 2c_{2,k}S_{\mathrm{base}}^{2}
\end{align*}
即目标对 $Q_g$、电压及换流器变量无直接二次项；无功与网损仅通过等式与不等式
间接进入目标。币种代码中未固定，单位记为``—''。

\subsubsection*{等式约束}
等式残差由共享潮流残差装配器与 OPF 自有项叠加而成
（\srcpath{src/optimal\_power\_flow/ac\_opf.cpp}:708--775，函数
\srcpath{assemble\_hybrid\_equalities}；潮流残差由
\srcpath{core::evaluate\_residual\_and\_jacobian} 在\textbf{剔除换流器}的
口径下计算，同文件:777--813）。

(1) 交流节点有功/无功平衡（$n_p+n_q$ 行，同文件:724--739）：
\begin{align*}
g^{P}_i &= P^{g}_i - P^{d}_i - P^{calc}_i(V,\theta) - \sum_{k\in\mathcal{C}(i)} P_{ac,k}
 & i &\in \text{非松弛母线}\\
g^{Q}_i &= Q^{g}_i - Q^{d}_i - Q^{calc}_i(V,\theta) - \sum_{k\in\mathcal{C}(i)} Q_{ac,k}
 & i &\in \text{PQ 母线}
\end{align*}
其中 $P^{calc}_i=V_i\sum_j V_j(G_{ij}\cos\theta_{ij}+B_{ij}\sin\theta_{ij})$、
$Q^{calc}_i=V_i\sum_j V_j(G_{ij}\sin\theta_{ij}-B_{ij}\cos\theta_{ij})$ 由 Ybus
直接求得（与 \srcpath{ac\_opf\_solver.hpp}:18--31 的文档式一致），
$\mathcal{C}(i)$ 为交流侧挂于母线 $i$ 的换流器集合。

(2) 直流节点平衡（$n_{dc}$ 行，同文件:741--755）：
\begin{align*}
g^{dc}_k &= \frac{P^{d}_{dc,k}}{S_{\mathrm{base}}}
  + V_{dc,k}\sum_{m} G^{dc}_{km}V_{dc,m} - \sum_{k'\in\mathcal{C}_{dc}(k)} P_{dc,k'}
\end{align*}
其中 $G^{dc}$ 为直流电导矩阵 \fld{gdc}，第一项为直流负荷，第二项为直流网络
注入 $V_k(G^{dc}V)_k$，末项扣除换流器直流口功率。

(3) 换流器交直流耦合（$n_c$ 行，同文件:757--772）：
\begin{align*}
g^{conv}_k &= P_{ac,k} + P_{dc,k} + P_{loss,k} = 0\\
P_{loss,k} &= a_k + b_k I_{ac,k} + c_k I_{ac,k}^{2}\\
I_{ac,k} &= \frac{\sqrt{\max(\varepsilon,\;P_{ac,k}^{2}+Q_{ac,k}^{2}+\varepsilon)}}
                    {\max(V_{m,k},\,10^{-6})}
\end{align*}
其中 $a_k=$~\fld{loss\_mw}$/S_{\mathrm{base}}$（固定损耗），
$b_k=$~\fld{loss\_percent}$/100$（比例损耗），$c_k=1-\eta$（\fld{eta} 效率，
导通损耗），数值软化常数 $\varepsilon=$~\fld{eps\_iac}=$10^{-6}$
（同文件:757--770）。符号约定为 $P_{ac}$ 交流口输出为正、
$P_{dc}=-(P_{ac}+P_{loss})$，即损耗由直流侧承担；该口径下\textbf{不计}
\fld{r\_conv\_ac\_pu} 交流侧导通损耗耦合（潮流模块才有该项），
\fld{converter\_model\_scope} 中 \fld{vsc\_ac\_conduction\_loss\_modelled}
保持 \fld{false}（同文件:3044--3050）。

\subsubsection*{非线性不等式（外点罚口径）}
紧致路径\textbf{不引入}不等式松弛与对应乘子，而是把支路热稳定界与换流器
容量圆写成 $h(x)\le 0$ 后以\textbf{外点二次罚}进入目标
（\srcpath{src/optimal\_power\_flow/ac\_opf.cpp}:853--966，函数
\srcpath{accumulate\_nonlinear\_limit\_terms}）：仅当 $h>0$ 时累加
\begin{align*}
\phi_{pen} &\mathrel{+}= \tfrac{1}{2}\rho\,h^{2}, &
\nabla\phi_{pen} &\mathrel{+}= \rho\,h\,\nabla h, &
\mathrm{diag}\,H_{pen} &\mathrel{+}= \rho\,(\nabla h)^{2}
\end{align*}
（同文件:865--885；Hessian 项为 Gauss--Newton 对角近似，舍去
$\rho h\nabla^{2}h$）。罚权重 $\rho$ 随迭代自适应，见
\ref{subsec:acopfnative:algorithm}~节。

(1) 支路视在功率平方界（两端各一条，同文件:887--952）。对在运且
\fld{rate\_a\_mva}$>0$ 的交流支路，以
$g_s=r/(r^{2}+x^{2})$、$b_s=-x/(r^{2}+x^{2})$、$b_c=b/2$、变比 $t$（0 回退 1）、
移相角 $\varphi=$~\fld{shift\_deg} 化弧度、
$\bar S^{2}=$(\fld{rate\_a\_mva}$/S_{\mathrm{base}})^{2}$
（\srcpath{build\_nonlinear\_limit\_workspace}，同文件:663--706）：
\begin{align*}
\theta &= \theta_i-\theta_j-\varphi\\
P_f &= \frac{g_s}{t^{2}}V_i^{2} - \frac{g_s\cos\theta+b_s\sin\theta}{t}V_iV_j\\
Q_f &= -\left(\frac{b_s}{t^{2}}+b_c\right)V_i^{2}
       - \frac{g_s\sin\theta-b_s\cos\theta}{t}V_iV_j\\
h_f &= P_f^{2}+Q_f^{2}-\bar S^{2}\le 0\\
P_t &= g_sV_j^{2} - \frac{g_s\cos\theta-b_s\sin\theta}{t}V_iV_j\\
Q_t &= -(b_s+b_c)V_j^{2} + \frac{g_s\sin\theta+b_s\cos\theta}{t}V_iV_j\\
h_t &= P_t^{2}+Q_t^{2}-\bar S^{2}\le 0
\end{align*}
（from 端：同文件:894--925；to 端：同文件:927--952；$\partial h/\partial x$
的四个分量在同文件:903--911、919--924、930--938、946--951 解析给出）。
串联阻抗为零（$r^{2}+x^{2}\le 0$）或未给 \fld{rate\_a\_mva} 的支路\textbf{跳过}
限值装配（同文件:667--678）。

(2) 换流器容量圆（每换流器一条，同文件:955--960）：
\begin{align*}
h^{conv}_k &= P_{ac,k}^{2}+Q_{ac,k}^{2}-\bar S^{2}_{conv,k}\le 0
\end{align*}
其中 $\bar S_{conv,k}$ 并非铭牌视在功率字段，而是从 P/Q 界推断的
$\max(|\,$\fld{pmax\_mw}$|,|\,$\fld{pmin\_mw}$|,|\,$\fld{qmax\_mvar}$|,|\,$
\fld{qmin\_mvar}$|)$，全部不可得时回退 \fld{kHugeBound}
（同文件:692--704）。

\subsubsection*{变量界的对数障碍}
盒式界经对数障碍进入目标（\srcpath{src/optimal\_power\_flow/ac\_opf.cpp}:
574--593，函数 \srcpath{accumulate\_log\_barrier}）：
\begin{align*}
\phi_{\mu} &= -\mu\sum_{i}\left[\ln(x_i-l_i)+\ln(u_i-x_i)\right]\\
\nabla\phi_{\mu,i} &= -\frac{\mu}{x_i-l_i}+\frac{\mu}{u_i-x_i} &
H_{\mu,ii} &= \frac{\mu}{(x_i-l_i)^{2}}+\frac{\mu}{(u_i-x_i)^{2}}
\end{align*}
任一分量触及或越出界时函数返回失败，主循环据以终止并报
``left the interior of variable bounds''（同文件:3641--3646）。

\subsection{原生内点算法全流程}\label{subsec:acopfnative:algorithm}

紧致路径的主循环为``外层障碍参数 $\mu$ 递减 + 内层障碍--罚 merit 牛顿法''
的双层结构（外层 \srcpath{src/optimal\_power\_flow/ac\_opf.cpp}:3609--4098，
内层 :3613--4081）。动态量初始化：
$\mu=\max($\fld{barrier\_mu0}$,$\fld{barrier\_mu\_min}$)$，
正则化 $\delta=\max($\fld{regularization}$,\,10^{-9})$，
罚权重 $\rho_{m}=\max($\fld{merit\_penalty}$,\,10)$，
目标权重 $w_f=0.05$（同文件:3596--3607）。

\subsubsection*{障碍--罚混合子问题}
每个内层迭代在当前 $(x,\mu,\rho_m,w_f)$ 下最小化
（同文件:3635--3678）：
\begin{align*}
F(x) = w_f f(x) + \phi_{\mu}(x) + \phi_{pen}(x;\rho_{nl})
       + \tfrac{1}{2}\rho_m\left(\lVert g(x)\rVert^{2}+\lVert h^{+}(x)\rVert^{2}\right)
\end{align*}
其中前四项进入梯度和 Hessian 对角（目标与障碍项精确，罚项取
Gauss--Newton 近似），末项仅以 merit 形式参与线搜索判据；非线性罚权重取
$\rho_{nl}=100\,\rho_m\cdot\mathrm{clamp}(10^{-2}/\max(\lVert g\rVert_{\infty},
10^{-12}),\,10^{-3},1)$，即可行域附近满罚、远离可行域时降罚以防梯度淹没
（同文件:3647--3650）。目标权重 $w_f$ 初值仅 0.05，在接近可行后逐步放大至
上限 1（同文件:3601、4078--4080）——这意味着迭代前中期求解的是一个
``可行性优先''的软化问题，判敛口径也随之针对障碍--罚子问题而非原始 KKT 系统。

\subsubsection*{KKT 牛顿系统组装与稀疏求解}
每次信赖尝试装配对称增广系统
（结构 \srcpath{build\_kkt\_workspace}，\srcpath{src/optimal\_power\_flow/ac\_opf.cpp}:1323--1399；
填值 \srcpath{update\_kkt\_values}，同文件:1401--1434）：
\begin{align*}
\begin{bmatrix} H+\delta I & J^{T}\\ J & -\delta I\end{bmatrix}
\begin{bmatrix}\Delta x\\ \Delta\lambda\end{bmatrix}
=
\begin{bmatrix}-\nabla F\\ -g\end{bmatrix}
\end{align*}
其中 $H$ 为目标+障碍+罚的对角 Hessian，$J=\partial g/\partial x$ 为等式雅可比，
右端 $\nabla F$ 不含 $J^{T}\lambda$ 项——对偶量 $\lambda$ 不跨迭代保持
（每次接受步后置零，同文件:4059），故该系统实质是``瞬时原始对偶''牛顿步；
解出的 $\Delta\lambda$ 仅随本步 merit 评估短暂留存，不作为对偶迭代量累积。
雅可比填值规则：潮流块取共享装配器输出的\textbf{负值}
（因 $g=$ 设定 $-$ 计算，注释见同文件:1236--1238）；发电机注入 $+1$、换流器
交流注入 $-1$（同文件:1241--1260）；直流块非对角元 $G^{dc}_{km}V_{dc,k}$、
对角元 $(G^{dc}V)_k+G^{dc}_{kk}V_{dc,k}$、$\partial g^{dc}/\partial P_{dc}=-1$
（同文件:1262--1283）；换流器耦合行
$\partial g/\partial P_{ac}=1+\partial P_{loss}/\partial P$、
$\partial g/\partial Q_{ac}=\partial P_{loss}/\partial Q$、
$\partial g/\partial P_{dc}=1$、$\partial g/\partial V_m=\partial P_{loss}/\partial V_m$，
其中 $\partial I/\partial P=P/(S_{ac}V_m)$ 等链式项解析给出
（同文件:1285--1320）。雅可比与 KKT 矩阵均采用``结构一次装配、非零元按索引
映射原地填值''的工作区复用设计（\srcpath{JWorkspace}:144--166、
\srcpath{KKTWorkspace}:168--176、\srcpath{build\_jacobian\_workspace}:968--1219）。
线性求解器由 \srcpath{core::make\_default\_sparse\_solver} 创建
（同文件:1397），后端按编译可用性取 KLU$\to$UMFPACK$\to$MKL
Pardiso$\to$SuperLU$\to$Eigen SparseLU 的优先级（
\srcpath{MIPSolvers/src/engine/kernel/linear\_algebra/linear\_solver.cpp}:230--242），
实际后端名记入 \fld{profiling.linear\_solver\_backend}（同文件:3551）。
稀疏模式分析只做一次并缓存，分解或回代失败即将本轮正则化乘 10 重试
（同文件:3694--3713、3835--3841）。

\subsubsection*{信赖尝试、步长截断与分数边界}
单个内层迭代最多做 \fld{kTrustAttempts}=6 次信赖尝试，每次乘以 10 递增正则化；
停滞计数 $\ge 3$ 且不可行度高于 $\max(10^{-3},\,10\,\varepsilon_{feas})$ 时跳过
信赖尝试直接进入恢复阶段（同文件:3686--3691）。牛顿步先做三层截断
（同文件:3715--3734）：
\begin{align*}
|\Delta x_i| &\le 0.25\,(u_i-l_i) && \text{逐变量}\\
\lVert\Delta x\rVert_{\infty} &\le 0.25 && \text{整体（\fld{kMaxDxInf}）}\\
|\Delta\lambda_r| &\le 10^{3} && \text{（\fld{kMaxDlambdaAbs}）}
\end{align*}
随后按分数边界规则取最大可行步长
（\srcpath{max\_feasible\_step}，同文件:1436--1454）：
\begin{align*}
\alpha_0 = \min\left(1,\;\min_{i:\Delta x_i>0}\frac{\tau(u_i-x_i)}{\Delta x_i},\;
                         \min_{i:\Delta x_i<0}\frac{\tau(x_i-l_i)}{-\Delta x_i}\right),
\qquad \tau=\mathrm{clamp}(\text{\fld{interior\_fraction}},\,0.5,\,0.9999)
\end{align*}
默认 $\tau=0.995$。

\subsubsection*{线搜索与接受判据}
在 $\alpha\leftarrow\alpha\cdot$\fld{step\_backoff}（默认 0.5）的回退链上最多试
\fld{max\_line\_search\_steps}（默认 20）个试验点，每个试验点要求严格内点、
重估等式与罚项（同文件:3742--3815）。接受判据为不可行度充分下降
\textbf{或}近可行时 merit 下降（同文件:3803--3812）：
\begin{align*}
\text{接受}\iff
\underbrace{\varphi^{trial}_{inf} < (1-5\times10^{-3}\alpha)\,\varphi^{cur}_{inf}}_{\text{不可行度下降}}
\;\lor\;
\big(\varphi^{cur}_{inf}\le\max(10^{-3},10\varepsilon_{feas})
 \;\land\; \varphi^{trial}_{merit}<\varphi^{cur}_{merit}\big)
\end{align*}
其中 $\varphi_{inf}=\sqrt{\lVert g\rVert^{2}+\lVert h^{+}\rVert^{2}}$。
若回退链上没有严格满足判据的点，启用``最优 $\alpha$ 兜底''：全程最佳试验点
在不可行度改善 $0.5\%$、或近可行且 merit 改善、或``近中性''（双指标恶化均
不超过 $0.5\%$）时仍可接受，近中性情形步长钳到 $5\times10^{-3}$
（同文件:3817--3833）。接受后正则化回收为
$\max($\fld{regularization}$,\delta_{trial}/2)$（同文件:3835--3837）。

\subsubsection*{恢复阶段与梯度回退}
信赖尝试全部失败时进入\textbf{恢复阶段}：最多 4 次尝试，Hessian 对角置零、
右端仅取 $-g$（纯可行性牛顿方向），步长上限收紧为
\fld{kMaxRestoreDxInf}=0.12 与逐变量 $0.10(u_i-l_i)$，接受判据放宽为
$(1-10^{-3}\alpha)$ 倍不可行度下降（同文件:3843--3958）。恢复再失败则取
最速可行性方向 $\Delta x=-J^{T}g$，上限 \fld{kMaxGradFallbackInf}=0.08，
判据 $(1-5\times10^{-4}\alpha)$ 或最佳点改善 $0.1\%$（同文件:3960--4035）。
三级机制均失败即终止，报 ``globalization failed (no acceptable
primal-dual/restoration step)''（同文件:4037--4041）。

\subsubsection*{步后更新与参数自适应}
接受步后：$x$ 压回内点（受约束变量边距比例 $10^{-3}$、自由变量 $10^{-8}$，
同文件:4043--4057）；$\lambda$ 置零（:4059）；按进展调整动态参数
（同文件:4061--4080）：
\begin{itemize}
  \item 恢复步：$\rho_m\leftarrow\min(10^{8},1.2\rho_m)$，
        $\delta\leftarrow\min(1,1.5\delta)$，$w_f\leftarrow\max(10^{-3},0.9w_f)$；
  \item 连续 4 步不可行度改善不足 $1\%$（停滞）：$\rho_m\times2$、$\delta\times5$、
        $w_f\times0.8$，停滞计数清零；
  \item 不可行度改善超 $20\%$：$\rho_m\leftarrow\max($\fld{merit\_penalty}$,
        0.95\rho_m)$（回落）；
  \item 不可行度进入 $10^{-2}$ 以内：$w_f\leftarrow\min(1,1.15w_f)$（目标权重爬坡）。
\end{itemize}

\subsubsection*{外层障碍参数更新与收敛判据}
内层收敛于当前 $\mu$（不可行度与平稳性同时达标，见下）或外层已``近可行''
（$\varphi_{inf}\le\max(10^{-4},\,50\varepsilon_{feas})$）时，外层把障碍参数
压向 $\mu\leftarrow\max(\mu\cdot$\fld{barrier\_mu\_reduction}$,\,$
\fld{barrier\_mu\_min}$)$（默认缩减率 0.2、下限 $10^{-8}$）；否则改为加重
$\rho_m\times1.5$、$\delta\times2$（同文件:4083--4098）。

收敛判据（同文件:3667--3674）：当
\begin{align*}
\max(\lVert g\rVert_{\infty},\;\max h^{+}) &\le \text{\fld{feasibility\_tol}}
&&\text{且} &
s_{dual} &\le \text{\fld{stationarity\_tol}}
\end{align*}
时认为当前 $\mu$ 的子问题收敛；若此时 $\mu\le1.01\,$\fld{barrier\_mu\_min}
则宣布整体收敛。平稳性度量 $s_{dual}$ 取 $\lVert\nabla F\rVert_{\infty}$；
当 $\varphi_{inf}\le10^{-3}$ 时改用其与最小二乘对偶估计
$\lVert\nabla F+J^{T}\lambda_{ls}\rVert_{\infty}$ 的较小者，$\lambda_{ls}$
由 $(JJ^{T}+10^{-8}I)\lambda_{ls}=-J\nabla F$ 经 \fld{SimplicialLDLT} 解得
（同文件:3655--3661、599--627）。

\subsubsection*{失败处理、诊断与回退}
终态汇总（\srcpath{finalize} 段，同文件:4100--4208）重估全部残差，按
$P/Q/DC/Conv$ 四族与 $\max h^{+}$ 分别报告，未收敛时 \fld{status} 携带各族
最大残差与最差平稳性变量的具名信息
（\srcpath{variable\_name}:1467--1499，状态串拼装:4196--4200）。失败原因枚举：
迭代耗尽（默认 ``maximum iterations reached''，:3597）、离开界内部（:3643）、
全局化失败（:4038）、初值非严格内点（:3558--3562）、维度退化或无在运发电机
（:3530--3539）。若 \fld{allow\_fallback} 且算例\textbf{不含} DC/VSC/LCC 组件，
未收敛时回退``经济调度 + AC 潮流''：先以 $\lambda$ 对分法做经济调度
（\srcpath{solve\_economic\_dispatch}:1501--1615，80 次对分、含截断退化的
最小平衡修正），再以牛顿潮流（80 次迭代、容差 $10^{-8}$）投影回交流可行域并
按母线均分 P/Q 到各机（\srcpath{try\_dispatch\_pf\_fallback}:1764--1867）；
混合算例则显式抑制该回退并写入提示（同文件:3266--3278、4189--4191）。
显式 \fld{EconomicDispatch} 后端直接使用同一快速路径，同样仅限纯 AC
（同文件:3501--3526）。

\subsubsection*{关键常数汇总}
\begin{center}
\footnotesize
\begin{tabular}{llll}
\toprule
常数 & 数值 & 出处（ac\_opf.cpp） & 含义\\
\midrule
\fld{kMinInteriorWidth} & $10^{-6}$ & :43 & 变量界最小宽度\\
\fld{kHugeBound} & $10^{4}$ & :44 & 自由变量数值界（标幺）\\
内点边距 & $\min(\max(10^{-8},10^{-3}w),0.49w)$ & :221 & 初值压入内点\\
$w_f$ 初值 & 0.05 & :3601 & 目标动态权重初值\\
$V_{dc}$ 界 & $[0.8,\,1.2]$ / DC\_V $\pm0.05$ & :418--423 & 直流电压硬编码界\\
\fld{eps\_iac} & $10^{-6}$ & :757 & 换流器电流数值软化\\
罚闸值 & $10^{-2}/\lVert g\rVert_\infty$ 钳于 $[10^{-3},1]$ & :3647 & 不等式罚自适应门\\
\fld{kTrustAttempts} & 6 & :3686 & 单迭代信赖尝试上限\\
\fld{kMaxDxInf} & 0.25 & :3687 & 原始步 $\infty$ 范数上限\\
\fld{kMaxDlambdaAbs} & $10^{3}$ & :3688 & 对偶步幅值上限\\
线搜索接受 & $(1-5\times10^{-3}\alpha)$ & :3804 & 不可行度充分下降系数\\
\fld{kRestorationAttempts} & 4 & :3844 & 恢复尝试上限\\
\fld{kMaxRestoreDxInf} & 0.12 & :3845 & 恢复步 $\infty$ 范数上限\\
\fld{kMaxGradFallbackInf} & 0.08 & :3963 & 梯度回退步上限\\
$JJ^{T}$ 正则 & $10^{-8}$ & :611 & 最小二乘对偶估计\\
ED 对分 & 80 次 / $10^{-8}$ & :1573--1581 & 经济调度 $\lambda$ 对分\\
PF 回退 & 80 次 / $10^{-8}$ & :1772--1773 & 回退潮流迭代与容差\\
\bottomrule
\end{tabular}
\end{center}

\subsection{Native parity IPM 核与求解编排}\label{subsec:acopfnative:kernel}

\subsubsection*{接口与数据结构}
作用于 Parity 全空间问题的自研原始对偶内点核声明于
\srcpath{include/hacdcpf/optimal\_power\_flow/native\_ipm\_solver.hpp}（命名空间
\fld{hacdcpf::opf::parity}，与 \fld{parity\_ipm\_solver.hpp} 共享结果类型）。
文档化的数学形式为（同文件:61--89）：
\begin{align*}
\min_x\; f(x)\quad \text{s.t.}\quad g(x)=0,\quad h(x)\le 0,\quad x_{min}\le x\le x_{max}
\end{align*}
不等式与盒式界引入松弛后，对摄动 KKT 系统
$\nabla f+J_g^{T}\lambda+J_h^{T}\mu+z^{+}-z^{-}=0$、$g=0$、$h+s=0$、$S\mu=\tau e$
取牛顿步，步长采用分数边界控制；返回对偶按 \fld{Problem::cidx} 的行序排列。
\fld{IPMOptions}（同文件:17--32）给出迭代上限 \fld{max\_iter}=400、三项容差
\fld{tol\_primal}/\fld{tol\_dual}/\fld{tol\_complementarity}=$10^{-6}$、正则化
$10^{-8}$、步长上限 \fld{alpha\_max}=0.95，以及可选的原始/等式对偶/不等式对偶/
松弛暖启动指针（仅在全部维数与正性检查通过时采用）。\fld{IPMResult}
（同文件:34--59）返回收敛标志、三项 KKT 残差及其初值、暖启动使用标记、
分解/回代/接受/拒绝/符号分析/后端升级/缩放重建计数、具体线性代数路径字符串
（如 \fld{dense\_lu}、\fld{sparse\_umfpack}、\fld{sparse\_klu}、
\fld{sparse\_eigen\_lu}）与原始-对偶解向量。核内部实现位于
\srcpath{src/optimal\_power\_flow/parity\_ipm.cpp}，细节见第~\ref{sec:parityipm}~章。
此外 \srcpath{homotopy\_tangent}（同文件:92--113）给出目标同伦 $P(t)=\min t\,f(x)$
在当前路径点的 Davidenko 切向：对增广系统右端取 $(-\nabla f/t,\,0,\,0)$ 复用
IPM 自身的 LDL$^{T}$ 机构一次分解回代，得 $(dx/dt,dz/dt,d\lambda/dt,d\mu/dt)$，
$t\approx0$、分解失败或非有限时返回 \fld{false}。

\subsubsection*{选项映射与暖启动编排}
\srcpath{solve\_with\_parity\_ipm}（\srcpath{src/optimal\_power\_flow/ac\_opf.cpp}:2486--3020）
把 \fld{ACOPFOptions} 翻译为核选项（同文件:2544--2551）：
\fld{max\_iter} 取内外层迭代上限之积；\fld{tol\_primal} 取
$\max($\fld{feasibility\_tol}$,10^{-10})$；\fld{tol\_dual} 取
$\max($\fld{stationarity\_tol}$,10^{-10})$；\fld{tol\_complementarity} 取
$\max($\fld{barrier\_mu\_min}$,10^{-10})$；\fld{regularization} 取
$\max($\fld{regularization}$,10^{-12})$；\fld{alpha\_max} 取
\fld{interior\_fraction} 钳于 $[0.5,0.9999]$。建模侧
\fld{load\_shedding} 仅当目标为 \fld{Economic} 时开启且 \fld{voll} 显式置 0，
四个 \fld{enforce\_*} 开关原样下传（同文件:2526--2540）。

暖启动有两条独立来源（同文件:2557--2602）：
\begin{itemize}
  \item \fld{ac\_pf\_warm\_start}（默认关）：先以混合潮流
        \srcpath{powerflow::solve\_hybrid}（开启换流器模式切换、PV-PQ 转换、
        自动摆节点选择）解出耦合可行点，再把电压、换流器功率、DC/DC 功率经
        \srcpath{build\_hybrid\_power\_flow\_warm\_start} 与
        \srcpath{recover\_power\_flow\_generator\_dispatch} 映射进 OPF 变量空间，
        其中发电机分派由潮流平衡残差在母线内按松弛优先重新分摊
        （同文件:2057--2265）；
  \item \fld{warm\_start}：把上一次兼容求解的 \fld{ACOPFResult} 逐变量族映射，
        优先采用不透明的 \fld{ipm\_primal\_state}（尺寸匹配且有限时），并以
        $10^{-8}$ 边距钳入界内；对偶与松弛状态仅在映射成功时一并传入
        （\srcpath{build\_parity\_primal\_warm\_start}:1879--2008）。
\end{itemize}

\subsubsection*{内层求解器选择与 Ipopt 回退}
内层引擎按 \fld{ParityInnerSolver}（枚举于同文件:1869--1873）与选项组合
（同文件:2604--2664）：
\begin{itemize}
  \item \fld{Auto} 且 \fld{ac\_pf\_warm\_start} 且非同伦：先用 Ipopt（若编译入），
        其终止于 acceptable 带而未严格收敛时，把最末原始点连同对偶/松弛交给
        自研核精化；Ipopt 不可用时直接自研核；
  \item \fld{Auto} 其余情形：先自研核，未收敛且 \fld{allow\_fallback} 时回退 Ipopt；
  \item \fld{NativeIPM}：仅自研核；\fld{Ipopt}：仅 Ipopt，未编译时回退自研核。
\end{itemize}
Ipopt 适配（\srcpath{solve\_parity\_with\_ipopt}:2384--2484，仅在
\srcpath{HACDCPF\_HAVE\_IPOPT} 下编译）把 Parity 问题的目标/梯度/等式/不等式/
雅可比/拉格朗日 Hessian 回调一一挂到 \fld{engine::NLPModel}
（\srcpath{build\_parity\_nlp\_input}:2275--2372），容差映射为：主容差取三项
最小、acceptable 带取三项最大、acceptable 互补下限 $\max(10^{-2},100\,
$\fld{tol\_complementarity}$)$（同文件:2321--2338）。结果接受口径除
\fld{success} 外还接受``acceptable residuals''：原始残差
$\le$\fld{tol\_primal}、对偶 $\le1000\,$\fld{tol\_dual}、互补
$\le\max(10^{-2},1000\,$\fld{tol\_complementarity}$)$，状态串保留 Ipopt 原始终止码
（同文件:2412--2428、2472--2475）。对偶重建时非线性不等式与盒式界乘子以
$2\times10^{-12}$ 为下限拼回 \fld{mu}/\fld{z}（同文件:2435--2471）。

\subsubsection*{目标同伦延续}
\fld{objective\_homotopy} 开启且内层非 Ipopt 时，进入两阶段延续
（同文件:3352--3482；理论引用 docs/numerical\_methods.md §12--13）：沿
$P(t)=\min\,t\,f$ 从 $t=0$（纯可行性）到 $t=1$（全成本），成本系数与 RPO 目标
权重同步乘 $t$（同文件:3371--3385）；最多 120 步，步长初值
\fld{homotopy\_dt0}（默认 0.10），快收敛（$\le20$ 次迭代）步长加倍至上限 0.5、
失败减半，$dt<0.005$ 判定进入最小违反域并停止（同文件:3367--3368、3395--3398、
3437--3443）。相邻步之间用 Davidenko 切向做预测
$w_{pred}=w+h\,w'$，预测步长受参数度量信赖半径 $\kappa=0.05$ 约束
$h=\min(dt,\kappa/\lVert dx/dt\rVert_\infty)$，松弛与不等式对偶以 $10^{-8}$ 为底
（同文件:3399--3436）。$t=1$ 未达且 \fld{allow\_fallback} 时，从全新耦合潮流点
对全目标做最后一次独立求解，成功则标记
``converged (full objective after incomplete objective homotopy)''；否则诚实报告
``not converged: objective homotopy stopped at t=...''（同文件:3451--3478）。
显式 Ipopt 请求与同伦互斥：同伦标志被清除后做单次独立求解（同文件:3483--3488）。

\subsubsection*{结果提取、LMP 与诊断量}
求解返回后（同文件:2666--3019）：电压、发电、甩负荷、换流器、DER、储能、
DC/DC、柔性负荷、能量路由器端口按块从原始向量提取并乘回 MW/MVAr
（:2716--2852）；在运 LCC 在最优电压状态下重估准稳态工作点
（$\alpha$/$\gamma$/$U_{d0}$/$I_d$/越限审计），触发角或熄弧角越窗时写入
\fld{model\_limitations}（:2854--2926）；LMP 从等式对偶提取
$\lambda^{P}_i\cdot$\fld{scale\_p}$/S_{\mathrm{base}}$ 与
$\lambda^{Q}_i\cdot$\fld{scale\_q}$/S_{\mathrm{base}}$，仅当对偶向量完整覆盖
平衡行时置 \fld{lmp\_valid=true}，否则写明原因（Ipopt 适配器未返回约束乘子或
对偶不完整，:2930--2952）；终态按平衡族报告 pu 残差并以
$\max(\lVert g\rVert_\infty,\max h^{+},\text{界违反})$ 汇总
\fld{max\_constraint\_violation}（:2954--2996）；原始/对偶/松弛终态存入
\fld{ipm\_*\_state} 供续解（:2684--2695）；同伦切向存入 \fld{ipm\_tangent\_*}
（:2697--2714）。最后按 \fld{bus\_merge\_map} 反投影到用户母线序并按投影证书
裁剪内部直流母线（:3011--3017；\srcpath{unproject\_per\_bus\_ac\_opf\_result}
与 \srcpath{trim\_internal\_dc\_bus\_results}，同文件:46--77）。计数口径：
\fld{iterations} 报告为 IPM 迭代数 $+1$、\fld{outer\_iterations} 恒为 1
（:2667--2668）。

\subsection{公共接口与门面}\label{subsec:acopfnative:api}

本引擎对外的公开 API 均为\textbf{自由函数}（当前代码检出处不存在名为
\fld{NativeIpmSolver} 或 \fld{AcOpfSolver} 的类）：
\begin{itemize}
  \item \texttt{opf::solve\_ac\_opf(sys, opt)}：主入口
        （\srcpath{include/hacdcpf/optimal\_power\_flow/ac\_opf\_solver.hpp}:42--43）；
  \item \texttt{opf::compute\_jacobian\_diagnostics(sys, opt)}：雅可比诊断声明
        （同文件:54--55）。\textbf{注意}：该函数在当前检出处\textbf{只有声明、
        没有定义}；\srcpath{ac\_opf.cpp}:4211 实现的是
        \srcpath{check\_ac\_opf\_jacobian\_fd}，且无任何调用方，两者均未接入
        构建产物（详见 \ref{subsec:acopfnative:verification}~节）；
  \item \texttt{parity::solve\_primal\_dual\_ipm(prob, opt)} 与
        \srcpath{parity::homotopy\_tangent(...)}：Native IPM 核与同伦切向
        （\srcpath{native\_ipm\_solver.hpp}:90--113）；
  \item \fld{hacdcpf::solve\_ac\_opf} 全局门面
        （\srcpath{src/api/hacdcpf.cpp}:1865--1890）：调用 \fld{opf::solve\_ac\_opf}
        后，对收敛且含开关/断路器的算例补算开关类设备端子潮流
        （\fld{ac\_switch\_flows}/\fld{ac\_circuit\_breaker\_flows}）；
  \item \fld{IOPFSolver}/\fld{DefaultOPFSolver} 抽象
        （\srcpath{include/hacdcpf/optimal\_power\_flow/opf\_solver\_interface.hpp}:24--42）：
        把 \fld{OPFProblem} 路由到 \fld{hacdcpf::solve\_ac\_opf} 或
        \fld{solve\_dc\_opf}，\fld{sys} 为空或 \fld{is\_ac}/\fld{is\_dc} 不唯一时
        抛 \fld{std::invalid\_argument}。
\end{itemize}

\subsection{选项字段表}\label{subsec:acopfnative:options}

下表给出本章引擎直接消费的 \fld{ACOPFOptions} 字段（定义于
\srcpath{include/hacdcpf/optimal\_power\_flow/opf\_options.hpp}:49--123；
公共结果字段 \fld{converged}/\fld{objective}/\fld{status} 等的通用口径见
第~\ref{sec:overview}~章）。``$\dagger$''标记仅 Parity 路径消费的字段。

\begin{paramtable}
\fld{ac\_solver\_backend} & — & — & Auto/\allowbreak ParityIPM/\allowbreak Ipopt/\allowbreak EconomicDispatch & Auto & 后端选择；翻译规则见 \ref{subsec:acopfnative:role}~节\\
\fld{objective} & — & — & Economic/\allowbreak VoltageDeviation/\allowbreak ActiveLoss/\allowbreak VoltageDeviation\allowbreak AndLoss & Economic & 目标类别；非经济目标经 Parity 建模生效$\dagger$\\
\fld{voltage\_target\_pu} & $V^{ref}$ & pu & 0.95--1.05（经验值） & 1.0 & 电压偏差目标$\dagger$\\
\fld{voltage\_deviation\_weight} & $w_V$ & — & $>0$ & 1.0 & 电压偏差目标权重$\dagger$\\
\fld{active\_loss\_weight} & $w_{loss}$ & — & $>0$ & 1.0 & 网损目标权重$\dagger$\\
\fld{max\_inner\_iterations} & — & — & $\ge1$ & 80 & 内层迭代上限；$\le0$ 时重置 80\\
\fld{max\_outer\_iterations} & — & — & $\ge1$ & 8 & 外层 $\mu$ 迭代上限；$\le0$ 时重置 8\\
\fld{max\_line\_search\_steps} & — & — & $\ge1$ & 20 & 线搜索回退上限\\
\fld{feasibility\_tol} & $\varepsilon_{feas}$ & pu & $10^{-8}$--$10^{-4}$ & $10^{-6}$ & 等式/不等式违反判敛容差\\
\fld{stationarity\_tol} & $\varepsilon_{stat}$ & — & $10^{-8}$--$10^{-4}$ & $10^{-6}$ & 平稳性判敛容差\\
\fld{barrier\_mu0} & $\mu_0$ & — & $10^{-4}$--1 & $10^{-2}$ & 初始障碍参数\\
\fld{barrier\_mu\_reduction} & $\sigma_\mu$ & — & 0.1--0.5 & 0.2 & 外层 $\mu$ 缩减率\\
\fld{barrier\_mu\_min} & $\mu_{min}$ & — & $10^{-10}$--$10^{-6}$ & $10^{-8}$ & 障碍参数下限（收敛门槛）\\
\fld{merit\_penalty} & $\rho_m$ & — & 1--$10^{3}$ & 10.0 & merit 等式罚权重初值\\
\fld{regularization} & $\delta$ & — & $10^{-9}$--$10^{-4}$ & $10^{-6}$ & KKT 正则化初值\\
\fld{step\_backoff} & $\beta$ & — & 0.3--0.7 & 0.5 & 线搜索回退因子\\
\fld{interior\_fraction} & $\tau$ & — & 0.9--0.9999 & 0.995 & 分数边界比例\\
\fld{ac\_eval\_threads} & — & — & $\ge1$ & 1 & 交流残差并行线程数\\
\fld{enable\_primal\_dual} & — & — & true/false & false & legacy 标志：启用内点路径（Auto 下混合算例自动置位）\\
\fld{use\_parity\_ipm} & — & — & true/false & false & legacy 标志：改用 Parity 全空间建模\\
\fld{allow\_fallback} & — & — & true/false & true & 允许 Ipopt/ED+PF 回退与同伦末段挽救\\
\fld{verbose} & — & — & true/false & false & 迭代日志\\
\fld{enable\_phase\_one} & — & — & true/false & true & Parity Native IPM 的有界 primal/dual warm-start 构造$\dagger$\\
\fld{phase\_one\_time\_limit\_ms} & — & ms & $>0$ 协作截止；$\le0$ 不限 & 5000 & Phase I 墙钟预算；单次稀疏分解不可中断$\dagger$\\
\fld{phase\_one\_max\_iterations} & — & — & $\ge0$ & 12 & Phase I Newton 迭代硬上限$\dagger$\\
\fld{phase\_one\_max\_factorizations} & — & — & $\ge0$ & 14 & primal/dual 共用稀疏分解硬上限$\dagger$\\
\fld{phase\_one\_max\_backtracks} & — & — & $\ge0$ & 12 & 每次 Phase I 步的回溯上限$\dagger$\\
\fld{phase\_one\_dispatch\_dual\_predictor} & — & — & true/false & false & 实验性零分解连通分量 dispatch-price dual 候选$\dagger$\\
\fld{phase\_one\_dispatch\_dual\_min\_improvement} & — & — & 0--1 & 0.05 & raw 与 normalized dual residual 必须同时达到的改善比例$\dagger$\\
\fld{warm\_start} & — & — & 指针/null & null & 上次兼容结果的非拥有指针，求解期间须存活\\
\fld{ac\_pf\_warm\_start} & — & — & true/false & false & 先做混合潮流并映射为暖启动（默认关，注释称良态算例物理初值更快）\\
\fld{compute\_homotopy\_tangent} & — & — & true/false & false & 返回点计算 Davidenko 切向$\dagger$\\
\fld{homotopy\_t} & $t$ & — & $(0,1]$ & 1.0 & 所解问题的成本标度（切向右端 $\nabla f/t$）$\dagger$\\
\fld{objective\_homotopy} & — & — & true/false & false & 两阶段目标同伦延续；与 Ipopt 内层互斥\\
\fld{homotopy\_dt0} & $dt_0$ & — & 0.01--0.5 & 0.10 & 同伦初始步长；$\le0$ 回退 0.10\\
\fld{enforce\_branch\_limits} & — & — & true/false & true & 支路限值族开关$\dagger$\\
\fld{enforce\_converter\_capacity} & — & — & true/false & true & 换流器容量圆开关$\dagger$\\
\fld{enforce\_converter\_current\_limits} & — & — & true/false & true & 换流器电流限值开关$\dagger$\\
\fld{enforce\_converter\_modulation\_limits} & — & — & true/false & true & 调制度/占空比窗口开关$\dagger$\\
\end{paramtable}

\subsection{结果字段表}\label{subsec:acopfnative:results}

\fld{ACOPFResult}（\srcpath{include/hacdcpf/optimal\_power\_flow/opf\_result.hpp}:72--169）
字段众多；下表仅列与本章引擎直接相关者，派生设备族（DER/储能/DC-DC/柔性负荷/
能量路由器，仅 Parity 路径填充）与通用字段的完整口径见第~\ref{sec:overview}~章。

\begin{paramtable}
\fld{vm} / \fld{va} & $V,\theta$ & pu, rad & 约 0.9--1.1 / $\pm\pi$ & 1.0 / 0 & 交流母线电压（反投影回用户母线序）\\
\fld{vdc} & $V_{dc}$ & pu & 0.8--1.2 & 空 & 直流母线电压（按投影证书裁剪）\\
\fld{pg\_mw} / \fld{qg\_mvar} & $P_g,Q_g$ & MW, MVAr & 机界内 & 输入值 & 发电机调度（已剔除合成外部电网机）\\
\fld{external\_grid\_p\_mw} / \fld{external\_grid\_q\_mvar} & $P_{eg},Q_{eg}$ & MW, MVAr & $\pm$合成界 & 0 & 外部电网分派（由合成机拆回）\\
\fld{dpd\_mw} / \fld{dqd\_mvar} & $\Delta P_d,\Delta Q_d$ & MW, MVAr & $\ge0$ & 空 & 逐母线甩负荷（仅 Parity 路径）\\
\fld{pac\_mw} / \fld{qac\_mvar} & $P_{ac},Q_{ac}$ & MW, MVAr & 换流器界内 & 空 & VSC 交流口功率（仅 Parity 路径填充）\\
\fld{lcc\_transfers} & — & — & — & 空 & 最优电压点重估的 LCC 准稳态工作点（仅 Parity 路径）\\
\fld{lmp\_p} / \fld{lmp\_q} & $\lambda^P,\lambda^Q$ & — & — & 空 & 节点价格，$\lambda\cdot$scale$/S_{\mathrm{base}}$；币种未定\\
\fld{lmp\_valid} / \fld{lmp\_validity\_reason} & — & — & true/false & false/空 & LMP 是否来自认证等式乘子及原因串\\
\fld{ipm\_primal\_state} 等四元 & — & — & — & 空 & 不透明续解状态（原始/等式对偶/不等式对偶/松弛），配合 \fld{warm\_start} 使用\\
\fld{ipm\_tangent\_primal} 等四元 & $dw/dt$ & — & — & 空 & 同伦切向（\fld{compute\_homotopy\_tangent} 时填充）\\
\fld{gen\_map} & — & — & — & — & 发电机行到 \fld{index} 的稳定归因映射\\
\fld{iterations} / \fld{outer\_iterations} & — & — & $\ge0$ & 0 & 迭代计数（Parity 路径为 IPM 数+1 与恒 1）\\
\fld{max\_constraint\_violation} & $\varphi_{inf}$ & pu & — & 0 & 终态最大约束违反（等式、$h^+$、界三者取大）\\
\fld{max\_stationarity} & $s_{dual}$ & — & — & 0 & 终态平稳性度量（障碍-罚口径，近可行时取最小二乘估计）\\
\fld{solver\_path} & — & — & NativeAC/\allowbreak ParityIPM & Unknown & 实际求解路径标记\\
\fld{profiling.linear\_solver\_backend} & — & — & 后端名字符串 & 空 & KKT 线性代数路径\\
\fld{profiling.\allowbreak factorization\_calls} 等计数 & — & — & $\ge0$ & 0 & 分解/回代/接受/拒绝/分析/升级/缩放重建计数\\
\fld{profiling.phase\_one\_*} & — & pu/ms/计数 & — & inf/0/not-run & Phase I 前后违反量、dual-fit、预算消耗、终止/线性后端及 Phase II 接受状态\\
\fld{profiling.dispatch\_dual\_predictor\_*} & — & ms/残差/布尔 & — & inf/0/not-attempted & 零分解 dual predictor 的尝试、接受、成本及 exact residual 审计\\
\fld{profiling.max\_*\_residual\_pu} & — & pu & — & 0 & 分族终态残差（P/Q/DC/Conv/其他/非线性不等式）\\
\fld{profiling.final\_barrier\_mu} & $\mu$ & — & — & 0 & 终态障碍参数（Parity 路径记录的是互补度）\\
\fld{infeasibility\_hints} & — & — & — & 空 & 不可行/不收敛的可能原因\\
\fld{model\_limitations} & — & — & — & 空 & 诚实口径声明（非凸局部解、储能单时段等）\\
\fld{converter\_model\_scope} & — & — & scope 标签 & 空 & 换流器建模覆盖声明（NativeAC 为 \texttt{ac-opf-\allowbreak native:\allowbreak vsc-free-\allowbreak pq+\allowbreak capacity-\allowbreak circle}）\\
\fld{audit} & — & — & — & 空 & 求解后独立可行性审计（由 \fld{verify\_opf\_result} 填充，见第~\ref{sec:verification}~章）\\
\fld{ac\_switch\_flows} / \fld{ac\_circuit\_breaker\_flows} & — & MW 等 & — & 空 & 开关/断路器端子潮流（全局门面补算）\\
\end{paramtable}

\subsection{验证与校核}\label{subsec:acopfnative:verification}

\begin{itemize}
  \item \textbf{解析雅可比的有限差分自检}：
        \srcpath{check\_ac\_opf\_jacobian\_fd}（\srcpath{src/optimal\_power\_flow/ac\_opf.cpp}
        第 4211--4409 行）对全部或均匀抽样的变量列做中心差分
        $h=\min($\fld{fd\_eps}$\cdot\max(1,|x|),\;0.49\min(x-l,u-x))$，
        \fld{fd\_eps} 缺省 $10^{-6}$，逐元比较解析列与差分列，通过阈值为
        最大绝对误差 $\le5\times10^{-4}$ 或最大相对误差 $\le5\times10^{-3}$
        （同文件:4319--4407）。\textbf{如实说明}：该函数在当前检出处没有任何
        调用方；公共头声明的 \srcpath{compute\_jacobian\_diagnostics}
        （\srcpath{ac\_opf\_solver.hpp}:54--55）在 \srcpath{src/} 内无定义，
        两者名字不一致且均未接入默认测试目标，属已登记的存疑点。
  \item \textbf{后端选择与混合算例收敛}：\srcpath{tests/test\_opf\_solver\_backends.cpp}
        覆盖内部混合算例 \fld{case300\_acdc}/\fld{case2000\_acdc} 的 ParityIPM
        收敛（:86 起）、Ipopt 环境变量门控，以及 \fld{case2869pegase} 大算例的
        稀疏 KKT 路径收敛（:1154 起）。
  \item \textbf{三方法交叉验证}：\srcpath{tests/test\_acopf\_dcopf\_crossval.cpp}
        在 MATPOWER 小算例（case5/9/14/30/39/57）上对 AC OPF、DC OPF 与潮流
        三者做调度一致性互查（:126、:508、:637 调用 \fld{solve\_ac\_opf}）。
  \item \textbf{回归与特殊场景}：\srcpath{tests/test\_hacdcpf.cpp}:390 的
        IEEE14 AC/DC 混合 OPF；\srcpath{tests/test\_nighttime\_opf.cpp}:58 起的
        夜间无光照 OPF 回归；\srcpath{tests/test\_graph\_roundtrip.cpp}:249 的
        图聚合往返后 OPF 一致性。
  \item \textbf{手算校核数据}：\srcpath{external\_data/opf\_hand\_cases/} 提供
        单母线 merit-order、两母线拥塞、甩负荷、二次成本线性化四个手算可复算
        算例（见其 README）；当前检出处未见直接消费该目录的 C++ 测试目标，
        如实标注为数据资产。
  \item \textbf{大网实测声明（代码注释口径）}：\fld{ac\_pf\_warm\_start} 与
        \fld{objective\_homotopy} 的字段注释记载 PEGASE 大网实测——
        case13659pegase 经暖启动以 189 次迭代收敛于参考最优、同伦延续认证
        2869 节点精确解与 9241 节点 $-0.013\%$ 偏差
        （\srcpath{opf\_options.hpp}:88--114）。这些数字为注释记载、非本章实测复核。
  \item \textbf{万节点结构化启动实测}：Release/KLU、一次预热加三次计时的固定协议下，
        case9241pegase 从普通 AC-PF 启动的 $12.657\,\mathrm{s}$/77 次迭代/76 次
        Phase-II 分解，降至 compact DCOPF dispatch seed 的
        $11.088\,\mathrm{s}$/56/55；端到端中位数下降 $12.4\%$，分解下降
        $27.6\%$。两者目标分别为 315911.571480 与 315911.571406，终态
        primal/dual 均低于 $5\times10^{-7}$。case13659pegase 的 baseline dual
        为 1.970，预筛选零成本跳过 DC，保持 138 次迭代、137 次分解、目标
        386116.837520 及 $2.05\times10^{-7}/5.89\times10^{-9}$ 的最终证书。
        该路径按真实 Parity 初始指标选择 dispatch 点；9241 的 dual
        $838.166\to236.753$，但最大 primal metric
        $1.601\times10^{-3}\to1.506\times10^{-1}$，故它是 dual-dispatch
        warm start 而非 central-state 证书，最终 KKT 仍由 Phase II 独立认证。
        DC Phase I 的 $2000\,\mathrm{ms}$ 默认墙钟预算覆盖投影、formulation、
        结构初值、symbolic 与 numeric 工作；本例耗时 $1638\,\mathrm{ms}$、一次
        DC symbolic analyze，且 baseline/candidate 只共享构建一次的 Parity
        formulation。稀疏分解不可中断，故允许一个在途分解造成的可诊断超时。
  \item \textbf{两万节点级边界实测}：25,000-bus ACTIVSg25k（4,834 台发电机、
        32,230 条交流支路）在同一 Release/KLU 主机上按一次预热加三次计时协议，
        普通 AC-PF 启动的样本为 $183.967/182.383/183.808\,\mathrm{s}$，中位数
        $183.808\,\mathrm{s}$；三次均以 66 次迭代、65 次 Phase-II 分解收敛，
        目标为 6252070.20694，primal/dual 为
        $3.239\times10^{-7}/1.004\times10^{-7}$。默认 $2\,\mathrm{s}$ DC Phase I
        的中位数为 $186.588\,\mathrm{s}$（慢 $1.51\%$）：一次 symbolic、两次
        迭代后以 $2064\,\mathrm{ms}$ 超时（overshoot $64\,\mathrm{ms}$），未形成
        合格候选并精确回退，故 Phase II 工作量与最终 KKT 不变。
        将 DC 预算扩大到 $30\,\mathrm{s}$ 虽把 DC 残差降至
        $8.767\times10^{-4}$，但真实 Parity primal 从
        $3.929\times10^{-4}$ 恶化到 $9.786\times10^{-2}$，dual 仅改善
        $4.15\%$，仍被固定 admission gate 拒绝，总时长反而增至
        $215.012\,\mathrm{s}$。因此当前实现已验证能求解该 25k ACOPF，但不能
        声称结构化 DC Phase I 在该算例上加速；继续提高辅助 DC-IPM 精度不是
        合理方向。
  \item \textbf{零分解 dual predictor 负结果}：实验模式按交流导通分量从未贴界
        机组 stationarity 中位数构造共同 $P/Q$ 平衡乘子，不增加 symbolic/numeric
        分解，并要求 raw 与 normalized full dual residual 均改善至少 $5\%$。
        ACTIVSg25k 单次 Release/KLU 探针仅以 $1.978\,\mathrm{ms}$ 将二者从
        $47278.840/468.10733$ 改为 $47276.850/468.08762$（约 $0.0042\%$），
        因而拒绝；总时长 $184.344\,\mathrm{s}$，Phase II 仍为 66 次迭代、
        65 次分解。该选项默认关闭，不能作为 25k 加速声明。
  \item \textbf{25k prepared-session 重复求解}：同一 ACTIVSg25k 在一个
        \srcpath{PreparedACOPFSession} 内连续求解两次，首轮 AC-PF 启动耗时
        $184.645\,\mathrm{s}$、66 次迭代、65 次 numeric 分解；第二轮复用
        formulation/maps、symbolic ordering 与完整 continuation state，耗时
        $8.666\,\mathrm{s}$、4 次迭代、3 次新 numeric 分解、0 次 symbolic
        analyze，墙钟降低 $95.31\%$。第二轮目标 6252069.93711，primal/dual
        为 $3.389\times10^{-8}/2.092\times10^{-8}$；峰值 RSS 为
        $4709.6\,\mathrm{MiB}$。因此 25k 重复求解加速已验证，但内存效率仍是
        更大批处理的边界；该结果不表示复用了旧 numeric factors。
\end{itemize}

\subsection{边界与局限}\label{subsec:acopfnative:limitations}

以下条目均按代码实际行为（含其自述注释与 \fld{model\_limitations} 文案）列出：
\begin{enumerate}
  \item \textbf{非凸局部解}：AC OPF 为非凸非线性规划，收敛仅认证局部 KKT 点而非
        全局最优——该语句作为 \fld{model\_limitations} 固定写入每次返回
        （\srcpath{src/optimal\_power\_flow/ac\_opf.cpp}:3205--3206）。
  \item \textbf{legacy 定位}：紧致 NativeAC 路径自述为旧式``外部障碍/罚循环''，
        结果中建议生产环境改用 ParityIPM（同文件:3219--3222）。
  \item \textbf{不等式非硬约束}：支路热稳定界与容量圆仅经外点二次罚进入目标，
        终解允许受罚权重控制的残余违反；罚 Hessian 为 Gauss--Newton 对角近似，
        且罚权重在远离可行域时按门值降权（同文件:853--966、3647--3650）。
  \item \textbf{无持续对偶迭代}：$\lambda$ 每次接受步后置零（同文件:4059），
        平稳性判据针对障碍--罚软化子问题，近可行时才切换到最小二乘对偶估计；
        目标权重 $w_f$ 初值 0.05、随可行性爬坡，前中期迭代并非在解原始成本问题
        （同文件:3601、3655--3661、4066--4080）。
  \item \textbf{换流器控制模式无关}：紧致路径把 VSC 处理为自由 $P_{ac}/Q_{ac}/P_{dc}$
        盒式变量，不锚定 VDC\_Q 直流电压、不实施电流/调制度窗口，
        \fld{vsc\_vdc\_control\_modelled=false}，scope 标签固定为
        \fld{ac-opf-native:vsc-free-pq+capacity-circle}（同文件:3041--3050）；
        亦不计 \fld{r\_conv\_ac\_pu} 交流侧导通损耗。
  \item \textbf{推断界与硬编码界}：换流器容量圆半径与 $P_{dc}$ 界由 P/Q 界推断
        而非铭牌视在功率（同文件:443--447、692--704）；$V_{dc}$ 界硬编码
        $[0.8,1.2]$、DC\_V 型收缩为设定值 $\pm0.05$（:415--428）；$\theta$ 界
        硬编码 $\pm\pi$（:378）；无显式 SLACK 时取第 0 号母线（:178--185）。
  \item \textbf{LCC 不进入紧致路径}：无 $V_{dc}$ 耦合列，凡在运 LCC 一律强制
        Parity/Ipopt 路径（同文件:3343--3350）；Parity 路径中 LCC 的
        $P/I/\alpha/\gamma$ 指令与换流变分接头为固定输入而非调度变量，额定电流
        饱和纳入模型而 $\alpha/\gamma$ 限值仅为求解后审计，换相重叠迭代、
        分接头控制与有功损耗未建模，等式二阶导用解析雅可比的局部有限差分且
        额定电流转折不可微（同文件:2514--2523、2854--2926）。
  \item \textbf{设备覆盖边界}：储能按单时段功率注入优化、不含跨时段 SOC 动态；
        能量路由器端口平衡按无损处理（同文件:3207--3214）；支路限值仅取
        \fld{rate\_a\_mva}（:667）；甩负荷仅 Economic 目标开启且 \fld{voll}
        置 0，其价格口径由 Parity 建模层决定（:2527--2531，见
        第~\ref{sec:parityform}~章）。
  \item \textbf{外部电网合成界}：短路容量并非进口限值，仅作有限尺度；该近似已在
        代码注释中自述（同文件:3173--3189）。
  \item \textbf{回退口径}：ED+PF 回退对混合算例显式抑制（同文件:3266--3278、
        4189--4191）；回退成功时 \fld{max\_stationarity} 置 0、目标按钳制后出力
        重算，属可行性投影而非优化解（:1832--1845）。
  \item \textbf{门控依赖}：Ipopt 路径仅在 \srcpath{HACDCPF\_HAVE\_IPOPT} 编译时
        存在，且可被 \srcpath{HACDCPF\_DISABLE\_IPOPT\_OPF} 运行时关闭
        （同文件:33--35、2388--2396）；未编译时 \fld{Auto} 的全部回退自动失效。
  \item \textbf{已知不一致}：\srcpath{compute\_jacobian\_diagnostics} 有声明无定义、
        \srcpath{check\_ac\_opf\_jacobian\_fd} 有定义无调用方；
        \srcpath{rebalance\_pg\_within\_bus} 标记 \fld{[[maybe\_unused]]} 为死代码
        （同文件:1617）；回退用 \srcpath{compute\_ac\_injections} 为稠密 Ybus
        双重循环，仅小算例回退可接受（:1738--1762）。
\end{enumerate}

% 直流潮流近似最优潮流（DC OPF）
% 本章文件由 \input 引入，不含导言区。
\section{直流潮流近似最优潮流（DC OPF）}\label{sec:dcopf}

DC OPF 是平台 OPF 家族中的线性化快速通道：把交流网络按直流潮流近似
（无损、无电压幅值变量、无无功方程）化为线性规划（LP）或凸二次规划（QP），
用于经济调度、节点边际电价（LMP）估计、可靠性评估中的大规模故障状态筛选等
对速度敏感、对电压/无功精度不敏感的场景。求解器本体为
\srcpath{src/optimal\_power\_flow/dc\_opf.cpp}（1394 行），公共入口声明于
\srcpath{include/hacdcpf/optimal\_power\_flow/dc\_opf\_solver.hpp}。与 AC OPF
（第\ref{sec:acopfnative}章、第\ref{sec:parityipm}章）不同，本求解器只覆盖
交流侧有功--相角子问题：直流网络、VSC/DC-DC/LCC 换流器物理一律不建模
（详见本章\ref{subsec:dcopf-dcgrid}节）。

\subsection{功能定位与公共接口}

\compmeta{solve\_dc\_opf / check\_dc\_opf\_feasibility}{\srcpath{include/hacdcpf/optimal\_power\_flow/dc\_opf\_solver.hpp}}

\subsubsection*{入口函数}
公共接口仅两个函数
（\srcpath{include/hacdcpf/optimal\_power\_flow/dc\_opf\_solver.hpp:solve\_dc\_opf、check\_dc\_opf\_feasibility}）：
\begin{align*}
&\text{\srcpath{solve\_dc\_opf}(\fld{sys}, \fld{opts})} \;\longrightarrow\;
  \text{\fld{DCOPFResult}}\\
&\text{\srcpath{check\_dc\_opf\_feasibility}(\fld{sys}, \fld{result}, \fld{tol})}
  \;\longrightarrow\; (\text{feasible},\ \text{max\_violation},\ \text{message})
\end{align*}
前者组装并求解 DC OPF；后者对一份已求解结果做独立的可行性复查
（发电机界、支路界、含松弛节点在内的逐节点功率平衡），容差 \fld{tol} 默认
$10^{-6}$，判据按 $\fld{tol}\times S_{\mathrm{base}}$ 折算为 MW
（\srcpath{src/optimal\_power\_flow/dc\_opf.cpp:check\_dc\_opf\_feasibility}，函数
\srcpath{check\_dc\_opf\_feasibility}；判据在 :check\_dc\_opf\_feasibility）。公共门面
\srcpath{src/api/hacdcpf.cpp:solve\_dc\_opf} 仅是透传封装，不附加语义。

\subsubsection*{内部装配结构}
建模产物集中在匿名命名空间的结构体 \fld{DCOPFFormulation}
（\srcpath{src/optimal\_power\_flow/dc\_opf.cpp:DCOPFFormulation}）：成员 \fld{lp} 与
\fld{qp} 分别承载线性化成本 LP 模型与真二次成本 QP 模型（约束结构共用），
\fld{gen\_map}/\fld{branch\_map} 记录活跃发电机/支路到原始组件向量的映射，
\fld{bus\_map} 记录母线稳定编号（\fld{ACBus::index}）到 0-based 位置的映射
（构建于 :build\_bus\_map，函数 \srcpath{build\_bus\_map}；发现重复母线编号时抛出
\fld{std::invalid\_argument}，:build\_bus\_map）。基准功率取
$\fld{base\_mva}=\max(\fld{sys.ac.base\_mva},\,1.0)$，默认 100~MVA
（:DCOPFFormulation、:build\_dc\_opf\_lp）。全章未注明处功率变量均为 $S_{\mathrm{base}}$ 标幺值。

\subsection{DC 线性化假设}\label{subsec:dcopf-assumptions}

代码中的近似假设以注释与结果声明为准，逐条如下。

\begin{enumerate}
  \item \textbf{忽略电阻、以电纳 $b=1/x$ 传输}。建模注释原文：
        ``In DC power flow: $P_{f,ij}=(\theta_i-\theta_j)/x_{ij}$ (ignore r,
        assume $x \gg r$)''（\srcpath{src/optimal\_power\_flow/dc\_opf.cpp:build\_dc\_opf\_lp}，
        函数 \srcpath{build\_dc\_opf\_lp} 的注释块 :build\_dc\_opf\_lp）。支路模型只读取
        \fld{Branch::x\_pu}，不读取电阻、分接头变比、移相角与充电电纳
        （:build\_dc\_opf\_lp）。
  \item \textbf{无电压幅值变量}。决策向量 $x=[\theta;\ P_g;\ P_f;\ \Delta P_d;\
        \lambda]$ 中只有相角 $\theta$（变量布局注释 :DCOPFFormulation），结果结构
        \fld{DCOPFResult} 也只有 \fld{va} 没有 \fld{vm}
        （\srcpath{include/hacdcpf/optimal\_power\_flow/opf\_result.hpp:DCOPFResult}），
        即经典 DC 近似中 $V_m\equiv 1$~pu 的内涵在代码里体现为“根本不存在
        电压幅值未知量”。
  \item \textbf{无无功方程}。负荷侧仅聚合有功 \fld{pd\_mw}/\fld{p\_mw}
        （:build\_dc\_opf\_lp），结果中没有任何无功字段；\fld{Generator::cost\_*} 之外的
        \fld{qmin\_mvar}/\fld{qmax\_mvar} 不进入模型。
  \item \textbf{无损}。结果对象固定写入模型局限声明：``DC OPF uses a lossless
        linearized AC network and does not model AC voltage magnitude, reactive
        power, or converter physics.''（:extract\_dc\_opf\_result）。
  \item \textbf{小角度假设未显式检验}。非松弛节点相角仅受宽界
        $-\pi\le\theta_i\le\pi$ 约束（:build\_dc\_opf\_lp），代码不做
        $|\theta_i-\theta_j|$ 越限告警；小角度是 DC 近似的理论前提而非运行期
        校验项，角度张角大的工况误差由使用者自担。
\end{enumerate}

\subsection{数学模型}\label{subsec:dcopf-model}

\subsubsection*{决策变量与布局}
设活跃母线 $n_b$、活跃发电机 $n_g$、活跃支路 $n_l$（均按 \fld{in\_service}
过滤后计数，:build\_dc\_opf\_lp），决策向量分段布局为
（:DCOPFFormulation 注释，偏移量赋值 :build\_dc\_opf\_lp）
\begin{align*}
x=[\ \underbrace{\theta}_{n_b}\ ;\ \underbrace{P_g}_{n_g}\ ;\
\underbrace{P_f}_{n_l\ \text{（可选）}}\ ;\
\underbrace{\Delta P_d}_{n_b\ \text{（可选）}}\ ;\
\underbrace{P^{gc}}_{n_b\ \text{（可选）}}\ ;\
\underbrace{\lambda}_{\text{PWL 断点}}\ ]
\end{align*}
（\srcpath{src/optimal\_power\_flow/dc\_opf.cpp:build\_dc\_opf\_lp}，函数
\srcpath{build\_dc\_opf\_lp}）。其中 $P_f$ 变量仅当
\fld{include\_branch\_limits}=\texttt{true} 时引入（:build\_dc\_opf\_lp），$\Delta P_d$
削减变量仅当 \fld{load\_shedding}=\texttt{true} 时每母线一个（:build\_dc\_opf\_lp），$P^{gc}$
仅当 \fld{allow\_fixed\_generation\_curtailment}=\texttt{true} 时每母线一个，
$\lambda$ 为凸二次成本发电机（$\fld{cost\_c2}>10^{-12}$ 且
$\fld{pmax\_mw}>\fld{pmin\_mw}$，:build\_dc\_opf\_lp）的分段线性凸组合权重，段数取
$\mathrm{clamp}(\fld{pwl\_segments},1,1000)$（:build\_dc\_opf\_lp）。

\subsubsection*{目标函数}
目标为发电成本最小化，存在 LP 与 QP 两种口径，由后端分派决定实际生效者
（见\ref{subsec:dcopf-solve}节）。

\textbf{（a）LP 口径。} 无线性化段的发电机按边际成本计价；成本数据全零时
取默认边际成本 1（成本单位/MWh），避免目标恒零
（\srcpath{src/optimal\_power\_flow/dc\_opf.cpp:build\_dc\_opf\_lp}）。由于 $P_g$ 以
标幺值进入 LP，成本系数乘以 $S_{\mathrm{base}}$ 使报告目标值与等式行对偶
保持 MW 工程单位（注释 :build\_dc\_opf\_lp）：
\begin{align*}
\min\ \sum_{k\in\mathcal{G}_{\mathrm{lin}}} c_{1,k}\,S_{\mathrm{base}}\,P_{g,k}
\;+\;\sum_{i} \mathrm{VOLL}\cdot S_{\mathrm{base}}\,\Delta P_{d,i}
\;+\;\sum_i\gamma_i S_{\mathrm{base}}P^{gc}_i
\;+\;\sum_{k\in\mathcal{G}_{\mathrm{pwl}}}\sum_{p}
f_k(\hat P_{k,p})\,\lambda_{k,p}
\end{align*}
（:build\_dc\_opf\_lp）。凸二次成本的发电机以 $\lambda$ 凸组合
插值逼近：断点
\begin{align*}
\hat P_{k,p} &= \fld{pmin\_mw}_k + \alpha_p\bigl(\fld{pmax\_mw}_k-\fld{pmin\_mw}_k\bigr),
& \alpha_p &= \frac{p}{N_{\mathrm{seg}}},\quad p=0,\dots,N_{\mathrm{seg}}\\
f_k(\hat P_{k,p}) &= c_{2,k}\hat P_{k,p}^{2}+c_{1,k}\hat P_{k,p}+c_{0,k}
\end{align*}
（:build\_dc\_opf\_lp），$\lambda_{k,p}\in[0,1]$（:build\_dc\_opf\_lp）。注释说明：凸性加
最小化使下包络自动选取相邻断点，无需二进制或 SOS2 变量（:build\_dc\_opf\_lp）。

\textbf{（b）QP 口径。} 真二次成本写成标准形
$\min \tfrac12 x^{\top}Qx+c^{\top}x$（\srcpath{src/optimal\_power\_flow/dc\_opf.cpp:build\_dc\_opf\_qp}，
函数 \srcpath{build\_dc\_opf\_qp}）：
\begin{align*}
Q_{P_{g,k},P_{g,k}} &= 2\,c_{2,k}\,S_{\mathrm{base}}^{2},
& c_{P_{g,k}} &= c_{1,k}\,S_{\mathrm{base}}
\end{align*}
（标幺换算推导见注释 :build\_dc\_opf\_qp；赋值 :build\_dc\_opf\_qp）。$c_2$、$c_1$ 同时为零时
同样回退默认边际成本 1（:build\_dc\_opf\_qp）；削减惩罚从 LP 模型原样拷贝（:build\_dc\_opf\_qp）。
常数项 $c_0$ 不进 QP 矩阵，仅在报告目标值时补加
（:compute\_qp\_objective；$c_0$ 累加亦在 :compute\_qp\_objective）。

\subsubsection*{节点功率平衡}
所有母线（含松弛节点）各有一行平衡方程
（\srcpath{src/optimal\_power\_flow/dc\_opf.cpp:build\_dc\_opf\_lp}）。显式 $P_f$ 口径下
\begin{align*}
\sum_{g\in\mathcal{G}(i)} P_{g,g} + \Delta P_{d,i}-P^{gc}_i
- \sum_{l:\,\mathrm{from}(l)=i} P_{f,l} + \sum_{l:\,\mathrm{to}(l)=i} P_{f,l}
= P^{\rm gross}_{d,i}-P^{\rm fix}_{g,i}
\end{align*}
（发电机系数 $+1$：:build\_dc\_opf\_lp；削减变量系数 $+1$ 的符号推导注释 :build\_dc\_opf\_lp；
$P_f$ 在 from 侧取 $-1$、to 侧取 $+1$：:build\_dc\_opf\_lp；右端项 $P_d$：:build\_dc\_opf\_lp）。
实现先分别聚合毛需求 $P^{\rm gross}_{d,i}$ 与固定正注入 $P^{\rm fix}_{g,i}$，再形成净右端项。
母线 \fld{pd\_mw} 和 \fld{Load::p\_mw}$\times$\fld{scaling} 的正值进入毛需求，负值进入固定注入；
\fld{StaticGenerator}、\fld{RenewableGen}、\fld{PVSystem} 与 \fld{Storage::p\_mw} 的正值进入固定注入，
负值进入毛需求。因而 $0\le\Delta P_{d,i}\le P^{\rm gross}_{d,i}$，启用 $P^{gc}$ 时
$0\le P^{gc}_i\le P^{\rm fix}_{g,i}$，不会用净负荷错误压缩可削负荷上界。

无显式 $P_f$ 变量时（\fld{include\_branch\_limits}=\texttt{false}），支路潮流
直接代入平衡行，形成经典 $B\theta$ 形式（:build\_dc\_opf\_lp）：
\begin{align*}
\sum_{j} B_{ij}\,\theta_j &= P_{g,i}+\Delta P_{d,i}-P_{d,i},
& B_{ii} &= \textstyle\sum_{j} b_{ij},\quad B_{ij}=-b_{ij},\quad b_{ij}=1/x_{ij}
\end{align*}
（$B$ 矩阵装配 :build\_dc\_opf\_lp；写入等式行 :build\_dc\_opf\_lp）。该路径曾漏掉整个松弛节点
平衡行，现仅跳过 $\theta$ 列中松弛节点的零值项而保留其平衡行（BUG-3 修复
注释 :build\_dc\_opf\_lp；对角元 :build\_dc\_opf\_lp；非对角元 :build\_dc\_opf\_lp）。

\subsubsection*{支路功率定义与输电界}
每条活跃支路一行潮流定义方程（\srcpath{src/optimal\_power\_flow/dc\_opf.cpp:build\_dc\_opf\_lp}）：
\begin{align*}
P_{f,k} - b_k\,(\theta_{\mathrm{from}(k)}-\theta_{\mathrm{to}(k)}) = 0,
\qquad b_k = 1/x_k
\end{align*}
其中 $x_k$ 取 \fld{Branch::x\_pu}，$|x_k|<10^{-12}$ 时兜底为 $10^{-6}$ 防除零
（:build\_dc\_opf\_lp；$B\theta$ 路径同兜底 :build\_dc\_opf\_lp；结果回算同兜底 :extract\_dc\_opf\_result）。
输电能力以变量盒界实现而非显式不等式（:build\_dc\_opf\_lp 注释）：
\begin{align*}
|P_{f,k}| \;\le\; \fld{rate\_a\_mva}_k\cdot\fld{branch\_limit\_margin}\,/\,S_{\mathrm{base}}
\end{align*}
（:build\_dc\_opf\_lp；\fld{rate\_a\_mva}$<10^{-9}$ 视为不限，界放宽至
\fld{kHugeBound}$=10^{20}$，:build\_dc\_opf\_lp、:kHugeBound）。注意仅长期额定 \fld{rate\_a\_mva}
生效，\fld{rate\_b}/\fld{rate\_c} 不进入模型。

\subsubsection*{发电机出力界与相角参考}
发电机界直接折算标幺盒界（:build\_dc\_opf\_lp）：
\begin{align*}
\fld{pmin\_mw}_k/S_{\mathrm{base}} \;\le\; P_{g,k} \;\le\;
\fld{pmax\_mw}_k/S_{\mathrm{base}}
\end{align*}
并以 $\pm 10^{20}$ 截断防无界（:build\_dc\_opf\_lp）。相角参考按连通分量逐分量锚定：
\srcpath{find\_slack\_bus} 先选出第一台 \fld{BusType::SLACK} 型母线，没有则
退化为第 0 号母线（:find\_slack\_bus）；再以在运支路做并查集连通分量划分，含松弛母线
的分量以松弛母线为参考，其余每个分量各自以分量内首母线为参考（:root（lambda））。
各参考相角以盒界 $[0,0]$ 钉死，其余节点自由于 $[-\pi,\pi]$（:build\_dc\_opf\_lp）。
$B\theta$ 路径保留参考列于等式约束中（其盒界已将其钉零，:build\_dc\_opf\_lp）。

\subsubsection*{负荷削减变量与 VOLL}
开启 \fld{load\_shedding} 时每母线一个削减变量
（\srcpath{src/optimal\_power\_flow/dc\_opf.cpp:build\_dc\_opf\_lp}）：
\begin{align*}
0 \;\le\; \Delta P_{d,i} \;\le\; \max\!\bigl(P_{d,i},\ 0.01/S_{\mathrm{base}}\bigr)
\end{align*}
（:build\_dc\_opf\_lp；只能削减正需求）。失负荷价值 \fld{voll}$\le 0$ 时自动取
（:build\_dc\_opf\_lp）
\begin{align*}
\mathrm{VOLL}_{\mathrm{eff}} = \max\Bigl(100,\;
1.1\max_{k}\bigl(c_{1,k}+2c_{2,k}\,\fld{pmax\_mw}_k\bigr)\Bigr)
\end{align*}
即略高于最贵发电机的边际成本上限，保证“能发不削”的字典序。

\subsubsection*{外部电网升格与固定注入}
在求解入口，每台在役 \fld{ExternalGrid} 被升格为一台挂在原母线上的松弛发电机，
界 $\pm 10^{5}$~MW/MVAr，成本系数照搬 \fld{ExternalGrid::cost\_c2/c1/c0}
（\srcpath{src/optimal\_power\_flow/dc\_opf.cpp:solve\_dc\_opf}，函数
\srcpath{solve\_dc\_opf}）；求解后其出力拆回
\fld{external\_grid\_p\_mw}，并从 \fld{pg\_mw} 中裁掉（:separate\_external\_grid\_dispatch（lambda），定义于
\srcpath{solve\_dc\_opf}）。

\subsection{求解流程与后端分派}\label{subsec:dcopf-solve}

\subsubsection*{求解主流程}
\srcpath{solve\_dc\_opf}（\srcpath{src/optimal\_power\_flow/dc\_opf.cpp:solve\_dc\_opf}）
的总体流程如图~\ref{fig:dcopf-flow}：先做 LCC 拒绝与规范投影，再做图论孤岛
预检，然后一次性装配 LP 约束结构并派生 QP 模型，最后按后端枚举分派求解并
回收对偶。

\begin{figure}[H]
\centering
\begin{tikzpicture}[
  node distance=0.55cm,
  box/.style={draw, rectangle, align=center, inner sep=3.5pt, font=\small},
  dec/.style={draw, diamond, aspect=2.1, align=center, inner sep=1.5pt, font=\small},
  arr/.style={-{Stealth[length=2.2mm]}, thick}]
\node[box] (in) {输入 \fld{HybridPowerSystem}};
\node[dec, below=of in] (lcc) {含在役 LCC？};
\node[box, right=1.6cm of lcc] (rej) {拒绝求解\\ \fld{dc-opf:rejected-lcc}};
\node[box, below=of lcc] (proj) {规范投影（开关合并母线）\\ + 外网升格为松弛发电机};
\node[dec, below=of proj] (isl) {孤岛预检\\ 存在无松弛孤岛？};
\node[box, right=1.6cm of isl] (shed) {预削减孤岛负荷\\ 无有效孤岛则报不可行};
\node[box, below=of isl] (build) {装配 LP 约束结构\\ 派生 QP 成本模型};
\node[box, below=of build] (disp) {后端分派（回退链）\\ Gurobi / LCQP / HiGHS / 对偶单纯形};
\node[box, below=of disp] (dual) {支撑 LP 回收对偶（LMP）\\ \fld{compute\_lmp} 门控};
\node[box, below=of dual] (out) {结果归因：反投影回原始母线/支路\\ 填写 \fld{model\_limitations} 等诚实字段};
\draw[arr] (in) -- (lcc);
\draw[arr] (lcc) -- node[above, font=\small]{是} (rej);
\draw[arr] (lcc) -- node[right, font=\small]{否} (proj);
\draw[arr] (proj) -- (isl);
\draw[arr] (isl) -- node[above, font=\small]{是} (shed);
\draw[arr] (isl) -- node[right, font=\small]{否/剪枝后} (build);
\draw[arr] (shed.south) |- (build.east);
\draw[arr] (build) -- (disp);
\draw[arr] (disp) -- (dual);
\draw[arr] (dual) -- (out);
\end{tikzpicture}
\caption{DC OPF 求解主流程（对应
\srcpath{src/optimal\_power\_flow/dc\_opf.cpp:solve\_dc\_opf}）}
\label{fig:dcopf-flow}
\end{figure}

\textbf{LCC 拒绝}：任何在役 \fld{LCCConverter} 都使求解直接返回，状态写明
“this AC-only linear formulation does not model LCC AC/DC coupling”，并提示改用
ParityIPM 或 Ipopt 的 AC OPF（:solve\_dc\_opf；见第\ref{sec:parityipm}章）。

\textbf{规范投影}：与 AC OPF/潮流共用同一份 canonical 拓扑——闭合开关与
断路器合并母线（\fld{strip\_dead\_islands}=\texttt{false}，:solve\_dc\_opf）。求解后
\fld{va}/\fld{lmp} 按密集量、\fld{load\_shedding\_mw} 按广延量反投影回原始
母线，\fld{pf\_mw} 按 \fld{branch\_orig\_to\_proj} 映射回原始支路（:separate\_external\_grid\_dispatch（lambda））。
失败解不回灌投影，以免用尺寸异常掩盖原始不可行性（注释 :separate\_external\_grid\_dispatch（lambda））。

\textbf{孤岛预检}：用图模块 \srcpath{build\_power\_system\_graph} 与
\srcpath{analyze\_topology} 在昂贵的 LP 装配前识别死岛——\fld{IsolatedLoad}
孤岛，或 \fld{NoSlack} 且岛内无在役可调度发电机的孤岛（:solve\_dc\_opf）；
\fld{NoSlack} 但含可调度电源的孤岛判活并正常建模。死岛负荷（母线 \fld{pd\_mw} 与
\fld{ac.loads} 两源相加，口径与建模一致，P1a 注释 :solve\_dc\_opf）计入直接削减
（:solve\_dc\_opf）；全部孤岛均无效时直接返回
\texttt{"Infeasible: no island with slack bus"}（:solve\_dc\_opf）；否则把死岛上的
负荷、发电机、静态发电机、新能源、光伏、储能与相连支路在本地副本中置为
退出再建模（:solve\_dc\_opf）。

\textbf{装配}：\srcpath{build\_dc\_opf\_lp} 构建约束结构（:solve\_dc\_opf 调用；
重复母线编号等拓扑错误在此捕获并转状态，:solve\_dc\_opf）；空系统或无发电机
直接返回 \texttt{"Empty system or no generators"}（:solve\_dc\_opf）；随后
\srcpath{build\_dc\_opf\_qp} 派生真二次成本模型（:solve\_dc\_opf）。

\subsubsection*{后端选择与回退链}
后端枚举 \fld{DCOPFSolverBackend} 取值为 \fld{Auto}/\fld{Native}/
\fld{NativeQP}/\fld{HiGHS}/\fld{Gurobi}
（\srcpath{include/hacdcpf/optimal\_power\_flow/opf\_options.hpp:DCOPFSolverBackend}）。
各候选通道（均定义于 \srcpath{src/optimal\_power\_flow/dc\_opf.cpp}）：
\begin{itemize}
  \item \textbf{Gurobi QP}（:try\_gurobi\_qp（lambda），定义于 \srcpath{solve\_dc\_opf}）：
        经 engine 适配器 \fld{GurobiAdapter} 求解 QP；仅在
        \fld{gurobi.available()} 探测到安装/许可证时尝试，否则记
        \fld{unavailable}；
  \item \textbf{NativeLCQP}（:try\_native\_qp（lambda），定义于 \srcpath{solve\_dc\_opf}）：
        自研线性约束 QP 内点求解器（MIPSolvers engine 的
        \srcpath{lcqp\_solver}，经转发头
        \srcpath{include/hacdcpf/engine/kernel/ipm/lcqp\_solver.hpp} 引入），
        求解真二次成本 QP；
  \item \textbf{HiGHS 通道}（:try\_highs（lambda），定义于 \srcpath{solve\_dc\_opf}）：
        HiGHS 本身无原生 QP 接口，先经 \fld{SolverEngine} 以
        \fld{preferred\_solver}=\fld{"HiGHS"} 路由一次 QP（实际求解者可能是
        Gurobi 或 LCQP，\fld{solver\_name} 如实反映，注释 :try\_highs（lambda））；QP
        失败后回退 \fld{HighsAdapter::solve\_lp} 求 LP，此时目标与对偶只反映
        线性成本模型，求解器名记为 \fld{"HiGHS-LP(QP-fallback)"}
        （:try\_highs（lambda））；
  \item \textbf{NativeDualSimplex}（:try\_native\_simplex（lambda），定义于
        \srcpath{solve\_dc\_opf}）：自研对偶单纯形
        \srcpath{solve\_lp\_with\_basis}（同样来自 MIPSolvers engine，经转发头
        \srcpath{include/hacdcpf/engine/kernel/lp\_kernel/dual\_simplex.hpp}），
        求解线性化成本 LP。
\end{itemize}
回退顺序（:solve\_dc\_opf）：\fld{Auto} 依次试 Gurobi-QP → NativeLCQP → HiGHS
通道 → NativeDualSimplex（:solve\_dc\_opf）；\fld{Gurobi} 同序但从 Gurobi 起
（:solve\_dc\_opf）；\fld{NativeQP} 失败仅回退 NativeDualSimplex（:solve\_dc\_opf）；
\fld{HiGHS} 不可用时直接失败返回（:solve\_dc\_opf）；\fld{Native} 只用对偶单纯形
（:solve\_dc\_opf）。每一步尝试都记入 \fld{solver\_chain}，格式
\fld{<backend>:<ok|fail(detail)>}，供事后审计真实回退路径（:record\_chain（lambda）、
:solve\_dc\_opf）。

\textbf{目标语义消歧}：LP 回退成功时必须把 \fld{use\_qp} 复位，否则会用 QP
口径重算 LP 可行点，造成目标/对偶语义不一致（注释 :solve\_dc\_opf）。最终按收敛路径重算或保留目标值并填写
\fld{objective\_model}：QP 收敛为 \fld{"QP"}（真二次成本含 $c_0$，由
\srcpath{compute\_qp\_objective} 重算，:solve\_dc\_opf），否则为 \fld{"LP-PWL"}
或 \fld{"LP"}，并以 \fld{pwl\_segments\_effective} 记录实际分段数（QP 时置 0，
:solve\_dc\_opf）。

\subsubsection*{对偶恢复与支撑 LP}
QP/LP 主解的对偶若不完整（\srcpath{has\_complete\_duals} 校验行对偶与盒对偶
数量，:has\_complete\_duals），且 \fld{compute\_lmp}=\texttt{true} 时，追加一次支撑 LP：
以原 LP 模型（QP 情形把成本换成梯度 $Qx+c$）用对偶单纯形重解，取回全部行
对偶与盒对偶（:populate\_missing\_duals\_from\_supporting\_lp；QP 梯度替换
:populate\_missing\_duals\_from\_supporting\_lp；门控 :solve\_dc\_opf）。代码注释如实声明该通道的已知缺陷：自研单纯形对
有界变量 LP 的盒对偶恢复不正确，\fld{include\_branch\_limits}=\texttt{true}
且支路界真正起作用时 \fld{branch\_mu\_lower}/\fld{branch\_mu\_upper}
\textbf{不可靠}；LMP（节点平衡行对偶）仍有效可用于经济调度
（NOTE 注释块 :populate\_missing\_duals\_from\_supporting\_lp）。因此结果中 \fld{branch\_mu\_valid} 在 KKT 级提取
实现之前\textbf{恒为 false}（:extract\_dc\_opf\_result）。

\subsubsection*{与 HiGHS/AML 建模层的分派关系}
本求解器只经 engine 适配器调用 HiGHS/Gurobi（转发头
\srcpath{include/hacdcpf/engine/solver/external/adapters.hpp}），自身不经过
AML 建模层。平台另有一套基于 MIPSolvers AML 的 DC OPF 实现
\srcpath{power\_models::solve\_dc\_opf}
（\srcpath{src/power\_models/dc\_opf\_builder.cpp:solve\_dc\_opf}），以集合/参数/变量的
代数建模语言重构同一数学模型并交由 AML 后端求解，用于交叉验证；其建模细节
仅见独立的 \srcpath{docs/modules/power\_models/power\_models\_manual.tex}。
两套实现相互独立：本章所述的孤岛预检、外网
升格、负荷削减与支撑 LP 对偶恢复只存在于引擎版。

\subsection{直流电网与换流器覆盖情况}\label{subsec:dcopf-dcgrid}

\textbf{不覆盖。} 具体口径：
\begin{itemize}
  \item 直流母线、直流支路、VSC 换流器、DC/DC 变换器、能量路由器均不进入
        模型——建模仅遍历 \fld{sys.ac.*} 容器（
        \srcpath{src/optimal\_power\_flow/dc\_opf.cpp:build\_dc\_opf\_lp}），结果结构也
        无任何直流侧字段
        （\srcpath{include/hacdcpf/optimal\_power\_flow/opf\_result.hpp:DCOPFResult}）；
  \item 结果固定声明 \fld{model\_scope}=\fld{"dc-opf:linear-no-converter-model"}
        （:extract\_dc\_opf\_result），\fld{converter\_model\_scope} 各换流器标志位恒为
        \texttt{false}（注释见
        \srcpath{include/hacdcpf/optimal\_power\_flow/opf\_result.hpp:DCOPFResult}）；
  \item 含在役 LCC 的系统被整体拒绝而非近似处理（:solve\_dc\_opf）；
  \item 直流侧负荷由此形成已知盲区：可靠性评估调用方对此有显式告警
        （“DC 负荷中断不计入 OPF，EENS/LOLE 可能偏乐观”，
        \srcpath{src/reliability/reliability\_assessment.cpp:evaluate\_state}）。
\end{itemize}
需要交直流联合优化的场景应使用 Parity IPM 全空间 OPF
（第\ref{sec:parityipm}章）或 AML 层的 acdcopf builder；后者见独立的
\srcpath{docs/modules/power\_models/power\_models\_manual.tex}。

\subsection{NativeLCQP 的结构化初值}

对每个在役 AC 支路连通分量 $\mathcal C$，先在发电上下界内调整 $P_g$，不足
部分再由允许的 $\Delta P_d$ 补足，使
\begin{equation}
\sum_{i\in\mathcal C}(P_{g,i}+\Delta P_{d,i}-P_{d,i})=0.
\end{equation}
每个分量选择一个参考角并删除其行列，求约化 Laplacian
$B_{\mathcal C,r}\theta_r=p_{\mathcal C,r}$；随后由
$P_{f,ij}=(\theta_i-\theta_j)/x_{ij}$ 回填显式支路流，并补齐 PWL convexity/
interpolation 变量。因此连通分量功率和是 Laplacian 可解性的兼容条件，而非
额外优化问题。生成点经过变量盒界与 $\|A_{eq}x_0-b_{eq}\|_\infty$ 审计后，
只通过 \fld{QPModel::x0} 交给 NativeLCQP；其他后端不消费该初值，构建失败则
精确回到确定性冷初始化。该步骤复杂度为图遍历线性工作加各分量一次约化
Laplacian 稀疏分解，不运行非线性 Phase I。

作为 ACOPF warm-start-only 子问题时，
\fld{phase\_one\_time\_limit\_ms} 从 \srcpath{solve\_dc\_opf} 入口开始计时，
因此 canonical projection、formulation、上述结构初值、symbolic analyze 与
numeric iteration 都计入总预算。NativeLCQP 在迭代/分解边界协作式检查截止；
一个已经开始的稀疏分解不可中断，允许完成后报告
\fld{phase\_one\_budget\_overshoot\_ms}。超时迭代仍须通过显式原坐标残差门槛
才可作为 warm start，且不声明 DCOPF 最优；有限预算耗尽后不再运行 simplex
fallback。

\subsection{选项：DCOPFOptions}

\compmeta{DCOPFOptions}{\srcpath{include/hacdcpf/optimal\_power\_flow/opf\_options.hpp}}

全部字段如下（定义 :DCOPFOptions）。

\begin{paramtable}
\fld{feasibility\_tol} & $\varepsilon_{\mathrm{f}}$ & — & $10^{-9}$--$10^{-4}$（经验值） & $10^{-6}$ & 原始可行性容差；传入 LCQP（\fld{tol\_primal}）与对偶单纯形\\
\fld{optimality\_tol} & $\varepsilon_{\mathrm{o}}$ & — & $10^{-9}$--$10^{-4}$（经验值） & $10^{-6}$ & 对偶/最优性容差；传入 LCQP（\fld{tol\_dual}/\fld{tol\_gap}）与对偶单纯形\\
\fld{max\_iterations} & — & — & $10^{2}$--$10^{5}$（经验值） & 10000 & 单纯形/LCQP 最大迭代数\\
\fld{verbose} & — & — & true/false & false & spdlog 打印规模与收敛信息（:solve\_dc\_opf）\\
\fld{solver} & — & — & Auto/Native/\allowbreak NativeQP/\allowbreak HiGHS/\allowbreak Gurobi & Auto & 后端选择与回退链入口（:DCOPFOptions）\\
\fld{pwl\_segments} & $N_{\mathrm{seg}}$ & — & 1--1000（clamp，:build\_dc\_opf\_lp） & 4 & 凸二次成本的 LP 分段线性插值段数；QP 路径不使用\\
\fld{include\_branch\_limits} & — & — & true/false & true & 是否引入显式 $P_f$ 变量及 \fld{rate\_a\_mva} 盒界；false 时退回纯 $B\theta$ 形式且无支路约束\\
\fld{branch\_limit\_margin} & — & — & 0.9--1.0（经验值） & 1.0 & 支路界缩放系数，乘在 \fld{rate\_a\_mva} 上（:DCOPFOptions）\\
\fld{load\_shedding} & — & — & true/false & true & 是否引入每母线削减变量 $\Delta P_d$\\
\fld{voll} & VOLL & —/MWh & $\ge 0$ & 0.0 & 失负荷价值；$\le 0$ 时按 $\max(100,\,1.1\times$最高边际成本$)$ 自动取值（:DCOPFOptions）\\
\fld{lexicographic\_load\_shedding} & — & — & true/false & false & 两阶段精确优先级：先最小总削减，再固定最优削减量并最小发电成本；HiGHS 第二阶段使用 PWL 成本 LP\\
\fld{full\_redispatch\_from\_zero} & — & — & true/false & false & HL-I/HL-II 充裕度再调度语义，在线机组统一采用 $0\le P_g\le P_g^{\max}$\\
\fld{allow\_fixed\_generation\_curtailment} & — & — & true/false & false & 是否引入逐母线 $P^{gc}$；与全负荷削减、零机组下界共同给出任意拓扑的构造性可行点\\
\fld{fixed\_generation\_curtailment\_cost} & $\gamma$ & 成本/MWh & $\ge0$ & 1.0 & 第二级目标中的固定正注入削减代价；第一级最小切负荷不使用\\
\fld{compute\_lmp} & — & — & true/false & true & 主解后是否追加支撑 LP 回收对偶；false 时 \fld{lmp} 清空并注明原因（:solve\_dc\_opf）\\
\fld{structural\_warm\_start} & — & — & true/false & true & 为 NativeLCQP 构造逐分量功率配平与约化 Laplacian 初值；LP/外部后端忽略\\
\fld{phase\_one\_linear\_relaxation} & — & — & true/false & false & 仅 HiGHS 后端的 Phase-I 凸松弛：直接解 LP-PWL 而非把二次模型路由到 QP 后端；返回点为 warm start，非二次成本 DCOPF 最优证书\\
\fld{compact\_quadratic\_model} & — & — & true/false & false & Phase-I QP 省略仅供 LP-PWL 使用的插值变量与行\\
\fld{accept\_phase\_one\_iterate} & — & — & true/false & false & 允许非最优 NativeLCQP 末迭代经显式残差门槛成为 warm-start-only 结果\\
\fld{phase\_one\_iterate\_tolerance} & $\varepsilon_I$ & pu & $\ge0$ & 0.1 & warm-start-only 末迭代的 equality/box 最大残差门槛\\
\fld{phase\_one\_time\_limit\_ms} & $T_I$ & ms & $\le0$ 或 $>0$ & 0 & warm-start-only 端到端协作式墙钟预算；$\le0$ 不限时\\
\end{paramtable}

\subsection{结果字段：DCOPFResult}

\compmeta{DCOPFResult}{\srcpath{include/hacdcpf/optimal\_power\_flow/opf\_result.hpp}}

全部字段如下（定义 :DCOPFResult）。向量字段的索引空间：\fld{va}/\fld{lmp}/
\fld{load\_shedding\_mw}/\fld{fixed\_generation\_curtailment\_mw} 按原始母线向量位置（经反投影）；\fld{pg\_mw} 按
\fld{sys.ac.generators} 位置；\fld{pf\_mw}/\fld{branch\_mu\_*} 按
\fld{sys.ac.branches} 位置；\fld{external\_grid\_p\_mw} 按
\fld{sys.ac.external\_grids} 位置。

\begin{paramtable}
\fld{va} & $\theta_i$ & rad & $(-\pi,\pi)$ & 空 & 母线电压相角（各连通分量参考母线处恒 0）\\
\fld{pg\_mw} & $P_g$ & MW & \fld{pmin\_mw}--\fld{pmax\_mw} & 空 & 发电机有功调度（由标幺乘 $S_{\mathrm{base}}$ 还原，:extract\_dc\_opf\_result）\\
\fld{external\_grid\_p\_mw} & $P_{\mathrm{ext}}$ & MW & $\pm 10^{5}$ & 空 & 外网（升格发电机）有功（:DCOPFResult）\\
\fld{pf\_mw} & $P_f$ & MW & $\pm$\fld{rate\_a\_mva} & 空 & 支路有功，from$\to$to 为正；无显式 $P_f$ 时由 $(\theta_f-\theta_t)/x$ 回算（:DCOPFResult）\\
\fld{converged} & — & — & true/false & false & 后端成功且规范坐标原始可行性证书通过时为 true\\
\fld{iterations} & — & — & — & 0 & 最终后端迭代数\\
\fld{objective} & $f$ & —/h & — & 0.0 & 目标值；口径见 \fld{objective\_model}（QP 含 $c_0$，LP 为线性/PWL 成本）\\
\fld{status} & — & — & — & 空 & 后端状态串或拒绝/失败原因（如 \texttt{"HiGHS not available"}）\\
\fld{solver\_name} & — & — & — & 空 & 实际收敛后端名（可能含 \fld{(QP-fallback)} 后缀）\\
\fld{runtime\_sec} & — & s & — & 0.0 & 含投影与预检的端到端耗时（:DCOPFResult）\\
\fld{structural\_warm\_start\_*} & — & pu/计数/布尔 & — & inf/0/false & 请求、构建、实际消费、分量/分解数、原坐标 equality residual 与状态原因\\
\fld{solver\_initial\_primal\_residual} & — & 缩放坐标 & $\ge0$ & inf & NativeLCQP 第 0 次迭代的 primal residual；含严格正 bound-slack 残差\\
\fld{phase\_one\_warm\_start\_only} & — & — & true/false & false & 末迭代只获准供上层 warm start 使用，不是 DCOPF 最优证书\\
\fld{phase\_one\_budget\_exhausted} & — & — & true/false & false & formulation 或 LCQP 阶段达到有限墙钟预算\\
\fld{phase\_one\_budget\_overshoot\_ms} & — & ms & $\ge0$ & 0 & 不可中断在途工作造成的预算超时\\
\fld{native\_qp\_symbolic\_analyze\_calls} & — & 次 & 0--1 & 0 & 本次 NativeLCQP KKT 稀疏模式分析次数；每轮 numeric factorization 复用该模式\\
\fld{lmp} & $\lambda_i$ & —/MWh & $\ge 0$（无罚项时） & 空 & 节点边际电价：平衡行对偶除以 $S_{\mathrm{base}}$（:DCOPFResult）\\
\fld{branch\_mu\_lower} & $\mu^{-}$ & —/MWh & — & 空 & 支路下界盒对偶/$S_{\mathrm{base}}$；当前\textbf{不可靠}（:DCOPFResult）\\
\fld{branch\_mu\_upper} & $\mu^{+}$ & —/MWh & — & 空 & 支路上界盒对偶/$S_{\mathrm{base}}$；当前\textbf{不可靠}（:DCOPFResult）\\
\fld{branch\_mu\_valid} & — & — & true/false & false & 恒 false，直至 KKT 级盒对偶提取实现（:extract\_dc\_opf\_result）\\
\fld{branch\_mu\_validity\_reason} & — & — & — & 空 & 上述无效性的人类可读说明\\
\fld{lmp\_valid} & — & — & true/false & false & 行对偶证书完整时为 true（:DCOPFResult）\\
\fld{lmp\_validity\_reason} & — & — & — & 空 & LMP 来源/缺失原因说明\\
\fld{load\_shedding\_mw} & $\Delta P_d$ & MW & $\ge 0$ & 空 & 逐母线模型内生削减量；$<10^{-6}$ 数值噪声清零\\
\fld{total\_load\_shedding\_mw} & — & MW & $\ge 0$ & 0.0 & 削减总量\\
\fld{fixed\_generation\_curtailment\_mw} & $P^{gc}$ & MW & $\ge0$ & 空 & 逐母线固定正注入削减量\\
\fld{total\_fixed\_generation\_curtailment\_mw} & — & MW & $\ge0$ & 0.0 & 固定正注入削减总量\\
\fld{full\_redispatch\_from\_zero} & — & — & true/false & false & 回显本次有效机组下界语义\\
\fld{primal\_feasibility\_certified} & — & — & true/false & false & 反投影前在实际规范 LP/QP 坐标中认证全部等式、不等式和盒界\\
\fld{maximum\_primal\_violation\_mw} & — & MW & $\ge0$ & inf & 上述规范坐标证书的最大违反量\\
\fld{primal\_feasibility\_reason} & — & — & — & 空 & 证书通过或拒绝原因\\
\fld{solver\_chain} & — & — & — & 空 & 依序记录各后端尝试结果，供审计回退路径（:DCOPFResult）\\
\fld{objective\_model} & — & — & QP/LP-PWL/LP & 空 & 报告目标值对应的成本口径（:DCOPFResult）\\
\fld{pwl\_segments\_effective} & — & — & 0--1000 & 0 & LP-PWL 实际段数；0 表示 QP 真二次成本\\
\fld{converter\_model\_scope} & — & — & — & 全 false & 换流器建模覆盖声明（DC OPF 不建模任何换流器）\\
\fld{model\_limitations} & — & — & — & 空 & 本次求解的模型局限声明（无损线性化、储能单时段等，:extract\_dc\_opf\_result）\\
\end{paramtable}

\subsection{典型使用场景}

\begin{itemize}
  \item \textbf{可靠性评估（FMEA/状态评估）}：每个故障状态开启毛负荷削减、固定正注入削减和
        零机组下界；先解纯最小切负荷 LP，再加入最优切负荷等式并最小化 PWL 发电与 $P^{gc}$ 成本，
        不使用有限 VOLL 冒充字典序
        （\srcpath{src/reliability/reliability\_assessment.cpp:evaluate\_state} 调用，
        选项设置 :run\_nonsequential\_mc；AC-only 适用范围声明 :evaluate\_state）。
  \item \textbf{EV--交通耦合联合优化}：按时段把 EV 充电负荷注入系统后逐时段
        求解 DC OPF，累加发电成本并取 LMP 构成联合社会福利目标
        （\srcpath{src/ev\_power\_traffic/joint\_social\_welfare.cpp:solve\_joint\_social\_welfare}；
        CTM 变体 \srcpath{src/ev\_power\_traffic/ctm\_joint\_welfare.cpp:simulate\_ev\_power\_traffic\_ctm\_joint}）。
  \item \textbf{SPPT 语义保持验证}：MR3d 用 DC OPF 节点电价校验
        rich$\to$canonical 投影前后的语义不变性
        （\srcpath{src/sppt/metamorphic.cpp:mr3\_semantic\_preservation\_opf\_dual}）。
  \item \textbf{HTTP 后端}：v1 分析接口在 \fld{analysis}=\fld{"optimal\_power\_flow"}
        且 \fld{solver}=\fld{"dc"} 时调用本求解器并序列化结果
        （\srcpath{src/server/runtime\_api\_v1.cpp:execute\_analysis}，函数
        \srcpath{execute\_analysis}）。
  \item \textbf{公共门面}：\srcpath{src/api/hacdcpf.cpp:solve\_dc\_opf} 透传封装。
\end{itemize}
如实说明两处“不调用”：\textbf{time\_series} 多时段/年度流水线挂接的是
AC OPF（\srcpath{src/time\_series/time\_series\_pf.cpp} 引入
\srcpath{ac\_opf\_solver.hpp}），不使用本求解器；\textbf{market} 模块的
SCUC/SCED 为其自研 LP 体系（\srcpath{src/market/market\_simulation.cpp:build\_pricing\_model}
附近），亦不复用 \srcpath{solve\_dc\_opf}。

\subsection{验证与校核}

\begin{itemize}
  \item \textbf{MATPOWER 标准算例}：case9 收敛且目标有限、\fld{pg\_mw} 尺寸
        正确（\srcpath{tests/test\_matpower\_cases.cpp:TEST\_CASE("DC OPF: case9 solves and objective is finite")}）；case14 调度
        逐机检界（同文件 :TEST\_CASE("DC OPF: case14 dispatch within generator limits")）。
  \item \textbf{AC/DC/PF 三方交叉验证}：小算例集上 DC OPF 收敛后以 AC 潮流
        回代其调度检验可行性；case57 上“DC 调度$\to$AC 潮流”不保证收敛，
        仅告警不断言（\srcpath{tests/test\_acopf\_dcopf\_crossval.cpp:TEST\_CASE("ACOPF/DCOPF/PF cross-validation on small cases")}，
        豁免逻辑 :TEST\_CASE("ACOPF/DCOPF/PF cross-validation on small cases")）。
  \item \textbf{可行性自检}：\srcpath{check\_dc\_opf\_feasibility} 对 case14
        解复检通过（同文件 :TEST\_CASE("DC OPF solution feasibility")）。
  \item \textbf{孤岛回归}：孤立负荷岛逐母线削减记账（:TEST\_CASE("DC OPF isolated-island shedding is reported per bus")）、无 authored
        松弛但含可调度电源的孤岛\textbf{被保留}、岛内机组出力承担岛内负荷而不做削减
        （:TEST\_CASE("DC OPF preserves a source-capable island without an authored slack")）、松弛节点平衡行纳入复检（:TEST\_CASE("DC OPF feasibility check includes slack-bus balance")）。
  \item \textbf{调度质量}：case30 上与 AC OPF（NativeAC）对比，$\sum|\Delta
        P_g|/\sum P_g$ 容忍阈值为 1.0——注释明言“DC OPF is a crude
        approximation”（同文件 :TEST\_CASE("DC OPF dispatch quality")，阈值同该用例）。
  \item \textbf{LMP 口径}：无支路界时 LMP 非负且不超过最高边际成本的数量级
        （防 $S_{\mathrm{base}}$ 换算 bug，同文件 :TEST\_CASE("DC OPF LMP: sign and scale")）；无拥塞时全网
        LMP 一致（:TEST\_CASE("DC OPF LMP: uncongested uniformity")）。
\end{itemize}
验证体系的总体组织见第\ref{sec:verification}章。

\subsection{边界与局限}

\begin{enumerate}
  \item \textbf{DC 近似的固有盲区}：无损、无电压幅值、无无功（结果
        \fld{model\_limitations} 首条，\srcpath{src/optimal\_power\_flow/dc\_opf.cpp:extract\_dc\_opf\_result}）。
        电压越限、无功不足、网损引起的附加阻塞均不可见；重载、大角度张角或
        低 $x/r$ 比配电网场景下误差显著，代码不做任何运行期误差估计。
  \item \textbf{支路模型简化}：仅取 \fld{x\_pu}，忽略电阻、分接头、移相角与
        充电电纳（:build\_dc\_opf\_lp）；限值只用 \fld{rate\_a\_mva}（:build\_dc\_opf\_lp）；零电抗支路
        以 $x=10^{-6}$ 兜底而非合并（:build\_dc\_opf\_lp），可能引入数值
        病态。
  \item \textbf{支路拥塞对偶不可用}：\fld{branch\_mu\_valid} 恒 false
        （:extract\_dc\_opf\_result）；需要拥塞租金分解的研究只能使用 LMP 总量，不能拆分
        能量/拥塞分量（NOTE 注释 :populate\_missing\_duals\_from\_supporting\_lp）。
  \item \textbf{LMP 的适用范围}：LMP 取自支撑 LP 平衡行对偶，标幺换算已按
        $1/S_{\mathrm{base}}$ 还原（:extract\_dc\_opf\_result）；\fld{include\_branch\_limits}
        =\texttt{true} 时 LMP 数值仍有效，但注意测试约定 LMP 专项用例是在
        \fld{include\_branch\_limits}=\texttt{false} 下跑的
        （\srcpath{tests/test\_acopf\_dcopf\_crossval.cpp:TEST\_CASE("DC OPF LMP: sign and scale")} 注释）。
  \item \textbf{无直流网络与换流器}：见本章\ref{subsec:dcopf-dcgrid}节；
        直流负荷对结果完全不可见。
  \item \textbf{储能为固定注入}：\fld{Storage::p\_mw} 作为不可调注入扣除，
        无跨时段 SOC 动态（:extract\_dc\_opf\_result 如实声明）。
  \item \textbf{外网升格为宽界松弛发电机}：$\pm 10^{5}$~MW（:solve\_dc\_opf），
        多外网并存时各外网出力由成本排序决定，不模拟实际互联调度协议。
  \item \textbf{后端可用性依赖环境}：Gurobi 需本机安装与许可证
        （:try\_gurobi\_qp（lambda））；\fld{HiGHS} 后端在库不可用时直接失败
        （:solve\_dc\_opf）；QP 经 \fld{SolverEngine} 路由时实际求解者以
        \fld{solver\_name}/\fld{solver\_chain} 为准（:try\_highs（lambda））。
  \item \textbf{LP-PWL 插值误差}：凸二次成本走 LP 路径时为分段线性近似，
        精度由 \fld{pwl\_segments}（默认 4）控制；报告口径以
        \fld{objective\_model}/\fld{pwl\_segments\_effective} 区分
        （:solve\_dc\_opf）。
  \item \textbf{相角界为宽界}：$[-\pi,\pi]$ 仅防无界（:build\_dc\_opf\_lp），不构成
        静态安全（角差）约束；需要角差约束的研究应改用
        AC OPF 或 AML 建模层（见
        \srcpath{docs/modules/power\_models/power\_models\_manual.tex}）。
\end{enumerate}

% 交直流混合最优潮流：Parity 建模层（全空间 NLP 的变量/约束布局、目标组装与 KKT 求值）
% 本章文件由 \input 引入，不含导言区。
\section{交直流混合最优潮流：Parity 建模层}\label{sec:parityform}

本章与第~\ref{sec:parityipm}~章合起来给出平台的\textbf{交直流混合最优潮流}
（AC/DC hybrid OPF）本体：统一入口 \srcpath{opf::solve\_ac\_opf} 在 Auto 后端下
检测到任何 DC/VSC/LCC 组件时，即把算例分派到本章建模 + 下章求解核的路线
（分派规则见第~\ref{sec:acopfnative}~章；纯交流算例是本模型的直流块为零的
退化情形，亦可走第~\ref{sec:acopfnative}~章的紧致路径）。本章给出该建模层的
完整口径：它把富语义 \fld{HybridPowerSystem} 经 canonical 投影装配为一个
全空间（full-space）非线性规划（NLP），并暴露目标值、目标梯度/对角 Hessian、
等式残差、非线性不等式、两套稀疏雅可比与拉格朗日 Hessian 共七类 KKT 求值函数。
该建模层是 Parity IPM 求解核（第~\ref{sec:parityipm}~节）与嵌入式 Ipopt 路径共用
的同一份问题定义；本章只覆盖建模与求值，迭代算法本身见
第~\ref{sec:parityipm}~节。全部论述以
\srcpath{include/hacdcpf/optimal\_power\_flow/formulation.hpp}（447 行）
与 \srcpath{src/optimal\_power\_flow/parity\_formulation.cpp}（2831 行）
为准。

\subsection{建模层定位与求解路径}

\subsubsection*{全空间 NLP 声明}
Parity 建模层装配的 NLP 形如（
\srcpath{include/hacdcpf/optimal\_power\_flow/formulation.hpp}:8--14）：
\begin{align*}
\min_{x}\; f(x)\quad
\text{s.t.}\quad g(x)=0,\quad h(x)\le 0,\quad
x_{\min}\le x\le x_{\max}
\end{align*}
与 Native 精简空间路径（第~\ref{sec:acopfnative}~节）不同，这里所有交流
母线相角、电压幅值、发电机 $P/Q$、换流器三端口功率、可控资源与可选切负荷
变量都保留为显式原变量，约束全部写成显式行块；头文件注释明确说明这一
布局"刻意便于审查"：\fld{VarIndex} 给出每个变量块、\fld{ConstraintIndex}
给出每个等式/不等式行块的起止（同文件:17--21）。变量界
$x_{\min}\le x\le x_{\max}$ 以盒式约束单独处理，不计入 $h(x)$。

\subsubsection*{命名与对齐对象（如实考证）}
"Parity"（对齐）一名在当前代码注释中未再展开定义，现存表述只有
"Parity full-space OPF formulation"（同文件:5）与"与 Native 精简空间路径
相对照的全空间路径"（\srcpath{docs/technical\_notebook/\allowbreak
sections/11c\_native\_ipm\_vs\_parity\_ipm.tex}）。git 历史显示该文件最初
由团队内部 Julia 原型逐方程移植而来：初始提交（\texttt{de00140}，
2026-05-12）中 \srcpath{clamp\_interior} 与直流预热启动处留有
"matching Julia"、"Julia :clamp\_1pct" 等注释；当前检出处这些注释已删除，
但 \srcpath{src/optimal\_power\_flow/parity\_ipm.cpp}:1184 仍留有
"the reference formulation stores dg as n x neq" 的对照痕迹。因此本章对
"对齐对象"的表述以可核实材料为限：Parity 建模层是\textbf{与参考实现
（内部 Julia 原型）逐方程对齐的全空间 AC OPF 建模}，其当前工程含义是
统一交直流的显式全空间 NLP；它与潮流方程的口径一致性由
第~\ref{subsec:parity-pf-consistency}~小节逐条说明。同一份
\fld{parity::Problem} 同时服务两个求解后端：自研 Parity IPM 与嵌入式
Ipopt（"applied to the same parity formulation"，
\srcpath{include/hacdcpf/optimal\_power\_flow/opf\_options.hpp}:19--25
注释）。

\subsubsection*{与潮流共享 canonical 数据}
\begin{sloppypar}
装配入口 \srcpath{build\_problem} 的第一步调用
\fld{core::make\_solver\_data(sys,\allowbreak}\allowbreak\ \fld{LossModelType::Linear)}，
完成富系统到 canonical \fld{SolverData} 的投影（
\srcpath{src/optimal\_power\_flow/parity\_formulation.cpp}:109--112；
投影与装配机制见第~\ref{sec:overview}~节）。因此 Parity NLP 的
$Y_{\mathrm{bus}}$、直流电导矩阵 $G^{dc}$、基准容量 \fld{base\_mva} 与
混合潮流完全同源；\fld{LossModelType} 仅作为字段存入 \fld{SolverData}
（\srcpath{src/power\_flow/solver\_data.cpp}:324），Parity 换流器耦合行
不读取该字段，而是自行按 $a+bI+cI^2$ 计损耗（见后文换流器耦合约束）。
随后 \srcpath{build\_problem} 依次完成：逐母线负荷/ZIP 系数预取
（:117--161）、各组件族的变量映射筛选（:163--327）、交流岛相角参考与
DC\_V 电压参考选择（:329--395）、变量块与约束行块计数及偏移
（:397--528）、行缩放因子（:530--537）、$Y_{\mathrm{bus}}$ 对角元与
$G^{dc}$ 稠密副本缓存（:540--549）、不可观测直流母线锚点检测
（:551--561）、固定注入汇总（:563--625）。
\end{sloppypar}

\subsubsection*{装配入口与数据流}
建模层对外暴露的函数链为 \srcpath{build\_problem} $\to$
\srcpath{build\_variable\_bounds} $\to$ \srcpath{build\_initial\_point}
$\to$ 七个 KKT 求值函数（声明见
\srcpath{include/hacdcpf/optimal\_power\_flow/formulation.hpp}:234--270、
287--311、354--367、391--402、419--445）。数据流如图~\ref{fig:parity-flow}
所示。

\begin{figure}[H]
\centering
\resizebox{\textwidth}{!}{%
\begin{tikzpicture}[
  font=\small,
  box/.style={draw, rectangle, minimum height=0.9cm, inner sep=4pt, align=center},
  arr/.style={-{Stealth[length=2.5mm]}, thick}]
  \node[box] (rich) {富语义系统\\\fld{HybridPowerSystem}};
  \node[box, right=0.9cm of rich] (sd) {canonical 投影\\\srcpath{make\_solver\_data}};
  \node[box, right=0.9cm of sd] (prob) {\fld{parity::Problem}\\\srcpath{build\_problem}};
  \node[box, right=0.9cm of prob] (kkt) {KKT 求值族\\objective / $g$ / $h$ / $J_g$ / $J_h$ / $\nabla^2 L$};
  \node[box, right=0.9cm of kkt] (ipm) {求解核\\Parity IPM / Ipopt};
  \draw[arr] (rich) -- (sd);
  \draw[arr] (sd) -- (prob);
  \draw[arr] (prob) -- node[above, font=\footnotesize]{bounds, $x_0$} (kkt);
  \draw[arr] (kkt) -- (ipm);
\end{tikzpicture}}
\caption{Parity 建模层数据流：富系统经 canonical 投影装配为
\fld{parity::Problem}，向 Parity IPM（第~\ref{sec:parityipm}~节）或
嵌入式 Ipopt 提供统一 KKT 求值接口}
\label{fig:parity-flow}
\end{figure}

\subsubsection*{运行期选项}
\fld{ParityOptions} 控制切负荷、目标类型与四个约束族开关（定义于
\srcpath{include/hacdcpf/optimal\_power\_flow/formulation.hpp}:35--49）。
约束族开关默认全部打开（always-enforce），调用方（如 GUI）可逐族关闭以
松弛问题；关闭只是不生成对应行块或把容量圆半径放到 $10^{12}$，不改变
其余方程。

\begin{paramtable}
\fld{load\_shedding} & — & — & true/false & true & 是否为每交流母线分配切负荷变量 $\Delta P^d,\Delta Q^d$（false 时两块长度为 0）\\
\fld{voll} & VOLL & — & $>0$ & 0.0 & 失负荷价值；0 表示按最大发电边际成本自动取值（见目标函数节）\\
\fld{eps\_iac} & $\varepsilon_I$ & pu & $10^{-12}$--$10^{-4}$（经验值） & 1e-6 & 换流器电流 $|I_{ac}|=\sqrt{P^2+Q^2+\varepsilon_I}/V_m$ 的平滑项\\
\fld{objective} & — & — & Economic/\allowbreak VoltageDeviation/\allowbreak ActiveLoss/\allowbreak VoltageDeviation\allowbreak AndLoss & Economic & 连续目标类型（\fld{ACOPFObjective} 枚举，\srcpath{opf\_options.hpp}:39--44）；后三者服务 RPO（第~\ref{sec:rpo}~节）\\
\fld{voltage\_target\_pu} & $V^{\mathrm{ref}}$ & pu & 0.95--1.05（经验值） & 1.0 & 电压偏差目标的参考电压\\
\fld{voltage\_deviation\_weight} & $w_V$ & — & $>0$ & 1.0 & 电压偏差项权重\\
\fld{active\_loss\_weight} & $w_L$ & — & $>0$ & 1.0 & 有功网损项权重\\
\fld{enforce\_branch\_limits} & — & — & true/false & true & 是否生成 AC/DC 支路容量不等式行\\
\fld{enforce\_converter\_capacity} & — & — & true/false & true & 是否按容量圆约束换流器视在功率（false 时半径取 $10^{12}$）\\
\fld{enforce\_converter\_current\_limits} & — & — & true/false & true & 是否生成换流器交流电流不等式行\\
\fld{enforce\_converter\_modulation\_limits} & — & — & true/false & true & 是否生成 VSC 调制窗口与 DC/DC 占空比线性不等式行（同一开关门控两族）\\
\end{paramtable}

\subsection{决策变量向量布局}\label{subsec:parity-varlayout}

\subsubsection*{块序与偏移}
原变量向量按 17 个连续块排列（文档注释见
\srcpath{include/hacdcpf/optimal\_power\_flow/formulation.hpp}:51--65；
块长赋值见 \srcpath{src/optimal\_power\_flow/parity\_formulation.cpp}:397--423，
起始偏移累加见同文件:425--464）：
\begin{align*}
x = [\,&\theta,\; V,\; P^g,\; Q^g,\; V^{dc},\; P^{ac}_c,\; Q^{ac}_c,\; P^{dc}_c,\;
\Delta P^d,\; \Delta Q^d,\; P^{ren},\; Q^{ren},\\
&P^{stor},\; Q^{stor},\; P^{stor,dc},\; P^{dcdc},\; P^{flex},\; P^{er},\; Q^{er}\,]
\end{align*}
能量路由器端口块 $P^{er},Q^{er}$ 在头文件公式注释之后补入
\fld{VarIndex}（同文件:73--77、86：\fld{n\_erp}/\fld{n\_erq} 与
\fld{i\_erp}/\fld{i\_erq}），实现中位于 $P^{flex}$ 之后
（parity\_formulation.cpp:460--463）。空组件族块长为 0，但仍保留等于
下一块起点的有效偏移（formulation.hpp:63--65）。各块长度、含义与
组件筛选口径如下表（筛选循环见 parity\_formulation.cpp:163--327）：

\begin{center}
\footnotesize
\begin{tabular}{@{}l l L{4.6cm} L{6.2cm}@{}}
\toprule
块符号 & 长度字段 & 块内容（单位均为标幺） & 组件筛选口径\\
\midrule
$\theta$ & \fld{n\_va}$=n_b$ & 交流母线电压相角（rad） & 全部交流母线\\
$V$ & \fld{n\_vm}$=n_b$ & 交流母线电压幅值 & 全部交流母线\\
$P^g,Q^g$ & \fld{n\_pg} & 发电机有功/无功 & 全部母线有效的发电机（含退运，退运机由界钉定）\\
$V^{dc}$ & \fld{n\_vdc}$=n_{dc}$ & 直流母线电压 & 全部直流母线\\
$P^{ac}_c,Q^{ac}_c,P^{dc}_c$ & \fld{n\_pac} & VSC 交流有功/无功、直流有功 & 在役且两端母线有效的换流器\\
$\Delta P^d,\Delta Q^d$ & \fld{n\_dpd} & 切负荷量 & \fld{load\_shedding}=true 时每交流母线一个，否则 0\\
$P^{ren},Q^{ren}$ & \fld{n\_pren} & 可控可再生有功/无功 & 可削减 \fld{RenewableGen}（\fld{curtailable}）与可控 \fld{PVSystem}（\fld{controllable}），以 \fld{ren\_source}=0/1 区分\\
$P^{stor},Q^{stor}$ & \fld{n\_pstor} & 交流储能有功/无功 & 在役 \fld{StorageUnit}\\
$P^{stor,dc}$ & \fld{n\_pstordc} & 直流储能有功 & 在役 \fld{DCStorage}\\
$P^{dcdc}$ & \fld{n\_pdcdc} & DC/DC 输入侧有功 $p_{in}$ & 在役且两端口母线有效的 DC/DC\\
$P^{flex}$ & \fld{n\_pflex} & 柔性负荷实际需求 & 在役且可控 \fld{FlexibleLoad}\\
$P^{er},Q^{er}$ & \fld{n\_erp},\fld{n\_erq} & 能量路由器端口有功/无功 & 在役路由器的在役端口；仅交流端口有无功变量\\
\bottomrule
\end{tabular}
\end{center}

\subsubsection*{索引映射}
\fld{Problem} 保存从紧凑变量位置回到 canonical 组件数组的全部映射
（\srcpath{include/hacdcpf/optimal\_power\_flow/formulation.hpp}:131--192）：
\fld{gen\_var\_to\_data}\allowbreak/\fld{gen\_bus}、
\fld{conv\_var\_to\_data}\allowbreak/\fld{conv\_ac\_bus}\allowbreak/\fld{conv\_dc\_bus}、
\fld{ren\_var\_to\_data}/\fld{ren\_source}/\fld{ren\_bus}、
\fld{stor\_var\_to\_data}、\fld{stor\_dc\_var\_to\_data}、
\fld{dcdc\_var\_to\_data}/\fld{dcdc\_bus\_in}/\fld{dcdc\_bus\_out}、
\fld{flex\_var\_to\_data}/\fld{flex\_bus}、
\fld{er\_ports}（\fld{ERPortInfo}：路由器号、端口号、母线、交/直标志、
块内列偏移，同文件:181--189）以及受限组件清单
\fld{branch\_limited}、\fld{dc\_branch\_limited}、
\fld{conv\_iac\_limited}、\fld{conv\_mmax\_limited}、
\fld{conv\_mmin\_limited}、\fld{dcdc\_duty\_limited}。结果归因回原始组件
由求解驱动层使用这些映射完成（第~\ref{sec:parityipm}~节）。

\subsubsection*{固定注入与负荷预取}
不进入决策向量的注入在装配期汇总为逐母线常量
\fld{p\_fixed\_inj}/\fld{q\_fixed\_inj}（
\srcpath{src/optimal\_power\_flow/parity\_formulation.cpp}:563--625）：
\fld{StaticGenerator}（恒固定，:568--574）、不可削减
\fld{RenewableGen}（:577--585）、不可控 \fld{PVSystem}（有功经
\srcpath{compute\_pv\_power\_mw} 曲线折算，:587--597）、\fld{VPP}
（:600--606）、并网 \fld{Microgrid} 交换功率（:609--615）与就地
\fld{MobileStorage}（:618--625）。交流负荷需求优先取
\fld{SolverData::pd\_pu}/\fld{qd\_pu}（分量负荷存在时），否则回退母线
\fld{pd\_mw}/\fld{qd\_mvar} 字段（:117--132）；ZIP 权重按母线加权限值
或全局权重填入 \fld{zip\_pp}$\sim$\fld{zip\_zq} 六个向量，缺省为恒功率
（:134--161）。

\subsection{目标函数组装}

\subsubsection*{经济目标（默认）}
经济目标为发电机二次成本、VSC 无功设定点正则、切负荷惩罚与可再生削减
成本之和（\srcpath{src/optimal\_power\_flow/parity\_formulation.cpp}:1204--1306，
函数 \srcpath{objective}）：
\begin{align*}
f(x)=\;&\sum_{g}\Big(c_{2g}\,(S_{\mathrm{base}}\,p_g)^2
       +c_{1g}\,(S_{\mathrm{base}}\,p_g)+c_{0g}\Big)
 +k_{Q}\sum_{c}\Big(S_{\mathrm{base}}\,q^{ac}_c-q^{set}_c\Big)^2\\
 &+\mathrm{VOLL}\sum_{i}S_{\mathrm{base}}\,(\Delta p^d_i+\Delta q^d_i)
 +\sum_{r}\Big(c_{1r}\,S_{\mathrm{base}}\,p^{ren}_r
 +c^{curt}_r\,(p^{rated}_r-S_{\mathrm{base}}\,p^{ren}_r)^+\Big)
 +C_{\mathrm{const}}
\end{align*}
其中 $p_g$、$q^{ac}_c$、$p^{ren}_r$ 为标幺变量，成本系数
\fld{cost\_c2}/\fld{cost\_c1}/\fld{cost\_c0} 以工程单位（MW、MVAr）存储，
装配时乘 $S_{\mathrm{base}}$=\fld{base\_mva} 化为有名值再代入
（:1253--1257；头文件说明见 formulation.hpp:272--286）。币种未在代码中
固定，目标值量纲随成本系数币种而定。

各分项出处：发电机二次成本 :1253--1257；VSC 无功设定点正则
$k_{Q}=10^{-3}$（常量 \srcpath{kReactiveSetpointRegularization}，:29；
装配 :1258--1265），其动机按注释是消除无功无直接成本造成的"平方向"，
防止不受约束的 VSC 吸收任意无功（:24--28 注释）；切负荷惩罚 :1283--1290；
可再生削减成本仅对 \fld{RenewableGen} 源生效
（\fld{cost\_curtail\_mwh}$>0$ 且存在削减量时，:1291--1299），光伏源只有
线性 $c_1$ 项（:1300--1303）。$C_{\mathrm{const}}$ 为固定注入的线性成本
常数（在役 \fld{StaticGenerator}、直流静态机、不可削减可再生、不可控
光伏、直流光伏，:1266--1282）——它是常量，不影响极小点，只影响报告
目标值。

\subsubsection*{标幺换算与梯度/对角 Hessian}
IPM 需要对标幺变量的导数。对发电机成本
$c_2(S_{\mathrm{base}}p_g)^2+c_1(S_{\mathrm{base}}p_g)+c_0$（
\srcpath{include/hacdcpf/optimal\_power\_flow/formulation.hpp}:289--306）：
\begin{align*}
\frac{\partial f}{\partial p_g}
 =2c_{2g}S_{\mathrm{base}}^2\,p_g+c_{1g}S_{\mathrm{base}},\qquad
\frac{\partial^2 f}{\partial p_g^2}=2c_{2g}S_{\mathrm{base}}^2
\end{align*}
（实现于 \srcpath{src/optimal\_power\_flow/parity\_formulation.cpp}:1356--1363，
函数 \srcpath{objective\_gradient\_hessian\_diag}）。同理 VSC 正则项
梯度 $2k_{Q}S_{\mathrm{base}}(S_{\mathrm{base}}q^{ac}_c-q^{set}_c)$、
对角曲率 $2k_{Q}S_{\mathrm{base}}^2$（:1364--1374）；切负荷梯度
$\mathrm{VOLL}\cdot S_{\mathrm{base}}$（:1375--1383）；可再生梯度
$(c_{1r}-\max(0,c^{curt}_r))S_{\mathrm{base}}$ 或 $c_1 S_{\mathrm{base}}$
（:1384--1394）。该函数只填可分目标项，其余分量梯度为零
（formulation.hpp:291--292 注释）。

\subsubsection*{VOLL 自动取值}
\fld{voll}$=0$ 时按下式自动取值（
\srcpath{src/optimal\_power\_flow/parity\_formulation.cpp}:65--84，函数
\srcpath{resolve\_voll\_auto}）：
\begin{align*}
\mathrm{VOLL}=\max\Big\{100,\;\max_{l}\,\fld{cost\_mw}_l,\;
1.1\max_{g}\big|c_{1g}+2c_{2g}\max(P^{\max}_g,P^{\min}_g)\big|\Big\}
\end{align*}
即不低于 100、不低于任何负荷申报的 \fld{cost\_mw}、并高于最高发电边际
成本斜率的 1.1 倍，保证切负荷恒劣于发电调度。结果存入
\fld{voll\_effective}（:538）。

\subsubsection*{RPO 目标族}
\fld{objective} 非 Economic 时（服务无功优化，见第~\ref{sec:rpo}~节），
目标换为电压偏差与/或可变部分有功网损（:1204--1251）：
\begin{align*}
f_{V}&=w_V\sum_{i}\big(V_i-V^{\mathrm{ref}}\big)^2\\
f_{L}&=w_L\,S_{\mathrm{base}}\Big(\sum p_g+\sum p^{ren}+\sum p^{stor}
+\sum p^{stor,dc}-\sum p^{flex}+\sum \Delta p^d\Big)
\end{align*}
（电压偏差 :1215--1220；网损账本 :1221--1237）。注释明确：固定注入与
需求为常量、不影响极小点，RPO 报告口径的网损在目标之外补回常数
（:1200--1203 注释）。非经济目标还叠加 $10^{-6}$ 量级的经济平局项以
选定唯一调度并改善 KKT 条件数（常量 \fld{kRpoEconomicTieBreak}，
:1207、1238--1251），随后提前返回，不再计入 VOLL 与削减成本。

\subsection{等式约束}

\subsubsection*{行块布局}
等式残差 $g(x)$ 按 8 个行块排列（文档注释见
\srcpath{include/hacdcpf/optimal\_power\_flow/formulation.hpp}\allowbreak :89--98；
计数与偏移见 \srcpath{src/optimal\_power\_flow/parity\_formulation.cpp}\allowbreak :467--488）：
\begin{align*}
g(x)=\big[\,g^P_{ac},\; g^Q_{ac},\; g^P_{dc},\; g^{vsc},\;
g^{dcdc},\; g^{er},\; g^{\theta\text{-ref}},\; g^{vdc\text{-ref}}\,\big]
\end{align*}
其中 DC/DC 平衡块 \fld{n\_dcdc\_bal} 恒为 0：$p_{in}$ 直接折入两个直流
节点平衡行，注释说明"分配了行块却不填充会留下空行使 KKT 结构奇异"
（:472--475）。残差求值入口为 \srcpath{equality\_constraints}
（:1401--1666），约定 AC 行用"网络注入加需求减可控供给"符号、DC 行用
"负荷加电导潮流减注入"符号（:1397--1400 注释）。

\subsubsection*{交流节点功率平衡}
交流网络注入取极坐标潮流方程（初始化对角项 :1433--1438，非对角遍历
$Y_{\mathrm{bus}}$ 非零元 :1439--1455；公式注释见
formulation.hpp:328--336）：
\begin{align*}
P^{net}_i&=V_i\sum_{j}V_j\big(G_{ij}\cos\theta_{ij}+B_{ij}\sin\theta_{ij}\big)\\
Q^{net}_i&=V_i\sum_{j}V_j\big(G_{ij}\sin\theta_{ij}-B_{ij}\cos\theta_{ij}\big)
\end{align*}
负荷按 ZIP 电压依赖模型取值，与潮流求解器口径一致（:1496--1498）：
\begin{align*}
P^d_i(V_i)&=P^d_i\,\big(w^{P}_0+w^{P}_1 V_i+w^{P}_2 V_i^2\big)\\
Q^d_i(V_i)&=Q^d_i\,\big(w^{Q}_0+w^{Q}_1 V_i+w^{Q}_2 V_i^2\big)
\end{align*}
平衡行（:1492--1507；符号约定注释见 formulation.hpp:313--325）：
\begin{align*}
g^P_i&=s_P\Big(P^{net}_i-P^g_i-P^{fix}_i-P^{ren}_i-P^{stor}_i
      +P^d_i(V_i)-\Delta P^d_i+P^{flex}_i-P^{ac}_i\Big)\\
g^Q_i&=s_Q\Big(Q^{net}_i-Q^g_i-Q^{fix}_i-Q^{ren}_i-Q^{stor}_i
      +Q^d_i(V_i)-\Delta Q^d_i-Q^{ac}_i\Big)
\end{align*}
式中 $P^g_i,P^{ren}_i,P^{stor}_i,P^{ac}_i$ 等为按母线累加的分量注入
（:1457--1488）；$s_P=1/\max(1,\max_i|P^d_i|)$、$s_Q$ 同理，为逐行常数
缩放（:530--537、1505--1506），用以平衡残差量级；DC 行与换流器行
不缩放。

\subsubsection*{直流节点功率平衡}
直流侧按电导矩阵取潮流，与潮流定义逐字一致（注释
"Match the PF definition exactly"，:1509--1512）：
\begin{align*}
g^{P,dc}_k=P^{d,dc}_k+V^{dc}_k\sum_{m}G^{dc}_{km}V^{dc}_m-P^{dc,conv}_k
\end{align*}
（:1509--1524）。$G^{dc}$ 对角为正、非对角为负，故高压发送端注入为正；
$P^{d,dc}_k$ 在无分量 \fld{DCLoad} 时取母线 \fld{pd\_mw}$/S_{\mathrm{base}}$，
否则逐负荷以 \srcpath{effective\_load\_p\_mw}（含 \fld{scaling}）累加
（:1519--1522、1528--1537）。直流静态机、直流光伏、直流储能分别以
固定负注入与变量注入进入（:1538--1554）。

\subsubsection*{VSC 换流器耦合约束}
每台 VSC 贡献一行能量平衡（:1604--1618；公式注释见
formulation.hpp:347--352）：
\begin{align*}
g^{vsc}_c=P^{ac}_c+P^{dc}_c+a_c+b_c I^{ac}_c+c_c\,(I^{ac}_c)^2=0,\qquad
I^{ac}_c=\frac{\sqrt{(P^{ac}_c)^2+(Q^{ac}_c)^2+\varepsilon_I}}{V_{m,ac}}
\end{align*}
其中 $a_c=$\fld{loss\_mw}$/S_{\mathrm{base}}$（固定损耗）、
$b_c=$\fld{loss\_percent}$/100$（比例损耗）、$c_c=1-\eta$（导通损耗），
$\varepsilon_I$=\fld{eps\_iac} 为 $|I_{ac}|$ 在原点处的可微平滑；
$V_{m,ac}$ 取下限 $10^{-8}$（:1608）。两侧功率均为"注入为正"约定，
故无损耗时 $P^{dc}_c=-P^{ac}_c$。

\subsubsection*{DC/DC 变换器折叠}
DC/DC 不设独立平衡行，输入侧变量 $p_{in}$ 直接折入两个直流平衡行
（:1578--1602，注释 :1620--1626）：输入母线行加 $p_{in}$（抽取），输出
母线行减 $p_{out}$（注入），即
\begin{align*}
g^{P,dc}_{in}&\;\leftarrow\;g^{P,dc}_{in}+p_{in},
&g^{P,dc}_{out}&\;\leftarrow\;g^{P,dc}_{out}-p_{out}\\
p_{out}&=\begin{cases}p_{in}\,\eta,&p_{in}\ge 0\\ p_{in}/\eta,&p_{in}<0\end{cases}
&p_{out}&\;\leftarrow\;p_{out}-r_{eq}\,\frac{p_{in}^2}{\max(V^{dc}_{in},\,0.1)^2}
\end{align*}
正方向约定为 $p_{in}>0$ 表示正向传输（$\mathrm{bus}_{in}\to\mathrm{bus}_{out}$）；效率
$\eta$ 被钳到 $[0.01,1]$；$r_{eq}>0$ 时叠加 $I^2R$ 导通损耗
（:1594--1600）。注意 OPF 变量是\textbf{输入侧}口径，而潮流
\fld{p\_ref\_mw} 是输出侧口径，二者经 $\eta$ 换算（注释 :1579--1583；
预热启动换算 :1070--1078）。

\subsubsection*{LCC 换流站注入}
在役 LCC 以"定阶、状态依赖"注入进入交直流平衡：不产生 OPF 决策变量，
直接在当前 $(V_m,V_{dc})$ 迭代点上调用与混合潮流同一套准稳态特性
\srcpath{lcc\_ac\_injection}/\srcpath{lcc\_dc\_injection} 求值并从平衡行
减去（:1556--1576；换相母线由 \srcpath{lcc\_commutation\_ac\_bus} 判定，
:1562--1563）。换流变分接头因此在 Parity NLP 中不可优化。

\subsubsection*{能量路由器}
每个在役端口的有功（交流端口另有无功）为决策变量，按注入为正从母线
平衡行减去（交流行与 DC 行的缩放口径分别对齐，:1628--1644）；每台在役
路由器再配一行有功守恒（:1645--1654）：
\begin{align*}
g^{er}_r=\sum_{p\in r}P^{er}_p=0
\end{align*}
即设备只"路由"不"产生"功率。当前为\textbf{无损}守恒，注释声明
\fld{loss\_percent} 精确化留作后续（"refinement deferred"，:1651--1652）。

\subsubsection*{参考行}
每个导电交流岛锚定一个相角参考（DFS 划岛，优先 \fld{BusType::SLACK}
母线，其次在役 \fld{is\_slack} 发电机所在母线，否则取岛首母线，
:329--385），参考行 $g^{\theta\text{-ref}}_r=\theta_{\mathrm{bus}}-\theta^{set}$
（:1655--1660）；每个在役 \fld{DCBusType::DC\_V} 直流母线锚定电压
$g^{vdc\text{-ref}}_r=V^{dc}_{\mathrm{bus}}-V^{set}_{dc}$（:387--395、1661--1665）。
交流岛间无相角耦合，必须逐岛接地（:329--331 注释）。

\subsection{变量界（盒式约束）}

$x_{\min}\le x\le x_{\max}$ 由 \srcpath{build\_variable\_bounds} 装配
（\srcpath{src/optimal\_power\_flow/parity\_formulation.cpp}:630--815；
声明注释见 formulation.hpp:236--247）。初始为 $\pm\infty$，逐块覆盖：

\begin{center}
\footnotesize
\begin{tabular}{lll}
\toprule
变量块 & 界（标幺，MW 系量除 $S_{\mathrm{base}}$） & 出处\\
\midrule
$\theta_i$ & $[-\pi,\;\pi]$ & :635--638\\
$V_i$ & $[$\fld{vmin\_pu}$,$\fld{vmax\_pu}$]$ & :639--643\\
$P^g_k$ & $[$\fld{pmin\_mw}$, $\fld{pmax\_mw}$]/S_{\mathrm{base}}$；退运机钉定 \fld{pg\_mw}$\pm10^{-8}$ & :645--660\\
$Q^g_k$ & $[$\fld{qmin\_mvar}$, $\fld{qmax\_mvar}$]/S_{\mathrm{base}}$；退运同上 & :651--660\\
$V^{dc}_k$ & $[$\fld{vmin\_pu}$,$\fld{vmax\_pu}$]$，非法时回退 $[0.8,1.2]$ & :663--673\\
$P^{ac}_c,Q^{ac}_c$ & \fld{pmin}/\fld{pmax}、\fld{qmin}/\fld{qmax} 换算 & :690--695\\
$P^{dc}_c$ & $\pm S^{\max}_c$（容量圆半径，见下节） & :696--698\\
$\Delta P^d_i,\Delta Q^d_i$ & $[0,\max(P^d_i,\,0.01/S_{\mathrm{base}})]$（无功取绝对值） & :701--708\\
$P^{ren}_k$ & $[0,\min(P^{rated},\,P^{avail})]$，$P^{avail}$ 取日程 \fld{p\_mw} & :712--732\\
$Q^{ren}_k$ & \fld{qmin}/\fld{qmax} 换算 & :722--723、730--731\\
$P^{stor}_k$ & 钉定于日程 \fld{p\_mw}$\pm 10^{-4}/S_{\mathrm{base}}$ 并夹到 \fld{pmin}/\fld{pmax} & :740--758\\
$Q^{stor}_k$ & \fld{qmin}/\fld{qmax} 换算 & :756--757\\
$P^{stor,dc}_k$ & 同交流储能钉定口径 & :760--778\\
$P^{dcdc}_k$ & $[$\fld{pmin\_mw}$,$\fld{pmax\_mw}$]$；\fld{pmax} 缺省回退 \fld{sn\_mva}，\fld{pmin} 回退 $-p_{\max}$ & :780--788\\
$P^{flex}_k$ & $[$\fld{p\_mw}$-$\fld{flex\_down\_mw}$, $\fld{p\_mw}$+$\fld{flex\_up\_mw}$]$ & :790--795\\
$P^{er},Q^{er}$ & 端口 \fld{pmin}/\fld{pmax}、\fld{qmin}/\fld{qmax}，非法回退 $\pm10^{6}$ & :797--814\\
\bottomrule
\end{tabular}
\end{center}

三点需要说明。其一，\textbf{发电机无功仅为盒式界}：当前实现不含无功
能力曲线（$P$--$Q$ 能力图）约束，$Q^g$ 上下界直接取自
\fld{qmin\_mvar}/\fld{qmax\_mvar}（:651--652）。其二，储能钉定源于单
时段口径：逐时段 OPF 无跨时段 SOC 状态转移，放开储能会每个时段都以最
大功率放电、无视 UC 日程，故以 $\pm 10^{-4}/S_{\mathrm{base}}$ 容差带
钉定（:735--739 注释）；\fld{cap\_charging\_strategy}=\texttt{"static"}
的承载力静态储能完全固定于 $-$\fld{cap\_static\_charging\_mw}
（:743--752、767--772）。其三，\textbf{不可观测直流母线}（无任何导电
直流支路相连，:551--561 检测）通过把盒式界收紧到 \fld{vm\_pu} 的
$\pm\varepsilon$ 邻域钉定（:678--688），注释明确禁止改用"电压偏差项
并入有功平衡"的做法（那会等效为虚构源，:674--677）；头文件
formulation.hpp:173--176 关于"软并联锚定项"的表述与当前实现不符，
以紧界实现为准。

\subsection{非线性不等式约束}

\subsubsection*{行块布局与门控}
不等式统一采用 $h(x)\le 0$ 约定，行序为：AC 支路 from 端、AC 支路 to
端、VSC 视在功率、DC 支路功率、VSC 交流电流、调制上限、调制下限、
DC/DC 占空比（文档注释见
\srcpath{include/hacdcpf/optimal\_power\_flow/formulation.hpp}:100--103；
计数见 \srcpath{src/optimal\_power\_flow/parity\_formulation.cpp}:489--527；
求值见 \srcpath{nonlinear\_inequality\_constraints}，:1958--2059）。
各族的生成门控：AC/DC 支路要求 \fld{enforce\_branch\_limits} 且在役、
\fld{rate\_a\_mva}$>0$（:197--217）；VSC 电流要求
\fld{enforce\_converter\_current\_limits} 且 \fld{i\_ac\_max\_pu}$>0$
（:499--501）；调制窗口要求
\fld{enforce\_converter\_modulation\_limits} 且
\fld{k\_m\_modulation}、\fld{vn\_ac\_kv}、\fld{vn\_dc\_kv} 均正，
\fld{m\_max}/\fld{m\_min}$>0$ 分别生成上/下行（:503--506）；DC/DC 占空
比要求非 Generic 拓扑且 $0\le d_{\min}<d_{\max}\le 1$（:511--525）。
未声明限值的组件不产生行（"unconstrained converters add no rows"，
formulation.hpp:115--117）。

\subsubsection*{AC 支路两端容量}
对每条受限支路，以 $\pi$ 等值计算两端潮流（:1970--1995）。记串联导纳
$y_s=1/(r+jx)=g_s+jb_s$、半充电 $b_c=b/2$、固定变比 $t$ 与固定移相角
$\varphi$（\fld{tap}、\fld{shift\_deg} 为\textbf{参数}而非变量，$|t|<10^{-12}$
时回退 $t=1$，:1974--1980），$\theta=\theta_i-\theta_j-\varphi$：
\begin{align*}
P^{f}&=\frac{g_s}{t^2}V_i^2-\frac{g_s\cos\theta+b_s\sin\theta}{t}V_iV_j
&Q^{f}&=-\Big(\frac{b_s}{t^2}+b_c\Big)V_i^2-\frac{g_s\sin\theta-b_s\cos\theta}{t}V_iV_j\\
P^{t}&=g_sV_j^2-\frac{g_s\cos\theta-b_s\sin\theta}{t}V_iV_j
&Q^{t}&=-(b_s+b_c)V_j^2+\frac{g_s\sin\theta+b_s\cos\theta}{t}V_iV_j
\end{align*}
容量约束取视在功率平方口径（$S_{\max}=$\fld{rate\_a\_mva}$/S_{\mathrm{base}}$，
:1985--1986）：
\begin{align*}
h^{Sf}&=(P^{f})^2+(Q^{f})^2-S_{\max}^2\le 0
&h^{St}&=(P^{t})^2+(Q^{t})^2-S_{\max}^2\le 0
\end{align*}
即 from/to 两端各一行、共用同一 $S_{\max}$（\fld{n\_sf}$=$\fld{n\_st}，
:489--490）。变压器分接头与移相器只以固定 $t$、$\varphi$ 进入本式与
$Y_{\mathrm{bus}}$，Parity NLP 不优化档位；离散档位优化属于 RPO 层
（第~\ref{sec:rpo}~节）。

\subsubsection*{VSC 容量圆与交流电流}
容量圆（:1996--2001）：
\begin{align*}
h^{Sconv}_c=(P^{ac}_c)^2+(Q^{ac}_c)^2-(S^{\max}_c)^2\le 0
\end{align*}
半径 $S^{\max}_c=$\fld{p\_rated\_mw}$/S_{\mathrm{base}}$；额定缺失或非正
时回退 $\max(|P_{\max}|,|P_{\min}|,|Q_{\max}|,|Q_{\min}|,10^{-6})/
S_{\mathrm{base}}$；\fld{enforce\_converter\_capacity}=false 时半径取
$10^{12}$ 使约束恒松弛（函数 \srcpath{converter\_smax\_pu}，:86--102）。交流电流限
（:2015--2024，"multi-converter model §3.1.5"）：
\begin{align*}
h^{Iac}_c=(P^{ac}_c)^2+(Q^{ac}_c)^2-I^{2}_{ac,\max}\,V_{m,ac}^2\le 0
\end{align*}
即 $I_{ac}=|S_{ac}|/V_m\le$\fld{i\_ac\_max\_pu} 的平方等价形（系统基准
标幺）。

\subsubsection*{DC 支路功率限}
\begin{align*}
h^{Sdc}_{km}=\Big[g_{km}V^{dc}_k\big(V^{dc}_k-V^{dc}_m\big)\Big]^2-P_{\max}^2\le 0
\end{align*}
其中 $g_{km}=1/$\fld{r\_pu}（$|r|<10^{-14}$ 时取 0），
$P_{\max}=$\fld{rate\_a\_mva}$/S_{\mathrm{base}}$（:2003--2013）。
注意这是\textbf{功率}平方口径，且只按 from 端计一行。

\subsubsection*{VSC 调制窗口（线性不等式）}
调制度 $m=\dfrac{V_{m,ac}\,V_{N,ac}}{K_m\,V_{dc}\,V_{N,dc}}$ 的两个界
关于 $(V_{m,ac},V_{dc})$ 均为线性，以两行放入同一非线性不等式向量，
使 IPM 只面对一个不等式接口（:2025--2042；线性行与 Jacobian 常系数见
:2161--2181）：
\begin{align*}
h^{mmax}_c&=V_{m,ac}\,V_{N,ac}-m_{\max}K_mV_{N,dc}\,V^{dc}\le 0\\
h^{mmin}_c&=m_{\min}K_mV_{N,dc}\,V^{dc}-V_{m,ac}\,V_{N,ac}\le 0
\end{align*}
（$V_{N,ac}$=\fld{vn\_ac\_kv}，$V_{N,dc}$=\fld{vn\_dc\_kv}，
$K_m$=\fld{k\_m\_modulation}；门控 :503--506）。潮流后校验口径的同一
窗口见《元件模型技术手册》VSC 节；此处为优化内嵌形式。

\subsubsection*{DC/DC 占空比窗口（线性不等式）}
非 Generic 拓扑的 DC/DC 由理想 CCM 电压变换律把占空比 $D$ 表为两端口
直流电压的线性关系，$d_{\min}\le D\le d_{\max}$ 化作两行线性不等式
（\srcpath{dcdc\_duty\_rows}，:1931--1952；行生成 :2043--2058）：
\begin{align*}
\text{Buck:}\;&{-d_{\max}}V_{N,in}V_{in}+V_{N,out}V_{out}\le 0,\quad
 d_{\min}V_{N,in}V_{in}-V_{N,out}V_{out}\le 0\\
\text{Boost:}\;&{-V_{N,in}V_{in}}+(1-d_{\max})V_{N,out}V_{out}\le 0,\quad
 V_{N,in}V_{in}-(1-d_{\min})V_{N,out}V_{out}\le 0\\
\text{BuckBoost:}\;&{-d_{\max}V_{N,in}V_{in}}+(1-d_{\max})V_{N,out}V_{out}\le 0,\quad
 d_{\min}V_{N,in}V_{in}-(1-d_{\min})V_{N,out}V_{out}\le 0\\
\text{Isolated:}\;&{-d_{\max}nV_{N,in}V_{in}}+V_{N,out}V_{out}\le 0,\quad
 d_{\min}nV_{N,in}V_{in}-V_{N,out}V_{out}\le 0
\end{align*}
对应电压律依次为 $V_{out}=DV_{in}$、$V_{out}=V_{in}/(1-D)$、
$V_{out}=\tfrac{D}{1-D}V_{in}$、$V_{out}=DnV_{in}$（$n$=\fld{n\_ratio}，
缺省 1）；端口额定 $V_N$ 取 \fld{vn\_in\_kv}/\fld{vn\_out\_kv}，缺省回退
母线 \fld{base\_kv}、再回退 1.0（\srcpath{dcdc\_port\_nominal\_kv}，
:1925--1929）。Generic 拓扑不产生行（:518--523）。

\subsection{初值构造}

\srcpath{build\_initial\_point} 生成严格内点暖启动（
\srcpath{src/optimal\_power\_flow/parity\_formulation.cpp}\allowbreak :985--1198；
声明注释见 formulation.hpp:249--270），分五步：

\begin{enumerate}
\item \textbf{物理初值}：$\theta$ 取母线 \fld{va\_deg} 化弧度、$V$ 取
  \fld{vm\_pu}（:991--996）；发电机取 \fld{pg\_mw}/\fld{qg\_mvar}
  （:997--1001）；当大量发电机 $P_g=0$（MATPOWER CDF 换算算例常见）时，
  把功率缺额按 $P_{\max}$ 比例摊到零值机上，给直流预热启动一个现实的
  分布式调度（"smart dispatch balancing"，:1003--1032）。$V^{dc}$ 取
  直流母线 \fld{vm\_pu}；VSC 取 $P^{ac}=$\fld{p\_set\_mw}、
  $Q^{ac}=$\fld{q\_set\_mvar}、$P^{dc}\approx-P^{ac}$（:1034--1044）；
  DER 各块取当前运行点（:1046--1095）。
\item \textbf{直流暖启动}：在每个交流岛接地后组装降阶 $B'$（取自
  $Y_{\mathrm{bus}}$ 虚部），解 $B'_{\mathrm{red}}\theta=-P_{\mathrm{inj}}$ 得与有功调度
  一致的相角（\srcpath{dc\_warm\_start}，:817--923）。采用稀疏 LU
  （COLAMD 排序）而非稠密 LU——注释记录了 13659 母线规模下稠密暖启动
  曾占全部求解时间 90\% 以上的教训（:867--871、900--912）。
\item \textbf{无功暖启动}：由 $Q^{net}(V,\theta)+Q^d(V)$ 估算母线无功
  需求，按母线上发电机数均分后夹到界内（\srcpath{qg\_warm\_start}，
  :925--983）。
\item \textbf{切负荷暖启动}：令 $\Delta P^d,\Delta Q^d$ 吸收逐母线残差
  使 $g^P=g^Q=0$，再夹到界内（:1103--1161）。
\item \textbf{最小二乘可行性修正}：解正规方程
  $(J_gJ_g^{\mathsf{T}})v=-g$、$\Delta x=J_g^{\mathsf{T}}v$，以 0.8 阻尼
  修正 $x_0$；分解失败时保留未修正点（:1163--1192）。最后全向量经
  \srcpath{clamp\_interior} 钳到界的内部（:1194--1197）。
\end{enumerate}

\srcpath{clamp\_interior}（:37--55）把 $x$ 钳入
$[l+\varepsilon,\,u-\varepsilon]$，$\varepsilon=\min(\max(10^{-9},0.01w),
0.49w)$、$w=u-l$；$l>u$ 时先交换，零宽度区间取中点。这保证 IPM 初始点
不立即违反界松弛（formulation.hpp:249--261 注释）。

\subsection{雅可比与稀疏结构组装}

\subsubsection*{等式雅可比}
\srcpath{equality\_jacobian}（
\srcpath{src/optimal\_power\_flow/parity\_formulation.cpp}:1668--1914）
以三元组累加、\fld{setFromTriplets} 汇总重复项成形
$n_{eq}\times n$ 稀疏矩阵（:1686--1688、1911--1913）。前置条件：工作区
\fld{ws.p\_calc}/\fld{ws.q\_calc} 须已由同一点的
\srcpath{equality\_constraints} 调用填好，本函数不重算（:1677--1678
注释）。

交流块对角元利用恒等式（:1690--1708）：
\begin{align*}
\frac{\partial g^P_i}{\partial\theta_i}&=s_P\big({-Q^{net}_i}-B_{ii}V_i^2\big)
&\frac{\partial g^Q_i}{\partial\theta_i}&=s_Q\big(P^{net}_i-G_{ii}V_i^2\big)\\
\frac{\partial g^P_i}{\partial V_i}&=s_P\Big(\frac{P^{net}_i}{V_i}+G_{ii}V_i
 +P^d_i\big(w^{P}_1+2w^{P}_2V_i\big)\Big)
&\frac{\partial g^Q_i}{\partial V_i}&=s_Q\Big(\frac{Q^{net}_i}{V_i}-B_{ii}V_i
 +Q^d_i\big(w^{Q}_1+2w^{Q}_2V_i\big)\Big)
\end{align*}
非对角元遍历 $Y_{\mathrm{bus}}$（:1710--1735）：
\begin{align*}
\frac{\partial P^{net}_i}{\partial\theta_j}&=V_iV_j\big(G_{ij}\sin\theta_{ij}-B_{ij}\cos\theta_{ij}\big)
&\frac{\partial P^{net}_i}{\partial V_j}&=V_i\big(G_{ij}\cos\theta_{ij}+B_{ij}\sin\theta_{ij}\big)\\
\frac{\partial Q^{net}_i}{\partial\theta_j}&=-V_iV_j\big(G_{ij}\cos\theta_{ij}+B_{ij}\sin\theta_{ij}\big)
&\frac{\partial Q^{net}_i}{\partial V_j}&=V_i\big(G_{ij}\sin\theta_{ij}-B_{ij}\cos\theta_{ij}\big)
\end{align*}
分量列（发电机、VSC、切负荷、DER、能量路由器）为 $-s_P$/$-s_Q$/$-1$
常数元（:1737--1754、1831--1899）；直流块
$\partial g^{P,dc}_k/\partial V^{dc}_m=V^{dc}_kG^{dc}_{km}
+\delta_{km}\sum_lG^{dc}_{kl}V^{dc}_l$（:1756--1765）；VSC 耦合行的损耗
导数（:1807--1829）：
\begin{align*}
\frac{\partial g^{vsc}_c}{\partial P^{ac}_c}=1+(b_c+2c_cI^{ac}_c)\frac{\partial I}{\partial P},\quad
\frac{\partial I}{\partial P}=\frac{P^{ac}_c}{S^{ac}_cV_m},\quad
\frac{\partial I}{\partial V_m}=-\frac{S^{ac}_c}{V_m^2}
\end{align*}
（$S^{ac}_c=\sqrt{(P^{ac}_c)^2+(Q^{ac}_c)^2+\varepsilon_I}$）；DC/DC 折入
导数含 $\eta$ 方向切换与 $I^2R$ 项（:1852--1875）；LCC 行的解析一阶导由
\srcpath{lcc\_ac\_dc\_jacobian} 提供并冠以负注入号（:1772--1805）；
参考行为单位元（:1900--1909）。

\subsubsection*{不等式雅可比}
\srcpath{nonlinear\_inequality\_jacobian}（:2061--2205）按同一行序填
$h(x)$ 的导数：支路容量行取 $2(P\,\partial P+Q\,\partial Q)$ 链式形式，
from/to 两端各 4 个非零元（:2074--2124）；VSC 容量圆 $2P,2Q$
（:2126--2130）；DC 支路 $2P^{flow}\partial P^{flow}$（:2132--2147）；
VSC 电流行 $\partial h/\partial V_m=-2I_{ac,\max}^2V_m$（:2149--2160）；
调制与占空比两族为常系数线性行（:2161--2200）。行序与求值函数严格
一致，使 IPM 对偶 $\mu$ 可按 \fld{ConstraintIndex} 计数解释
（formulation.hpp:395--399 注释）。

\subsubsection*{拉格朗日 Hessian}
\srcpath{lagrangian\_hessian}（:2221--2829）组装
$\nabla^2_{xx}L=\nabla^2f+\sum_i\lambda_i\nabla^2g_i+\sum_j\nu_j\nabla^2h_j$：
对称三plet 累加（同一非零元同时写入 $(r,c)$ 与 $(c,r)$，
:2250--2260），最终 \fld{setFromTriplets} 求和成形 $n\times n$ 对称
稀疏矩阵（:2825--2828）。分块构成：

\begin{itemize}
\item \textbf{目标曲率}：仅对角——发电机成本 $2c_{2g}S_{\mathrm{base}}^2$
  与 VSC 无功正则 $2k_QS_{\mathrm{base}}^2$（:2262--2272）。
\item \textbf{交流平衡曲率}：遍历 $Y_{\mathrm{bus}}$ 非对角元累积
  $\theta$--$\theta$、$V$--$V$、$\theta$--$V$ 交叉项，对角元含
  $2G_{ii}$/$-2B_{ii}$ 与 ZIP 二阶项 $2P^d_iw^{P}_2$、$2Q^d_iw^{Q}_2$；
  每行乘 $\lambda_i$ 并带上行缩放 $s_P$/$s_Q$（:2280--2378）。
\item \textbf{直流平衡曲率}：$\lambda_k\big(2G^{dc}_{kk}\big)$ 对角与
  $\lambda_kG^{dc}_{km}$ 交叉（:2380--2398）。
\item \textbf{LCC 曲率（有限差分）}：LCC 等式曲率不是解析式——对解析
  一阶雅可比做中心差分，步长 $10^{-5}\max(1,|v|)$，并在额定电流切换
  拐点处只取留在当前活跃集内的扰动方向，避免人为 $O(1/h)$ 曲率拖住
  KKT 步（:2400--2533，注释 :2400--2408；活跃集守护
  :2485--2496）。这是建模层中唯一数值微分部分。
\item \textbf{VSC 损耗曲率}：$I^{ac}$ 平滑后的全部二阶导
  （:2535--2586）。
\item \textbf{DC/DC $I^2R$ 曲率}：仅 $r_{eq}>0$ 时非零，
  $\partial^2/\partial p_{in}^2=2r/V^2$ 等（:2588--2620）。
\item \textbf{不等式曲率}：支路容量行完整二阶展开（:2626--2768）；
  VSC 容量圆 $2\nu$ 对角（:2770--2777）；DC 支路（:2780--2802）；VSC
  电流（:2804--2816）；调制与占空比为线性行、曲率为零，仅推进对偶
  游标保持对齐（:2817--2822）。
\end{itemize}

稠密版本 \srcpath{lagrangian\_hessian\_dense} 只是稀疏结果的转换包装，
供诊断与稠密实验使用，生产路径应走稀疏重载（:2207--2216；
formulation.hpp:404--418 注释）。参数 \fld{drop\_tol} 目前保留未实现，
仅为 API 兼容（formulation.hpp:438--439；实现中形参未命名，:2226）。

\subsection{与潮流方程口径的一致性}\label{subsec:parity-pf-consistency}

Parity 建模层与混合潮流共享同一份 canonical 数据与网络方程，一致性
可逐条核实：

\begin{itemize}
\item \textbf{同源装配}：\srcpath{build\_problem} 与潮流一样经
  \srcpath{make\_solver\_data} 投影，$Y_{\mathrm{bus}}$、$G^{dc}$、
  \fld{base\_mva} 完全同源（:111）。
\item \textbf{交流注入}：$P^{net}$/$Q^{net}$ 极坐标方程与潮流失配式
  同形（:1433--1455）；ZIP 负荷注释明确"consistent with PF solver"
  （:1496）。
\item \textbf{直流平衡}：注释明确"Match the PF definition exactly:
  $Pcalc_k=V_k(G_{dc}V)_k$"（:1510--1512）；DC/DC 折入注释说明镜像
  \srcpath{residual\_evaluator.cpp} 的做法（:1526--1527）。
\item \textbf{VSC 损耗}：$a+bI+cI^2$ 与潮流换流器损耗分段口径同形，
  系数 $a=$\fld{loss\_mw}$/S_{\mathrm{base}}$、$b=$\fld{loss\_percent}$/100$、
  $c=1-\eta$ 一致（:1613--1616；潮流侧见
  \srcpath{src/power\_flow/converter\_model.cpp}）。
\item \textbf{LCC}：直接复用潮流模块的
  \srcpath{lcc\_ac\_injection}/\srcpath{lcc\_dc\_injection}/
  \srcpath{lcc\_ac\_dc\_jacobian}（:1569--1572、1784--1785；
  \srcpath{include/hacdcpf/power\_flow/lcc\_model.hpp}:120--163），
  模型本体只有一份。
\item \textbf{符号约定差异（显式声明）}：DC/DC 的 OPF 变量为输入侧
  $p_{in}$，潮流 \fld{p\_ref\_mw} 为输出侧，二者经 $\eta$ 换算
  （:1579--1583）；能量路由器端口"注入为正"与潮流
  \fld{p\_spec}/\fld{pdc\_spec} 约定一致（:1628--1631）。
\item \textbf{解后回代}：Parity OPF 收敛后驱动层在调度点重跑一次
  潮流并返回 \fld{post\_pf} 支路潮流/换流器量（公共路由行为，见
  第~\ref{sec:overview}~节），交叉验证优化点的潮流可行性；系统化
  数值校核见第~\ref{sec:verification}~节。
\end{itemize}

\subsection{边界与局限}

按诚实口径列明当前实现的近似与未覆盖项（出处见前文各节）：

\begin{enumerate}
\item \textbf{单时段口径}：无跨时段 SOC，储能变量以 $\pm10^{-4}/
  S_{\mathrm{base}}$ 容差钉定于日程出力（:735--758）；静态承载力储能
  完全固定。多时段能量调度需上层 UC/日程给定。
\item \textbf{无功能力曲线未实现}：发电机 $Q$ 仅为盒式界
  （:651--652），不随 $P$ 收缩；VSC 亦只有盒式 $P$/$Q$ 界加容量圆。
\item \textbf{分接头/移相角不可优化}：$t$、$\varphi$ 在支路容量式与
  $Y_{\mathrm{bus}}$ 中为固定参数（:1974--1980）；LCC 换流变 tap 同样
  不进决策向量（:1556--1558）。离散档位优化在 RPO 层外环完成
  （第~\ref{sec:rpo}~节）。
\item \textbf{LCC 为准稳态固定阶注入}：不占决策变量，其功率随
  $(V_m,V_{dc})$ 被动变化；Hessian 贡献为有限差分而非解析
  （:2400--2533），拐点处依赖活跃集守护。
\item \textbf{能量路由器无损}：端口有功严格守恒，\fld{loss\_percent}
  损耗精确化注释声明留作后续（:1651--1652）。
\item \textbf{近似与平滑}：$|I_{ac}|$ 以 $\varepsilon_I$ 平滑
  （\fld{eps\_iac} 默认 $10^{-6}$）；$V_m$ 在换流器行取下限 $10^{-8}$、
  DC/DC 损耗取下限 0.1（:1608、1597）；容量圆半径在额定缺失时回退
  四界最大值（:94--100）。
\item \textbf{行缩放的对偶口径}：AC 平衡行乘 $s_P$/$s_Q$
  （:530--537），故对应拉格朗日乘子需除以缩放因子才是功率量纲的
  影子价格；Hessian 组装已带同一缩放（:2321--2322）。
\item \textbf{文档/实现不一致两处}：头文件关于不可观测直流母线"软
  并联锚定项"的表述（formulation.hpp\allowbreak :173--176）与紧界实现
  （:678--688）不符；\fld{drop\_tol} 参数保留未实现
  （formulation.hpp:438--439）。
\item \textbf{目标的常数项}：固定注入的线性成本与 RPO 网损账本中的
  常量不影响极小点，只影响报告值（:1200--1203、1266--1282）。
\item \textbf{暖启动不保证可行}：最小二乘修正分解失败时保留未修正
  点（:1189--1191）；可行性由 IPM 迭代自行恢复
  （第~\ref{sec:parityipm}~节）。
\end{enumerate}

% 交直流混合最优潮流：Parity IPM 求解核
% 本章文件由 \input 引入，不含导言区。
\section{交直流混合最优潮流：Parity IPM 求解核}\label{sec:parityipm}

本章给出交直流混合 OPF（第~\ref{sec:parityform}~章建模）的自研原始对偶内点
求解核（下称 Parity IPM 核）的数值口径：它以 Parity 建模层装配出的
\srcpath{parity::Problem} 为输入，在同一个牛顿框架内同时推进原始变量、
松弛变量与全部乘子，输出收敛证书、状态向量与诊断计数。全章符号锚点均指
\srcpath{src/optimal\_power\_flow/parity\_ipm.cpp} 与
\srcpath{include/hacdcpf/optimal\_power\_flow/native\_ipm\_solver.hpp}
（IPMOptions/IPMResult 类型定义），驱动层锚点指
\srcpath{src/optimal\_power\_flow/ac\_opf.cpp}。核内全部方程为标幺值
（$S_{\mathrm{base}}=100$~MVA，见 \ref{sec:parityform}~节）；成本币种未定。

\subsection{求解核定位与调用链}

\subsubsection*{功能定位}
Parity IPM 核是平台 AC OPF 三条非线性求解路线之一（另两条为
\ref{sec:acopfnative}~节的 Native AC 外部障碍环路与本章
\ref{subsec:pipm-ipopt}~节的嵌入式 Ipopt 适配）。与 Native AC 路径的
“外层障碍参数 $+$ 内层牛顿”两层结构不同，本核在 Parity 建模层给出的
全空间问题上直接运行 Mehrotra 预测--校正原始对偶内点法：所有变量块与
约束行在 \srcpath{Problem::vidx}/\srcpath{Problem::cidx} 中显式编号
（\srcpath{include/hacdcpf/optimal\_power\_flow/formulation.hpp:VarIndex、:ConstraintIndex}），
核不再做任何设备级聚合或惩罚逼近，因此它是混合 AC/DC 统一建模的
首选审查路径（\srcpath{include/hacdcpf/optimal\_power\_flow/parity\_ipm\_solver.hpp:solve\_ac\_opf\_parity}
头文件注释）。核的公开入口为
\srcpath{solve\_primal\_dual\_ipm}（声明
\srcpath{native\_ipm\_solver.hpp:solve\_primal\_dual\_ipm}，定义 :solve\_primal\_dual\_ipm），配套的同伦切线
例程 \srcpath{homotopy\_tangent} 复用同一套 KKT 机器为 Davidenko 目标
同伦提供路径切向量（声明 \srcpath{native\_ipm\_solver.hpp:homotopy\_tangent}，
定义 :homotopy\_tangent）。

\subsubsection*{调用链与驱动层}
头文件 \srcpath{parity\_ipm\_solver.hpp:solve\_ac\_opf\_parity} 声明了高层驱动
\srcpath{solve\_ac\_opf\_parity}，但该声明在本检出处\textbf{没有对应定义}
（全仓检索仅命中声明本身），属悬置声明；实际驱动是
\srcpath{ac\_opf.cpp} 中的文件静态函数 \srcpath{solve\_with\_parity\_ipm}
（src/optimal\_power\_flow/ac\_opf.cpp:solve\_with\_parity\_ipm），由公共门面 \srcpath{solve\_ac\_opf}
（src/optimal\_power\_flow/ac\_opf.cpp:solve\_ac\_opf）按 \fld{ACOPFOptions::ac\_solver\_backend} 分派：
\begin{itemize}
  \item \fld{ACOPFSolverBackend::ParityIPM}：强制本核
        （内层枚举 \fld{ParityInnerSolver::NativeIPM}，src/optimal\_power\_flow/ac\_opf.cpp:solve\_ac\_opf\_impl）；
  \item \fld{ACOPFSolverBackend::Ipopt}：强制嵌入式 Ipopt
        （src/optimal\_power\_flow/ac\_opf.cpp:solve\_ac\_opf\_impl）；
  \item \fld{Auto}（默认）：本核先行，未收敛且 \fld{allow\_fallback}
        为真时回退 Ipopt（src/optimal\_power\_flow/ac\_opf.cpp:solve\_with\_parity\_ipm）；凡检测到 DC/VSC/LCC
        子系统即自动启用本核所在的 Parity 路径（src/optimal\_power\_flow/ac\_opf.cpp:solve\_ac\_opf\_impl）。
\end{itemize}
内层求解器三态枚举 \fld{ParityInnerSolver} 定义于 src/optimal\_power\_flow/ac\_opf.cpp:ParityInnerSolver。
驱动层还负责把面向用户的 \fld{ACOPFOptions} 折算为核参数
\fld{IPMOptions}（src/optimal\_power\_flow/ac\_opf.cpp:solve\_with\_parity\_ipm，映射关系见
\ref{subsec:pipm-params}~节参数表说明列），并在求解后把 \fld{IPMResult}
归因回 \fld{ACOPFResult}（src/optimal\_power\_flow/ac\_opf.cpp:solve\_with\_parity\_ipm，见
\ref{subsec:pipm-errors}~节）。整体调用链与门控位置见图~\ref{fig:pipm-chain}。

\begin{figure}[H]
\centering
\resizebox{\textwidth}{!}{%
\begin{tikzpicture}[
  font=\small,
  box/.style={draw, rectangle, align=center, minimum height=0.85cm,
              inner sep=4pt},
  lab/.style={font=\footnotesize, align=center},
  arr/.style={-{Stealth[length=2.2mm]}, thick},
  node distance=0.55cm and 0.9cm]
  \node[box] (api) {\srcpath{solve\_ac\_opf}\\[-2pt]{\footnotesize src/optimal\_power\_flow/ac\_opf.cpp:solve\_ac\_opf}};
  \node[box, below=of api] (drv) {\srcpath{solve\_with\_parity\_ipm}\\[-2pt]
        {\footnotesize src/optimal\_power\_flow/ac\_opf.cpp:solve\_with\_parity\_ipm（后端分派/选项折算/结果归因）}};
  \node[box, below left=0.7cm and 0.2cm of drv] (prob) {Parity 建模层\\
        \srcpath{build\_problem}\\[-2pt]{\footnotesize 见 \ref{sec:parityform}~节}};
  \node[box, below right=0.7cm and 0.2cm of drv] (ipopt) {\srcpath{solve\_parity\_with\_ipopt}\\[-2pt]
        {\footnotesize src/optimal\_power\_flow/ac\_opf.cpp:solve\_parity\_with\_ipopt（\srcpath{HACDCPF\_HAVE\_IPOPT} 门控）}};
  \node[box, below=2.0cm of drv] (core) {Parity IPM 核\\
        \srcpath{solve\_primal\_dual\_ipm}\\[-2pt]{\footnotesize src/optimal\_power\_flow/parity\_ipm.cpp:solve\_primal\_dual\_ipm}};
  \node[box, below left=0.7cm and 0.2cm of core] (kkt) {KKT 线性代数\\
        稠密 LU / UMFPACK / KLU\\/ Eigen SparseLU（MUMPS 被门控）\\[-2pt]
        {\footnotesize src/optimal\_power\_flow/parity\_ipm.cpp:backend\_order}};
  \node[box, below right=0.7cm and 0.2cm of core] (res) {\fld{IPMResult}\\
        收敛标志/残差/状态串\\[-2pt]{\footnotesize include/hacdcpf/optimal\_power\_flow/native\_ipm\_solver.hpp:IPMResult}};
  \draw[arr] (api) -- (drv);
  \draw[arr] (drv) -- (prob);
  \draw[arr] (drv) -- (ipopt);
  \draw[arr] (prob) -- (core);
  \draw[arr] (drv) -- node[lab, right=2pt]{NativeIPM/Auto} (core);
  \draw[arr] (core) -- (kkt);
  \draw[arr] (core) -- (res);
  \draw[arr, dashed] (ipopt) -- node[lab, above=1pt]{Auto 回退} (res);
\end{tikzpicture}}
\caption{Parity IPM 核的调用链与门控位置。虚线为仅当本核未收敛且
\fld{allow\_fallback} 时启用的 Ipopt 替代路径（\ref{subsec:pipm-ipopt}~节）。}
\label{fig:pipm-chain}
\end{figure}

\subsubsection*{选项与结果结构体}
\compmeta{IPMOptions}{\srcpath{include/hacdcpf/optimal\_power\_flow/native\_ipm\_solver.hpp:IPMOptions}}
\fld{IPMOptions} 承载迭代上限、三类容忍度、正则化基数、fraction-to-boundary
因子 $\tau$ 与可选的完整热启动状态（原始点、等式/不等式乘子、松弛）；
全字段见 \ref{subsec:pipm-params}~节参数表。如实说明两点：其一，
\fld{tol\_complementarity} 与 \fld{regularization} 两个字段在本核实现中
\textbf{从未被读取}（:solve\_primal\_dual\_ipm 内仅出现 \fld{tol\_primal}、\fld{tol\_dual}、
\fld{alpha\_max}、\fld{max\_iter} 与 \fld{verbose}），互补容忍度由
\fld{tol\_primal} 派生、正则化梯级为硬编码常数（见
\ref{subsec:pipm-conv}、\ref{subsec:pipm-kkt}~节）；其二，\fld{tol\_dual}
不进入收敛判据本身，仅用于滤波器线搜索的平稳性守卫（:solve\_primal\_dual\_ipm）。

\compmeta{IPMResult}{\srcpath{include/hacdcpf/optimal\_power\_flow/native\_ipm\_solver.hpp:IPMResult}}
\fld{IPMResult} 返回收敛标志、迭代数、三类归一化残差
（\fld{primal\_inf}/\fld{dual\_inf}/\fld{complementarity}）、初值残差、
热启动标记、线性代数计数（\fld{factorization\_calls}、
\fld{linear\_solve\_calls}、\fld{accepted\_steps}、\fld{rejected\_steps}、
\fld{symbolic\_analyze\_calls}、\fld{backend\_escalations}、
\fld{scaling\_rebuilds}）、英文状态串 \fld{status}、实际 KKT 后端标签
\fld{linear\_solver}（\fld{dense\_lu}/\fld{sparse\_umfpack}/
\fld{sparse\_klu}/\fld{sparse\_eigen\_lu}/\fld{sparse}，:solve\_primal\_dual\_ipm）
、Phase I 前后违反量/dual-fit/预算/终止原因、Phase II 接受标志，
与最终原始-对偶状态 $(x,\lambda,\mu,z)$。状态串的可取值全集见
\ref{subsec:pipm-errors}~节的状态串表。

\subsection{障碍问题形式与中心路径参数}\label{subsec:pipm-barrier}

\subsubsection*{障碍问题与松弛化}
核求解的问题是 Parity 建模层装配的非线性规划（头文件问题声明
\srcpath{native\_ipm\_solver.hpp:solve\_primal\_dual\_ipm}）：
\begin{align*}
\min_{x}\; & f(x)\\
\text{s.t.}\; & g(x) = 0,\quad h(x) \le 0,\quad x_{\min} \le x \le x_{\max}
\end{align*}
其中 $x\in\mathbb{R}^{n}$ 为 \srcpath{Problem::vidx} 定义的全空间变量
（$n=$~\fld{n\_total}），$g:\mathbb{R}^n\to\mathbb{R}^{m_{eq}}$ 为等式行
（$m_{eq}=$~\fld{n\_eq\_total}，含 AC/DC 功率平衡、换流器平衡、参考角与
直流参考电压），$h:\mathbb{R}^n\to\mathbb{R}^{m_{nl}}$ 为非线性不等式
（$m_{nl}=$~\fld{n\_ineq\_nonlin}，线路容量、VSC 视在功率/电流/调制窗口等）。
核在入口处把变量上下界改写为两行一组的一次不等式并入 $h$：对每个有限下界
列 $c$ 追加 $x_{\min,c}-x_c\le 0$（雅可比行仅含 $-1$），对每个有限上界列
追加 $x_c-x_{\max,c}\le 0$（仅含 $+1$）（\srcpath{assemble\_inequalities}，
界列扫描 :solve\_primal\_dual\_ipm）。合并后的不等式总数记
$m=n_{iq}=m_{nl}+n_{lb}+n_{ub}$；若 $n_{iq}=0$ 直接抛异常
（:solve\_primal\_dual\_ipm）。引入松弛 $z>0$ 与不等式乘子 $\mu>0$ 后，核迭代求解
扰动 KKT 系统
\begin{align*}
\nabla f(x) + J_g(x)^{\top}\lambda + J_h(x)^{\top}\mu &= 0\\
g(x) &= 0\\
h(x) + z &= 0\\
Z\mu &= \gamma\, e
\end{align*}
（\srcpath{native\_ipm\_solver.hpp:solve\_primal\_dual\_ipm} 文档注释与 :solve\_primal\_dual\_ipm 的
右端组装一致；$Z=\mathrm{diag}(z)$，$e$ 为全 1 向量，$\gamma>0$ 为当前
障碍目标）。建模层对目标梯度作工程缩放：当
$\lVert\nabla f\rVert_\infty>100$ 时取
$\omega_f=\min(1,\,100/\lVert\nabla f\rVert_\infty)$ 统一缩小目标及其
对角 Hessian，把拉氏梯度与最优乘子压回 $O(1)$ 量级（:solve\_primal\_dual\_ipm、:apply\_obj\_scale\_hessian（lambda），参考
Ipopt \fld{nlp\_scaling\_method} 的做法）；下文 $\tilde f=\omega_f f$，
$\tilde\nabla f$、$\nabla^2\tilde f$ 同义。

\subsubsection*{初始点与初值规则}
原始初值优先取调用方热启动点：逐变量投影到界内，并保留极小内点余量
（:solve\_primal\_dual\_ipm）$\varepsilon=\min(\,10^{-6}\max(1,\mathrm{width}),$
$0.49\,\mathrm{width}\,)$；无热启动时由建模层 \srcpath{build\_initial\_point}
给点
（\srcpath{formulation.hpp:build\_initial\_point}）。冷启动松弛取
$z_i=\max(|h_i(x_0)|,\,1)$（$h_i\ge0$）或 $\max(-h_i,\,10^{-2})$
（$h_i<0$），乘子取 $\mu_i=\max(1/z_i,\,10^{-2})$ 以维持
$\mu_i z_i\approx1$；等式乘子 $\lambda=0$（:solve\_primal\_dual\_ipm）。热启动四元组
仅在维度、有限性与正性检查全部通过时整体采用（:solve\_primal\_dual\_ipm）。
$\tau$ 取 \fld{alpha\_max}（:solve\_primal\_dual\_ipm）。

\subsubsection*{有界 Phase I 原--对偶起点构造}
在没有兼容完整四元组且 \fld{enable\_phase\_one=true} 时，核先运行局部约束
Newton 恢复。标准 AC 潮流分区把电压作为 state/basic 列；发电机母线以一个
$Q_g$ 替代对应 $V_m$，参考母线再以一个 $P_g$ 完成方阵。Hybrid AC/DC 的
$V_{dc}$、换流器/DC/DC/能量路由器平衡列按同一规则补齐。若当前没有违反的
非线性不等式且方阵可分解，求
\begin{equation}
 J_B(x_k)\Delta x_B=-g(x_k),\qquad \Delta x_F=0;
\end{equation}
否则把所有 $h_i(x_k)>\varepsilon_p$ 行追加到残差，对未固定在界上的列做一次
变量区间缩放后的 SparseQR。该构造对应 Nocedal--Wright（2006）第 7.5、
11.1 与 16.5 节的缩放约束 Newton/基变量思想；线搜索按 Deuflhard（2011）
第 2.2--2.3 节只接受原 pu 坐标最大违反量严格下降的步。

Hybrid AC/DC 在上述全模型恢复前先利用 DC 电导块做一次结构步。对换流器行
$g_c$ 与 DC 平衡行 $g_d$ 的线性化写成
\begin{align}
J_{dv}\Delta V_{dc}+J_{dp}\Delta P_{dc}&=-g_d,\\
J_{cp}\Delta P_{dc}+J_{ca}\Delta P_{ac}&=-g_c.
\end{align}
先由第二式消去局部 $\Delta P_{dc}$，得到只含
$[\Delta V_{dc},\Delta P_{ac}]$ 的 Schur 系统；DC reference 行直接追加。
这一步只分解约 $n_{dc}+n_{conv}$ 的稀疏块，不形成 OPF Hessian/KKT。接受准则
要求 DC balance/converter/reference 子块无穷范数严格下降，且全模型原坐标
违反量不增；全局 admission 仍然独立生效，因此局部 DC 改善不能强迫远离
可行域的 case300 进入昂贵全模型 QR。诊断字段
\fld{phase\_one\_structural\_*} 与 \fld{phase\_one\_structure} 区分该步、
后续 state/basic Newton 和普通回退。

Phase I 与 Phase II 的默认耦合参数为
$\mu_0=0.1$、$\eta_a=1$、$\eta_p=0.1$、$\eta_c=0.5$。只有初始违反量满足
$\|r_p(x_0)\|_\infty\le\eta_a\mu_0$ 才进入局部恢复；否则以
\fld{outside-phase-one-admission} 零分解退出，且不改变 Phase II 冷起点。
恢复在 $\epsilon_p=\max(\mathrm{tol}_p,\eta_p\mu_0)$ 停止，而不是追到最终
精度。令 $v=\max(0,\max_i h_i)$，松弛下限取
$s_{\min}=\max(2\times10^{-10},(\epsilon_p-v)/2)$，于是
$h_i+s_i\le(\epsilon_p+v)/2\le\epsilon_p$；再置
$\nu_i=\mu_0/s_i$，避免仅用机器下限造成巨大乘子，同时精确保持中心乘积。
移交还要求 $\max_i|s_i\nu_i/\mu_0-1|\le\eta_c$。

Phase I 只有迭代、稀疏分解、回溯与协作式墙钟预算，不含障碍序列、filter、
Hessian、二阶校正或弹性变量。原始点认证后只做一次 $J_g^T\lambda$ 对偶拟合，
\fld{phase\_one\_dual\_fit\_residual} 报告 basic 列 stationarity（或 SparseQR
正规残差）的缩放范数，保留自由控制列的 reduced gradient 给 Phase II。
完整 continuation state 直接绕过 Phase I。\fld{phase\_one\_primal\_feasible}
仍专指达到最终容差；\fld{phase\_one\_in\_handoff\_corridor} 才表示达到上述
有限精度移交邻域。若 primal/dual/centrality 证书任一失败，
Phase II 精确恢复进入 Phase I 之前的原始点以及普通冷启动 slack/dual；因此
\fld{phase\_two\_start\_accepted=false} 的试探轨迹不会改变最优性求解。
这一区分使 Phase I 保持 warm-start producer，而不是第二套 infeasible-start IPM。

\subsubsection*{Prepared OPF session 与完整状态延续}
公开 RAII 类型 \srcpath{PreparedACOPFSession} 为重复 Parity ACOPF 保留
formulation、组件/行映射、稀疏 pattern 的 symbolic ordering，以及上一轮完整
$(x,\lambda,\mu,z)$。兼容键不是维数检查，而是覆盖变量块、等式/不等式块、
有限 bound-row 顺序、reference row、受限设备集合、energy-router port 顺序与
$Y_{bus}/G_{dc}$ 非零 pattern 的 64 位签名；LCC 的稳定 ID、AC/DC 端点、解析后的
换流母线以及 energy-router port 稳定 ID 同样进入语义映射审计。参数变化先重新投影 canonical numeric
\srcpath{SolverData}；若结构审计通过，只刷新负荷/ZIP、固定注入、目标/界与矩阵
数值，保留 OPF maps 和 symbolic ordering。拓扑、服务状态、PV/PQ/reference、
constraint-family 或 row membership 改变时完整 rebuild，并在该次求解禁止 opaque
continuation state。linear backend 还对压缩 KKT 的全部行列坐标计算 pattern 指纹；
维数和非零元个数相同但坐标不同仍强制重新 symbolic analyze。

四块状态必须同时存在、签名相同、维数正确、数值有限且 slack/inequality dual
为正，随后仍通过本节 Phase-II residual/centrality 审计；任一条件失败时只允许
公共物理变量映射，不允许“维数碰巧相同”的 dual row 复用。当前 SuiteSparse
adapter 与 MIPSolvers MUMPS wrapper 只支持复用 symbolic 后重新 numeric factorize，
没有可按矩阵漂移、scaled backward residual 与 refinement 独立审计的跨求解
numeric-factor 对象。因此 \fld{prepared\_numeric\_refactor} 默认关闭；即使开启也
返回 \fld{unsupported} 并做新 numeric factorization，不虚报 numeric reuse。
AC-PF adjoint dual seed 未实现，因为 reduced PV/PQ Jacobian 与 full-space OPF
equality/bound multiplier 布局不兼容，且未保留 PF factor 时仍需额外分解。

\subsubsection*{预测--校正与中心路径参数更新}
每次迭代先对 $z,\mu$ 施加 $10^{-12}$ 下限保护并构成障碍比
$\Sigma=\mathrm{diag}(\mu_i/z_i)$（:solve\_primal\_dual\_ipm），随后执行
Mehrotra 预测--校正。预测步以 $\gamma=0$ 解牛顿系统得仿射方向
$(\Delta x^{\mathrm{aff}},\Delta\lambda^{\mathrm{aff}},
\Delta z^{\mathrm{aff}},\Delta\mu^{\mathrm{aff}})$，并以 $\tau=1$ 探测
仿射步长（凝集与增广形式均见 :solve\_primal\_dual\_ipm）：
\begin{align*}
\alpha_p^{\mathrm{aff}} &= \min\{\mathrm{ftb}(z,\Delta z^{\mathrm{aff}},1),\,1\},
&
\alpha_d^{\mathrm{aff}} &= \min\{\mathrm{ftb}(\mu,\Delta\mu^{\mathrm{aff}},1),\,1\}
\end{align*}
对偶间隙与仿射间隙为（凝集与增广均见 :solve\_primal\_dual\_ipm）
\begin{align*}
\bar\mu &= \frac{z^{\top}\mu}{n_{iq}},
&
\bar\mu_{\mathrm{aff}} &=
\frac{(z+\alpha_p^{\mathrm{aff}}\Delta z^{\mathrm{aff}})^{\top}
      (\mu+\alpha_d^{\mathrm{aff}}\Delta\mu^{\mathrm{aff}})}{n_{iq}}
\end{align*}
定心参数取立方律并截断（凝集与增广均见 :solve\_primal\_dual\_ipm）：
\begin{align*}
\sigma &= \left(\mathrm{clamp}\!\left(
  \frac{\bar\mu_{\mathrm{aff}}}{\bar\mu+10^{-16}},\,0,\,1\right)\right)^{3},
&
\gamma &= \mathrm{clamp}\!\left(\sigma\,\bar\mu,\;10^{-14},\,10^{4}\right)
\end{align*}
校正步右端补入 Mehrotra 二阶交叉项
$c_i=\mathrm{clamp}(\Delta\mu_i^{\mathrm{aff}}\Delta z_i^{\mathrm{aff}},
\,\pm 0.1\gamma)$（:solve\_primal\_dual\_ipm）。凝集形式的校正右端为
（:solve\_primal\_dual\_ipm）
\begin{align*}
N_{\mathrm{corr}} &= \nabla\tilde f + J_g^{\top}\lambda + J_h^{\top}\mu
+ J_h^{\top}\!\left((\mu\circ h + \gamma e - c)\oslash z\right)
\end{align*}
其中 $\circ$、$\oslash$ 为逐分量乘/除；校正方向的松弛与乘子恢复为
（:solve\_primal\_dual\_ipm）
\begin{align*}
\Delta z &= -(h+z) - J_h\,\Delta x,
&
\Delta\mu &= -\left(\mu + \Sigma\,\Delta z + (c-\gamma e)\oslash z\right)
\end{align*}
此后可追加至多 2 次 Gondzio 额外校正：对试步互补对
$\tilde z_i\tilde\mu_i$ 低于 $0.1\gamma$ 的分量以零等式残差右端再解一次
同分解方程并累加方向（凝集与增广均见 :solve\_primal\_dual\_ipm，阈值
$10^{-12}\gamma$）。增广形式另有两条防线：当 fraction-to-boundary 步长
被压到 $10^{-9}$ 以下时判定方向数值爆炸，把 $\gamma$ 提升到
$\max(\gamma,\,0.5\bar\mu)$ 并清零交叉项后重解校正，至多 2 次
（重新定心，:solve\_primal\_dual\_ipm）；滤波器全拒时还可以 $\gamma\leftarrow\max(\gamma,\bar\mu)$
的纯定心步整体重跑一次线搜索（中心性恢复重试，:solve\_primal\_dual\_ipm）。

环境变量 \srcpath{HACDCPF\_OPF\_MU\_STRATEGY=abo}（仅增广形式生效，
:solve\_primal\_dual\_ipm）把立方律换成 Armand--Benoist--Orban 动态 $\mu$：以
$\mu^{+}=\mu_{\mathrm{dyn}}^{2}/\mu_0$ 为更新函数、
\begin{align*}
\gamma &= \mathrm{clamp}\!\left(
\mu_{\mathrm{dyn}} + \alpha_{\mathrm{aff}}(\mu^{+}-\mu_{\mathrm{dyn}}),
\;10^{-14},\,10^{4}\right)
\end{align*}
其中 $\alpha_{\mathrm{aff}}=\min(\alpha_p^{\mathrm{aff}},
\alpha_d^{\mathrm{aff}},1)$，$\mu_0$ 为首个迭代的 $\bar\mu$
（:solve\_primal\_dual\_ipm）；实际采用的 $\gamma$（含重新定心修正）回写为下一迭代的
$\mu_{\mathrm{dyn}}$（:solve\_primal\_dual\_ipm）。默认 \fld{mehrotra}。

\subsection{KKT 系统组装与线性求解}\label{subsec:pipm-kkt}

\subsubsection*{两种牛顿系统形式}
核提供两种 KKT 组装形式（枚举 \fld{KktForm}；环境变量
\srcpath{HACDCPF\_OPF\_KKT\_FORM}，默认 \fld{condensed}，:kkt\_form\_preference）：
\begin{itemize}
  \item \textbf{凝集式（condensed，生产默认）}：把不等式乘子消入
        $(1,1)$ 块，
        \begin{align*}
        W &= \nabla^{2}_{xx}\tilde L + J_h^{\top}\Sigma J_h,
        &
        \begin{bmatrix} W+\delta I & J_g^{\top}\\ J_g & -\delta I\end{bmatrix}
        \begin{bmatrix}\Delta x\\ \Delta\lambda\end{bmatrix}
        &= -\begin{bmatrix} N\\ g\end{bmatrix}
        \end{align*}
        $W$ 的三重积按 $\mathrm{diag}(\Sigma)\cdot J_h$ 先行缩放再乘
        装配（:solve\_primal\_dual\_ipm），KKT 三元组装配见
        \srcpath{factor\_kkt\_sparse}（:factor\_kkt\_sparse）；$N$ 在预测步取
        $\tilde L_x + J_h^{\top}(\mu\circ h\oslash z)$（:solve\_primal\_dual\_ipm），
        校正步取 $N_{\mathrm{corr}}$。头文件注释指出凝集式条件数被
        平方（$\kappa\sim\sigma^2$，$\sigma=\mu/z$ 跨 $10^{\pm10}$），
        是 PEGASE 算例滤波器停滞的根因之一（:KktForm）。
  \item \textbf{增广式（augmented）}：保留不等式乘子为未知量，
        \begin{align*}
        \begin{bmatrix}
        \nabla^{2}_{xx}\tilde L+\delta_W I & J_g^{\top} & J_h^{\top}\\
        J_g & -\delta_C I & 0\\
        J_h & 0 & -ZM^{-1}
        \end{bmatrix}
        \begin{bmatrix}\Delta x\\ \Delta\lambda\\ \Delta\mu\end{bmatrix}
        = \begin{bmatrix}-\tilde L_x\\ -g\\ r_3\end{bmatrix}
        \end{align*}
        （块结构注释、装配与 $(3,3)$ 块 $-z_i/\mu_i$ 均见
        \srcpath{factor\_kkt\_sparse\_augmented}）。条件数对障碍比为
        线性，且对称不定分解可精确核对惯性：约化 Hessian 正定时惯性应为
        $(n,\;m_{eq}+n_{iq},\;0)$。
\end{itemize}
稠密路径（下文）与 $n_{iq}=0$ 时一律回落凝集式（:solve\_primal\_dual\_ipm）。

\subsubsection*{正则化与惯性控制}
凝集式采用渐进正则化：依次尝试
$\delta\in\{0,\;10^{-8}\lVert W\rVert,\;10^{-6}\lVert W\rVert,
\;10^{-4}\lVert W\rVert\}$（下限 $10^{-12}$），先分解再回代，任一失败
即升级，四级全败报 \texttt{"KKT factorization failed"}（:solve\_primal\_dual\_ipm）。
增广式实现 Wächter--Biegler $\delta_W$ 惯性校正：仅当 LDL$^{\top}$ 后端
（MUMPS）报告负主元数且不等于 $m_{eq}+n_{iq}$ 时，以
$\delta_W=10^{-8}\lVert\nabla^2_{xx}\tilde L\rVert$ 起步、每次 $\times8$、
上限 $10^{-2}\lVert\nabla^2_{xx}\tilde L\rVert$ 重分解，至多 8 次，末次
封顶后按 best-effort 接受（:solve\_primal\_dual\_ipm）；LU 类后端不报惯性
（\fld{negs}$=-1$）则\textbf{不加判别地接受未正则化分解}（:solve\_primal\_dual\_ipm）。
$\delta_W$ 热启动取 $\max(10^{-8}\lVert L_{xx}\rVert,\,0.5\,\delta_W^{prev})$
且只对 MUMPS 步启用，避免泄漏进 KLU 轨迹（:solve\_primal\_dual\_ipm）。$(2,2)$ 块正则
$\delta_C$ 默认 $10^{-8}\lVert L_{xx}\rVert$，可用
\srcpath{HACDCPF\_OPF\_DELTA\_C} 覆盖（:solve\_primal\_dual\_ipm）。如实说明：由于
MUMPS 在本检出处不可用（见下），$\delta_W$ 校正在实际构建中被静默绕过，
增广式的惯性保证随之失效。

\subsubsection*{平衡、分解与后端选择}
稀疏 KKT 缓存 \fld{SparseKKTCache}（:SparseKKTCache）按“模式不变”假设复用符号
分析：每次求解先对 KKT 做对称 Ruiz 平衡——迭代至多 5 轮、每轮把
$D_i$ 除以当前行 $\infty$-范数的平方根、行失衡偏差 $|\log r_{\max}|<10^{-3}$
即停（:compute\_ruiz\_scaling）——得到 $D\!KD$，此后求解按
$(D\!KD)\,y=D\,b$、$x=D\,y$ 进行（:kkt\_solve\_scaled）。平衡矩阵默认整个
求解过程只算一次并冻结（\srcpath{HACDCPF\_OPF\_RESCALE\_PER\_ITER=1}
恢复逐迭代重算，:factor\_assembled\_kkt）；注释记录了冻结缩放把 IEEE24 三区域扩展
算例从 640 迭代停滞修复为 36 迭代收敛的教训（:factor\_assembled\_kkt）。
回代自带至多 2 步迭代精化，且仅当校正\textbf{确实减小}残差时才采纳
（防数值奇异 KKT 上精化发散，:kkt\_solve\_scaled）。

线性后端按结构感知顺序选取（\srcpath{backend\_order}）：
\begin{itemize}
  \item 纯 AC 问题（无 DC 母线/支路、VSC、DC/DC、能量路由器，判定
        :solve\_primal\_dual\_ipm）：KLU $\to$ MUMPS $\to$ UMFPACK $\to$ Eigen
        SparseLU（:backend\_order）。KLU 的块三角分解匹配电网 KKT 的近块三角
        结构，注释给出 case2869 分解 0.006~s 对 MUMPS 0.015~s、端到端
        5.97~s 对 12.05~s 且目标值一致的实测（:backend\_order）；
  \item 混合 AC/DC 问题：MUMPS $\to$ UMFPACK $\to$ KLU $\to$ Eigen
        SparseLU（:backend\_order），理由是 $\delta_W$ 校正需要 LDL$^{\top}$
        惯性（:backend\_order）。
\end{itemize}
活动后端具有粘性；仅当 Eigen SparseLU 回代残差超过
$10^{-6}\lVert b\rVert_\infty$ 时置 \fld{solve\_degraded}，下次分解沿候选
序列升级一档（:kkt\_solve\_scaled、:factor\_assembled\_kkt）；UMFPACK/KLU 的告警残差不触发降级
（注释 :kkt\_solve\_scaled）。每个后端按 $(\dim,\mathrm{nnz})$ 缓存各自的符号分析，
模式不变即复用（:SparseKKTCache、:factor\_assembled\_kkt）。$n+m_{eq}\le512$ 的小系统改走
稠密 \fld{Eigen::PartialPivLU}（阈值 \fld{kDenseAutoKktDim}，:solve\_primal\_dual\_ipm；
稠密装配/分解 :factor\_kkt\_dense，同样带 2 步精化 :kkt\_solve\_dense）；$0\times0$
KKT 按 $\mathbb R^0$ 上唯一空解直接闭合，不调用 Eigen 分解。环境变量
\srcpath{HACDCPF\_OPF\_LINEAR\_SOLVER}（旧拼写
\srcpath{HACDCPF\_OPF\_SPARSE\_SOLVER} 兼容）可强制
\fld{dense}/\fld{mumps}/\fld{umfpack}/\fld{klu}/\fld{eigen}/\fld{auto}
（:linear\_solver\_preference\_env）。

\subsubsection*{稀疏库的编译期门控与回退}
UMFPACK/KLU 支持由 \srcpath{HACDCPF\_OPF\_HAVE\_UMFPACK} 与
\srcpath{HACDCPF\_OPF\_HAVE\_KLU} 门控（:SparseKKTCache、:backend\_order）：CMake 在
\srcpath{HACDCPF\_USE\_SUITESPARSE=ON}（默认 ON，\srcpath{CMakeLists.txt:HACDCPF\_USE\_SUITESPARSE}）
下优先链接 MIPSolvers 内置的 \fld{umfpack\_vendored}/\fld{klu\_vendored}
目标并定义两宏（\srcpath{CMakeLists.txt:HACDCPF\_OPF\_HAVE\_UMFPACK、:HACDCPF\_OPF\_HAVE\_KLU}），否则回退系统 SuiteSparse 探测
（\srcpath{CMakeLists.txt:HACDCPF\_SUITESPARSE\_INCLUDE\_DIR}）；两级都找不到时两个宏不定义，稀疏路径只剩
Eigen \fld{SparseLU}（COLAMD 排序，:SparseKKTCache、:sparse\_factorize\_active）。MUMPS 分支由
\srcpath{HACDCPF\_HAVE\_MUMPS} 门控（:SparseKKTCache、:backend\_order），文件头注释称该宏
“继承自 mipsolvers 目标的公共编译定义”；但本检出处 MIPSolvers 的 CMake
从不定义该宏（全仓检索仅在 parity\_ipm.cpp 内部出现），且其所引用的
\srcpath{mipsolvers::engine::MumpsSolver} 类在当前 MIPSolvers 头文件树中
不存在。因此\textbf{本检出处 MUMPS 恒被编译排除}：混合 AC/DC 问题的实际
后端序列为 UMFPACK $\to$ KLU $\to$ Eigen SparseLU，上文的 $\delta_W$
惯性校正与 case 注释中的 MUMPS 实测数字对当前构建均不适用。

\subsection{步长规则与全局化}\label{subsec:pipm-step}

\subsubsection*{fraction-to-boundary 与空步判定}
原始/对偶步长分别由 fraction-to-boundary 规则给出
（\srcpath{ftb}）：
\begin{align*}
\alpha(v,\Delta v,\tau) &= \min\!\left\{1,\;
\min_{i:\,\Delta v_i<0}\frac{-\tau\, v_i}{\Delta v_i}\right\}
\end{align*}
最终取 $\alpha_p=\max(\min(\mathrm{ftb}(z,\Delta z,\tau),1),10^{-10})$、
$\alpha_d$ 同理（:solve\_primal\_dual\_ipm）。当 $\min(\alpha_p,\alpha_d)<10^{-7}$
时判定该方向为空步（vacuous step）：滤波器的进步余量
$\eta\alpha$、$\gamma_\theta\theta$ 已数值失效，直接放弃线搜索进入
异常路径（:solve\_primal\_dual\_ipm；注释记录了 case13659 在 $\alpha\sim10^{-9}$
“被接受”步上随机游走数百迭代的事故）。

\subsubsection*{回溯线搜索与 $(\theta,\varphi)$ 滤波器}
线搜索对试步 $(x+\alpha_p\Delta x,\;z+\alpha_p\Delta z,\;
\mu+\alpha_d\Delta\mu,\;\lambda+\alpha_d\Delta\lambda)$ 重估全部非线性
残差并做接受测试，拒绝则 $\alpha\leftarrow0.5\alpha$，至多 14 次回溯
（\fld{kMaxBacktracks}/\fld{kBacktrack}，:solve\_primal\_dual\_ipm；主循环
\srcpath{run\_backtracks}（lambda））。接受准则有两层：
\begin{enumerate}
  \item \textbf{遗留三轴准则}：可行轴要求
        $\theta_t\le(1-\eta\bar\alpha_p)\theta$ 且对偶两轴不放大超 100 倍；
        平稳轴与互补轴分别在可行/平稳守卫
        （$\theta_t\le\max(10\varepsilon_p,1.02\theta)$、
        $\|L_x\|_t\le\max(10\varepsilon_d,1.02\|L_x\|)$）下要求对应量
        按 $(1-\eta\bar\alpha_d)$ 下降，$\eta=10^{-4}$（:eval\_and\_test\_trial（lambda））。
        注释指出互补轴守卫宽松（每步 2\%，复利）曾使 case2869pegase
        目标与平稳性虚胀约 3 倍，可用 \srcpath{HACDCPF\_OPF\_NO\_LEGACY=1}
        关闭（:solve\_primal\_dual\_ipm）；
  \item \textbf{$(\theta,\varphi)$ 滤波器}（Wächter--Biegler §3.2，Ipopt
        默认常数）：$\varphi_\gamma(x,z)=\tilde f(x)-\gamma\sum_i\ln z_i$，
        试点须相对当前对与\textbf{全部历史对}在
        $\theta$ 上占优（$\theta_t\le(1-\gamma_\theta)\theta$）或在
        $\varphi$ 上至少降 $\gamma_\varphi\theta$（支配记忆见
        :eval\_and\_test\_trial（lambda），历史入栈 :run\_backtracks（lambda））；近可行且方向为真实 $\varphi$ 下降
        时追加 Armijo 条件
        $\varphi_t\le\varphi_k+\eta_\varphi\bar\alpha\, d\varphi$
        （switching 条件与 Armijo 均见 :eval\_and\_test\_trial（lambda））。常数
        $\gamma_\theta=\gamma_\varphi=10^{-5}$、$\delta_s=1.0$、
        $s_\theta=1.1$、$\eta_\varphi=10^{-8}$、$\theta_{\min}=10^{-4}$
        （:solve\_primal\_dual\_ipm）；另有绝对天花板
        $\theta_{\max}=10^{4}\max(1,\theta_0)$（:solve\_primal\_dual\_ipm）。
\end{enumerate}
滤波器机制本身即非单调策略：接受只要求相对历史包络的进步，允许单轴
短期恶化，因此目标与残差沿轨迹均不单调；核另以 best-iterate 档案兜底
（:solve\_primal\_dual\_ipm）。如实说明：只有当 $n+m_{eq}\ge512$ 或
\srcpath{HACDCPF\_OPF\_FILTER\_ALL=1} 时线搜索才真正执行（:solve\_primal\_dual\_ipm）；
小 KKT 系统在 \srcpath{eval\_and\_test\_trial} 内无条件接受首个有限试步
（:eval\_and\_test\_trial（lambda）），即小算例走的是无阻尼全步牛顿轨迹，发散风险由后续异常路径
承担。诊断开关 \srcpath{HACDCPF\_OPF\_STRICT\_THETA=1} 把接受准则收紧为
纯可行下降（:solve\_primal\_dual\_ipm、:eval\_and\_test\_trial（lambda））。

\subsubsection*{二阶校正、中心性恢复与恢复阶段}
三个逐级升级的全局化补救（均仅增广式；凝集式失败后直接进入异常路径）：
\begin{enumerate}
  \item \textbf{二阶校正（SOC）}：首轮试步因可行性恶化被拒时，用同一
        分解把等式块右端换成 $-(\alpha_p g+g(x+\alpha_p\Delta x))$ 再解
        一次并尝试复合步，抵消 Maratos 效应的一阶残差上升，每迭代至多
        一次（:solve\_primal\_dual\_ipm 注释，实现 :run\_backtracks（lambda））；
  \item \textbf{中心性恢复重试}：滤波器全拒但首轮可行性尚可时，以
        $\gamma\leftarrow\max(\gamma,\bar\mu)$、交叉项清零重解校正方向
        并再跑一次线搜索（Gondzio--Grothey 式定心步，:solve\_primal\_dual\_ipm）；
  \item \textbf{恢复阶段}：以上均失败时，把 $(1,1)$ 块换成
        $\rho I$（$\rho=10^{-8}\lVert L_{xx}\rVert$）做一步 Gauss-Newton
        型可行性步，$\alpha$ 折半至多 8 次，接受首个
        $\theta$ 下降且对偶轴有界（$\le100$ 倍）的试步；下一主迭代用真实
        $L_{xx}$ 重新分解，等效于 Ipopt 恢复模式的有界版本
        （:solve\_primal\_dual\_ipm）。
\end{enumerate}

\subsection{收敛判据}\label{subsec:pipm-conv}

\subsubsection*{归一化残差度量}
核以四个归一化量度量收敛（\srcpath{compute\_convergence}）：
\begin{align*}
\theta &= \frac{\max(\lVert g\rVert_\infty,\,\max_i h_i^{+})}
               {1+\max(\lVert x\rVert_\infty,\lVert z\rVert_\infty)},
&
\theta_{\mathrm{raw}} &= \max(\lVert g\rVert_\infty,\,\max_i h_i^{+})\\
\kappa_d &= \frac{\lVert \nabla\tilde f+J_g^{\top}\lambda
            +J_h^{\top}\mu\rVert_\infty}
            {1+\max(\lVert\mu\rVert_\infty,\lVert\lambda\rVert_\infty)},
&
\kappa_c &= \frac{z^{\top}\mu/n_{iq}}{1+\lVert x\rVert_\infty}
\end{align*}
另有成本条件 $\kappa_f=|\tilde f-\tilde f_0|/(1+|\tilde f_0|)$（:compute\_convergence（lambda）），
但如实说明：$\kappa_f$ 仅用于日志输出，\textbf{不进入任何收敛判据}。

\subsubsection*{判据全表}
下表“典型范围”列填判据的不等式形式，“默认值”列填默认容忍度，“说明”列
末注源码行号（均为 \srcpath{parity\_ipm.cpp}）。
\begin{paramtable}
原始可行（归一化） & $\theta$ & — & $<\varepsilon_p=\fld{tol\_primal}$ & $10^{-6}$ & 驱动层取 $\max(\fld{feasibility\_tol},10^{-10})$（src/optimal\_power\_flow/ac\_opf.cpp:solve\_with\_parity\_ipm；:compute\_convergence（lambda）、:solve\_primal\_dual\_ipm）\\
原始可行（原始量纲） & $\theta_{\mathrm{raw}}$ & — & $<\varepsilon_p$ & $10^{-6}$ & 与归一化值同时要求，防大 $\lVert x\rVert$ 稀释（:compute\_convergence（lambda）、:solve\_primal\_dual\_ipm）\\
对偶可行（平稳性） & $\kappa_d$ & — & $<\varepsilon_p$ & $10^{-6}$ & 注意：比较的是 $\varepsilon_p$ 而非 $\varepsilon_d$；$\fld{tol\_dual}$ 只进滤波器守卫（:compute\_convergence（lambda）、:solve\_primal\_dual\_ipm、:eval\_and\_test\_trial（lambda））\\
互补间隙 & $\kappa_c$ & — & $<\max(10\varepsilon_p,\,2\times10^{-3})$ & $2\times10^{-3}$ & 大算例障碍病态下平均 $z\mu/n_{iq}$ 常停在 $10^{-3}$ 级，故设工程下限；\fld{tol\_complementarity} 字段未被读取（:solve\_primal\_dual\_ipm）\\
迭代上限 & — & — & $\fld{max\_iter}$ & 640 & 驱动层折算为内、外上限乘积 $80\times8$（src/optimal\_power\_flow/ac\_opf.cpp:solve\_with\_parity\_ipm；:solve\_primal\_dual\_ipm）\\
可接受级收敛 & — & — & $\theta,\theta_{\mathrm{raw}},\kappa_d<100\,\varepsilon_p$ 且 $\kappa_c$ 达标 & $100\times$ & 取当前点或 best-iterate 中更优者，状态串如实标注 acceptable（Ipopt \fld{Solved\_To\_Acceptable\_Level} 惯例）（:solve\_primal\_dual\_ipm）\\
目标停滞早停 & — & — & $\theta_{\mathrm{raw}}<10^{-4}$、$\kappa_c$ 达标、$|\Delta \tilde f|\le10^{-6}(1+|\tilde f|)$ 连续 20 次 & 20 次 & 早停后状态为 \texttt{objective-stagnant (dual-polish certification pending)}，交由对偶抛光认证（:solve\_primal\_dual\_ipm）\\
对偶最小二乘认证 & — & — & 终点 $\theta,\theta_{\mathrm{raw}},\kappa_c$ 达 acceptable 且抛光后 $\kappa_d<100\,\varepsilon_p$ & $100\times$ & 解 $\min_{\lambda,\mu}\tfrac12\lVert\tilde\nabla f+J_g^{\top}\lambda+J_h^{\top}\mu\rVert^2$ 的 KKT 系统（$(1,1)$ 块取 $I$），$\mu$ 截断非负；状态串 \texttt{converged (dual least-squares certified)}（:solve\_primal\_dual\_ipm）\\
\end{paramtable}

严格判据为 $\theta<\varepsilon_p$、$\theta_{\mathrm{raw}}<\varepsilon_p$、
$\kappa_d<\varepsilon_p$、$\kappa_c<\kappa_c^{tol}$ 四者同时成立
（:solve\_primal\_dual\_ipm 初判与迭代内判定，置 \fld{converged} 与状态
\fld{"converged"}）。对偶抛光只在 $\kappa_d$ 不劣于轨迹自身乘子时采纳
抛光乘子（:solve\_primal\_dual\_ipm）。

\subsection{异常路径与结果标记}\label{subsec:pipm-errors}

\subsubsection*{核内状态串全集}
\fld{IPMResult::status} 的全部可取值及其触发位置如下（默认初值
\texttt{"maximum iterations reached"}，:solve\_primal\_dual\_ipm；“说明”列末注源码行号）：
\begin{paramtable}
\fld{converged} & — & — & — & — & 严格判据满足（:solve\_primal\_dual\_ipm）\\
\texttt{converged (acceptable tolerance, best iterate)} & — & — & — & — & best-iterate 达 $100\times$ 容忍度（:solve\_primal\_dual\_ipm）\\
\texttt{converged (acceptable tolerance, final iterate)} & — & — & — & — & 终点达 $100\times$ 容忍度（:solve\_primal\_dual\_ipm）\\
\texttt{converged (dual least-squares certified)} & — & — & — & — & 对偶抛光认证通过（:solve\_primal\_dual\_ipm）\\
\texttt{objective-stagnant (dual-polish certification pending)} & — & — & — & — & 目标停滞窗口触发；随后若抛光认证失败则保持未收敛（:solve\_primal\_dual\_ipm）\\
\texttt{maximum iterations reached} & — & — & — & — & 迭代耗尽（:solve\_primal\_dual\_ipm 默认值）\\
\texttt{KKT factorization failed} & — & — & — & — & 正则化梯级/惯性校正全败（:solve\_primal\_dual\_ipm 凝集与增广两支）\\
\texttt{predictor back-solve failed} & — & — & — & — & 预测步回代失败（:solve\_primal\_dual\_ipm，仅增广）\\
\texttt{NaN in corrector step} & — & — & — & — & 校正方向出现非有限值（:solve\_primal\_dual\_ipm）\\
\texttt{NaN in corrector step} & — & — & — & — & 校正方向出现非有限值（:build\_kkt\_workspace、:solve\_economic\_dispatch）\\
\texttt{filter line search failed} & — & — & — & — & 回溯耗尽且恢复阶段失败（:solve\_primal\_dual\_ipm）\\
\end{paramtable}
两条硬异常在入口处抛出 \fld{std::runtime\_error}：问题规模非法
（$n\le0$ 或 $m_{eq}<0$，:solve\_primal\_dual\_ipm）与装配后无不等式（:solve\_primal\_dual\_ipm）。

\subsubsection*{向 ACOPFResult 的映射（诚实口径）}
驱动层把核结果翻译为公共结果（src/optimal\_power\_flow/ac\_opf.cpp:solve\_with\_parity\_ipm）：
\begin{itemize}
  \item \fld{converged} 与 \fld{status}：收敛时
        \texttt{"converged (<backend>)"}（backend 如
        \fld{parity\_ipm:sparse\_klu}、\fld{ipopt\_filter\_linesearch}），
        Ipopt 可接受残差另加 \texttt{[acceptable residuals]}；未收敛时
        \texttt{"not converged: <IPMResult::status>"}，核的 acceptable/抛光
        语义透过状态串保留（src/optimal\_power\_flow/ac\_opf.cpp:solve\_with\_parity\_ipm）；
  \item 计数与残差进 \fld{ACOPFProfiling}：线性后端标签、分解/回代/
        接受/拒绝步数、符号分析次数、后端升级次数、缩放重建次数、热启动
        标记与初值残差（src/optimal\_power\_flow/ac\_opf.cpp:solve\_with\_parity\_ipm；结构体
        \srcpath{opf\_result.hpp:ACOPFProfiling}）；
  \item 原始解不存在或含非有限值时 \fld{objective} 与
        \fld{max\_constraint\_violation} 置 $+\infty$，跳过全部状态提取
        （src/optimal\_power\_flow/ac\_opf.cpp:solve\_with\_parity\_ipm）；
  \item 求解后在终态重估各族残差写入 \fld{profiling} 的
        \fld{max\_ac\_p\_balance\_residual\_pu} 等审计字段
        （src/optimal\_power\_flow/ac\_opf.cpp:solve\_with\_parity\_ipm）；未收敛时由 \srcpath{solve\_ac\_opf}
        汇总 \fld{infeasibility\_hints}（状态串、孤岛无松弛等，
        src/optimal\_power\_flow/ac\_opf.cpp:solve\_ac\_opf\_impl）；
  \item \fld{lmp\_valid} 仅在等式乘子长度覆盖 AC 平衡行时置真并注明
        来源，否则写明不可用原因（Ipopt 适配器不返回约束乘子时 LMP
        不可得，src/optimal\_power\_flow/ac\_opf.cpp:solve\_with\_parity\_ipm）；
  \item 模型覆盖声明走 \fld{converter\_model\_scope}/\fld{ValidityFlags}
        与 \fld{model\_limitations}（路径标签按实际启用的约束族拼接，
        LCC 准稳态、触发角/熄弧角求解后审计越限等逐条列出，
        src/optimal\_power\_flow/ac\_opf.cpp:solve\_with\_parity\_ipm），口径与 \ref{sec:overview}~节
        的诚实结果约定一致。
\end{itemize}

\subsection{嵌入式 Ipopt 替代路径}\label{subsec:pipm-ipopt}

\subsubsection*{编译期与运行期门控}
Parity 问题同时可交由嵌入式 Ipopt 求解（适配器
\srcpath{solve\_parity\_with\_ipopt}，src/optimal\_power\_flow/ac\_opf.cpp:solve\_parity\_with\_ipopt）：Parity
\fld{Problem} 暴露的目标、梯度、等式/不等式残差与雅可比、拉氏 Hessian
回调与 \fld{engine::NLPModel} 一一对应，界约束直接作为变量界传入
（src/optimal\_power\_flow/ac\_opf.cpp:build\_parity\_nlp\_input 注释）。门控分两层：
\begin{itemize}
  \item 编译期 \srcpath{HACDCPF\_HAVE\_IPOPT}（src/optimal\_power\_flow/ac\_opf.cpp:solve\_parity\_with\_ipopt）：
        平台选项 \srcpath{HACDCPF\_ENABLE\_IPOPT} 在 Apple 与 Windows 默认 ON、
        其余平台默认 OFF（\srcpath{CMakeLists.txt:HACDCPF\_ENABLE\_IPOPT}，任何平台均可显式开启）；
        Windows 默认选择 MIPSolvers 的 \fld{pardisomkl} 后端并要求本地静态
        sequential oneMKL 包完整，
        并强制透传为 \fld{MIPSOLVERS\_BUILD\_LOCAL\_IPOPT}；MIPSolvers 侧仅在以内嵌
        Ipopt 源码构建出 \fld{ipopt\_local} 时置 \fld{MIPSOLVERS\_HAVE\_IPOPT} 并透传
        \srcpath{HACDCPF\_HAVE\_IPOPT=1}（\srcpath{MIPSolvers/cmake/Dependencies.cmake:MIPSOLVERS\_HAVE\_IPOPT}、
        \srcpath{MIPSolvers/CMakeLists.txt:HACDCPF\_HAVE\_IPOPT}）——开启但内嵌
        \texttt{ipopt/} 源码缺失时在 MIPSolvers 侧 \fld{FATAL\_ERROR}（必须用仓库内
        Ipopt 源码，不用系统 Ipopt）。未编译入时适配器返回
        \fld{available=false}、状态 \texttt{"Ipopt not compiled in"}
        （src/optimal\_power\_flow/ac\_opf.cpp:solve\_parity\_with\_ipopt）；
  \item 运行期 \srcpath{HACDCPF\_DISABLE\_IPOPT\_OPF}：设置即强制不可用，
        用于后端隔离与诊断（src/optimal\_power\_flow/ac\_opf.cpp:solve\_parity\_with\_ipopt）。
\end{itemize}

\subsubsection*{与核的衔接方式}
内层求解器为 \fld{Auto} 且关闭目标同伦时，若启用了混合潮流热启动
（\fld{ac\_pf\_warm\_start}），先用 Ipopt 解、其未收敛但给出有限原始点时
再以完整四元组热启动本核精化（src/optimal\_power\_flow/ac\_opf.cpp:solve\_with\_parity\_ipm）；默认 \fld{Auto}
分支为本核先行、未收敛且 \fld{allow\_fallback} 时回退 Ipopt
（src/optimal\_power\_flow/ac\_opf.cpp:solve\_with\_parity\_ipm）；显式 \fld{Ipopt} 后端在 Ipopt 不可用时回落
本核并在后端标签中注明 \texttt{(ipopt unavailable)}（src/optimal\_power\_flow/ac\_opf.cpp:solve\_with\_parity\_ipm）。
如实说明两条口径差异：其一，Ipopt 报
\fld{Maximum\_Iterations\_Exceeded} 但原始可行严格达标、对偶与互补落在
$1000\times$ 容忍带（互补另取 $10^{-2}$ 工程下限）时仍按收敛接受，状态
串保留原终止码（src/optimal\_power\_flow/ac\_opf.cpp:solve\_parity\_with\_ipopt）；其二，适配器不返回
界乘子，驱动层以 $2\times10^{-12}$ 内点下限重构 $\mu,z$ 以满足核状态
格式（src/optimal\_power\_flow/ac\_opf.cpp:solve\_parity\_with\_ipopt），因此 Ipopt 路径的 LMP 不可用
（src/optimal\_power\_flow/ac\_opf.cpp:solve\_with\_parity\_ipm）。

\subsection{调谐参数表}\label{subsec:pipm-params}

\subsubsection*{IPMOptions 字段}
\begin{paramtable}
\fld{verbose} & — & — & true/false & false & 逐迭代调试日志（:IPMOptions，:solve\_primal\_dual\_ipm）\\
\fld{tol\_primal} & $\varepsilon_p$ & — & $10^{-8}$--$10^{-5}$（经验值） & $10^{-6}$ & 原始可行与平稳性判据容忍度；驱动层取 $\max(\fld{feasibility\_tol},10^{-10})$（\srcpath{native\_ipm\_solver.hpp}:IPMOptions，src/optimal\_power\_flow/ac\_opf.cpp:solve\_with\_parity\_ipm）\\
\fld{tol\_dual} & $\varepsilon_d$ & — & $10^{-8}$--$10^{-5}$（经验值） & $10^{-6}$ & 不进收敛判据；仅滤波器平稳守卫 $\max(10\varepsilon_d,1.02\kappa_d)$（\srcpath{native\_ipm\_solver.hpp}:IPMOptions，src/optimal\_power\_flow/parity\_ipm.cpp:solve\_primal\_dual\_ipm）\\
\fld{tol\_complementarity} & $\varepsilon_c$ & — & — & $10^{-6}$ & \textbf{本核未读取}；互补阈值固定为 $\max(10\varepsilon_p,2\times10^{-3})$（\srcpath{native\_ipm\_solver.hpp}:IPMOptions，src/optimal\_power\_flow/parity\_ipm.cpp:solve\_primal\_dual\_ipm）\\
\fld{regularization} & $\delta_0$ & — & — & $10^{-8}$ & \textbf{本核未读取}；正则化梯级硬编码（\srcpath{native\_ipm\_solver.hpp}:IPMOptions，src/optimal\_power\_flow/parity\_ipm.cpp:solve\_primal\_dual\_ipm）\\
\fld{alpha\_max} & $\tau$ & — & 0.9--0.9999（经验值） & 0.95 & fraction-to-boundary 因子；驱动层把 \fld{interior\_fraction} 经 clamp 夹到 $[0.5,\,0.9999]$，默认 0.995（\srcpath{native\_ipm\_solver.hpp}:IPMOptions，src/optimal\_power\_flow/ac\_opf.cpp:solve\_with\_parity\_ipm）\\
\fld{verbose} & — & — & true/false & false & 逐迭代调试日志（\srcpath{native\_ipm\_solver.hpp}:IPMOptions，src/optimal\_power\_flow/parity\_ipm.cpp:solve\_primal\_dual\_ipm）\\
\fld{enable\_phase\_one} & — & — & true/false & true & 无兼容完整状态时运行有界 primal/dual 起点构造\\
\fld{phase\_one\_time\_limit\_ms} & — & ms & $>0$ 协作截止；$\le0$ 不限 & 5000 & 在稀疏分解之间检查，单次分解不可中断\\
\fld{phase\_one\_max\_iterations} & — & — & $\ge0$ & 12 & 约束 Newton 迭代硬上限\\
\fld{phase\_one\_max\_factorizations} & — & — & $\ge0$ & 14 & primal 恢复与一次 dual fit 共用的数值分解硬上限\\
\fld{phase\_one\_max\_backtracks} & — & — & $\ge0$ & 12 & 每次 Newton 步的回溯上限\\
\fld{phase\_one\_dispatch\_dual\_min\_improvement} & — & — & 0--1 & 0.05 & raw 与 normalized dual residual 必须同时达到的改善比例（:IPMOptions，:solve\_primal\_dual\_ipm）\\
\fld{phase\_one\_admission\_mu\_factor} & — & — & $\ge0$ & 1.0 & Phase I 准入要求 violation $\le$ 因子$\times\mu_I$，否则记 \texttt{outside-phase-one-admission}；非有限回退 1.0（:inf\_norm，:elapsed\_ms（lambda））\\
\fld{phase\_one\_primal\_mu\_factor} & — & — & $\ge0$ & 0.1 & 交接 primal 容差 $\max(\varepsilon_p,\text{因子}\times\mu_I)$；非有限回退 0.1（:inf\_norm，:restore\_phase\_one\_primal）\\
\fld{phase\_one\_centrality\_tolerance} & $\eta_c$ & — & $\ge0$ & 0.5 & 中心性门槛（$\max_i|z_i\mu_i/\mu_0-1|$ 判据）；非有限回退 0.5（:ftb，:initialize\_phase\_one\_duals）\\
\fld{phase\_one\_dispatch\_dual\_predictor} & — & — & true/false & false & 实验性零分解连通分量 dispatch-price dual 候选（:ftb，:apply\_obj\_scale\_hessian（lambda））\\
\fld{phase\_one\_dispatch\_dual\_min\_improvement} & — & — & 0--1 & 0.05 & raw 与 normalized dual residual 必须同时达到的改善比例（:ftb，:apply\_obj\_scale\_hessian（lambda））\\
\fld{prepared\_state} & — & — & 指针/null & null & 跨求解保有 KKT 符号分析的所有者；\fld{layout\_signature} 变化时调用方须 \texttt{reset}（:IPMOptions，:solve\_primal\_dual\_ipm）\\
\fld{prepared\_numeric\_refactor} & — & — & true/false & false & 实验性数值重分解请求；当前后端无独立可审计 refactor 对象，fail-closed 报 \texttt{unsupported} 并总是新算数值因子（:IPMOptions，:solve\_primal\_dual\_ipm）\\
\fld{prepared\_numeric\_max\_relative\_drift} & — & — & $\ge0$ & $10^{-3}$ & 数值重分解允许的矩阵值相对漂移上限（:IPMOptions；当前后端不消费，见上行）\\
\fld{prepared\_numeric\_backward\_error\_tolerance} & — & — & $\ge0$ & $10^{-8}$ & 数值重分解的后向误差审计门槛（:IPMOptions；当前后端不消费，见上行）\\
\fld{slack\_start} & $z_0$ & — & — & nullptr & 热启动松弛，须严格为正（:IPMOptions，:solve\_primal\_dual\_ipm）\\
\fld{equality\_dual\_start} & $\lambda_0$ & — & — & nullptr & 热启动等式乘子，与下两项同时通过检查才启用（\srcpath{native\_ipm\_solver.hpp}:IPMOptions，src/optimal\_power\_flow/parity\_ipm.cpp:solve\_primal\_dual\_ipm）\\
\fld{inequality\_dual\_start} & $\mu_0$ & — & — & nullptr & 热启动不等式乘子，须严格为正（\srcpath{native\_ipm\_solver.hpp}:IPMOptions，src/optimal\_power\_flow/parity\_ipm.cpp:solve\_primal\_dual\_ipm）\\
\fld{slack\_start} & $z_0$ & — & — & nullptr & 热启动松弛，须严格为正（\srcpath{native\_ipm\_solver.hpp}:IPMOptions，src/optimal\_power\_flow/parity\_ipm.cpp:solve\_primal\_dual\_ipm）\\
\end{paramtable}

\subsubsection*{环境变量开关}
\begin{paramtable}
\srcpath{HACDCPF\_OPF\_LINEAR\_SOLVER} & — & — & dense/mumps/\allowbreak umfpack/klu/\allowbreak eigen/auto & auto & 强制 KKT 线性后端（旧拼写 \srcpath{HACDCPF\_OPF\_SPARSE\_SOLVER} 兼容）（:linear\_solver\_preference\_env）\\
\srcpath{HACDCPF\_OPF\_KKT\_FORM} & — & — & condensed/\allowbreak augmented & condensed & 牛顿系统形式；condensed 为生产默认，augmented 供刚性模型（:kkt\_form\_preference）\\
\srcpath{HACDCPF\_OPF\_MU\_STRATEGY} & — & — & mehrotra/abo & mehrotra & 中心路径参数策略（ABO 仅增广式）（:solve\_primal\_dual\_ipm）\\
\srcpath{HACDCPF\_OPF\_DELTA\_C} & $\delta_C$ & — & $10^{-10}$--$10^{-6}$（经验值） & $10^{-8}\lVert L_{xx}\rVert$ & 增广式 $(2,2)$ 块正则（绝对值）（:solve\_primal\_dual\_ipm）\\
\srcpath{HACDCPF\_OPF\_INERTIA\_GATE} & $N_{ig}$ & — & 0--20 & 0（关） & 纯 AC KLU 路径每 $N$ 次分解用 MUMPS 复检惯性；默认关闭（效果算例相关，注释 :solve\_primal\_dual\_ipm），且依赖本检出处不可用的 MUMPS（:solve\_primal\_dual\_ipm）\\
\srcpath{HACDCPF\_OPF\_RESCALE\_PER\_ITER} & — & — & 置位/未置位 & 未置位 & 置 1 恢复逐迭代 Ruiz 重平衡（默认冻结）（:factor\_assembled\_kkt）\\
\srcpath{HACDCPF\_OPF\_FILTER\_ALL} & — & — & 置位/未置位 & 未置位 & 对任意规模强制启用残差滤波器（默认仅 $n+m_{eq}\ge512$）（:solve\_primal\_dual\_ipm）\\
\srcpath{HACDCPF\_OPF\_STRICT\_THETA} & — & — & 置位/未置位 & 未置位 & 诊断：仅按纯可行下降接受（:solve\_primal\_dual\_ipm、:eval\_and\_test\_trial（lambda））\\
\srcpath{HACDCPF\_OPF\_NO\_LEGACY} & — & — & 置位/未置位 & 未置位 & 诊断：关闭遗留三轴接受路径（:solve\_primal\_dual\_ipm、:eval\_and\_test\_trial（lambda））\\
\srcpath{HACDCPF\_OPF\_PROF} & — & — & 置位/未置位 & 未置位 & 进程级分解/回代/装配耗时剖析，退出时打印（:IpmProf）\\
\srcpath{HACDCPF\_DISABLE\_IPOPT\_OPF} & — & — & 置位/未置位 & 未置位 & 运行期强制禁用嵌入式 Ipopt（src/optimal\_power\_flow/ac\_opf.cpp:solve\_parity\_with\_ipopt）\\
\end{paramtable}

\subsubsection*{内部硬编码常数}
\begin{paramtable}
\fld{kDenseAutoKktDim} & — & — & — & 512 & $n+m_{eq}$ 不大于此值走稠密 PartialPivLU（:solve\_primal\_dual\_ipm）\\
\fld{kLargeKktFilterThreshold} & — & — & — & 512 & 启用残差滤波器的 KKT 维数下限（:solve\_primal\_dual\_ipm）\\
\fld{kMaxBacktracks} & — & — & — & 14 & 主线搜索最大回溯次数（:solve\_primal\_dual\_ipm）\\
\fld{kBacktrack} & — & — & — & 0.5 & 回溯折半因子（:solve\_primal\_dual\_ipm）\\
\fld{kFilterEta} & $\eta$ & — & — & $10^{-4}$ & 遗留三轴进步余量（:solve\_primal\_dual\_ipm、:eval\_and\_test\_trial（lambda））\\
\fld{kGammaTheta}/\fld{kGammaPhi} & $\gamma_\theta$/$\gamma_\varphi$ & — & — & $10^{-5}$ & 滤波器占优余量（:solve\_primal\_dual\_ipm）\\
\fld{kDelta}\allowbreak /\fld{kSTheta}\allowbreak /\fld{kEtaPhi}\allowbreak /\fld{kThetaMin} & $\delta_s$/$s_\theta$/$\eta_\varphi$/$\theta_{\min}$ & — & — & 1.0/\allowbreak 1.1/\allowbreak $10^{-8}$/\allowbreak $10^{-4}$ & switching 与 Armijo 常数（:solve\_primal\_dual\_ipm）\\
\fld{kRegLevels} & $\delta$ & — & — & 见说明 & 凝集式渐进正则化梯级 $\{0,\,10^{-8},\,10^{-6},\,10^{-4}\}\lVert W\rVert$（:solve\_primal\_dual\_ipm）\\
— & $\delta_W$ & — & — & $10^{-8}\lVert L_{xx}\rVert$ 起、$\times8$、封顶 $10^{-2}\lVert L_{xx}\rVert$ & 增广式惯性校正（仅 MUMPS；本检出处不生效）（:solve\_primal\_dual\_ipm）\\
— & — & — & — & $10^{-12}$ & $z,\mu$ 下限保护（:solve\_primal\_dual\_ipm）与步长下限 $10^{-10}$（:solve\_primal\_dual\_ipm）\\
— & — & — & — & $10^{-7}$ & 空步判定阈值（:solve\_primal\_dual\_ipm）\\
\fld{kBetaGC} & $\beta_{gc}$ & — & — & 0.1 & Gondzio 额外校正目标 $0.1\gamma$（:solve\_primal\_dual\_ipm）\\
\fld{kStagnationWindow} & — & — & — & 20 & 目标停滞早停窗口（:solve\_primal\_dual\_ipm）\\
— & — & — & — & $100\times$ & 可接受级收敛的容忍度倍率（:solve\_primal\_dual\_ipm）\\
— & — & — & — & 5 & Ruiz 平衡最大轮数、$10^{-3}$ 停则（:compute\_ruiz\_scaling）\\
— & — & — & — & 2 & KKT 迭代精化步数（:kkt\_solve\_scaled、:kkt\_solve\_dense）；Gondzio 校正次数上限（:solve\_primal\_dual\_ipm）\\
— & — & — & — & 100 & 目标梯度缩放参考范数 $\lVert\nabla f\rVert_\infty$（:solve\_primal\_dual\_ipm）\\
\end{paramtable}

\subsection{验证与校核}\label{subsec:pipm-verification}

\begin{itemize}
  \item \textbf{单元/回归测试}：\srcpath{tests/test\_opf\_solver\_backends.cpp}
        以 \fld{ACOPFSolverBackend::ParityIPM} 覆盖后端分派、
        \fld{solver\_path} 标记、热启动状态回灌、不可行诊断与
        \fld{profiling} 计数；另覆盖纯 AC case30 和 Hybrid microgrid 的
        primal/dual 证书、零分解预算、兼容状态绕过，以及 case300/case2000
        的单调有界非认证路径。case300 还断言未接受 Phase I 时 Phase II 的
        初始残差与关闭 Phase I 完全一致；
        \srcpath{tests/test\_multiscale\_comprehensive.cpp} 与
        \srcpath{tests/test\_nighttime\_opf.cpp} 亦引用本路径；
  \item \textbf{基准算例}：源码注释记录了 case2000\_acdc（约 2~s）、
        case2869pegase（Release 54 迭代）、case13659 等 PEGASE/大电网
        算例的实测轨迹（:kkt\_form\_preference、:backend\_order），手算 OPF 算例与跨引擎
        对照通道见 \ref{sec:verification}~节；
  \item \textbf{雅可比正确性}：Parity 建模层的解析导数经有限差分核对
        （见 \ref{sec:parityform}~节），本核的 KKT 右端与恢复公式与其
        共享同一组回调（:solve\_primal\_dual\_ipm、:eval\_and\_test\_trial（lambda））。
  \item \textbf{性能口径}：\srcpath{phase\_one\_restoration\_benchmark}
        固定两次预热、五次计时。默认参数下 case30 从
        $3.1061\times10^{-2}$ 恢复到 $4.0305\times10^{-4}$，Phase II 从
        17/16 次迭代/分解降到 16/15，端到端中位数
        $4.458\to4.225$~ms（$1.055\times$）。Hybrid microgrid 的初始违反量
        $0.144>\mu_0$，case300 为 $1.129>\mu_0$，二者均由 admission gate
        以零分解跳过；Phase II 的迭代、分解与最终残差和关闭 Phase I 完全一致，
        结构探针中位成本分别约 $0.008$ 与 $0.124$~ms；对应端到端开/关
        中位数为 $0.235/0.228$~ms 与 $230.058/230.462$~ms。降低 $\mu_0$ 到
        0.03/0.01 会使 case30 也在 admission 外，因此默认 0.1 保留。
\end{itemize}

\subsection{边界与局限}\label{subsec:pipm-limits}

\begin{enumerate}
  \item \textbf{MUMPS 分支为死代码}：\srcpath{HACDCPF\_HAVE\_MUMPS} 在本
        检出处无任何定义点，\srcpath{MumpsSolver} 类不存在于当前
        MIPSolvers 头文件树；增广式的 $\delta_W$ 惯性校正、惯性门复检
        与注释中的 MUMPS 性能数字对当前构建均不适用，混合 AC/DC 问题实际
        由 UMFPACK/KLU/Eigen SparseLU 承接，且在 LU 类后端下惯性校验被
        静默跳过（:solve\_primal\_dual\_ipm）。
  \item \textbf{选项字段未生效}：\fld{IPMOptions::tol\_complementarity}
        与 \fld{regularization} 从未被读取；\fld{tol\_dual} 不进收敛
        判据（判据统一用 $\varepsilon_p$）。调用方按字段语义调参不会
        改变核行为。
  \item \textbf{小系统无全局化}：$n+m_{eq}<512$ 且未置
        \srcpath{HACDCPF\_OPF\_FILTER\_ALL} 时，线搜索无条件接受首个
        有限试步（:eval\_and\_test\_trial（lambda）），发散只能由 \texttt{KKT factorization failed}
        等异常路径兜底。
  \item \textbf{互补容忍度宽松}：互补判据下限 $2\times10^{-3}$（:solve\_primal\_dual\_ipm），
        故“收敛”不保证紧互补；严格 KKT 点需依赖对偶最小二乘认证或
        acceptable 级判据的对偶分量。
  \item \textbf{单线程}：核为单线程实现，剖析器按此假设设计
        （:IpmProf 注释）。
  \item \textbf{凝集式条件数}：默认凝集形式的 $\kappa(W)$ 随障碍比平方
        增长，刚性混合算例可能出现“精确求解但方向错误”的停滞；增广式
        可缓解但需自行开启（:KktForm、:kkt\_form\_preference）。
  \item \textbf{Ipopt 路径的平台限制与口径差异}：嵌入式 Ipopt 的可用性取决于
        MIPSolvers 检出是否含可构建的内嵌 \texttt{ipopt/} 源码（见本章门控节，
        默认仅 Apple 开启）；其 acceptable 判据（$1000\times$）
        比本核（$100\times$）宽松，且不返回约束乘子，LMP 在该路径不可用
        （src/optimal\_power\_flow/ac\_opf.cpp:solve\_parity\_with\_ipopt）。
  \item \textbf{悬置声明}：\srcpath{parity\_ipm\_solver.hpp:solve\_ac\_opf\_parity} 的
        \srcpath{solve\_ac\_opf\_parity} 无定义，调用入口实为
        \srcpath{solve\_ac\_opf} 内的静态驱动（src/optimal\_power\_flow/ac\_opf.cpp:solve\_with\_parity\_ipm）。
  \item 恢复阶段与中心性恢复仅存在于增广形式；默认凝集形式在滤波器
        全拒时直接报 \texttt{filter line search failed}（:solve\_primal\_dual\_ipm）。
\end{enumerate}

% 无功优化与 OLTC 离散档位（RPO）
% 本章文件由 \input 引入，不含导言区。
\section{无功优化与 OLTC 离散档位（RPO）}\label{sec:rpo}

本章给出无功优化（Reactive Power Optimization, RPO）模块的技术口径：功能定位与
交流 OPF 的分工、双层数学模型、离散档位搜索算法、公开接口的全部选项与结果字段、
验证与校核通道及已知局限。内容依据
\srcpath{src/optimal\_power\_flow/reactive\_power\_opt.cpp}（1388 行）与
\srcpath{include/hacdcpf/optimal\_power\_flow/reactive\_power\_opt.hpp}（307 行）
逐行整理，验证素材取自
\srcpath{docs/archive/audits/reactive\_power\_optimization\_validation.md} 并与
\srcpath{tests/} 中的自动测试逐条核对。

\subsection{功能定位与适用范围}

\compmeta{RPOOptions / RPOResult / solve\_rpo}{\srcpath{include/hacdcpf/optimal\_power\_flow/reactive\_power\_opt.hpp}}

\subsubsection*{在 OPF 家族中的位置}
RPO 解决的是电压/无功协调优化问题：以母线电压偏差与系统有功网损为目标，协调
两类手段——\textbf{连续无功/电压手段}（发电机无功 $Q_g$、可控新能源与光伏逆变器
无功、储能无功、VSC 交流无功、母线电压幅值）与\textbf{离散调压设备}（两绕组
有载调压变压器 OLTC 档位、可投切并联补偿步数）。平台将该问题建模为混合整数
非线性规划（MINLP），并采用“外层离散坐标搜索 + 内层连续 AC/DC OPF 补全”的
双层结构求解（\srcpath{include/hacdcpf/optimal\_power\_flow/reactive\_power\_opt.hpp}:293--304，
函数 \srcpath{solve\_rpo} 的文档注释）。

与 AC OPF 各章的分工如下：RPO 不重新实现连续优化，而是把每一个离散配置
$(t,s)$ 的连续补全委托给通用 AC OPF 求解器 \srcpath{solve\_ac\_opf}
（\srcpath{src/optimal\_power\_flow/reactive\_power\_opt.cpp}:366，函数
\srcpath{solve\_opf\_inplace} 内调用）。内层默认后端为 Parity 全空间 IPM
（见第~\ref{sec:parityipm}~节），建模层为 Parity formulation
（见第~\ref{sec:parityform}~节），备选后端为嵌入式 Ipopt
（见第~\ref{sec:acopfnative}~节）；RPO 特有的部分只有离散变量层、灵敏度
排序坐标搜索与网损台账。因此本章只展开离散层与目标口径，连续层的潮流方程、
约束装配与 IPM 细节见前述各节。

\subsubsection*{最优性口径（诚实声明）}
求解器返回的是\textbf{经连续 OPF 补全的可行局部最优点}，不是全局 MINLP 证书
（\srcpath{include/hacdcpf/optimal\_power\_flow/reactive\_power\_opt.hpp}:301--304）。
结果中 \fld{globally\_certified} 恒为 false、\fld{optimality\_gap\_available}
恒为 false、\fld{gap} 恒为占位值 1.0
（\srcpath{include/hacdcpf/optimal\_power\_flow/reactive\_power\_opt.hpp}:246--255；
\srcpath{src/optimal\_power\_flow/reactive\_power\_opt.cpp}:1335），
并且每次求解都在 \fld{model\_limitations} 中显式写入该声明
（\srcpath{src/optimal\_power\_flow/reactive\_power\_opt.cpp}:1301--1302）。
算法标识字符串固定为
\srcpath{discrete\_coordinate\_search\_with\_ac\_opf}
（\srcpath{include/hacdcpf/optimal\_power\_flow/reactive\_power\_opt.hpp}:253）。

\subsubsection*{元件利用审计}
按验证文档的口径，系统元件在 RPO 中分四类处理
（\srcpath{docs/archive/audits/reactive\_power\_optimization\_validation.md} 第 3 节；
GUI 覆盖审计实现于 \srcpath{tests/run\_gui\_server.cpp}:19686--19799）：
\begin{enumerate}
  \item \textbf{直接离散优化}：两绕组 OLTC（档位）、可投切并联补偿（步数）；
  \item \textbf{连续协同优化}：机组 $P_g/Q_g$、可控新能源/光伏 $P/Q$、储能无功、
        柔性负荷、VSC/DC/DC/能量路由器端口等，全部由内层 AC/DC OPF 处理；
  \item \textbf{固定建模}：负荷、固定电源、固定变比变压器/固定并联元件及当前
        开关拓扑；
  \item \textbf{部分或未覆盖}：三绕组变压器分接头、独立调压控制器
        （\fld{regulator\_controls}）、单体三相不平衡 RPO，以及当前尚未完整
        回写到 authored-space PF 的富混合 AC/DC 换流器调度。
\end{enumerate}
需要强调：“模型中出现”不等于“作为决策变量优化”。平台没有独立的 SVG 元件
类型：静态无功补偿若以固定电纳给出，则作为网络参数固定建模；若以
\fld{Shunt::switchable} 可投切步进形式给出，则进入离散决策向量。连续无功源
（发电机、可控逆变器、储能、VSC）的无功上下限由元件数据给定，在内层 OPF
中作为界约束处理（见第~\ref{sec:parityform}~节）。

\subsection{数学模型}

\subsubsection*{双层结构}
记离散决策向量为
\begin{align*}
y &= (t_1,\ldots,t_{N_t},\,s_1,\ldots,s_{N_s})\in\mathbb{Z}^{N_t+N_s}
\end{align*}
其中 $t_k$ 为第 $k$ 台可选 OLTC 的档位，$s_j$ 为第 $j$ 台可投切并联补偿的
步数。给定 $y$ 后，连续运行点
\begin{align*}
x(y) &= \arg\min_{x\in\mathcal{F}(y)}\; f_{\mathrm{inner}}(x)
\end{align*}
由内层混合 AC/DC OPF 解出，$\mathcal{F}(y)$ 为该离散配置下的连续可行集
（潮流平衡与各类运行约束，见下文）。外层问题为
\begin{align*}
\min_{y\in\mathcal{Y}}\; f_{\mathrm{RPO}}\bigl(x(y)\bigr)
\end{align*}
其中 $\mathcal{Y}$ 为离散可行集（档位窗口与步数界）。代码按此结构组织：
外层是 \srcpath{solve\_rpo} 的坐标搜索
（\srcpath{src/optimal\_power\_flow/reactive\_power\_opt.cpp}:1015--1293），
内层是 \srcpath{solve\_opf\_inplace} 对 \srcpath{solve\_ac\_opf} 的封装
（\srcpath{src/optimal\_power\_flow/reactive\_power\_opt.cpp}:312--403）。

\subsubsection*{RPO 排名目标（外层评估口径）}
RPO 排名目标\textbf{始终由返回的潮流事后评估}，与内层 NLP 实际最小化的目标
解耦（\srcpath{include/hacdcpf/optimal\_power\_flow/reactive\_power\_opt.hpp}:19--24）。
评估函数 \srcpath{rpo\_objective} 按所选目标叠加两项
（\srcpath{src/optimal\_power\_flow/reactive\_power\_opt.cpp}:269--285）：
\begin{align*}
f_V &= w_V\sum_{i\in\mathcal{B}_{ac}} \bigl(V_{m,i}-V^{*}\bigr)^{2}
   && \text{（MinVoltageDeviation / Combined）}\\
f_L &= w_L\,P_{\mathrm{loss}}
   && \text{（MinActiveLoss / Combined）}\\
f_{\mathrm{RPO}} &=
   \begin{cases}
     f_V & \text{MinVoltageDeviation（默认）}\\
     f_L & \text{MinActiveLoss}\\
     f_V + f_L & \text{Combined}
   \end{cases}
\end{align*}
其中 $V_{m,i}$ 取内层 OPF 结果向量 \fld{vm}（标幺值），$V^{*}$ 为
\fld{v\_target}（默认 1.0 pu），$w_V$、$w_L$ 为 \fld{vdev\_weight} 与
\fld{loss\_weight}（默认均为 1.0）。目标枚举 \srcpath{RPOObjective} 定义于
\srcpath{include/hacdcpf/optimal\_power\_flow/reactive\_power\_opt.hpp}:13--17。
代码中\textbf{没有}“档位动作代价”项：离散动作不直接进入目标，仅通过改善
$f_V$/$f_L$ 体现价值；\fld{gap\_tol} 只是坐标下降判定“本轮有改进”的阈值
（见算法流程），并非动作成本。

网损 $P_{\mathrm{loss}}$ 采用\textbf{全系统电源与物理需求的有功差额}台账
（\srcpath{src/optimal\_power\_flow/reactive\_power\_opt.cpp}:189--266，函数
\srcpath{compute\_loss}）：
\begin{align*}
P_{\mathrm{loss}} &= \sum_{\text{全部电源}} P - \sum_{\text{全部需求}} P
\end{align*}
电源侧计列：常规机组 \fld{pg\_mw}、外部电网 \fld{external\_grid\_p\_mw}、
可弃新能源 \fld{pren\_mw}、储能放电 \fld{pstor\_mw}、在运静态发电机
（\fld{p\_mw}$\times$\fld{scaling}）、不可弃新能源、不可控光伏、直流侧
静态电源/光伏、移动储能、虚拟电厂，并减去微网交换功率
（\srcpath{src/optimal\_power\_flow/reactive\_power\_opt.cpp}:193--231）；
需求侧计列：母线挂接负荷、交流负荷、不对称负荷三相总量、柔性负荷
（优先取 OPF 调度 \fld{pflex\_mw}）、充电站与电动机折算有功、直流母线与
直流负荷，并减去切负荷量 \fld{dpd\_mw} 的非负部分
（\srcpath{src/optimal\_power\_flow/reactive\_power\_opt.cpp}:233--264）。
VSC、DC/DC 与能量路由器端口属\textbf{内部传输}，不计入电源台账
（同文件:compute\_loss 注释），其损耗已隐含在差额中。当
\fld{loss\_weight}$=1$ 时 MinActiveLoss 的 \fld{objective} 在数值上等于
物理网损 MW（单元测试断言此性质，
\srcpath{tests/test\_reactive\_power\_opt.cpp}:41）。

\subsubsection*{内层连续 NLP 目标（MatchRPO 与 LossEconomic）}
内层 NLP 目标由枚举 \srcpath{RPOInnerObjective} 选择
（\srcpath{include/hacdcpf/optimal\_power\_flow/reactive\_power\_opt.hpp}:25--34），
与外层排名目标解耦的原因写在头文件注释中：精简的物理目标（电压偏差/网损）
在大型僵硬电网上会使 Parity IPM 停滞，而单位边际成本经济调度继承经济路径的
稳健机制且收敛（同文件:RPOInnerObjective）。

\begin{itemize}
  \item \textbf{MatchRPO（默认）}：内层直接最小化 RPO 目标。映射关系为
        MinVoltageDeviation $\to$ \srcpath{ACOPFObjective::\allowbreak VoltageDeviation}，
        Combined $\to$ \srcpath{ACOPFObjective::\allowbreak VoltageDeviation\allowbreak AndLoss}
        （\srcpath{src/optimal\_power\_flow/reactive\_power\_opt.cpp}:344--348；
        目标枚举见 \srcpath{include/hacdcpf/optimal\_power\_flow/opf\_options.hpp}:39--43）。
        Parity 装配的目标为
        \begin{align*}
        f_{\mathrm{inner}} &= w_V\sum_{i}\bigl(V_{m,i}-V^{*}\bigr)^{2}
          + w_L\,P_{\mathrm{loss}}^{\mathrm{var}}
          + 10^{-6}\sum_{k}\bigl(c_{2,k}P_{g,k}^{2}+c_{1,k}P_{g,k}\bigr)
        \end{align*}
        其中 $P_{\mathrm{loss}}^{\mathrm{var}}$ 只含变量部分
        （$P_g$、$P_{\mathrm{ren}}$、$P_{\mathrm{stor}}$、直流储能、减柔性负荷、
        加切负荷），固定注入/需求为常数不影响极小解，报告网损时再补回
        （\srcpath{src/optimal\_power\_flow/parity\_formulation.cpp}:1204--1237）。
        末项是 $10^{-6}$ 权重的发电经济性\textbf{择优项}
        （\srcpath{kRpoEconomicTieBreak}），仅用于在等价有功调度方向中选定一个
        并改善 KKT 数值条件，不作用于 $V_m/Q$ 变量、不替代主 RPO 目标
        （同文件:objective）。
  \item \textbf{LossEconomic}：内层改为单位边际成本经济调度
        （\srcpath{ACOPFObjective::Economic}）。理论依据是关键等价关系：
        功率平衡固定需求下
        \begin{align*}
        P_{\mathrm{loss}} &= \sum P_{g,\text{全部}} - \sum P_{d,\text{固定}}
        \quad\Longrightarrow\quad
        \min P_{\mathrm{loss}} \equiv \min \sum P_g
        \end{align*}
        即全体常规机组成本改写为 $c_1=1,\;c_2=c_0=0$ 的经济调度
        （\srcpath{src/optimal\_power\_flow/reactive\_power\_opt.cpp}:687--711，
        函数 \srcpath{apply\_unit\_costs}；目标映射同文件:solve\_opf\_inplace）。
        报告网损仍由 \srcpath{compute\_loss} 按潮流重算，与成本数值无关。
        \textbf{MinActiveLoss 恒走此路径}；MinVoltageDeviation/Combined 在
        MatchRPO 基线不收敛且 \fld{robust\_inner\_on\_failure} 为真时自动升级
        到此路径（同文件:run\_baseline（lambda）），或交流母线数达到
        \fld{robust\_inner\_min\_buses}（默认 1000）时直接以该路径起步
        （同文件:apply\_unit\_costs（lambda））。
\end{itemize}

\subsubsection*{离散变量与档位约束}
OLTC 档位 $t_k$ 取整数值，铭牌范围为 $[t_k^{\min},t_k^{\max}]$
（字段 \fld{tap\_min}/\fld{tap\_max}），相对中性档 \fld{tap\_neutral}
（$t_{\mathrm{neutral}}$）的\textbf{相对变比}为
（\srcpath{src/optimal\_power\_flow/reactive\_power\_opt.cpp}:80--82，
函数 \srcpath{tap\_ratio}）：
\begin{align*}
r(t) &= 1 + (t - t_{\mathrm{neutral}})\cdot \frac{\Delta_{\%}}{100}
\end{align*}
其中 $\Delta_{\%}$ 为 \fld{tap\_step\_percent}（每档步长百分数）。
\fld{max\_tap\_move}（记 $m_{\mathrm{tap}}$）把每台已选 OLTC 的本次优化窗口
限制在当前档位 $t_k^{0}$ 上下 $m_{\mathrm{tap}}$ 档以内
（\srcpath{src/optimal\_power\_flow/reactive\_power\_opt.cpp}:541--547）：
\begin{align*}
t_k &\in [t_k^{\min},t_k^{\max}]\cap
        [\,t_k^{0}-m_{\mathrm{tap}},\;t_k^{0}+m_{\mathrm{tap}}\,]
\end{align*}
$m_{\mathrm{tap}}=-1$ 暴露完整铭牌范围（库默认值），$m_{\mathrm{tap}}=0$
保持所有 OLTC 固定；GUI 显式默认 $m_{\mathrm{tap}}=2$
（\srcpath{include/hacdcpf/optimal\_power\_flow/reactive\_power\_opt.hpp}:137--141；
\srcpath{tests/run\_gui\_server.cpp}:19410--19411）。
\fld{restrict\_tap\_indices} 与 \fld{enabled\_tap\_indices} 提供逐台参与
控制：开启后只有列入 vector position 的设备进入离散决策向量
（\srcpath{include/hacdcpf/optimal\_power\_flow/reactive\_power\_opt.hpp}:143--147；
\srcpath{src/optimal\_power\_flow/reactive\_power\_opt.cpp}:529--534）。

对带 \fld{source\_branch\_idx} 的变压器（MATPOWER 导入链），档位动作按
相对变比同步缩放原始交流支路的电气变比 $\tau$（字段 \fld{branch.tap}；
\srcpath{src/optimal\_power\_flow/reactive\_power\_opt.cpp}:94--104，函数
\srcpath{set\_transformer\_tap\_position}）：
\begin{align*}
\tau &\leftarrow \tau\cdot \frac{r(t^{\mathrm{new}})}{r(t^{\mathrm{old}})}
\end{align*}
从而保证 RPO 内层、独立 OPF 复算与 PF 回放使用同一调压动作；“相对变比”与
“实际支路 tap（电气变比）”在结果中分列（\fld{ratio\_after} 与
\fld{electrical\_tap\_after}，后者按
$\tau_{\mathrm{before}}\cdot r_{\mathrm{after}}/r_{\mathrm{before}}$
换算，同文件:solve\_rpo）。

可投切并联补偿步数 $s_j\in\{0,1,\ldots,n_j\}$（$n_j$ 为 \fld{n\_steps}），
当前电纳按（\srcpath{src/optimal\_power\_flow/reactive\_power\_opt.cpp}:302--307，
函数 \srcpath{apply\_settings}）：
\begin{align*}
B_{s,j} &= b_{\mathrm{step}}\cdot s_j
\end{align*}
其中 $b_{\mathrm{step}}$ 为 \fld{bs\_per\_step}；结果写入 \fld{bs\_mvar}
（MVAr）。

设备入选资格由排除原因函数给出（空串即入选）：变压器要求投运、铭牌上限
大于下限（\fld{tap\_max} $>$ \fld{tap\_min}）、\fld{tap\_step\_percent}
有限且为正、
当前档位在铭牌范围内（\srcpath{src/optimal\_power\_flow/reactive\_power\_opt.cpp}:108--116，
函数 \srcpath{tap\_exclusion\_reason}，排除原因串为
\fld{not\_in\_service}/\allowbreak\fld{fixed\_or\_invalid\_tap\_range}/\allowbreak
\fld{non\_positive\_tap\_step}/\allowbreak\fld{current\_tap\_out\_of\_range}）；
并联补偿要求投运、可投切（\fld{switchable}）、\fld{n\_steps} 大于 1、
\fld{bs\_per\_step} 的绝对值大于 $10^{-12}$、当前步数位于 0 至
\fld{n\_steps} 之间
（同文件:shunt\_exclusion\_reason，函数 \srcpath{shunt\_exclusion\_reason}）。
MATPOWER 支路只提供当前连续变比，不提供 OLTC 可控标志、档位上下限或每档
步长，因此导入器保留精确变比但\textbf{不}伪造 OLTC：导入设备的
\fld{tap\_min}=\fld{tap\_max}=\fld{tap\_pos}=0、\fld{tap\_step\_percent}=0，
按固定变比处理（单元测试断言，
\srcpath{tests/test\_reactive\_power\_opt.cpp}:184--211）。

\subsubsection*{连续约束与切负荷口径}
连续可行集 $\mathcal{F}(y)$ 由内层 AC/DC OPF 的约束族构成：交流/直流节点
功率平衡、母线电压幅值上下限、机组 $P/Q$ 上下限、支路热稳限值、VSC 容量圆、
交直流电流限值、调制比窗口与 DC/DC 占空比约束（装配细节见
第~\ref{sec:parityform}~节）。RPO 通过四个开关逐项启停后四类约束：
\fld{enforce\_branch\_limits}、\fld{enforce\_converter\_capacity}、
\fld{enforce\_converter\_current\_limits}、
\fld{enforce\_converter\_modulation\_limits}（默认全为 true，
\srcpath{include/hacdcpf/optimal\_power\_flow/reactive\_power\_opt.hpp}:151--154；
透传于 \srcpath{src/optimal\_power\_flow/reactive\_power\_opt.cpp}:358--363）。

切负荷口径需要精确说明：Parity 装配按目标类型决定是否生成切负荷变量——
仅当 \fld{objective} 为 \srcpath{ACOPFObjective::Economic} 时
\fld{load\_shedding} 取真
（\srcpath{src/optimal\_power\_flow/ac\_opf.cpp}:2526--2539，注释明确
“RPO objectives must not improve voltage/loss by dropping demand”）。
即电压偏差/网损/组合目标下\textbf{没有}切负荷变量，RPO 不能通过少供电
获得虚假的电压或网损改善；而 LossEconomic 内层（Economic 目标）保留
VOLL 兜底的切负荷作为可行性逃生阀。这与验证文档“RPO 模式关闭负荷削减”
的表述一致，但需注意该表述对经济内层例外。

\subsubsection*{电压界}
母线电压幅值上下限由内层 OPF 按元件数据（母线 \fld{vm\_min\_pu}/
\fld{vm\_max\_pu}）作为硬约束处理（见第~\ref{sec:parityform}~节），RPO
本身不另设电压界参数；\fld{v\_target} 只是偏差惩罚的目标值而非约束。
GUI 展示层固定以 $[0.95,1.05]$ pu 作为电压合格带的显示常量
（\srcpath{tests/run\_gui\_server.cpp}:19634--19635，输出字段
\fld{v\_min}/\fld{v\_max}），与求解约束无关。结果指标
\fld{max\_vdev\_before}/\fld{max\_vdev\_after} 为
$\max_i |V_{m,i}-V^{*}|$
（\srcpath{src/optimal\_power\_flow/reactive\_power\_opt.cpp}:287--291，
函数 \srcpath{max\_voltage\_deviation}）。

\subsection{算法流程}

\subsubsection*{总览}
求解入口 \srcpath{solve\_rpo}
（\srcpath{src/optimal\_power\_flow/reactive\_power\_opt.cpp}:621--1386）
按图~\ref{fig:rpo-flow} 组织：先做控制清单审计与基线连续 OPF，再执行
四个阶段的离散搜索，最后归因结果并写入局限声明。阶段划分与文件头注释
一致（同文件）。

\begin{figure}[H]
\centering
\begin{tikzpicture}[
  font=\footnotesize,
  box/.style={draw, rectangle, align=center, minimum height=0.85cm,
              inner sep=4pt},
  note/.style={draw, rectangle, dashed, align=center, inner sep=4pt},
  arr/.style={-{Stealth[length=2.2mm]}, thick}]
  \node[box] (input) {系统 + \fld{RPOOptions}};
  \node[box, below=0.5cm of input] (inv) {控制清单审计\\\srcpath{inspect\_rpo\_controls}};
  \node[box, below=0.5cm of inv] (base) {基线连续 OPF\\（经济种子/同伦/稳健升级）};
  \node[box, below=0.5cm of base] (p0) {Phase 0：$\pm1$ 扰动灵敏度排序};
  \node[box, below=0.5cm of p0] (p1) {Phase 1：逐变量整数三分搜索};
  \node[box, below=0.5cm of p1] (p2) {Phase 2：坐标下降（至多 3 轮）};
  \node[box, below=0.5cm of p2] (p3) {Phase 3：变量对 $\pm1$ 邻域精化};
  \node[box, below=0.5cm of p3] (out) {结果归因 + \fld{model\_limitations}};
  \node[note, right=1.4cm of p1] (eval) {每个候选 = 一次完整 AC/DC OPF\\解缓存、暖启动链、失败快速拒绝\\时间/评估预算只能在求解之间检查};
  \draw[arr] (input) -- (inv);
  \draw[arr] (inv) -- (base);
  \draw[arr] (base) -- (p0);
  \draw[arr] (p0) -- (p1);
  \draw[arr] (p1) -- (p2);
  \draw[arr] (p2) -- (p3);
  \draw[arr] (p3) -- (out);
  \draw[dashed] (eval.west) -- (p1.east);
\end{tikzpicture}
\caption{RPO 求解流程。离散搜索各阶段共享同一候选评估管线（内层 AC/DC OPF）。}
\label{fig:rpo-flow}
\end{figure}

\subsubsection*{基线求解与经济种子}
对交流母线数 $\ge 100$ 的系统，先以 Parity IPM 求一个高精度经济 OPF 种子：
迭代预算为 $\max(200,\min(K_{\mathrm{ipm}}^{\max},800))$（$K_{\mathrm{ipm}}^{\max}$
即 \fld{max\_ipm\_iter}），平稳性容差收紧到 \fld{stationarity\_tol} 与
$10^{-6}$ 的较小者；大电网（达到
\fld{robust\_inner\_min\_buses}）直接启用经济目标同伦，直接种子不收敛且
\fld{seed\_homotopy\_on\_failure} 为真时再同伦重解一次
（\srcpath{src/optimal\_power\_flow/reactive\_power\_opt.cpp}:737--782）。
随后在当前离散配置下解基线 OPF（“优化前”状态），以种子结果暖启动
（同文件:run\_baseline（lambda））。基线接受判据是\textbf{实测容差收敛}：
\begin{align*}
\text{接受} &\iff \mathrm{conv}\;\lor\;
  \bigl(\mathrm{viol}\le\epsilon_p\;\land\;\mathrm{stat}\le\epsilon_d\bigr)
\end{align*}
其中 $\mathrm{conv}$、$\mathrm{viol}$、$\mathrm{stat}$ 分别为内层结果的
\fld{converged}、\fld{max\_constraint\_violation} 与 \fld{max\_stationarity}，
$\epsilon_p$、$\epsilon_d$ 为 \fld{ipm\_tol} 与
\fld{stationarity\_tol}（同文件:run\_baseline（lambda））；按实测容差接受的基线在
\fld{status} 中附加说明（同文件:run\_baseline（lambda））。MatchRPO 基线失败时的稳健升级见内层目标小节
（同文件:run\_baseline（lambda））；升级成功同样写入 \fld{status}
（同文件:run\_baseline（lambda））。基线失败则直接返回，仅写入局限声明
（同文件:run\_baseline（lambda））。

\subsubsection*{候选评估管线}
所有搜索阶段共用评估函数 \srcpath{evaluate}
（\srcpath{src/optimal\_power\_flow/reactive\_power\_opt.cpp}:898--941），
要点如下。
\begin{enumerate}
  \item \textbf{解缓存}：以 $(t,s)$ 拼接向量为键的 \fld{eval\_cache} 避免
        重复 NLP 评估，基线先行入库（同文件:run\_baseline（lambda））。
  \item \textbf{预算检查}：每次求解前检查墙钟时间 \fld{time\_limit\_sec} 与
        累计内层求解次数 \fld{max\_nodes}；预算只能在\textbf{两次求解之间}
        检查，无法中断一次进行中的 IPM（同文件:budget\_ok（lambda）；
        另见 \srcpath{include/hacdcpf/optimal\_power\_flow/reactive\_power\_opt.hpp}:110--117
        注释）。
  \item \textbf{原位修改}：候选配置直接写入可变系统副本 \fld{mut\_sys}，
        求解后还原原始档位，避免深拷贝（
        \srcpath{src/optimal\_power\_flow/reactive\_power\_opt.cpp}:294--308、
        400--401）。
  \item \textbf{暖启动链}：候选以当前在位可行点 \fld{incumbent\_result}
        暖启动（同文件:evaluate（lambda））；实测 3000 节点电网上好的候选 2--7 次
        迭代即收敛（\srcpath{docs/archive/audits/reactive\_power\_optimization\_validation.md}
        第 5.3 节）。
  \item \textbf{失败快速拒绝}：纯交流电网上搜索评估的迭代上限取
        \fld{search\_ipm\_iter}（默认 200），超限候选按不可行拒绝
        （目标取 $10^{30}$），且不再启用第二后端；混合 AC/DC 电网保留完整
        预算与 Ipopt 第二腿——其电压偏差渐近行为合理需要它
        （\srcpath{src/optimal\_power\_flow/reactive\_power\_opt.cpp}:909--920；
        纯交流/混合判定同文件:solve\_rpo）。第二腿在另一后端上重解并按实测
        质量取优（同文件:solve\_opf\_inplace）。
  \item \textbf{并行波评估}：\fld{num\_threads}$>1$ 时一批独立候选按
        Jacobi 波形并行求解（动态工作窃取，每线程一次深拷贝），波间推进
        在位点（\srcpath{src/optimal\_power\_flow/reactive\_power\_opt.cpp}:418--486、
        952--1001）。混合 AC/DC 强制串行：MUMPS 路径由进程级互斥锁串行化，
        多线程只增开销；纯交流用 KLU，线程安全（同文件:solve\_rpo；
        \srcpath{include/hacdcpf/optimal\_power\_flow/reactive\_power\_opt.hpp}:84--96）。
\end{enumerate}

\subsubsection*{Phase 0：灵敏度排序}
每个离散变量在基线处做 $\pm1$ 扰动（每变量至多 2 次求解），取
$|\Delta f_{\mathrm{RPO}}|$ 作为灵敏度，按降序确定后续所有阶段的处理次序
——最有影响的变量先优化（\srcpath{src/optimal\_power\_flow/reactive\_power\_opt.cpp}:1015--1092）。
同时按有限差分梯度
$g_k=(f(y_k+d)-f(y_k))/d$ 记录一维局部灵敏度面（同文件:solve\_rpo）。并行实现为一次 Jacobi 波（同文件:evaluate\_wave（lambda）），串行实现为
Gauss--Seidel 顺序（同文件:solve\_rpo）。

\subsubsection*{Phase 1：逐变量整数三分搜索}
对每个变量（按灵敏度序），利用 RPO 目标对单变量的\textbf{近似单峰性}做整数
网格三分搜索：区间 $[lo,hi]$ 内取
$m_1=lo+\lfloor(hi-lo)/3\rfloor$、$m_2=hi-\lfloor(hi-lo)/3\rfloor$，
比较 $f(m_1)$ 与 $f(m_2)$ 后收缩区间，直到 $hi-lo\le 2$，再对残余小窗口
穷举（\srcpath{src/optimal\_power\_flow/reactive\_power\_opt.cpp}:1094--1167）。
单变量评估量由 $O(R)$ 降为 $O(\log R)$（$R$ 为档位数）。该阶段刻意保持
Gauss--Seidel 结构：跨变量 Jacobi 化实测会丢失定向优化能力、降低解质量
（同文件:solve\_rpo 注释）。近似单峰性是启发假设，非凸问题上不保证成立，
这也是本阶段之后仍需 Phase 2/3 的原因。

\subsubsection*{Phase 2：坐标下降}
对每个变量（仍按灵敏度序）穷举其全部可取值，遇更优即更新在位点；整轮
无超过 \fld{gap\_tol}（默认 $10^{-4}$）的改进即停止，至多
\srcpath{kMaxCDCycles}${}=3$ 轮
（\srcpath{src/optimal\_power\_flow/reactive\_power\_opt.cpp}:1171--1234，
轮数上限同文件:solve\_rpo）。局部灵敏度面在此阶段继续累积但\textbf{从不用于
剪枝}——注释明确：非凸 MINLP 上局部线性面不是全局有效下界
（同文件:LocalSensitivityPlane 注释；灵敏度面结构
\srcpath{LocalSensitivityPlane} 见同文件:LocalSensitivityPlane）。

\subsubsection*{Phase 3：变量对邻域精化}
对所有变量对 $(v_i,v_j)$ 检查围绕在位点的 $\pm1$ 联合扰动（每对至多 4 个
组合），捕捉坐标下降可能遗漏的双变量交互
（\srcpath{src/optimal\_power\_flow/reactive\_power\_opt.cpp}:1236--1293）。
并行为一次 Jacobi 精化波（同文件:evaluate\_wave（lambda）），串行为双重循环
（同文件:solve\_rpo）。

\subsubsection*{终止与结果归因}
搜索结束后按预算消耗置 \fld{terminated\_by\_time\_limit}/
\fld{terminated\_by\_evaluation\_limit}，并向 \fld{model\_limitations}
追加相应说明（\srcpath{src/optimal\_power\_flow/reactive\_power\_opt.cpp}:1295--1315）。
在位点可行（目标 $<10^{20}$）时 \fld{converged}=true，状态串区分
“Feasible local optimum (heuristic search)”/“Feasible incumbent at time
limit”/“Feasible incumbent at evaluation limit”
（同文件:solve\_rpo）；无可行配置时返回基线向量并标注
“No feasible configuration found”（同文件:solve\_rpo）。
逐设备明细（档位/步数、相对变比、电气变比、电纳前后值）于
同文件:solve\_rpo 装配。无可调离散设备时直接以基线为解并如实声明
“结果只是连续 OPF 基线，非全局 MINLP 证书”（同文件:run\_baseline（lambda））。
\fld{bc\_stats} 为遗留统计信封：\fld{nodes\_explored} 记候选评估数、
\fld{lp\_solves} 实际记 NLP 求解次数、\fld{cuts\_added} 记灵敏度面数量，
\textbf{不是}分支定界或全局间隙证书（同文件:solve\_rpo；
\srcpath{include/hacdcpf/optimal\_power\_flow/reactive\_power\_opt.hpp}:284--286）。
环境变量 \srcpath{HACDCPF\_RPO\_PROF}${}=1$ 可打开分阶段墙钟计时输出
（\srcpath{src/optimal\_power\_flow/reactive\_power\_opt.cpp}:42--54）。

\subsection{公开接口与参数}

\subsubsection*{接口函数}
模块对外暴露三个函数（命名空间 \srcpath{hacdcpf::opf}）：
\begin{itemize}
  \item \texttt{solve\_rpo(sys, opt)}：求解 RPO（MINLP），返回
        \srcpath{RPOResult}（声明与决策变量说明见
        \srcpath{include/hacdcpf/optimal\_power\_flow/reactive\_power\_opt.hpp}:293--305；
        实现 \srcpath{src/optimal\_power\_flow/reactive\_power\_opt.cpp}:621--1386）；
  \item \texttt{inspect\_rpo\_controls(sys, opt)}：逐设备审计全部候选离散
        控制，含被排除设备及确切排除原因；求解器与 GUI 输入预览共用这一
        入选合同（\srcpath{include/hacdcpf/optimal\_power\_flow/reactive\_power\_opt.hpp}:234--237；
        实现 \srcpath{src/optimal\_power\_flow/reactive\_power\_opt.cpp}:512--601）；
  \item \texttt{apply\_rpo\_discrete\_solution(sys, result)}：把解出的离散
        点写回 authored 系统（含联动 MATPOWER 支路变比与并联电纳），供独立
        OPF/PF 回放（\srcpath{include/hacdcpf/optimal\_power\_flow/reactive\_power\_opt.hpp}:239--242；
        实现 \srcpath{src/optimal\_power\_flow/reactive\_power\_opt.cpp}:603--619）。
\end{itemize}
GUI 通过 \srcpath{/api/session/rpo\_inputs} 与 \srcpath{/api/session/run\_rpo}
两个路由使用该接口；\fld{relax\_only} 请求把 \fld{max\_nodes} 压为 1，只做
基线（根松弛）评估（\srcpath{tests/run\_gui\_server.cpp}:19393--19430）。

\subsubsection*{求解选项（RPOOptions）}
\begin{paramtable}
\fld{objective} & — & — & MinVoltage\allowbreak Deviation/\allowbreak MinActiveLoss/\allowbreak Combined & MinVoltage\allowbreak Deviation & RPO 排名目标（外层评估口径；\srcpath{RPOObjective}）\\
\fld{inner\_nlp\_objective} & — & — & MatchRPO/\allowbreak LossEconomic & MatchRPO & 内层 NLP 目标；MinActiveLoss 恒用 LossEconomic（\srcpath{RPOInnerObjective}）\\
\fld{robust\_inner\_on\_failure} & — & — & true/false & true & MatchRPO 基线失败时自动升级 LossEconomic 并重解一次\\
\fld{robust\_inner\_min\_buses} & — & 母线条数 & 0--$10^{5}$（经验值） & 1000 & 交流母线数达到该值时 vdev/Combined 直接以 LossEconomic 起步；0=总先试 MatchRPO；对 MinActiveLoss 无影响\\
\fld{vdev\_weight} & $w_V$ & — & $10^{-4}$--$10^{2}$（经验值） & 1.0 & 电压偏差项权重（Combined/MinVoltageDeviation）\\
\fld{loss\_weight} & $w_L$ & — & $10^{-4}$--$10^{2}$（经验值） & 1.0 & 网损项权重（Combined/MinActiveLoss）；1 时目标量纲为 MW\\
\fld{v\_target} & $V^{*}$ & pu & 0.95--1.05（经验值） & 1.0 & 电压偏差惩罚目标值\\
\fld{max\_nodes} & $N_{\mathrm{eval}}^{\max}$ & 次 & $10^{2}$--$10^{5}$（经验值） & 50000 & 连续 OPF 评估次数上限（含基线与种子近似计数）\\
\fld{time\_limit\_sec} & $T^{\max}$ & s & 10--3600（经验值） & 120.0 & 离散搜索墙钟时间上限（只能在两次求解之间生效）\\
\fld{gap\_tol} & $\epsilon_{\mathrm{imp}}$ & — & $10^{-6}$--$10^{-2}$（经验值） & $10^{-4}$ & 坐标下降判定“本轮有改进”的最小目标改善量\\
\fld{int\_tol} & — & — & — & $10^{-5}$ & 保留字段（源码兼容），离散值以 int 表示，当前未使用\\
\fld{num\_threads} & — & 线程 & 0--硬件并发 & 0 & 离散候选并行评估线程数；0=自动，1=串行；混合 AC/DC 强制为 1；并行把搜索变为 Jacobi 波形，轨迹/局部解可能与串行略有差异\\
\fld{branching} & — & — & — & Pseudocost & 旧 B\&B 后端保留字段，当前局部搜索不消费\\
\fld{node\_sel} & — & — & — & Hybrid & 旧 B\&B 后端保留字段，当前局部搜索不消费\\
\fld{max\_ipm\_iter} & $K_{\mathrm{ipm}}^{\max}$ & 次 & 400--4000（经验值） & 2000 & 每次内层求解（基线）的 IPM 迭代预算上限\\
\fld{search\_ipm\_iter} & $K_{\mathrm{s}}^{\max}$ & 次 & 50--500（经验值） & 200 & 搜索评估的失败快速迭代上限（仅纯交流生效）；$\le 0$ 表示不设上限\\
\fld{ipm\_tol} & $\epsilon_p$ & pu 等 & $10^{-9}$--$10^{-4}$（经验值） & $10^{-6}$ & 内层可行性（最大约束违反）容差，兼作实测收敛判据\\
\fld{stationarity\_tol} & $\epsilon_d$ & — & $10^{-6}$--$10^{-2}$（经验值） & $10^{-3}$ & scaled KKT 平稳性容差，与可行性容差分列\\
\fld{inner\_solver\_backend} & — & — & ParityIPM/\allowbreak Ipopt/\allowbreak Auto/\allowbreak EconomicDispatch & ParityIPM & 全部内层 AC OPF 评估的后端（\srcpath{ACOPFSolverBackend}，见 \srcpath{include/hacdcpf/optimal\_power\_flow/opf\_options.hpp}:28--33）\\
\fld{seed\_homotopy\_on\_failure} & — & — & true/false & true & 大/僵硬电网经济种子不收敛时，沿经济目标同伦一次性兜底\\
\fld{max\_tap\_move} & $m_{\mathrm{tap}}$ & 档 & $-1$--$1000$ & $-1$ & 每台已选 OLTC 距当前档位的最大移动档数；$-1$=完整铭牌范围，0=全部固定；GUI 显式默认 2\\
\fld{restrict\_tap\_indices} & — & — & true/false & false & true 时仅 \fld{enabled\_tap\_indices} 列出的 vector position 进入离散决策向量\\
\fld{enabled\_tap\_indices} & — & — & — & 空 & 参与优化的变压器 vector position 列表（配合上一字段）\\
\fld{verbose} & — & — & true/false & false & 内层求解器详细输出\\
\fld{enforce\_branch\_limits} & — & — & true/false & true & 内层启用支路热稳限值\\
\fld{enforce\_converter\_capacity} & — & — & true/false & true & 内层启用 VSC 容量圆约束\\
\fld{enforce\_converter\_current\_limits} & — & — & true/false & true & 内层启用换流器电流限值\\
\fld{enforce\_converter\_modulation\_limits} & — & — & true/false & true & 内层启用调制比窗口与 DC/DC 占空比约束\\
\end{paramtable}

\subsubsection*{求解结果（RPOResult）}
\begin{paramtable}
\fld{converged} & — & — & true/false & false & 是否得到经 OPF 补全的可行在位点（局部意义，非全局证书）\\
\fld{objective} & $f_{\mathrm{RPO}}$ & —（量纲随目标） & — & 0.0 & 在位点的 RPO 排名目标值；MinActiveLoss 且 $w_L=1$ 时等于物理网损 MW\\
\fld{gap} & — & — & — & 1.0 & 占位值；不提供最优性间隙（无下界）\\
\fld{nodes\_explored} & — & 次 & — & 0 & 离散候选评估计数（含基线）\\
\fld{nlp\_solves} & — & 次 & — & 0 & 内层 NLP 求解次数（近似含种子/基线）\\
\fld{runtime\_sec} & — & s & — & 0.0 & 求解墙钟耗时\\
\fld{status} & — & — & — & 空 & 状态串（局部最优/时限可行/评估限可行/无可行配置及实测容差附注）\\
\fld{algorithm} & — & — & — & \fld{discrete\_coordinate\_search\_with\_ac\_opf} & 算法标识（固定字符串）\\
\fld{globally\_certified} & — & — & — & false & 恒 false：无全局 MINLP 证书\\
\fld{optimality\_gap\_available} & — & — & — & false & 恒 false：无最优性间隙\\
\fld{terminated\_by\_time\_limit} & — & — & true/false & false & 是否撞墙钟时限\\
\fld{terminated\_by\_evaluation\_limit} & — & — & true/false & false & 是否撞评估次数上限\\
\fld{model\_limitations} & — & — & — & 空 & 局限声明串列表（局部性、预算截断、内层目标升级等，逐条写入）\\
\fld{effective\_inner\_nlp\_objective} & — & — & MatchRPO/\allowbreak LossEconomic & MatchRPO & 实际生效的内层目标（可能因稳健升级不同于请求值）\\
\fld{vm\_before}/\allowbreak\fld{vm\_after} & $V_m$ & pu & — & 空 & 优化前/后交流母线电压幅值全向量（按 \fld{ac.buses} 位置序，与 GUI 回放写回口径一致）\\
\fld{va\_before}/\allowbreak\fld{va\_after} & $\theta$ & rad & — & 空 & 优化前/后交流母线电压相角（弧度；GUI 按 $\times 180/\pi$ 转度展示）\\
\fld{qg\_mvar\_before}/\allowbreak\fld{qg\_mvar\_after} & $Q_g$ & MVAr & — & 空 & 优化前/后机组无功（内层 OPF 调度）\\
\fld{pg\_mw\_before}/\allowbreak\fld{pg\_mw\_after} & $P_g$ & MW & — & 空 & 优化前/后机组有功（内层 OPF 调度）\\
\fld{taps} & — & — & — & 空 & 逐台 OLTC 明细（\srcpath{TapResult}，见下表）\\
\fld{shunts} & — & — & — & 空 & 逐台可投切并联明细（\srcpath{ShuntResult}，见下表）\\
\fld{total\_loss\_before\_mw}/\allowbreak\fld{total\_loss\_after\_mw} & $P_{\mathrm{loss}}$ & MW & — & 0.0 & 优化前/后全系统物理网损（\srcpath{compute\_loss} 台账口径）\\
\fld{max\_vdev\_before}/\allowbreak\fld{max\_vdev\_after} & $\Delta V_{\max}$ & pu & — & 0.0 & 优化前/后最大电压偏差 $\max_i|V_{m,i}-V^{*}|$\\
\fld{bc\_stats} & — & — & — & — & 遗留统计信封：记评估数/NLP 次数/灵敏度面数，非 B\&B 证书\\
\fld{baseline\_opf} & — & — & — & — & 基线内层 OPF 完整运行点（供独立 OPF/PF 回放）\\
\fld{optimized\_opf} & — & — & — & — & 在位点内层 OPF 完整运行点（供独立 OPF/PF 回放）\\
\end{paramtable}

\subsubsection*{逐设备结果（TapResult / ShuntResult）}
\begin{paramtable}
\fld{trafo\_index} & — & — & — & 0 & 变压器在 \fld{ac.transformers\_2w} 的 vector position\\
\fld{name} & — & — & — & 空 & 设备名（空名自动生成）\\
\fld{tap\_before}/\allowbreak\fld{tap\_after} & $t_k$ & 档 & 铭牌窗口内 & 0 & 优化前/后档位\\
\fld{ratio\_before}/\allowbreak\fld{ratio\_after} & $r$ & pu & 约 0.85--1.15（经验值） & 1.0 & 相对中性档的相对变比 $r(t)$\\
\fld{electrical\_tap\_before}/\allowbreak\fld{electrical\_tap\_after} & $\tau$ & pu & — & 1.0 & 联动 MATPOWER 支路的实际电气变比 \fld{branch.tap}（无联动时等于相对变比）\\
\midrule
\fld{shunt\_index} & — & — & — & 0 & 并联补偿在 \fld{ac.shunts} 的 vector position\\
\fld{name} & — & — & — & 空 & 设备名\\
\fld{step\_before}/\allowbreak\fld{step\_after} & $s_j$ & 步 & 0 至 \fld{n\_steps} & 0 & 优化前/后步数\\
\fld{bs\_mvar\_before}/\allowbreak\fld{bs\_mvar\_after} & $B_s$ & MVAr & — & 0.0 & 优化前/后电纳（等于 \fld{bs\_per\_step} 乘以步数 $s_j$）\\
\end{paramtable}
（结构体定义：
\srcpath{include/hacdcpf/optimal\_power\_flow/reactive\_power\_opt.hpp}:158--177；
结果装配：\srcpath{src/optimal\_power\_flow/reactive\_power\_opt.cpp}:1347--1371。
表中 \textbackslash midrule 上半为 \srcpath{TapResult}，下半为
\srcpath{ShuntResult}。）

\subsubsection*{控制清单审计（RPOTapControlInput / RPOShuntControlInput）}
\srcpath{inspect\_rpo\_controls} 返回 \srcpath{RPOControlInventory}
（\fld{taps} 与 \fld{shunts} 两个向量，
\srcpath{include/hacdcpf/optimal\_power\_flow/reactive\_power\_opt.hpp}:227--230），
是求解器与 GUI 输入预览共用的入选合同。OLTC 行字段如下
（同文件:RPOTapControlInput；装配实现
\srcpath{src/optimal\_power\_flow/reactive\_power\_opt.cpp}:512--577）：
\begin{paramtable}
\fld{trafo\_index} & — & — & — & 0 & 在 \fld{ac.transformers\_2w} 的 vector position\\
\fld{authored\_index} & — & — & — & 0 & 用户侧稳定元件编号 \fld{Transformer2W::index}\\
\fld{source\_branch\_idx} & — & — & — & 0 & 联动 MATPOWER 支路的 \fld{ACBranch::index}（0=无联动）\\
\fld{name} & — & — & — & 空 & 设备名\\
\fld{hv\_bus}/\allowbreak\fld{lv\_bus} & — & — & — & 0 & 高/低压侧母线编号\\
\fld{tap\_side} & — & — & 0/1 & 0 & 分接头所在侧（0=高压，1=低压）\\
\fld{in\_service} & — & — & true/false & false & 投运状态\\
\fld{adjustable} & — & — & true/false & false & 铭牌层面可调（排除原因为空）\\
\fld{selected\_for\_optimization} & — & — & true/false & false & 本次优化入选（可调且通过逐台选择且 \fld{max\_tap\_move} $\ne 0$）\\
\fld{exclusion\_reason} & — & — & — & 空 & 排除原因串（空=无排除；取值见数学模型小节）\\
\fld{tap\_pos} & $t_k^{0}$ & 档 & 铭牌范围内 & 0 & 当前档位\\
\fld{tap\_min}/\allowbreak\fld{tap\_max} & $t_k^{\min}$/$t_k^{\max}$ & 档 & — & 0 & 铭牌档位下/上限\\
\fld{tap\_neutral} & $t_{\mathrm{neutral}}$ & 档 & — & 0 & 中性档（相对变比基准）\\
\fld{position\_count} & — & 档 & — & 0 & 铭牌档位数 $t_k^{\max}-t_k^{\min}+1$（无效范围为 0）\\
\fld{optimization\_tap\_min}/\allowbreak\fld{optimization\_tap\_max} & — & 档 & — & 0 & 本次优化窗口下/上限（铭牌与 $\pm m_{\mathrm{tap}}$ 窗口之交）\\
\fld{optimization\_position\_count} & — & 档 & — & 0 & 本次可选位置数（未入选为 0）\\
\fld{tap\_step\_percent} & — & \% & 0.5--2.5（经验值） & 0.0 & 每档步长（百分数）\\
\fld{ratio\_current}/\allowbreak\fld{ratio\_min}/\allowbreak\fld{ratio\_max} & $r$ & pu & — & 1.0 & 当前/最小/最大相对变比 $r(t)$\\
\fld{electrical\_tap\_current}/\allowbreak\fld{electrical\_tap\_min}/\allowbreak\fld{electrical\_tap\_max} & $\tau$ & pu & — & 1.0 & 联动支路当前/最小/最大电气变比\\
\end{paramtable}
可投切并联行字段如下
（\srcpath{include/hacdcpf/optimal\_power\_flow/reactive\_power\_opt.hpp}:211--225；
装配实现 \srcpath{src/optimal\_power\_flow/reactive\_power\_opt.cpp}:579--599）：
\begin{paramtable}
\fld{shunt\_index} & — & — & — & 0 & 在 \fld{ac.shunts} 的 vector position\\
\fld{authored\_index} & — & — & — & 0 & 用户侧稳定元件编号 \fld{Shunt::index}\\
\fld{name} & — & — & — & 空 & 设备名\\
\fld{bus} & — & — & — & 0 & 挂接母线编号\\
\fld{in\_service} & — & — & true/false & false & 投运状态\\
\fld{switchable} & — & — & true/false & false & 可投切标志\\
\fld{adjustable} & — & — & true/false & false & 入选资格（排除原因为空）\\
\fld{exclusion\_reason} & — & — & — & 空 & 排除原因串（空=无排除）\\
\fld{current\_step} & $s_j^{0}$ & 步 & 0 至 \fld{n\_steps} & 0 & 当前步数\\
\fld{n\_steps} & $n_j$ & 步 & — & 0 & 最大步数\\
\fld{position\_count} & — & 步 & — & 0 & 可选位置数 $n_j+1$\\
\fld{bs\_per\_step\_mvar} & — & MVAr & — & 0.0 & 每步电纳\\
\fld{bs\_current\_mvar} & $B_s$ & MVAr & — & 0.0 & 当前电纳\\
\end{paramtable}

\subsection{验证与校核}

\subsubsection*{全过程交叉验证（GUI 服务层）}
验证文档第 4 节规定的全过程交叉验证在 GUI 服务层固定执行
（\srcpath{tests/run\_gui\_server.cpp}:19432--19684），步骤为：
\begin{enumerate}
  \item 保存 RPO 内层最优连续点与离散档位（\fld{optimized\_opf}）；
  \item 用 \srcpath{apply\_rpo\_discrete\_solution} 把档位写入独立系统副本，
        按 \fld{effective\_inner\_nlp\_objective} 复刻单位成本改写后重解
        \textbf{独立 AC/DC OPF}（后端 Auto，非 RPO 内层管线，
        \srcpath{tests/run\_gui\_server.cpp}:19434--19484）；
  \item 比较两次 OPF 的目标值、最大 $|\Delta V_m|$、$|\Delta P_g|$、
        $|\Delta Q_g|$（同文件:run\_rpo（lambda））；
  \item 把优化调度（$P_g/Q_g$、新能源、储能、柔性负荷、$V_m/\theta$）写入
        PF 回放副本并运行非线性潮流（同文件:run\_rpo（lambda））；
  \item 按判据给结论：独立 OPF 一致性要求双方收敛、审计可行，目标相对差
        不超过 $\max(10^{-3},\,10\epsilon_d)$、电压差不超过
        $\max(5\times10^{-4},\,10\epsilon_p)$ pu（同文件:run\_rpo（lambda））——
        与验证文档第 4 节的表述逐字相符（本章已抽查核实）；纯交流的 PF
        电压差判据为 $10^{-3}$ pu（同文件:run\_rpo（lambda））；
  \item 富混合 AC/DC（含直流网/VSC/DC-DC/能量路由器）的 authored-space
        调度回放尚未完整映射时，\fld{pf\_check\_applicable}=false，结论
        标注“可用数值检查通过；PF 回放部分覆盖”，不宣称完整 PF 等价
        （同文件:run\_rpo（lambda））。
\end{enumerate}
非凸 OPF 可能存在目标等价但 $P/Q$ 控制分配不同的局部点，故
$|\Delta P_g|$/$|\Delta Q_g|$ 完整展示但不单独触发失败
（\srcpath{tests/run\_gui\_server.cpp}:19584--19587 注释）。

\subsubsection*{自动测试矩阵}
\begin{paramtable}
\fld{tests/test\_reactive\_power\_opt.cpp}:16--42 & case9（MATPOWER）& MinActiveLoss & — & — & 断言收敛、\fld{globally\_certified}/\fld{optimality\_gap\_available} 为 false、\fld{model\_limitations} 非空、\fld{objective} 等于物理网损且 $0<P_{\mathrm{loss}}<50$ MW\\
\fld{tests/test\_reactive\_power\_opt.cpp}:44--82 & case9 + 2 步电容器组 & MinVoltage\allowbreak Deviation & — & — & 断言离散步数动作落在 $[0,2]$、优化目标不差于基线、\fld{gap}=1.0（不报虚假 gap=0）\\
\fld{tests/test\_reactive\_power\_opt.cpp}:84--110 & \srcpath{build\_case300\_acdc} & MinVoltage\allowbreak Deviation & — & — & 种子求解返回完整展示向量、已选 OLTC 数与清单一致；标记 \srcpath{[.performance]}，默认 ctest 不运行\\
\fld{tests/test\_reactive\_power\_opt.cpp}:112--182 & 手工 4 台变压器 + 2 台并联 & 清单审计 & — & — & 逐条断言排除原因串、位置计数与 $\pm m_{\mathrm{tap}}$ 窗口、逐台选择开关\\
\fld{tests/test\_reactive\_power\_opt.cpp}:184--211 & case300.m 导入链 & 清单审计 & — & — & MATPOWER 导入不虚构 4001 档 OLTC；\srcpath{build\_case300\_acdc} 恰 3 台 $[-4,4]$/1.25\% 可调\\
\fld{tests/test\_reactive\_power\_opt.cpp}:213--243 & 手工联动支路 & 档位回写 & — & — & \srcpath{apply\_rpo\_discrete\_solution} 把档位写入变压器并按比率缩放 \fld{branch.tap}（$1.03\times1.025$）\\
\fld{tests/e2e/rpo\_gui\_e2e.mjs}:63--267 & case24/\allowbreak case9/\allowbreak case300\_acdc & 浏览器 E2E & — & — & case24 五个固定变比 $1.03/1.02$ 不被识别为 OLTC；case9 全过程交叉验证通过；case300 混合回归（62 台变压器/3 台可调、\fld{pf\_check\_applicable}=false 部分覆盖）\\
\end{paramtable}
（说明栏为本章概括，非表字段；该表借用 paramtable 版式。）
测试注册：\srcpath{tests/CMakeLists.txt}:34（单元测试）、
\srcpath{tests/CMakeLists.txt}:532--537（\srcpath{rpo\_gui\_e2e}，标签
\srcpath{gui;e2e;browser;rpo;cross\_validation}，需 chromium，超时 180 s）。

\subsubsection*{大规模可扩展性实测（文档存档）}
验证文档第 5 节记录了内层目标重构前后的大算例实测
（\srcpath{docs/archive/audits/reactive\_power\_optimization\_validation.md}）。关键数字：
\begin{itemize}
  \item \textbf{case2869（PEGASE，2869 节点 + 20 台可调 OLTC + 10 台可投切
        并联，authored 测试系统）}：基线（单位成本经济 + 同伦）74 次迭代 /
        25 s 收敛到参考目标 133999，约束违反 $1.1\times10^{-6}$；暖启动链
        （tap$\pm1$/并联$\pm1$）全部收敛，2--7 次迭代、平均约 0.46 s/次；
        完整 MinActiveLoss RPO 动作 22 台设备，物理网损 1561.8~MW
        $\to$ 1558.8~MW（约 $-3$ MW）；MinVoltageDeviation（稳健内层）收敛
        但 vdev 0.1392 $\to$ 0.1401 基本未降（机理见局限小节）。
  \item \textbf{IEEE24 三区域扩展回归}：vdev（MatchRPO）loss 36.45~MW
        $\to$ 35.34~MW、vdev 0.0180 $\to$ 0.0114，与旧版逐位一致；loss
        目标（新单位成本经济）基线 20.14~MW、最终 19.22~MW（优于旧 19.65）；
        combined：loss 29.85 $\to$ 20.02~MW、vdev 0.050 $\to$ 0.044。
\end{itemize}
\textbf{如实说明}：上述数字是验证文档记录的实测存档，本章在 \srcpath{tests/}
全库检索（133999、1561.8、19.22、36.45 等）未发现对应的自动断言——相关
自动测试只覆盖 case2869 的潮流/OPF 收敛性
（\srcpath{tests/test\_matpower\_cases.cpp}:336--337、
\srcpath{tests/test\_opf\_solver\_backends.cpp}:1154--1156），RPO 大算例
指标未纳入默认 ctest 回归，复现需按文档口径自建 authored 算例。

\subsection{边界与局限}

\begin{enumerate}
  \item \textbf{局部启发式，无全局证书}：离散层是灵敏度排序 + 坐标搜索 +
        成对邻域精化的局部启发式；\fld{globally\_certified}/
        \fld{optimality\_gap\_available} 恒为 false，\fld{gap} 恒为占位
        1.0，\fld{bc\_stats} 不是分支定界证书
        （\srcpath{src/optimal\_power\_flow/reactive\_power\_opt.cpp}:1301--1302、
        1335、1373--1383）。
  \item \textbf{启发假设}：Phase 1 依赖目标对单变量的近似单峰性，非凸问题
        不保证；局部灵敏度面不作为全局下界，从不用于剪枝
        （同文件:solve\_rpo、1171--1175）。
  \item \textbf{LossEconomic 内层不直接优化电压}：稳健内层按网损最优分配
        机组无功；实测 case2869 上 vdev RPO 虽收敛，离散设备未能抵消网损
        最优无功分配对电压的不利影响（vdev 0.1392 $\to$ 0.1401）。如需机组
        无功参与调压，仅中小电网用 MatchRPO
        （\srcpath{docs/archive/audits/reactive\_power\_optimization\_validation.md} 第 5.5 节；
        结果字段 \fld{effective\_inner\_nlp\_objective} 与相应
        \fld{model\_limitations} 条目如实披露该路径，
        \srcpath{src/optimal\_power\_flow/reactive\_power\_opt.cpp}:1311--1315）。
  \item \textbf{预算粒度}：时间与评估预算只能在两次内层求解之间检查，
        单次 IPM 可能超出剩余预算；\fld{search\_ipm\_iter} 是唯一的单次
        求解护栏，且只对纯交流电网生效
        （\srcpath{include/hacdcpf/optimal\_power\_flow/reactive\_power\_opt.hpp}:110--117）。
  \item \textbf{并行语义}：\fld{num\_threads}$>1$ 把搜索从 Gauss--Seidel
        变为 Jacobi（在位点波间推进），轨迹与所得局部解可能与串行不同；
        混合 AC/DC 强制串行（MUMPS 互斥锁）
        （\srcpath{include/hacdcpf/optimal\_power\_flow/reactive\_power\_opt.hpp}:84--96）。
  \item \textbf{未覆盖设备}：三绕组变压器分接头、独立调压控制器
        （\fld{regulator\_controls}）、单体三相不平衡 RPO 不进入离散决策
        向量；移动储能/虚拟电厂/微网为固定或投影注入
        （\srcpath{tests/run\_gui\_server.cpp}:19781--19796；
        \srcpath{docs/archive/audits/reactive\_power\_optimization\_validation.md} 第 3 节）。
  \item \textbf{混合 AC/DC 回放部分覆盖}：富混合元件的 authored-space PF
        调度回放尚未完整映射，交叉验证对此标注“部分覆盖”而非数值不一致
        （\srcpath{tests/run\_gui\_server.cpp}:19581--19583、19677--19684）。
  \item \textbf{MATPOWER 导入链}：导入支路变比一律按固定电气变比处理，
        只有算例或用户显式给出有效档位范围与步长后才进入 RPO
        （\srcpath{tests/test\_reactive\_power\_opt.cpp}:184--211）。
  \item \textbf{保留字段}：\fld{int\_tol}、\fld{branching}、\fld{node\_sel}
        为兼容旧 B\&B 构建而保留，当前局部搜索不消费
        （\srcpath{include/hacdcpf/optimal\_power\_flow/reactive\_power\_opt.hpp}:80--81、
        98--101）。
  \item \textbf{目标侧无档位动作代价}：频繁调档的机械寿命成本未建模；
        \fld{max\_tap\_move} 窗口与 \fld{gap\_tol} 改进阈值是仅有的动作
        抑制手段（本章数学模型小节）。
\end{enumerate}

% 三相混合最优潮流
% 本章文件由 \input 引入，不含导言区。
\section{三相混合最优潮流}\label{sec:threephase}

本章介绍平台的三相混合最优潮流模块（命名空间 \srcpath{hacdcpf::opf::phase\_hybrid}）。
该模块面向含单相/两相接入的不平衡配电系统与交直流混合系统，在\textbf{相域}
（phase domain）内直接建立并求解最优潮流：每个母线的每一相都是独立的电气节点，
支路以 $3\times 3$ 耦合导纳块建模，换流器按相端口接入并可选择相域控制律。
模块由三层组成——装配层 adapter 把 rich 工程模型 \fld{HybridPowerSystem} 装配为
相域 OPF 问题；OPF 本体把该问题作为单体非线性规划（NLP）用内点法求解；松弛层
给出同一问题的 SOCP 多面体外逼近与目标下界证书。与第~\ref{sec:acopfnative}~章的
单相 AC OPF 不同，本章模块不经过 canonical 投影与 \fld{SolverData} 装配，而是使用
专用的 \fld{ThreePhaseHybridOPFCase} 数据结构。

\textbf{成熟度声明（诚实口径）}：该模块是平台当前的活跃研发对象，接口、控制律覆盖
与验证深度仍在演进。本章全部结论以当前源代码为准；代码中的近似、占位、被拒绝的
输入与尚未生效的参数在~\ref{subsec:tp-limit}~节逐条列出，不作为已实现能力陈述。

\subsection{相域建模动机与适用范围}

\subsubsection*{动机}
输电网 OPF 的正序（平衡）等值在配电系统中失效：单相与两相分支、分相负荷、
分相分布式电源以及不平衡运行工况要求按相分别表示电压、电流与功率。代码层面
的对应设计如下：

\begin{itemize}
  \item \textbf{按相存在的节点}。装配层只为“母线在运且该相存在”的
        $(\text{母线},\varphi)$ 对分配相节点，\fld{phase\_mask} 不含某相时该相
        根本不进入问题，因而单相/两相分支自然得到少于 $3\times n_{\mathrm{bus}}$
        的节点集（\srcpath{src/optimal\_power\_flow/three\_phase\_hybrid\_adapter.cpp:build\_three\_phase\_hybrid\_model}）。
        节点的相归属记录在 \fld{ac\_phase\_index}（:build\_three\_phase\_hybrid\_model）。
  \item \textbf{分相负荷与分相电源}。母线三相负荷 \fld{pd\_a\_mw}/\fld{pd\_b\_mw}/%
        \fld{pd\_c\_mw}（无功同理）按相累加到对应节点的恒功率注入
        （:build\_three\_phase\_hybrid\_model）；三相发电机可按显式分相调度值或按相数均分
        接入单相节点（:build\_three\_phase\_hybrid\_model）。
  \item \textbf{相间耦合}。三相支路不以序分量解耦，而是以满 $3\times 3$ 复导纳块
        直接进入相域导纳矩阵 \fld{y\_ac}（:build\_three\_phase\_hybrid\_model），因此相间互感、
        不平衡潮流分布无需任何序网假设。
  \item \textbf{不平衡度约束}。对三相齐全的母线显式施加负序/正序电压幅值比
        （VUF）上限（\srcpath{src/optimal\_power\_flow/three\_phase\_hybrid\_opf.cpp:evaluate\_inequalities}），
        这是序分量在模块中的唯一显式用途（另见换流器正序电流控制）。
  \item \textbf{换流器相域控制}。VSC 换流器在相域内可选择均分相功率、
        跟网型正序电流、构网型下垂三种控制方程
        （\srcpath{include/hacdcpf/optimal\_power\_flow/three\_phase\_hybrid\_opf.hpp:PhaseVSCControlMode}），
        支撑不平衡接入下的逆变器调度研究。
\end{itemize}

\subsubsection*{适用范围（按装配层校验）}
装配入口 \srcpath{build\_three\_phase\_hybrid\_model} 对输入系统做硬性校验
（\srcpath{src/optimal\_power\_flow/three\_phase\_hybrid\_adapter.cpp:build\_three\_phase\_hybrid\_model}）：
必须含 \fld{three\_phase\_ac} 子系统、至少一个 DC 母线与一台 VSC 换流器；
\textbf{不支持} DC/DC 变换器与能量路由器，遇到即抛 \fld{std::invalid\_argument}。
此外还要求：至少一个在运相节点（:build\_three\_phase\_hybrid\_model）、至少一个交流参考源
（:build\_three\_phase\_hybrid\_model）、至少一个在运 DC 母线（:build\_three\_phase\_hybrid\_model）、DC 支路电阻
严格为正（:build\_three\_phase\_hybrid\_model）、每台 VSC 的交流三端子相 A/B/C 齐全
（:build\_three\_phase\_hybrid\_model）、至少一台两端均在运的 VSC（:build\_three\_phase\_hybrid\_model）。
可控 DC 静态发电机同样被拒绝（:build\_three\_phase\_hybrid\_model）。

\subsection{三层结构与代码组织}\label{subsec:tp-layers}

图~\ref{fig:tp-layers}~给出模块的三层组织与数据流。三层共享同一份问题数据
\fld{ThreePhaseHybridOPFCase}：装配层负责生成它，OPF 本体与松弛层各自独立地
消费它，二者不互相调用。

\begin{figure}[H]
\centering
\begin{tikzpicture}[
  box/.style={draw,rectangle,minimum width=5.2cm,minimum height=1.05cm,align=center,font=\footnotesize},
  arr/.style={-{Stealth},thick}]
  \node[box] (rich) {\fld{HybridPowerSystem}\\（rich 工程语义层）};
  \node[box,below=0.85cm of rich] (adapter) {装配层 adapter\\\srcpath{three\_phase\_hybrid\_adapter.cpp}};
  \node[box,below=0.85cm of adapter] (casedata) {\fld{ThreePhaseHybridOPFCase}\\（相域 OPF 问题数据）};
  \node[box,below left=1.0cm and -1.6cm of casedata] (opf) {相域 OPF 本体（NLP）\\\srcpath{three\_phase\_hybrid\_opf.cpp}};
  \node[box,below right=1.0cm and -1.6cm of casedata] (relax) {SOCP 松弛层（LP 序列）\\\srcpath{three\_phase\_hybrid\_relaxation.cpp}};
  \node[box,below=0.85cm of opf] (nlp) {\fld{engine::NLPModel}\\NativeIPM / Ipopt 后端};
  \node[box,below=0.85cm of relax] (lp) {\fld{engine::LPModel}\\HiGHS 后端};
  \draw[arr] (rich) -- (adapter);
  \draw[arr] (adapter) -- (casedata);
  \draw[arr] (casedata.south) -- (opf.north);
  \draw[arr] (casedata.south) -- (relax.north);
  \draw[arr] (opf) -- (nlp);
  \draw[arr] (relax) -- (lp);
\end{tikzpicture}
\caption{三相混合 OPF 模块的三层组织：adapter 负责 rich 模型到相域问题数据的装配，
OPF 本体与松弛层共享同一份 \fld{ThreePhaseHybridOPFCase}，分别下沉到
MIPSolvers 的 NLP 与 LP 求解通道}
\label{fig:tp-layers}
\end{figure}

\subsubsection*{装配层（adapter）}
\compmeta{ThreePhaseHybridModel}{\srcpath{include/hacdcpf/optimal\_power\_flow/three\_phase\_hybrid\_adapter.hpp}}
入口函数 \srcpath{build\_three\_phase\_hybrid\_model}
（头文件声明：\srcpath{include/hacdcpf/optimal\_power\_flow/three\_phase\_hybrid\_adapter.hpp:build\_three\_phase\_hybrid\_model}；
实现：\srcpath{src/optimal\_power\_flow/three\_phase\_hybrid\_adapter.cpp:build\_three\_phase\_hybrid\_model}）。
输出结构 \fld{ThreePhaseHybridModel}（\srcpath{include/hacdcpf/optimal\_power\_flow/three\_phase\_hybrid\_adapter.hpp:ThreePhaseHybridModel}）内嵌装配完成的
\fld{ThreePhaseHybridOPFCase}，并保留稳定组件 ID 与求解位置的映射：
\fld{ac\_bus\_ids}、\fld{bus\_phase\_nodes}（每母线三相的节点号，缺相为 $-1$）、
\fld{dc\_bus\_ids}、\fld{vsc\_component\_indices}、\fld{generator\_bindings}
（含组件 \fld{index}、母线、相号、名称、来源类型与爬坡参数）以及
\fld{model\_limitations} 字符串列表。查询接口 \texttt{phase\_node(ac\_bus\_id, phase)}
按 ID 反查节点号（实现：src/optimal\_power\_flow/three\_phase\_hybrid\_adapter.cpp:phase\_node）。装配选项
\fld{ThreePhaseHybridAdapterOptions} 的全部字段见~\ref{subsec:tp-adapter}~节末的
参数表。

\subsubsection*{相域 OPF 本体（opf）}
\compmeta{ThreePhaseHybridOPFCase}{\srcpath{include/hacdcpf/optimal\_power\_flow/three\_phase\_hybrid\_opf.hpp}}
公开接口（\srcpath{include/hacdcpf/optimal\_power\_flow/three\_phase\_hybrid\_opf.hpp:solve\_three\_phase\_hybrid\_opf、make\_three\_phase\_hybrid\_pf\_case、solve\_three\_phase\_hybrid\_opf\_sequence、solve\_three\_phase\_hybrid\_opf\_branches}）：单点求解 \srcpath{solve\_three\_phase\_hybrid\_opf}
（实现入口 src/optimal\_power\_flow/three\_phase\_hybrid\_opf.cpp:solve\_three\_phase\_hybrid\_opf，核心实现 \srcpath{solve\_three\_phase\_hybrid\_opf\_impl}
于 src/optimal\_power\_flow/three\_phase\_hybrid\_opf.cpp:solve\_three\_phase\_hybrid\_opf\_impl）、运行点转潮流算例 \srcpath{make\_three\_phase\_hybrid\_pf\_case}
（src/optimal\_power\_flow/three\_phase\_hybrid\_opf.cpp:make\_three\_phase\_hybrid\_pf\_case）、固定拓扑序列求解 \srcpath{solve\_three\_phase\_hybrid\_opf\_sequence}
（src/optimal\_power\_flow/three\_phase\_hybrid\_opf.cpp:solve\_three\_phase\_hybrid\_opf\_sequence）与参数摄动分支求解 \srcpath{solve\_three\_phase\_hybrid\_opf\_branches}
（src/optimal\_power\_flow/three\_phase\_hybrid\_opf.cpp:solve\_three\_phase\_hybrid\_opf\_branches）。内部组织为：数据校验 \srcpath{validate\_case}
（src/optimal\_power\_flow/three\_phase\_hybrid\_opf.cpp:validate\_case）$\to$ 模型数据 \fld{ModelData} 与变量布局 \fld{Layout}
装配 \srcpath{build\_model\_data}（src/optimal\_power\_flow/three\_phase\_hybrid\_opf.cpp:build\_model\_data）$\to$ NLP 回调装配
\srcpath{build\_nlp}（src/optimal\_power\_flow/three\_phase\_hybrid\_opf.cpp:build\_nlp）$\to$ 初值/对偶初始化 $\to$ 内点求解
$\to$ KKT 审计与结果归因。

\subsubsection*{松弛/凸化层（relaxation）}
\compmeta{ThreePhaseHybridRelaxationResult}{\srcpath{include/hacdcpf/optimal\_power\_flow/three\_phase\_hybrid\_relaxation.hpp}}
入口 \srcpath{solve\_three\_phase\_hybrid\_opf\_relaxation}
（\srcpath{include/hacdcpf/optimal\_power\_flow/three\_phase\_hybrid\_relaxation.hpp:solve\_three\_phase\_hybrid\_opf\_relaxation}；实现：\srcpath{src/optimal\_power\_flow/three\_phase\_hybrid\_relaxation.cpp:solve\_three\_phase\_hybrid\_opf\_relaxation}），
以及同选项序列求解 \srcpath{solve\_three\_phase\_hybrid\_opf\_relaxation\_sequence}
（src/optimal\_power\_flow/three\_phase\_hybrid\_relaxation.cpp:solve\_three\_phase\_hybrid\_opf\_relaxation\_sequence）。头文件注释明确了该层性质：对相域 AC/DC SOCP 松弛做
多面体外逼近，每一条锥割都是 SOCP 可行集的必要条件，因此每一轮 LP 都保持为原
非凸 OPF 的外松弛；凸二次发电成本用支撑切线低估（include/hacdcpf/optimal\_power\_flow/three\_phase\_hybrid\_relaxation.hpp:solve\_three\_phase\_hybrid\_opf\_relaxation）。
该层输出的是\textbf{目标下界证书}而非可行运行点。

\subsection{装配层：rich 模型到相域问题数据}\label{subsec:tp-adapter}

\subsubsection*{标幺化与问题规模}
基准功率取 \fld{three\_phase\_ac.base\_mva}，缺失时回退 \fld{system.base\_mva}，
再缺失取 100~MVA（src/optimal\_power\_flow/three\_phase\_hybrid\_adapter.cpp:build\_three\_phase\_hybrid\_model）；VUF 上限 \fld{vuf\_max} 被截断到
$[0,1]$（:build\_three\_phase\_hybrid\_model）。相节点编号按母线顺序、相号 A/B/C 顺序稠密分配，
仅含在运母线的存在相（:build\_three\_phase\_hybrid\_model）。以下功率、电压除注明外均为该基准下
的标幺值。

\subsubsection*{三相导纳矩阵与节点数据}
\fld{y\_ac} 由配网三相潮流模块的 \srcpath{analysis::build\_full\_ybus\_phase\_entries}
生成（声明于 \srcpath{include/hacdcpf/power\_flow/distribution\_power\_flow.hpp}，
调用处 src/optimal\_power\_flow/three\_phase\_hybrid\_adapter.cpp:build\_three\_phase\_hybrid\_model；该函数细节超出本章范围）。调用前装配层复制一份
相域网络并清空其 \fld{external\_grids}——参考电压与可调度电源在 OPF 中已有
显式表示，若再压入其诺顿等值将重复计数（src/optimal\_power\_flow/three\_phase\_hybrid\_adapter.cpp:build\_three\_phase\_hybrid\_model）。
\fld{include\_shunts} 选项控制并联补偿是否入矩阵（:build\_three\_phase\_hybrid\_model）。
节点默认电压窗口为 $[0.90,1.10]$~pu（:build\_three\_phase\_hybrid\_model），随后被母线自身的
\fld{vmin\_pu}/\fld{vmax\_pu}（非正时回退默认）覆盖（:build\_three\_phase\_hybrid\_model）；
初值电压由母线 \fld{vm\_*\_pu} 与 \fld{va\_*\_deg}（度转弧度）合成
（:build\_three\_phase\_hybrid\_model）。三相齐全的母线登记进 \fld{three\_phase\_bus\_nodes}，
作为 VUF 约束的作用集（:build\_three\_phase\_hybrid\_model）。\fld{i\_ac\_fixed}（固定电流注入）
在装配层恒置零（:build\_three\_phase\_hybrid\_model）。

\subsubsection*{负荷等值（诚实口径）}
OpenDSS 风格负荷若 \fld{connection} 为 delta，或带有恒阻抗/恒电流成分
（\fld{const\_z\_percent} 等六个百分比字段任一非零），装配层\textbf{不能}直接
表示其电压特性，只能采用导入时的分相恒功率运行点等值；当选项
\fld{allow\_constant\_power\_load\_equivalent} 为 false 时这类负荷直接抛错，
为 true 时继续装配但在 \fld{model\_limitations} 中登记该近似
（src/optimal\_power\_flow/three\_phase\_hybrid\_adapter.cpp:build\_three\_phase\_hybrid\_model，限制说明字符串见:build\_three\_phase\_hybrid\_model）。

\subsubsection*{发电机与外部电网}
三相发电机逐相生成 \fld{PhaseGenerator}：若给出了显式分相调度
（$|P_a|+|P_b|+|P_c|+|Q_a|+|Q_b|+|Q_c|>10^{-12}$）则取分相值，否则把
\fld{p\_mw}/\fld{q\_mvar} 按相数均分；功率上下界同样按相数分摊，未给出有效
上下界（\fld{pmax\_mw}$\le$\fld{pmin\_mw}）时退化为固定出力（该相上下界相等，
src/optimal\_power\_flow/three\_phase\_hybrid\_adapter.cpp:build\_three\_phase\_hybrid\_model）。成本系数通过名称或母线在 rich 模型中匹配同机
（\srcpath{matching\_rich\_generator}，src/optimal\_power\_flow/three\_phase\_hybrid\_adapter.cpp:matching\_rich\_generator），匹配不到时
$c_2=0$、$c_1=30$（:build\_three\_phase\_hybrid\_model）；匹配到时二次项按相数放大
（\fld{cost\_c2}$\,{}^*=c_2\cdot n_\varphi$）。每台机组的每相随带一条
\fld{generator\_bindings} 记录用于结果归因（:build\_three\_phase\_hybrid\_model）。

外部电网（\fld{external\_grids}）产生两类对象：其一，按
\fld{use\_phase\_voltage\_setpoint} 选择分相或对称（$\theta$、$\theta-120^\circ$、
$\theta+120^\circ$）电压设定，登记为参考节点（src/optimal\_power\_flow/three\_phase\_hybrid\_adapter.cpp:add\_reference（lambda），对称角度
构造见:add\_reference（lambda））；其二，每相一台界为 $\pm 5\times$ 总负荷的虚拟发电机作为
外部电网功率交换的松弛变量（:build\_three\_phase\_hybrid\_model），成本同样来自 rich 侧匹配
（$c_2^*=3c_2$，缺省 $c_1=30$，:build\_three\_phase\_hybrid\_model）。若没有任何外部电网参考，
回退取首台 \fld{BusType::SLACK} 母线（:build\_three\_phase\_hybrid\_model）；仍无则抛错
（:build\_three\_phase\_hybrid\_model）。

\subsubsection*{直流网络与换流器端口}
DC 支路电导 $g=\max(1,\fld{n\_parallel})/\fld{r\_pu}$ 压入拉普拉斯型电导矩阵
\fld{g\_dc}（src/optimal\_power\_flow/three\_phase\_hybrid\_adapter.cpp:build\_three\_phase\_hybrid\_model，电导计算:build\_three\_phase\_hybrid\_model）。DC 母线自带
\fld{pd\_mw} 与电压初值/窗口（:build\_three\_phase\_hybrid\_model）；\fld{DC\_V} 型母线登记为
直流电压参考（:build\_three\_phase\_hybrid\_model）。DC 负荷按 \fld{p\_mw}$\times$\fld{scaling}
计入（:build\_three\_phase\_hybrid\_model）；不可控 DC 静态发电机、光伏阵列与 DC 储能以负负荷口径计入
（:build\_three\_phase\_hybrid\_model）；可控 DC 静态发电机被拒绝（:build\_three\_phase\_hybrid\_model）。若全系统未声明
\fld{DC\_V} 母线，第一台 VSC 的直流端子锚定直流电压，并登记
\fld{model\_limitations}（:build\_three\_phase\_hybrid\_model）。

每台在运 VSC 生成一个 \fld{PhaseVSC}：交流三端子按相号映射到相节点，缺任一相
即抛错（src/optimal\_power\_flow/three\_phase\_hybrid\_adapter.cpp:build\_three\_phase\_hybrid\_model）；效率截断到 $[0.01,1]$（:build\_three\_phase\_hybrid\_model）；容量
\fld{s\_max\_pu} 取 \fld{p\_rated\_mw} 与各功率限幅绝对值的最大者除以基准，
选项 \fld{enforce\_converter\_capacity} 关闭时取 $10^6$（:build\_three\_phase\_hybrid\_model）；相电流上限
取 \fld{i\_ac\_max\_pu}，选项 \fld{enforce\_converter\_current} 关闭或非正时取
$10^6$（:build\_three\_phase\_hybrid\_model）；\fld{ac\_grid\_forming} 或 \fld{AC\_GRID\_FORMING} 模式映射
为构网型下垂，其余映射为均分相功率（:build\_three\_phase\_hybrid\_model）；虚拟阻抗回退值
$r_v=0.01$、$x_v=0.10$~pu（:build\_three\_phase\_hybrid\_model）；电压参考取 \fld{v\_ac\_set\_pu}，非正回退
1.0（:build\_three\_phase\_hybrid\_model）。装配结束时登记总括性限制说明：adapter 仅支持相域交流导纳、
电阻型 DC 支路、直接 VSC 端口与恒功率节点负荷方程（:build\_three\_phase\_hybrid\_model）。

\begin{paramtable}
\fld{include\_shunts} & — & — & true/false & true & 并联补偿是否压入 \fld{y\_ac}（src/optimal\_power\_flow/three\_phase\_hybrid\_adapter.cpp:build\_three\_phase\_hybrid\_model）\\
\fld{enforce\_converter\_capacity} & — & — & true/false & true & 是否施加 VSC 视在功率上限，false 时 \fld{s\_max\_pu} 取 $10^6$（src/optimal\_power\_flow/three\_phase\_hybrid\_adapter.cpp:build\_three\_phase\_hybrid\_model）\\
\fld{enforce\_converter\_current} & — & — & true/false & true & 是否施加 VSC 相电流上限，false 或限值非正时取 $10^6$（src/optimal\_power\_flow/three\_phase\_hybrid\_adapter.cpp:build\_three\_phase\_hybrid\_model）\\
\fld{allow\_constant\_power\_load\_equivalent} & — & — & true/false & true & 允许 delta/电压相关负荷以恒功率运行点等值装配；false 时遇此类负荷抛错（src/optimal\_power\_flow/three\_phase\_hybrid\_adapter.cpp:build\_three\_phase\_hybrid\_model）\\
\fld{vuf\_max} & $\varepsilon_{\mathrm{vuf}}$ & pu & $[0,1]$（装配时截断） & 0.03 & 三相母线负序/正序电压幅值比上限（src/optimal\_power\_flow/three\_phase\_hybrid\_adapter.cpp:build\_three\_phase\_hybrid\_model）\\
\end{paramtable}

\subsection{相域数学模型}\label{subsec:tp-model}

\subsubsection*{决策变量与布局}
OPF 本体采用\textbf{直角坐标}电压表示：相节点 $k$ 的电压
$V_k=e_k+\mathrm{j}f_k$（\srcpath{src/optimal\_power\_flow/three\_phase\_hybrid\_opf.cpp:model\_voltage}，
函数 \srcpath{model\_voltage}）。原始决策向量按
\begin{align*}
x=\bigl[e;\,f;\,p_g;\,q_g;\,u;\,p^{ac};\,q^{ac};\,e^{gfm};\,f^{gfm};\,p^{dc}\bigr]
\end{align*}
排布：$e,f\in\mathbb{R}^{n_v}$ 为（降阶后）相节点电压实/虚部，$p_g,q_g\in\mathbb{R}^{n_g}$
为分相发电机出力，$u\in\mathbb{R}^{n_{dc}}$ 为直流节点电压，$p^{ac},q^{ac}\in\mathbb{R}^{n_{cp}}$
为换流器分相交流功率，$e^{gfm},f^{gfm}\in\mathbb{R}^{n_{gfm}}$ 为构网型换流器 A 相
内电势，$p^{dc}\in\mathbb{R}^{n_c}$ 为换流器直流功率（布局结构 \fld{Layout} 与各段
偏移：src/optimal\_power\_flow/three\_phase\_hybrid\_opf.cpp:Layout、482--499）。维度计数为
（\srcpath{src/optimal\_power\_flow/three\_phase\_hybrid\_opf.cpp:build\_model\_data}）：
\begin{align*}
n_{\mathrm{var}} &= 2n_v+2n_g+n_{dc}+2n_{cp}+2n_{gfm}+n_c,\\
n_{\mathrm{eq}} &= 2n_v+n_{dc}+n_c+n_{\mathrm{dcref}}+n_{\mathrm{ctrl}}+2n_{\mathrm{ref}},\\
n_{\mathrm{ineq}} &= 2n_{\mathrm{full}}+n_{3\varphi}+n_c+n_{cp},
\end{align*}
其中 $n_{\mathrm{ctrl}}$ 为换流器控制等式数：构网型每台 $2n_{cp,c}$ 条，其余
每台 $2(n_{cp,c}-1)$ 条（src/optimal\_power\_flow/three\_phase\_hybrid\_opf.cpp:build\_model\_data）；$n_{\mathrm{full}}$ 为降阶前相节点数，
$n_{3\varphi}$ 为三相齐全母线数。所有变量均为连续量，默认界 $\pm10^{20}$，
随后按物理意义收紧：$e,f\in[-1.2,1.2]$，$p_g,q_g$ 取机组界，
$u\in[\fld{v\_dc\_min\_pu},\fld{v\_dc\_max\_pu}]$，$p^{ac},q^{ac}\in[-s_{\max},s_{\max}]$
（定单位功率因数时 $q^{ac}=0$），$e^{gfm},f^{gfm}\in[-1.5,1.5]$，
$p^{dc}\in[-s_{\max},s_{\max}]$（\srcpath{src/optimal\_power\_flow/three\_phase\_hybrid\_opf.cpp:build\_nlp}，
函数 \srcpath{build\_nlp}）。

\subsubsection*{三相支路模型与节点功率平衡}
交流网络以 $n\times n$ 复稀疏耦合导纳矩阵 \fld{y\_ac} 建模（$n$ 为相节点数，
三相支路对应 $3\times 3$ 满块）。节点电流与复功率注入为
\begin{align*}
I &= Y_{ac}\,V + i_{\mathrm{fix}}, &
S_k &= V_k\,\overline{I_k},
\end{align*}
其中 $i_{\mathrm{fix}}$ 为固定电流注入向量 \fld{i\_ac\_fixed}（装配层恒零；
\srcpath{src/optimal\_power\_flow/three\_phase\_hybrid\_opf.cpp:evaluate\_equalities}，
函数 \srcpath{evaluate\_equalities}）。每个相节点的有功/无功平衡等式为
\begin{align*}
g_k &= \operatorname{Re}\{S_k\} + p^{\mathrm{load}}_k
  - \sum_{g\in\mathcal{G}(k)} p_{g} - \sum_{c\in\mathcal{C}(k)} p^{ac}_{c,\varphi(k)} = 0,\\
g_{n_v+k} &= \operatorname{Im}\{S_k\} + q^{\mathrm{load}}_k
  - \sum_{g\in\mathcal{G}(k)} q_{g} - \sum_{c\in\mathcal{C}(k)} q^{ac}_{c,\varphi(k)} = 0,
\end{align*}
$\mathcal{G}(k)$、$\mathcal{C}(k)$ 分别为挂在节点 $k$ 的发电机与换流器相端口
（\srcpath{src/optimal\_power\_flow/three\_phase\_hybrid\_opf.cpp:evaluate\_equalities}；
发电机项:evaluate\_equalities，换流器项:evaluate\_equalities）。解析雅可比按 $S_k$ 对 $(e,f)$ 的偏导
展开装配（\srcpath{equality\_jacobian}，src/optimal\_power\_flow/three\_phase\_hybrid\_opf.cpp:equality\_jacobian，其中导纳块偏导:equality\_jacobian）。

\subsubsection*{直流网络与换流器功率耦合}
直流侧以电导矩阵 \fld{g\_dc} 建模，节点 $k$ 的平衡等式与换流器耦合等式为
\begin{align*}
u_k\,(G_{dc}\,u)_k + p^{\mathrm{dc,load}}_k - \sum_{c\in\mathcal{C}_{dc}(k)} p^{dc}_c &= 0,\\
p^{dc}_c + \eta_c \sum_{\varphi} p^{ac}_{c,\varphi} &= 0,
\end{align*}
即直流注入等于电阻网络功率加换流器抽取，换流器直流功率与三相交流有功之和
以效率 $\eta_c$ 闭合（\srcpath{src/optimal\_power\_flow/three\_phase\_hybrid\_opf.cpp:evaluate\_equalities}；
直流参考电压锚定 $u_t=u_{\mathrm{ref},t}$ 见:evaluate\_equalities）。该耦合为单效率近似，
不区分运行点的损耗曲线。

\subsubsection*{换流器相域控制方程}
每台换流器按 \fld{control\_mode} 选择一组等式（装配位置 src/optimal\_power\_flow/three\_phase\_hybrid\_opf.cpp:evaluate\_equalities，
类型枚举 \fld{PhaseVSCControlMode} 见 \srcpath{include/hacdcpf/optimal\_power\_flow/three\_phase\_hybrid\_opf.hpp:PhaseVSCControlMode}）：

（i）\textbf{均分相功率}（\fld{EqualPhasePower}）：以首相为基准，
\begin{align*}
p^{ac}_{\varphi} - p^{ac}_{a} = 0,\qquad
q^{ac}_{\varphi} - q^{ac}_{a} = 0,\qquad \varphi\neq a
\end{align*}
（\srcpath{src/optimal\_power\_flow/three\_phase\_hybrid\_opf.cpp:evaluate\_equalities}）。

（ii）\textbf{跟网型正序电流}（\fld{GridFollowingPLL}，别名
\fld{PositiveSequenceCurrent}）：记 $a=e^{\mathrm{j}2\pi/3}$，对
$\varphi\in\{b,c\}$ 施加
\begin{align*}
s_{\varphi}\,V_a - r_{\varphi}\,V_{\varphi}\,s_a = 0,\qquad
r_b=a,\; r_c=a^{2},
\end{align*}
其中 $s_\varphi=p^{ac}_\varphi+\mathrm{j}q^{ac}_\varphi=V_\varphi\overline{I_\varphi}$。
在电压非零时该式等价于 $I_\varphi=\overline{r_\varphi}\,I_a$，即换流器只输出
正序电流（\srcpath{src/optimal\_power\_flow/three\_phase\_hybrid\_opf.cpp:evaluate\_equalities}；
雅可比:equality\_jacobian；拉格朗日黑塞:lagrangian\_hessian）。

（iii）\textbf{构网型下垂}（\fld{GridFormingDroop}）：引入 A 相内电势
$E_a=e^{gfm}+\mathrm{j}f^{gfm}$，三相内电势取正序组
$E_\varphi\in\{E_a,\,a^{2}E_a,\,aE_a\}$，对每相施加
\begin{align*}
\overline{E_\varphi}\,V_\varphi - |V_\varphi|^{2} - \overline{z_v}\,s_\varphi = 0,
\qquad z_v = r_v + \mathrm{j}x_v,
\end{align*}
在 $V_\varphi\neq 0$ 时等价于诺顿等值 $E_\varphi = V_\varphi + z_v I_\varphi$
（\srcpath{src/optimal\_power\_flow/three\_phase\_hybrid\_opf.cpp:evaluate\_equalities}；
雅可比:equality\_jacobian）。校验要求构网型换流器虚拟阻抗非零（\srcpath{validate\_case}，
src/optimal\_power\_flow/three\_phase\_hybrid\_opf.cpp:validate\_case），序感知控制（后两种）要求三端子恰为相 A/B/C 且
\fld{ac\_phase\_index} 完整（src/optimal\_power\_flow/three\_phase\_hybrid\_opf.cpp:validate\_case）。

\textbf{诚实口径}：\fld{PhaseVSC} 中的 \fld{p\_droop\_pu}、\fld{q\_droop\_pu}、
\fld{voltage\_reference\_pu}、\fld{nominal\_frequency\_hz} 等字段被复制进内部
映射（src/optimal\_power\_flow/three\_phase\_hybrid\_opf.cpp:build\_model\_data）但\textbf{不进入任何等式/不等式}——所谓“下垂”在 OPF
中仅体现为上述诺顿耦合，功率-频率/无功-电压下垂律与电压设定值尚未构成约束；
\fld{pll\_kp}/\fld{pll\_ki} 与 \fld{voltage\_integral\_gain} 仅用于解后动态
平衡点恢复（src/optimal\_power\_flow/three\_phase\_hybrid\_opf.cpp:recover\_dynamic\_equilibria）。

\subsubsection*{不平衡约束（VUF）}
记三相齐全母线的节点电压为 $(V_a,V_b,V_c)$，正、负序分量按
\begin{align*}
V_1 = \tfrac{1}{3}\bigl(V_a + aV_b + a^{2}V_c\bigr),\qquad
V_2 = \tfrac{1}{3}\bigl(V_a + a^{2}V_b + aV_c\bigr)
\end{align*}
定义（\srcpath{src/optimal\_power\_flow/three\_phase\_hybrid\_opf.cpp:sequence\_components}，
函数 \srcpath{sequence\_components}），并对每座三相母线施加
\begin{align*}
|V_2|^{2} \;\le\; \varepsilon_{\mathrm{vuf}}^{2}\,|V_1|^{2},
\end{align*}
$\varepsilon_{\mathrm{vuf}}$ 即 \fld{vuf\_max}（\srcpath{src/optimal\_power\_flow/three\_phase\_hybrid\_opf.cpp:evaluate\_inequalities}，
函数 \srcpath{evaluate\_inequalities}；雅可比:inequality\_jacobian）。该约束只施加于
\fld{three\_phase\_bus\_nodes} 中登记的齐全三相母线，两相/单相分支不参与。

\subsubsection*{不等式约束汇总}
全部非线性不等式以 $h(x)\le 0$ 形式给出（求值
\srcpath{evaluate\_inequalities}，src/optimal\_power\_flow/three\_phase\_hybrid\_opf.cpp:evaluate\_inequalities；雅可比
\srcpath{inequality\_jacobian}，src/optimal\_power\_flow/three\_phase\_hybrid\_opf.cpp:inequality\_jacobian；拉格朗日黑塞的参数化模板
\srcpath{build\_inequality\_hessian\_template}，src/optimal\_power\_flow/three\_phase\_hybrid\_opf.cpp:build\_inequality\_hessian\_template）：
\begin{align*}
\underline{v}_k^{2} - |V_k|^{2} &\le 0, &
|V_k|^{2} - \overline{v}_k^{2} &\le 0, &
&\forall\,k,\\
|V_2|^{2} - \varepsilon_{\mathrm{vuf}}^{2}|V_1|^{2} &\le 0, &
\sum_{\varphi}\bigl((p^{ac}_{c,\varphi})^{2}+(q^{ac}_{c,\varphi})^{2}\bigr) - s_{\max,c}^{2} &\le 0, &
&\forall\,c,\\
(p^{ac}_{c,\varphi})^{2}+(q^{ac}_{c,\varphi})^{2}
  - i_{\max,c}^{2}\,|V_{c,\varphi}|^{2} &\le 0, &
& &\forall\,c,\varphi,
\end{align*}
依次为相节点电压幅值平方窗口（:evaluate\_inequalities）、VUF（:evaluate\_inequalities）、换流器三相总
视在功率（:evaluate\_inequalities）、换流器分相电流（:evaluate\_inequalities；由 $p^2+q^2\le i_{\max}^2|V|^2$
给出，$|V|$ 为恢复后的全节点电压）。变量界见决策变量一节。换流器相电流上限
未配置时回退 $s_{\max}/\sqrt{n_{\varphi}}$（src/optimal\_power\_flow/three\_phase\_hybrid\_opf.cpp:build\_model\_data）。

\subsubsection*{目标函数}
目标为发电成本加小幅正则化（\srcpath{src/optimal\_power\_flow/three\_phase\_hybrid\_opf.cpp:objective}，
函数 \srcpath{objective}）：
\begin{align*}
f = \sum_{g}\Bigl[c_{2,g}\bigl(S_{\mathrm{base}}p_g\bigr)^{2}
  + c_{1,g}\bigl(S_{\mathrm{base}}p_g\bigr)\Bigr]
  + 10^{-4}\sum_{g} q_g^{2}
  + 10^{-3}\sum_{c,\varphi}\bigl((p^{ac}_{c,\varphi})^{2}+(q^{ac}_{c,\varphi})^{2}\bigr)
  + 10^{-4}\sum_{c} (p^{dc}_c)^{2},
\end{align*}
其中成本以有名值 MW 计（$p$ 为标幺，$S_{\mathrm{base}}=$ \fld{base\_mva}），
三项正则化权重 $10^{-4}/10^{-3}/10^{-4}$ 为源码硬编码、不可配置
（梯度 \srcpath{objective\_gradient}，src/optimal\_power\_flow/three\_phase\_hybrid\_opf.cpp:objective\_gradient；黑塞
\srcpath{objective\_hessian}，src/optimal\_power\_flow/three\_phase\_hybrid\_opf.cpp:objective\_hessian）。成本币种代码中未固定。

\subsubsection*{图降阶与电压恢复}
\fld{ModelVariant::GraphReduced} 下，所有无注入的纯传输相节点经稀疏 Kron
（Schur 补）降阶消去：带负荷/固定电流注入、发电机、换流器端子、参考节点
一律保留（\srcpath{src/optimal\_power\_flow/three\_phase\_hybrid\_opf.cpp:build\_model\_data}，
函数 \srcpath{build\_model\_data}；降阶核为
\srcpath{graph::reduce\_sparse\_kron}，受 \fld{reduction\_options} 控制）。
\fld{ModelVariant::Full} 下保留全部节点，恢复算子取单位阵（:identity\_recovery）。
全节点电压由恢复算子线性还原
\begin{align*}
V_{\mathrm{full}} = R\,V_{\mathrm{model}}
\end{align*}
（\srcpath{src/optimal\_power\_flow/three\_phase\_hybrid\_opf.cpp:full\_voltage}，
函数 \srcpath{full\_voltage}；恢复行缓存:build\_model\_data）。任一全节点的恢复行为空
表明该节点与参考源电气脱开，此时抛运行错误并提示先移除无源组件
（:build\_model\_data）。消去节点数记入结果 \fld{eliminated\_phase\_nodes}（:solve\_three\_phase\_hybrid\_opf\_impl）。
注意电压窗口、VUF 等不等式始终定义在\textbf{全节点}上（
\srcpath{evaluate\_inequalities} 先恢复全电压再求值，src/optimal\_power\_flow/three\_phase\_hybrid\_opf.cpp:evaluate\_inequalities），
因此降阶不削弱约束覆盖面。

\subsubsection*{初值构造}
初值由 \srcpath{build\_initial\_point} 生成（src/optimal\_power\_flow/three\_phase\_hybrid\_opf.cpp:build\_initial\_point）：电压取
\fld{voltage\_start}（:build\_initial\_point），换流器功率按直流侧初始平衡反推并均分
（:build\_initial\_point）；随后进行“交流潮流初始化 + 换流器控制投影”三轮交替
（:build\_initial\_point），最后按节点平衡回算发电机出力并截断到上下界（:build\_initial\_point）。
交流潮流初始化 \srcpath{initialize\_ac\_power\_flow}（src/optimal\_power\_flow/three\_phase\_hybrid\_opf.cpp:initialize\_ac\_power\_flow）先尝试
直接牛顿法（至多 20 次迭代，线搜索要求全节点电压落在 $(0.5,1.3)$~pu，
:initialize\_ac\_power\_flow）；失败则回退 $x$ 并改用 Z-bus 不动点延拓（20 个负荷步、每步
至多 80 次迭代，电压幅值下限保护 0.2~pu，:initialize\_ac\_power\_flow）。存在非参考节点上的
发电机时初始化直接跳过（返回 false，:initialize\_ac\_power\_flow；调用方忽略该返回值并继续，
:build\_initial\_point——属已知粗糙点）。换流器控制投影
\srcpath{project\_converter\_control}（src/optimal\_power\_flow/three\_phase\_hybrid\_opf.cpp:project\_converter\_control）按控制律把当前总功率
重分配到各相：构网型反解 $E_a$ 并按 $I_\varphi=(E_\varphi-V_\varphi)/z_v$
重算相功率（:project\_converter\_control），跟网型按 $I_a=\overline{S_{\mathrm{tot}}}/(3V_1)$
正序重分配（:project\_converter\_control）。

\subsection{求解流程与求解器接口}\label{subsec:tp-solve}

\subsubsection*{主流程}
\srcpath{solve\_three\_phase\_hybrid\_opf\_impl}（src/optimal\_power\_flow/three\_phase\_hybrid\_opf.cpp:solve\_three\_phase\_hybrid\_opf\_impl）的步骤为：
数据校验与模型装配 $\to$ \fld{engine::NLPModel} 回调装配 $\to$ 初值构造
$\to$（NativeIPM 路径 Phase I）有预算的原始可行性恢复
$\to$ 对偶/松弛初始化 $\to$（Phase II）MIPSolvers Native IPM
$\to$ 结果归因与 KKT 审计。这里的 Phase I/II 是“可行原--对偶状态构造/最优性
求解”的接口分层，不是平衡 AC OPF 的 $t=0\to1$ 目标同伦。求解器容差统一取自
\fld{options.tolerance}，acceptable 级别放宽 10 倍（src/optimal\_power\_flow/three\_phase\_hybrid\_opf.cpp:solve\_three\_phase\_hybrid\_opf\_impl）。

\subsubsection*{后端：NativeIPM 与 Ipopt}
\fld{SolverBackend::NativeIPM} 走 MIPSolvers 原生内点核：滤波器全局化、关闭
问题缩放，Phase II 的 primal/dual/complementarity 始终使用
\fld{options.tolerance} 最终容差；Phase I 只构造 central corridor。适配器为
\srcpath{engine::NativeIPMAdapter}，求解调用:solve\_three\_phase\_hybrid\_opf\_impl；核头文件引入
见 src/optimal\_power\_flow/three\_phase\_hybrid\_opf.cpp:solve\_three\_phase\_hybrid\_opf\_impl）。该路径先做原始可行性恢复
\srcpath{restore\_primal\_feasibility}。若模型声明了自由控制列 $F$，其补集为
state/basic 列 $B$，每轮只解稀疏方阵 Newton 校正
\begin{equation}
 J_B(x_k)\,\Delta x_B=-g(x_k),\qquad \Delta x_F=0,
 \label{eq:phase-one-basic-newton}
\end{equation}
并用 SparseLU；分区不可用或方阵分解失败时，对矩形 $J$ 直接做 SparseQR。
该设计依据 Nocedal--Wright（2006）第 11.1 节的约束 Newton 校正，避免稠密
物化，也避免正规方程 $JJ^T$ 的条件数平方与额外填充。Deuflhard（2011）
第 2.2--2.3 节的单调阻尼规则只接受原 pu 坐标最大违反量严格下降的步；不引入
障碍参数序列、filter、二阶校正、Hessian 解或弹性变量，因而不复制 Phase II
的早期内点迭代。

Phase I 同时受迭代、稀疏数值分解和回溯三个硬预算约束。墙钟预算在分解之间
检查；单次稀疏分解不可协作中断，故它是软截止时间，可能被最后一个不可分割
分解略微越过，而总分解次数仍给出有限工作上界。有限墙钟预算
$\fld{phase\_one\_time\_limit\_ms}>0$ 时绝不调用 Ipopt；旧
\fld{warm\_start\_with\_ipopt} 只在显式无限预算（该字段 $\le0$）时保留为兼容
路径。
默认 $\mu_0=0.1$、$\eta_a=1$、$\eta_p=0.1$、$\eta_c=0.5$。初始违反量高于
$\eta_a\mu_0$ 时以零分解退出 Phase I；否则恢复到
$\epsilon_p=\max(\mathrm{tol},\eta_p\mu_0)$。slack floor 使用剩余 primal
corridor 的一半，并令 $\nu_i=\mu_0/s_i$，同时控制扰动 primal residual、乘子
尺度和中心性。Phase I 在原始 pu 坐标下重新审计，并请求 MIPSolvers 独立的
\fld{central\_warm\_start} 契约；它不同于最终可行的
\fld{primal\_feasible\_start}。只有完整 primal/dual/slack、原坐标与扰动 primal
residual、中心性及乘子尺度全部通过才保留 $x_0$，否则自动回到普通冷初始化。
审计请求/接受状态和实际稀疏线性求解
后端分别由 \fld{phase\_two\_start\_requested}、
\fld{phase\_two\_start\_accepted} 与
\fld{phase\_two\_linear\_solver\_backend} 返回。
\fld{SolverBackend::Ipopt} 直接把同一 NLP 交给
\srcpath{engine::IpoptAdapter}（src/optimal\_power\_flow/three\_phase\_hybrid\_opf.cpp:solve\_three\_phase\_hybrid\_opf\_impl）。\textbf{门控说明}：Ipopt
后端依赖嵌入式 Ipopt（宏 \srcpath{HACDCPF\_HAVE\_IPOPT}）；未编译时适配器不抛错，
而是返回 \fld{Unavailable:} 前缀的失败状态（
\srcpath{MIPSolvers/src/engine/solver/external/adapters.cpp:IpoptAdapter::available}），结果即
不收敛。默认后端为 Ipopt（\srcpath{include/hacdcpf/optimal\_power\_flow/three\_phase\_hybrid\_opf.hpp:ThreePhaseHybridOPFOptions}）。

\subsubsection*{对偶初始化与拉格朗日黑塞}
NativeIPM 路径的对偶初值由 \srcpath{initialize\_primal\_dual\_start} 构造：
仅当初点不超过上述 $\epsilon_p$ corridor 时启用；
松弛变量由 $s=-h(x_0)$ 截断生成（:initialize\_primal\_dual\_start），不等式乘子按障碍参数
$\mu/s$ 恢复（:initialize\_primal\_dual\_start）；等式乘子优先利用
\fld{equality\_free\_columns} 划分做基解（仅当 $n_{\mathrm{var}}-n_{\mathrm{eq}}=n_c$
时以每台换流器首相 $q^{ac}$ 为自由列），否则对稳态方程做一次 SparseQR
最小二乘恢复。求解器未返回等式对偶时（如 Ipopt 路径），在同一解点上重跑一次
该初始化并仅当对偶残差不超过 $10^{-6}$ 时采纳（src/optimal\_power\_flow/three\_phase\_hybrid\_opf.cpp:solve\_three\_phase\_hybrid\_opf\_impl）。

Phase I 的可行点通常不是最优点，因此完整 stationarity
$r_{\mathrm{stat}}=\nabla f+J_g^T\lambda+J_h^T\nu$ 应保留控制空间 reduced
gradient 给 Phase II，而不应被 Phase I 人为消去。对声明了 $B/F$ 分区的模型，
dual 质量证书为
\begin{equation}
 r_{\mathrm{dual,fit}}=
 \frac{\|r_{\mathrm{stat},B}\|_\infty}
 {1+\max(\|\lambda\|_\infty,\|\nu\|_\infty)};
 \label{eq:phase-one-dual-fit}
\end{equation}
SparseQR 回退则报告对应最小二乘正规残差
$\|J_g r_{\mathrm{stat}}\|_\infty$ 的同尺度值。结果中的
\fld{phase\_one\_dual\_fit\_residual} 用此口径判定 warm start 质量，
\fld{initial\_dual\_residual} 继续诚实报告完整 stationarity，二者不可互换。
拉格朗日黑塞 \srcpath{lagrangian\_hessian}（src/optimal\_power\_flow/three\_phase\_hybrid\_opf.cpp:lagrangian\_hessian）按目标黑塞、
潮流等式双线性项（:lagrangian\_hessian）、换流器控制等式项（:lagrangian\_hessian）与参数化
不等式模板（:lagrangian\_hessian）装配。

在 360 个三相母线、1080 个相节点、2175 个变量的稀疏回归上，初始违反量
$10^{-3}$ 已在默认 $10^{-2}$ corridor 内，因此 primal Newton 为 0 次，只做
一次 sparse-basis dual fit；dual-fit 为 $8.33\times10^{-15}$，五次计时中位数
为 5.438~ms。相较旧的严格 Phase I（两次 Newton、三次总分解），这里消除了
两次不必要的 primal 分解而不放宽 Phase II 最终证书。

\subsubsection*{KKT 审计与收敛口径}
求解器报成功后，模块在解点上\textbf{重新}评估物理等式与\textbf{全部}不等式
（含被预言机省略者），并施加审计阈值：原始/对偶/互补残差、电压越限、换流器
越限均不超过 $10^{-6}$ 才把 \fld{converged} 置真（
\srcpath{src/optimal\_power\_flow/three\_phase\_hybrid\_opf.cpp:solve\_three\_phase\_hybrid\_opf\_impl}，
审计条件:solve\_three\_phase\_hybrid\_opf\_impl）。即求解器自身成功只是必要条件。
\fld{verify\_derivatives} 为真时另做中心差分（步长 $10^{-5}$）校核等式/不等式
雅可比与拉格朗日黑塞，误差上界写入结果（:solve\_three\_phase\_hybrid\_opf\_impl；乘子取固定的伪随机模式
:solve\_three\_phase\_hybrid\_opf\_impl）。

\subsubsection*{约束预言机}
\fld{use\_constraint\_oracle} 为真时启用精确约束预言机
（src/optimal\_power\_flow/three\_phase\_hybrid\_opf.cpp:solve\_three\_phase\_hybrid\_opf\_impl）：以初值处 $h_j(x_0)\ge -\max(m_{\mathrm{init}},m_{\mathrm{act}})$
的行、全部换流器约束行与 \fld{oracle\_seed\_rows} 构成初始强制集
（:solve\_three\_phase\_hybrid\_opf\_impl）；每轮求解受限问题后在全量不等式上筛查，把
$h_j(x)\ge -m_{\mathrm{act}}$ 的行补入（:solve\_three\_phase\_hybrid\_opf\_impl）；
\fld{oracle\_probe\_iterations} 为正值时先跑“预测轮”（限迭代数），活跃集
稳定后再跑精确校正轮（:solve\_three\_phase\_hybrid\_opf\_impl）；相邻轮之间传递原始点与对偶
（:solve\_three\_phase\_hybrid\_opf\_impl）。达到 \fld{oracle\_max\_rounds} 仍有新增行时判不收敛
（:solve\_three\_phase\_hybrid\_opf\_impl）。被省略行的最大违反量记入 \fld{max\_omitted\_inequality}
（:solve\_three\_phase\_hybrid\_opf\_impl），作为近似口径的公开声明。

\subsubsection*{序列与分支求解}
\srcpath{solve\_three\_phase\_hybrid\_opf\_sequence}（src/optimal\_power\_flow/three\_phase\_hybrid\_opf.cpp:solve\_three\_phase\_hybrid\_opf\_sequence）要求各
时段网络与设备拓扑完全一致（否则抛错，:solve\_three\_phase\_hybrid\_opf\_sequence），用各时段最大负荷构造一次
共享 \fld{ModelData}（:solve\_three\_phase\_hybrid\_opf\_sequence），逐时段传递上一步的原始点、等式/不等式对偶
与松弛作为热启动，并把上一步强制行集作为预言机种子（:solve\_three\_phase\_hybrid\_opf\_sequence）。
\srcpath{solve\_three\_phase\_hybrid\_opf\_branches}（src/optimal\_power\_flow/three\_phase\_hybrid\_opf.cpp:solve\_three\_phase\_hybrid\_opf\_branches）从已认证的
基准点出发做同样的参数摄动分支求解，基准未收敛则抛错（:solve\_three\_phase\_hybrid\_opf\_branches）。

\subsubsection*{与单相 AC OPF 的接口关系}
本章模块与第~\ref{sec:acopfnative}~章的单相 AC OPF Native IPM 的关系是
“同核不同模”：两者共享 MIPSolvers 的 \fld{engine::NLPModel} 抽象与原生内点核
（及嵌入式 Ipopt 适配器），但三相模块不复用单相的 \fld{SolverData} 装配、
总线类型（PQ/PV/松弛）机制与投影回传，而是自建相域单体 NLP。结果侧的接口有
两个：其一，\srcpath{make\_three\_phase\_hybrid\_pf\_case}（src/optimal\_power\_flow/three\_phase\_hybrid\_opf.cpp:make\_three\_phase\_hybrid\_pf\_case）把
收敛运行点翻译为三相混合潮流算例 \fld{powerflow::ThreePhaseHybridPFCase}——
换流器总功率固定、按控制模式与是否直流电压参考映射到
\fld{EqualPhasePQ}\allowbreak /\fld{EqualPhaseVdcQ}\allowbreak /\fld{GridFollowingPQ}\allowbreak /\fld{GridFollowingVdcQ}\allowbreak /\fld{GridFormingVoltage}\allowbreak /\fld{GridFormingVdc} 六类潮流
控制（:make\_three\_phase\_hybrid\_pf\_case），供独立潮流复核 OPF 运行点（要求已收敛且维度匹配，
:make\_three\_phase\_hybrid\_pf\_case）；其二，上文序列/分支接口支撑时序与 N-1 类参数扫描。结果归因回
rich 组件依赖 adapter 保留的 \fld{generator\_bindings} 与
\fld{vsc\_component\_indices}。

\subsection{SOCP 松弛与目标下界证书}\label{subsec:tp-relax}

\subsubsection*{提升变量}
松弛层把电压双线性项提升为新变量：交流侧对每个相节点引入
$W_{kk}=|V_k|^2$，对每对相关联节点引入
$W_{ij}=V_i\overline{V_j}=W^{R}_{ij}+\mathrm{j}W^{I}_{ij}$（$i<j$）；
直流侧引入 $X_{kk}=u_k^2$ 与 $X_{ij}=u_i u_j$。关联节点对取自 \fld{y\_ac} 非零
模式、三相齐全母线的相两两组合、PSD 团（clique）候选与参考节点两两组合
（\srcpath{src/optimal\_power\_flow/three\_phase\_hybrid\_relaxation.cpp:solve\_three\_phase\_hybrid\_opf\_relaxation}；
团候选由导纳矩阵列邻居与三相母线生成并剔除被包含者:solve\_three\_phase\_hybrid\_opf\_relaxation）。
变量界为 $W_{kk}\in[\underline{v}_k^2,\overline{v}_k^2]$（参考节点固定为
$|V_{\mathrm{ref}}|^2$）、$W^{R/I}_{ij}\in[-\overline{v}_i\overline{v}_j,\,
\overline{v}_i\overline{v}_j]$（双参考对固定为 $V_i\overline{V_j}$ 的对应值）、
$X_{kk}\in[\underline{u}_k^2,\overline{u}_k^2]$、$X_{ij}\in[\underline{u}_i
\underline{u}_j,\overline{u}_i\overline{u}_j]$（:solve\_three\_phase\_hybrid\_opf\_relaxation）。发电机变量
$p_g,q_g$ 按原界引入，$c_2>10^{-14}$ 者附加成本上图变量 $t_g$（:solve\_three\_phase\_hybrid\_opf\_relaxation）；
换流器变量 $p^{ac},q^{ac}\in[-s_{\max},s_{\max}]$（定单位功率因数时
$q^{ac}=0$）与 $p^{dc}$（:solve\_three\_phase\_hybrid\_opf\_relaxation）。

\subsubsection*{线性化后的平衡方程与不等式}
提升后节点功率平衡成为线性等式
\begin{align*}
\sum_{j}\operatorname{Re}\{\overline{Y_{kj}}W_{kj}\} + p^{\mathrm{load}}_k
  - \sum_{g\in\mathcal{G}(k)} p_g - \sum_{c\in\mathcal{C}(k)} p^{ac}_{c,\varphi(k)} &= 0,\\
\sum_{j}\operatorname{Im}\{\overline{Y_{kj}}W_{kj}\} + q^{\mathrm{load}}_k
  - \sum_{g\in\mathcal{G}(k)} q_g - \sum_{c\in\mathcal{C}(k)} q^{ac}_{c,\varphi(k)} &= 0,\\
\sum_{j} (G_{dc})_{kj}X_{kj} + p^{\mathrm{dc,load}}_k
  - \sum_{c\in\mathcal{C}_{dc}(k)} p^{dc}_c &= 0,
\end{align*}
装配时按 $i<j$ 定向修正 $W^{I}$ 项符号（
\srcpath{src/optimal\_power\_flow/three\_phase\_hybrid\_relaxation.cpp:solve\_three\_phase\_hybrid\_opf\_relaxation}；
定向处理:solve\_three\_phase\_hybrid\_opf\_relaxation）。换流器耦合 $p^{dc}_c+\eta_c\sum_\varphi p^{ac}_{c,\varphi}=0$
与均分相功率等式按原样保留（:solve\_three\_phase\_hybrid\_opf\_relaxation）。VUF 约束对 $W=VV^{H}$ 是线性的：
$|V_2|^2-\varepsilon_{\mathrm{vuf}}^2|V_1|^2=\mathrm{tr}\bigl((n_2n_2^{H}-
\varepsilon_{\mathrm{vuf}}^2 n_1n_1^{H})W\bigr)\le 0$，其中 $n_1,n_2$ 为序分量
系数向量（:solve\_three\_phase\_hybrid\_opf\_relaxation）。若 \fld{y\_ac}（厄米特部最小特征值不小于
$-10^{-8}\cdot\max(1,\lVert Y\rVert)$）或 \fld{g\_dc} 通过无源性检查
（:ac\_admittance\_is\_passive），各加一条聚合网损非负割：发电与换流器注入之和不小于总负荷
（:solve\_three\_phase\_hybrid\_opf\_relaxation；未通过检查时放弃该割并登记 \fld{model\_limitations}）。

\subsubsection*{锥割与外逼近循环}
精确条件 $|W_{ij}|^2\le W_{ii}W_{jj}$（交流）、$X_{ij}^2\le X_{ii}X_{jj}$
（直流）、$\sum_\varphi(p^2+q^2)\le s_{\max}^2$ 与
$p_\varphi^2+q_\varphi^2\le i_{\max}^2 W_{nn}$ 均为旋转二阶锥约束；松弛层不调用
锥求解器，而是在每个 LP 解点上对违例锥生成支撑超平面割（
\srcpath{src/optimal\_power\_flow/three\_phase\_hybrid\_relaxation.cpp:violated\_cone\_cuts}，
函数 \srcpath{violated\_cone\_cuts}）：交流提升锥
$\lVert(2W^{R},2W^{I},W_{ii}-W_{jj})\rVert_2\le W_{ii}+W_{jj}$（:append\_if\_violated（lambda））、
PSD 团最小特征向量的秩一割（\srcpath{append\_ac\_psd\_cuts}）、
直流提升锥（:append\_if\_violated（lambda））、换流器视在功率锥（:append\_if\_violated（lambda））、换流器相电流锥
（:append\_if\_violated（lambda））。外逼近循环至多 \fld{max\_outer\_approximation\_rounds}$+1$ 轮，
每轮把累积割并入 LP 交 HiGHS 求解，无违例割时置
\fld{soc\_outer\_approximation\_converged} 并退出（:solve\_three\_phase\_hybrid\_opf\_relaxation；HiGHS 不可用直接
返回 :solve\_three\_phase\_hybrid\_opf\_relaxation）。

\subsubsection*{目标与对偶下界}
LP 目标为 $\sum_g[c_{1,g}S_{\mathrm{base}}p_g + t_g]$，二次成本
$c_2'(p)=c_2S_{\mathrm{base}}^2p^2$ 以 \fld{cost\_tangent\_points} 条支撑切线
$t_g\ge 2c_2'p_0p-c_2'p_0^2$ 在 $[p_{\min},p_{\max}]$ 上均匀布点低估
（\srcpath{src/optimal\_power\_flow/three\_phase\_hybrid\_relaxation.cpp:solve\_three\_phase\_hybrid\_opf\_relaxation}）。
无功正则项与换流器正则项从目标中省略（:solve\_three\_phase\_hybrid\_opf\_relaxation 登记的
\fld{model\_limitations}）。每轮 LP 的对偶按下式构造有效下界：不等式对偶截断到
非正锥后，在变量盒界上逐变量最小化拉格朗日函数（
\srcpath{evaluate\_lp\_dual\_bound}）；各轮取最大写入
\fld{round\_lower\_bounds} 与 \fld{objective\_lower\_bound}（:solve\_three\_phase\_hybrid\_opf\_relaxation）。
由于每条割都是必要条件、成本被低估，该下界对原非凸 OPF 成立
（\fld{outer\_relaxation\_valid} 标记，:solve\_three\_phase\_hybrid\_opf\_relaxation）。

\subsubsection*{输入拒绝与已知省略（诚实口径）}
以下输入被显式拒绝：选项非法（轮数负、切点数 $<2$、容差非正，:solve\_three\_phase\_hybrid\_opf\_relaxation）；
\fld{i\_ac\_fixed} 非零（固定电流注入需要增广电压提升，未实现，:solve\_three\_phase\_hybrid\_opf\_relaxation）；
负二次成本（非凸，:solve\_three\_phase\_hybrid\_opf\_relaxation）。以下模型成分被显式省略并登记：序感知换流器
（跟网/构网）控制等式为非凸，从下界模型中剔除（:solve\_three\_phase\_hybrid\_opf\_relaxation）；目标正则项省略
（:solve\_three\_phase\_hybrid\_opf\_relaxation）；无源性检查失败时聚合网损割不施加（:solve\_three\_phase\_hybrid\_opf\_relaxation）。
因此松弛层输出仅是下界证书与诊断，不给出可行运行点；其
\fld{generator\_active\_power\_pu} 等向量是 LP 解的对应分量，不保证交流可实施。

\subsection{公开选项与结果字段}\label{subsec:tp-fields}

\subsubsection*{问题数据 \fld{ThreePhaseHybridOPFCase}}
\begin{paramtable}
\fld{name} & — & — & — & 空 & 算例名称（日志用）\\
\fld{base\_mva} & $S_{\mathrm{base}}$ & MVA & $>0$ & 1.0 & 标幺基准功率\\
\fld{y\_ac} & $Y_{ac}$ & pu & — & 空 & $n\times n$ 相域复导纳矩阵（稀疏）\\
\fld{i\_ac\_fixed} & $i_{\mathrm{fix}}$ & pu & — & 空 & 固定电流注入；松弛层要求为零\\
\fld{p\_load\_pu} & $p^{\mathrm{load}}$ & pu & — & 空 & 相节点恒功率有功负荷\\
\fld{q\_load\_pu} & $q^{\mathrm{load}}$ & pu & — & 空 & 相节点恒功率无功负荷\\
\fld{v\_min\_pu} & $\underline{v}$ & pu & 0.85--1.0 & 空 & 相节点电压幅值下限\\
\fld{v\_max\_pu} & $\overline{v}$ & pu & 1.0--1.15 & 空 & 相节点电压幅值上限\\
\fld{voltage\_start} & $V^{(0)}$ & pu & — & 空 & 相节点复电压初值\\
\fld{ac\_phase\_index} & — & — & 0/1/2 & 空 & 各相节点所属相号；序感知控制必需\\
\fld{reference\_nodes} & — & — & — & 空 & 交流参考（松弛）相节点列表\\
\fld{reference\_voltage} & $V_{\mathrm{ref}}$ & pu & — & 空 & 参考节点复电压设定\\
\fld{three\_phase\_bus\_nodes} & — & — & — & 空 & 三相齐全母线的节点三元组，VUF 作用集\\
\fld{vuf\_max} & $\varepsilon_{\mathrm{vuf}}$ & pu & $[0,1]$ & 0.03 & 负序/正序电压幅值比上限\\
\fld{generators} & — & — & — & 空 & \fld{PhaseGenerator} 列表（分相机组）\\
\fld{converters} & — & — & — & 空 & \fld{PhaseVSC} 列表\\
\fld{g\_dc} & $G_{dc}$ & pu & — & 空 & 直流电导矩阵（拉普拉斯型）\\
\fld{p\_dc\_load\_pu} & $p^{\mathrm{dc,load}}$ & pu & — & 空 & 直流节点恒功率负荷\\
\fld{v\_dc\_start} & $u^{(0)}$ & pu & — & 空 & 直流电压初值\\
\fld{v\_dc\_min\_pu} & $\underline{u}$ & pu & 0.85--1.0 & 空 & 直流电压下限\\
\fld{v\_dc\_max\_pu} & $\overline{u}$ & pu & 1.0--1.15 & 空 & 直流电压上限\\
\fld{dc\_reference\_terminals} & — & — & — & 空 & 直流电压参考端子\\
\fld{dc\_reference\_voltage\_pu} & $u_{\mathrm{ref}}$ & pu & — & 空 & 直流参考电压设定\\
\end{paramtable}

\subsubsection*{设备结构 \fld{PhaseGenerator} 与 \fld{PhaseVSC}}
\begin{paramtable}
\fld{phase\_node} & — & — & $[0,n)$ & $-1$ & 机组挂接的相节点\\
\fld{p\_min\_pu} / \fld{p\_max\_pu} & $p_{\min}/p_{\max}$ & pu & — & 0 & 有功出力求解界\\
\fld{q\_min\_pu} / \fld{q\_max\_pu} & $q_{\min}/q_{\max}$ & pu & — & 0 & 无功出力求解界\\
\fld{cost\_c2} & $c_2$ & — & $\ge 0$ & 0 & 二次成本系数（按 MW 计；负值被松弛层拒绝）\\
\fld{cost\_c1} & $c_1$ & — & — & 0 & 一次成本系数（按 MW 计）\\
\end{paramtable}

\begin{paramtable}
\fld{phase\_nodes} & — & — & — & 空 & 换流器交流相端子节点（序感知控制须为 A/B/C 三相）\\
\fld{dc\_terminal} & — & — & $[0,n_{dc})$ & $-1$ & 直流端子编号\\
\fld{efficiency} & $\eta$ & pu & $[0.01,1]$（使用时截断） & 0.98 & 交直流功率耦合效率\\
\fld{s\_max\_pu} & $s_{\max}$ & pu & $\ge 0$ & 0 & 三相总视在功率上限\\
\fld{phase\_current\_max\_pu} & $i_{\max}$ & pu & $\ge 0$ & 0 & 相电流上限；非正时回退 $s_{\max}/\sqrt{n_\varphi}$\\
\fld{fixed\_unity\_power\_factor} & — & — & true/false & false & true 时强制 $q^{ac}=0$\\
\fld{control\_mode} & — & — & EqualPhase\allowbreak Power/\allowbreak GridFollowingPLL/\allowbreak GridFormingDroop & EqualPhase\allowbreak Power & 相域控制律；\fld{PositiveSequenceCurrent} 为 \fld{GridFollowingPLL} 别名\\
\fld{nominal\_frequency\_hz} & $f_N$ & Hz & — & 50.0 & 额定频率（当前仅存档，未入方程）\\
\fld{pll\_kp} / \fld{pll\_ki} & $k_p^{\mathrm{pll}}/k_i^{\mathrm{pll}}$ & — & — & 0.01 / 1.0 & PLL 增益；仅用于解后动态平衡点恢复\\
\fld{virtual\_r\_pu} / \fld{virtual\_x\_pu} & $r_v/x_v$ & pu & $>0$（构网型必需） & 0 / 0.10 & 构网诺顿虚拟阻抗\\
\fld{p\_droop\_pu} / \fld{q\_droop\_pu} & — & pu & — & 0.01 / 0.05 & 下垂系数（当前未进入 OPF 方程，仅存档）\\
\fld{voltage\_reference\_pu} & $V_{\mathrm{ref}}^{c}$ & pu & — & 1.0 & 构网电压参考（当前未进入 OPF 方程）\\
\fld{voltage\_integral\_gain} & $k_i^{V}$ & — & — & 10.0 & 电压积分增益；仅用于平衡点恢复\\
\end{paramtable}

\subsubsection*{求解选项 \fld{ThreePhaseHybridOPFOptions}}
\begin{paramtable}
\fld{variant} & — & — & Full/\allowbreak GraphReduced & Full & 全节点或稀疏 Kron 图降阶模型（src/optimal\_power\_flow/three\_phase\_hybrid\_opf.cpp:build\_model\_data）\\
\fld{backend} & — & — & Ipopt/\allowbreak NativeIPM & Ipopt & NLP 后端；Ipopt 受嵌入式编译门控\\
\fld{reduction\_options} & — & — & — & 见下 & 图降阶参数 \fld{SparseKronOptions}：\fld{max\_front}=12、\fld{max\_nnz\_ratio}=2.0、\fld{pivot\_tolerance}=$10^{-12}$、\fld{drop\_tolerance}=$10^{-13}$\\
\fld{max\_iterations} & — & — & $\ge 1$ & 300 & 内点最大迭代数\\
\fld{tolerance} & $\mathrm{tol}$ & — & $>0$ & $10^{-7}$ & 收敛容差；acceptable 级别放宽 10 倍\\
\fld{phase\_one\_time\_limit\_ms} & — & ms & — & 5000 & Phase I 协作式软截止；$\le0$ 表示无限墙钟预算\\
\fld{phase\_one\_max\_iterations} & — & — & $\ge0$ & 12 & Phase I Newton 迭代硬上限\\
\fld{phase\_one\_max\_factorizations} & — & — & $\ge0$ & 14 & primal 恢复与 dual 恢复共用的稀疏数值分解硬上限\\
\fld{phase\_one\_max\_backtracks} & — & — & $\ge0$ & 12 & 每次 Newton 全步之外允许的阻尼回溯数\\
\fld{phase\_one\_barrier\_mu} & $\mu_0$ & — & $(0,0.1]$ & 0.1 & Phase-I central target\\
\fld{phase\_one\_admission\_mu\_factor} & $\eta_a$ & — & $\ge0$ & 1.0 & 局部 Phase I 准入半径 $\eta_a\mu_0$\\
\fld{phase\_one\_primal\_mu\_factor} & $\eta_p$ & — & $\ge0$ & 0.1 & handoff primal corridor $\eta_p\mu_0$\\
\fld{phase\_one\_centrality\_tolerance} & $\eta_c$ & — & $\ge0$ & 0.5 & 最大相对中心乘积偏差\\
\fld{warm\_start\_with\_ipopt} & — & — & true/false & false & 仅无限墙钟预算下启用的旧 Ipopt 可行性热解兼容路径\\
\fld{verify\_derivatives} & — & — & true/false & false & 中心差分校核雅可比/黑塞（src/optimal\_power\_flow/three\_phase\_hybrid\_opf.cpp:solve\_three\_phase\_hybrid\_opf\_impl）\\
\fld{verbose} & — & — & true/false & false & 输出阶段日志与初始化诊断\\
\fld{use\_constraint\_oracle} & — & — & true/false & false & 启用精确约束预言机（src/optimal\_power\_flow/three\_phase\_hybrid\_opf.cpp:solve\_three\_phase\_hybrid\_opf\_impl）\\
\fld{oracle\_initial\_margin} & $m_{\mathrm{init}}$ & — & $\ge 0$ & $10^{-3}$ & 初始筛查裕度\\
\fld{oracle\_activation\_margin} & $m_{\mathrm{act}}$ & — & $\ge 0$ & $10^{-5}$ & 轮次激活裕度\\
\fld{oracle\_probe\_iterations} & — & — & $\ge 0$ & 0 & 预测轮迭代数；0 关闭预测轮\\
\fld{oracle\_max\_rounds} & — & — & $\ge 1$ & 8 & 预言机最大轮数\\
\fld{oracle\_seed\_rows} & — & — & — & 空 & 初始强制不等式行（热启动用）\\
\fld{enforced\_inequality\_rows} & — & — & — & 空 & 高级覆盖：显式强制行集；空=全部不等式（\srcpath{include/hacdcpf/optimal\_power\_flow/three\_phase\_hybrid\_opf.hpp:enforced\_inequality\_rows}）\\
\fld{primal\_start} & $x_0$ & — & — & 空 & 原始初值覆盖（维度匹配且有限才采纳）\\
\fld{equality\_dual\_start} & $\lambda_0$ & — & — & 空 & 等式对偶种子；NaN 行表示待恢复\\
\fld{nonlinear\_inequality\_dual\_start} & $\nu_0$ & — & — & 空 & 非线性不等式对偶种子（全量或强制行两种长度皆可）\\
\fld{nonlinear\_slack\_start} & $s_0$ & — & — & 空 & 非线性不等式松弛种子\\
\end{paramtable}

\subsubsection*{求解结果 \fld{ThreePhaseHybridOPFResult}}
\begin{paramtable}
\fld{converged} & — & — & true/false & false & 求解成功且通过 $10^{-6}$ KKT 审计（src/optimal\_power\_flow/three\_phase\_hybrid\_opf.cpp:solve\_three\_phase\_hybrid\_opf\_impl）\\
\fld{status} / \fld{solver} & — & — & — & 空 & 求解器返回状态串 / 求解器名\\
\fld{iterations} & — & — & — & 0 & 内点迭代数（预言机为累计）\\
\fld{variables} / \fld{equalities} / \fld{inequalities} & — & — & — & 0 & $n_{\mathrm{var}}/n_{\mathrm{eq}}/n_{\mathrm{ineq}}$\\
\fld{enforced\_inequalities} & — & — & — & 0 & 实际强制的非线性不等式行数\\
\fld{constraint\_oracle\_rounds} & — & — & — & 0 & 预言机轮数\\
\fld{constraint\_oracle\_added\_rows} & — & — & — & 0 & 预言机累计补入行数\\
\fld{equality\_jacobian\_nonzeros} / \fld{inequality\_jacobian\_nonzeros} & — & — & — & 0 & 初值处雅可比非零元统计（src/optimal\_power\_flow/three\_phase\_hybrid\_opf.cpp:solve\_three\_phase\_hybrid\_opf\_impl）\\
\fld{max\_omitted\_inequality} & — & — & — & $-\infty$ & 被省略行在解点的最大违反量（src/optimal\_power\_flow/three\_phase\_hybrid\_opf.cpp:solve\_three\_phase\_hybrid\_opf\_impl）\\
\fld{eliminated\_phase\_nodes} & — & — & — & 0 & 图降阶消去的相节点数\\
\fld{objective} & $f$ & — & — & 0 & 目标值（成本币种未定）\\
\fld{runtime\_ms} & — & ms & — & 0 & 求解耗时（序列求解首项含共享建模耗时）\\
\fld{initial\_primal\_residual} / \fld{initial\_dual\_residual} & — & — & — & 0 & 初值原始/对偶残差（对偶初始化后评估）\\
\fld{phase\_one\_initial\_violation} / \fld{phase\_one\_constraint\_violation} & — & pu & — & $+\infty$ & Phase I 前后原坐标最大约束/变量界违反量\\
\fld{phase\_one\_dual\_fit\_residual} & — & — & — & $+\infty$ & 式~\eqref{eq:phase-one-dual-fit} 的 basic-fit 或 SparseQR 正规残差；不是完整 stationarity\\
\fld{phase\_one\_primal\_feasible} / \fld{phase\_one\_in\_handoff\_corridor} & — & — & true/false & false & 最终容差可行 / 有限精度 central handoff 邻域；二者不可混用\\
\fld{phase\_one\_dual\_initialized} & — & — & true/false & false & dual-fit/正性/维度证书\\
\fld{phase\_one\_handoff\_primal\_tolerance} / \fld{phase\_one\_perturbed\_primal\_residual} & — & pu & — & 0 / $+\infty$ & central primal 门槛与实测值\\
\fld{phase\_one\_centrality} / \fld{phase\_one\_barrier\_mu} & — & — & — & $+\infty$ / 0 & 相对中心偏差与目标障碍参数\\
\fld{phase\_one\_budget\_exhausted} / \fld{phase\_one\_termination} & — & — & — & false / not-run & 预算耗尽标记与停止原因\\
\fld{phase\_one\_iterations} / \fld{phase\_one\_factorizations} / \fld{phase\_one\_backtracks} & — & — & — & 0 & Phase I 实际硬预算消耗\\
\fld{phase\_one\_runtime\_ms} / \fld{phase\_one\_linear\_solver} & — & ms / — & — & 0 / unselected & Phase I 耗时与 SparseLU/SparseQR 路径\\
\fld{phase\_two\_start\_requested} / \fld{phase\_two\_start\_accepted} & — & — & true/false & false & Phase II 是否请求并通过独立 central warm-start 审计保留起点\\
\fld{phase\_two\_start\_rejection\_reason} & — & — & — & 空 & central 合同未接受原因\\
\fld{phase\_two\_linear\_solver\_backend} & — & — & — & unselected & Phase II 实际 KKT 后端；零迭代时可保持 unselected\\
\fld{initial\_worst\_equality} / \fld{initial\_worst\_inequality} & — & — & — & $-1$ & 初值最差等式/不等式行号\\
\fld{primal\_residual} / \fld{dual\_residual} / \fld{complementarity} & — & — & — & 0 & 解点 KKT 残差（解点重评估口径）\\
\fld{max\_voltage\_violation} & — & pu & — & 0 & 电压幅值窗口最大越限\\
\fld{max\_vuf} & — & pu & — & 0 & 三相母线最大 $|V_2|/|V_1|$（src/optimal\_power\_flow/three\_phase\_hybrid\_opf.cpp:max\_vuf）\\
\fld{max\_converter\_current\_vuf} & — & pu & — & 0 & 换流器相电流负序/正序比最大值（src/optimal\_power\_flow/three\_phase\_hybrid\_opf.cpp:max\_converter\_current\_vuf）\\
\fld{max\_converter\_current\_loading} & — & pu & — & 0 & 换流器相电流负载率最大值（src/optimal\_power\_flow/three\_phase\_hybrid\_opf.cpp:max\_converter\_current\_loading）\\
\fld{max\_converter\_violation} & — & pu & — & 0 & 换流器容量/电流约束最大越限（src/optimal\_power\_flow/three\_phase\_hybrid\_opf.cpp:solve\_three\_phase\_hybrid\_opf\_impl）\\
\fld{max\_dynamic\_equilibrium\_residual} & — & — & — & 0 & 动态平衡点最大微分残差（src/optimal\_power\_flow/three\_phase\_hybrid\_opf.cpp:recover\_dynamic\_equilibria）\\
\fld{max\_equality\_jacobian\_error} / \fld{max\_inequality\_jacobian\_error} / \fld{max\_lagrangian\_hessian\_error} & — & — & — & 0 & 差分导数校核误差（\fld{verify\_derivatives} 开启时）\\
\fld{primal} & $x$ & — & — & 空 & 原始解向量（模型布局）\\
\fld{equality\_dual} / \fld{inequality\_dual} / \fld{inequality\_slack} & $\lambda/\nu/s$ & — & — & 空 & 对偶与松弛（不等式部分已扩展回全量行）\\
\fld{full\_voltage} & $V_{\mathrm{full}}$ & pu & — & 空 & 恢复后的全节点复电压\\
\fld{dc\_voltage} & $u$ & pu & — & 空 & 直流节点电压\\
\fld{generator\_active\_power\_pu} / \fld{generator\_reactive\_power\_pu} & $p_g/q_g$ & pu & — & 空 & 分相机组出力\\
\fld{converter\_phase\_power\_pu} & $s^{ac}$ & pu & — & 空 & 每台换流器各相复功率\\
\fld{converter\_dc\_power\_pu} & $p^{dc}$ & pu & — & 空 & 每台换流器直流功率\\
\fld{enforced\_inequality\_rows} & — & — & — & 空 & 实际强制不等式行集\\
\fld{constraint\_oracle\_trace} & — & — & — & 空 & 各轮诊断 \fld{ConstraintOracleRoundDiagnostic}（轮次/行数/新增/迭代/收敛/耗时/全量最大违反）\\
\fld{converter\_dynamic\_equilibria} & — & — & — & 空 & 每台换流器的动态平衡点 \fld{PhaseVSCDynamicEquilibrium}（PLL 角/积分器、内电势、功率参考、微分残差）\\
\fld{reduction} & — & — & — & — & 图降阶结果 \fld{SparseKronResult}（含恢复证书）\\
\end{paramtable}

\subsubsection*{松弛层选项与结果}
\begin{paramtable}
\fld{max\_outer\_approximation\_rounds} & — & — & $\ge 0$ & 12 & 外逼近最大轮数（0 表示仅基线 LP）\\
\fld{cost\_tangent\_points} & — & — & $\ge 2$ & 9 & 二次成本支撑切线条数\\
\fld{cone\_tolerance} & — & — & $>0$ & $10^{-6}$ & 锥违例判据\\
\fld{solver\_tolerance} & — & — & $>0$ & $10^{-8}$ & LP 求解容差（保留字段；当前实现未透传给 HiGHS 适配器）\\
\fld{verbose} & — & — & true/false & false & 日志开关（当前实现未使用）\\
\end{paramtable}

\begin{paramtable}
\fld{solved} & — & — & true/false & false & LP 循环正常完成\\
\fld{outer\_relaxation\_valid} & — & — & true/false & false & 外松弛有效性标记（求解完成即置真）\\
\fld{dual\_certificate\_available} & — & — & true/false & false & 对偶下界证书是否可用\\
\fld{soc\_outer\_approximation\_converged} & — & — & true/false & false & 末轮无违例锥割\\
\fld{ac\_passivity\_cut\_applied} / \fld{dc\_passivity\_cut\_applied} & — & — & true/false & false & 聚合网损割是否施加\\
\fld{status} / \fld{solver} & — & — & — & 空 & 状态串 / 求解器名\\
\fld{objective\_lower\_bound} & $\underline{f}$ & — & — & 0 & 原非凸 OPF 目标的有效下界（证书）\\
\fld{lp\_primal\_objective} & — & — & — & 0 & 末轮 LP 原始目标\\
\fld{lp\_primal\_dual\_gap} & — & — & — & 0 & 末轮 LP 原始目标与下界之差\\
\fld{runtime\_ms} & — & ms & — & 0 & 耗时\\
\fld{max\_soc\_violation} & — & — & — & 0 & 末轮最大锥违例\\
\fld{max\_ac\_lift\_violation} / \fld{max\_ac\_psd\_violation} / \fld{max\_dc\_lift\_violation} & — & — & — & 0 & 交流提升/PSD 团/直流提升分项违例\\
\fld{max\_converter\_apparent\_power\_violation} / \fld{max\_converter\_current\_violation} & — & — & — & 0 & 换流器视在功率/相电流锥分项违例\\
\fld{primal\_residual} / \fld{dual\_residual} & — & — & — & 0 & HiGHS 报告的原始/对偶可行度\\
\fld{rounds} / \fld{cuts\_added} & — & — & — & 0 & 实际轮数 / 累计割数\\
\fld{variables} / \fld{equalities} / \fld{inequalities} & — & — & — & 0 & 末轮 LP 规模\\
\fld{primal} & $x$ & — & — & 空 & 末轮 LP 解（提升空间，不保证 AC 可实施）\\
\fld{generator\_active\_power\_pu} / \fld{generator\_reactive\_power\_pu} & $p_g/q_g$ & pu & — & 空 & LP 解中的机组分量\\
\fld{round\_lower\_bounds} & — & — & — & 空 & 各轮下界序列（单调不减）\\
\fld{model\_limitations} & — & — & — & 空 & 省略/近似条目（诚实口径）\\
\end{paramtable}

\subsection{验证与校核}\label{subsec:tp-verify}

模块测试目标为 \fld{test\_three\_phase\_hybrid\_opf}（注册于
\srcpath{tests/CMakeLists.txt:test\_three\_phase\_hybrid\_opf}），全部算例共用一个三母线九节点合成混合系统
（\srcpath{tests/test\_three\_phase\_hybrid\_opf.cpp:make\_small\_hybrid\_case}，
\srcpath{make\_small\_hybrid\_case}）：三母线全相齐全，支路阻抗角一致的串联
导纳 $y=8-\mathrm{j}16$ 与 $y=6-\mathrm{j}12$（:make\_small\_hybrid\_case）；末母线分相负荷
$p=0.08+0.01\varphi$、$q=0.025+0.005\varphi$（:make\_small\_hybrid\_case）；电压窗口
$[0.85,1.10]$~pu（:make\_small\_hybrid\_case）；首母线为平衡参考 $1\angle\{0^\circ,-120^\circ,
+120^\circ\}$（:make\_small\_hybrid\_case）；$\varepsilon_{\mathrm{vuf}}=0.05$（:make\_small\_hybrid\_case）；首母线三相
各挂一台 $c_2=0.1$、$c_1=30$ 机组（:make\_small\_hybrid\_case）；直流两节点、电导 50、节点 1 负荷
0.10、节点 0 为直流参考（:make\_small\_hybrid\_case）；一台 VSC 接末母线三相与直流节点 0，
$\eta=0.98$、$s_{\max}=0.25$（:make\_small\_hybrid\_case）。除注明外各用例后端为 Ipopt、
容差 $10^{-7}$。

\begin{itemize}
  \item \textbf{Full 与图降阶一致性}（:TEST\_CASE("Monolithic phase hybrid OPF Full and GR recover the same solution")）：降阶恰消去 3 个相节点
        （:TEST\_CASE），目标相对偏差不超过 $10^{-6}$（:TEST\_CASE），全电压最大偏差
        小于 $10^{-6}$（:TEST\_CASE），电压越限小于 $10^{-8}$（:TEST\_CASE），
        VUF 不超过 $\varepsilon_{\mathrm{vuf}}+10^{-7}$（:TEST\_CASE），
        换流器电流不平衡度两模型一致到 $10^{-6}$（:TEST\_CASE）。
  \item \textbf{序列求解保留时变负荷节点}（:TEST\_CASE）：第二时段在降阶模型的
        被保留节点上改变负荷，序列解与独立求解目标相对偏差不超过 $10^{-7}$、
        电压偏差小于 $10^{-6}$（:TEST\_CASE）。
  \item \textbf{GFL/GFM 动态平衡点}（:TEST\_CASE）：跟网型正序电流模式下换流器
        电流负序比小于 $10^{-7}$、负载率不超过 $1+10^{-7}$、动态平衡残差小于
        $10^{-7}$（:TEST\_CASE），且解析雅可比/黑塞差分误差小于 $10^{-5}$
        （:TEST\_CASE）；把平衡点交给暂态 \fld{GridFollowingInverter} 后导数最大
        值小于 $10^{-7}$（:TEST\_CASE）。构网型下垂模式同样通过负载率/平衡残差/
        导数校核，内电势幅值大于 0.5~pu，\fld{GridFormingInverter} 导数小于
        $10^{-7}$（:TEST\_CASE）。
  \item \textbf{独立耦合潮流复核}（:TEST\_CASE）：
        \srcpath{make\_three\_phase\_hybrid\_pf\_case} 生成的潮流算例以
        $10^{-9}$ 容差收敛，潮流电压与 OPF 运行点交/直流电压偏差均小于
        $2\times10^{-6}$（:SECTION）。
  \item \textbf{NativeIPM 原-对偶搬运}（:SECTION）：先把图降阶解与对偶按
        保留节点映射搬回 Full 模型（含 NaN 标记的待恢复等式对偶行，
        :TEST\_CASE("Native Full certifies a graph-reduced primal-dual transport")），NativeIPM 从该点出发的初始原始/对偶残差不超过 $10^{-6}$，
        最终 KKT 三项残差均不超过 $10^{-6}$，目标与降阶解相对偏差不超过
        $10^{-6}$（:SECTION）。
  \item \textbf{约束预言机}（:TEST\_CASE）：强制行数严格少于全量且被省略行最大
        违反量低于 $-m_{\mathrm{act}}$（:TEST\_CASE），目标与全行求解相对偏差
        不超过 $10^{-7}$、电压偏差小于 $10^{-6}$、KKT 残差不超过 $10^{-6}$
        （:TEST\_CASE）；以预言机行集与解为种子重解时零新增行且行集不变
        （:TEST\_CASE）。
  \item \textbf{松弛下界证书}（:TEST\_CASE）：零轮基线 LP 的下界不超过非凸解
        目标 $+10^{-7}$（:TEST\_CASE）；10 轮外逼近下界不低于基线、不超过非凸
        解目标 $+10^{-7}$、各轮下界单调（容差 $10^{-9}$）、末轮锥违例不大于
        基线（:TEST\_CASE）。
  \item \textbf{固定电流拒绝}（:TEST\_CASE）：注入非零 \fld{i\_ac\_fixed} 时松弛
        层返回未求解、状态串含 \fld{fixed-current} 且登记
        \fld{model\_limitations}（:TEST\_CASE）。
\end{itemize}

另有两条模块级验证通道：其一，Sioux Falls 交通网 + IEEE 123 节点馈线联合算例
校验 adapter 装配规模与 ID 归因（$\ge 120$ 座三相母线、$\ge 250$ 个相节点、
4 条直流母线、4 台 VSC、$\ge 15$ 台分相机组，\fld{vuf\_max} 选项透传；
\srcpath{tests/test\_sioux\_falls\_ieee123\_case.cpp:TEST\_CASE("Sioux Falls--IEEE 123 case preserves public networks and assumptions")}、:TEST\_CASE("Sioux Falls--IEEE 123 stress controls preserve charging injection")）；其二，
论文复现实验工具 \fld{phase\_hybrid\_opf\_benchmark} 与
\fld{sioux\_falls\_ieee123\_experiment}（后者由 \srcpath{HACDCPF\_HAVE\_OPENDSS}
门控，\srcpath{CMakeLists.txt:phase\_hybrid\_opf\_benchmark、sioux\_falls\_ieee123\_experiment}）提供更大规模的基准与审计，
但不属于默认 ctest 套件。截至当前代码，模块尚无 OpenDSS 等外部工具的 OPF 级
对照算例（潮流层对照不属于本章范围）。

\subsection{边界与局限}\label{subsec:tp-limit}

\begin{itemize}
  \item \textbf{活跃研发}：模块接口与覆盖范围仍在演进；默认后端 Ipopt 受嵌入式
        编译门控，未编译时仅返回 \fld{Unavailable:} 状态（
        \srcpath{MIPSolvers/src/engine/solver/external/adapters.cpp:IpoptAdapter::available}）。
  \item \textbf{网络覆盖}：adapter 不支持 DC/DC 变换器、能量路由器与可控 DC 静态
        发电机（遇之抛错，src/optimal\_power\_flow/three\_phase\_hybrid\_adapter.cpp:build\_three\_phase\_hybrid\_model、391--395）；DC 支路仅电阻型；
        无线路潮流约束、无 OLTC/分接头与并联补偿调节变量；换流器损耗仅单效率
        近似。
  \item \textbf{负荷模型}：仅恒功率节点方程；delta 与电压相关负荷只能以导入
        运行点的恒功率等值进入，已在 \fld{model\_limitations} 声明
        （src/optimal\_power\_flow/three\_phase\_hybrid\_adapter.cpp:build\_three\_phase\_hybrid\_model、474--476）。
  \item \textbf{换流器控制}：序感知模式要求三相端子齐全；构网“下垂”在 OPF
        中仅有诺顿耦合等式，\fld{p\_droop\_pu}/\fld{q\_droop\_pu}/%
        \fld{voltage\_reference\_pu}/\fld{nominal\_frequency\_hz} 不进入任何
        约束（仅存档与平衡点恢复用，src/optimal\_power\_flow/three\_phase\_hybrid\_opf.cpp:build\_model\_data）；松弛层完全省略序感知
        控制等式（src/optimal\_power\_flow/three\_phase\_hybrid\_relaxation.cpp:solve\_three\_phase\_hybrid\_opf\_relaxation）。
  \item \textbf{外部电网}：以 $\pm5\times$ 总负荷的虚拟发电机近似（
        src/optimal\_power\_flow/three\_phase\_hybrid\_adapter.cpp:build\_three\_phase\_hybrid\_model），非无限大母线的阻抗模型。
  \item \textbf{目标函数}：三项正则化权重硬编码（src/optimal\_power\_flow/three\_phase\_hybrid\_opf.cpp:objective）；松弛层
        目标省略正则项且以切线低估二次成本，故其目标值是下界而非近似原目标。
  \item \textbf{初始化}：存在非参考节点机组时交流潮流初始化静默跳过
        （src/optimal\_power\_flow/three\_phase\_hybrid\_opf.cpp:initialize\_ac\_power\_flow、1602--1606），初值质量依赖 \fld{voltage\_start}。
  \item \textbf{松弛层}：拒绝非零固定电流注入与负二次成本
        （src/optimal\_power\_flow/three\_phase\_hybrid\_relaxation.cpp:solve\_three\_phase\_hybrid\_opf\_relaxation）；外逼近轮数耗尽时
        \fld{soc\_outer\_approximation\_converged} 为假，下界仍有效但松弛间隙
        未闭合；LP 解不保证交流可实施。
  \item \textbf{验证覆盖}：完整 OPF/KKT 功能测试仍以三母线合成算例为主；
        另有 360 个三相母线的 Phase I 稀疏规模/预算回归。IEEE 123 规模仅
        adapter 装配校验与实验工具；尚无外部工具 OPF 对照或大规模完整
        Phase II 收敛回归。
\end{itemize}

% 验证与校核、已知局限
% 本章文件由 \input 引入，不含导言区。
\section{验证与校核、已知局限}\label{sec:verification}

本章汇总 OPF 模块的全部验证证据与已知局限。前面各章（第~\ref{sec:acopfnative}~--~\ref{sec:threephase}~节）
给出各求解器的模型与算法，本章回答两个问题：\textbf{凭什么相信这些求解器的结果}，
以及\textbf{在哪些规模、工况与元件组合下结果尚未被充分验证}。写作口径与
《元件模型技术手册》一致：测试断言一律给出\texttt{文件:行号}，文档记录的实测数值
标注“文档实测”，被门控、被跳过、被降级为警告的条目如实列出，不以路线图代替现状。

\subsection{验证体系与证据分级}\label{subsec:ver:layers}

OPF 验证按四个层次组织（图~\ref{fig:ver:layers}）：

\begin{enumerate}
  \item \textbf{导数级}：Parity 建模层的解析等式 Jacobian、LCC 三变量一阶导数与局部
        Hessian、相域模型导数，全部与中心有限差分比对
        （\srcpath{tests/test\_opf\_solver\_backends.cpp}:274--308、310--422、424--598；
        \srcpath{tests/test\_three\_phase\_hybrid\_opf.cpp}:174--282）；
  \item \textbf{优化级}：收敛后检查 KKT 残差（\fld{max\_constraint\_violation}、
        \fld{max\_stationarity}）、目标值符号与量级、变量边界、实际生效的求解路径
        （\fld{solver\_path}）与线性代数后端（\fld{linear\_solver\_backend}）；
  \item \textbf{系统级}：把 OPF 调度写回系统副本运行独立 Newton 潮流回放（PF replay）、
        \srcpath{verify\_opf\_result} 物理审计、图归约/投影前后结果 round-trip；
  \item \textbf{外部基准}：手算闭式解、MATPOWER 标准算例、DSP（BPA）LCC 参考解、
        文档记录的参考目标值。
\end{enumerate}

\begin{figure}[H]
\centering
\begin{tikzpicture}[
  font=\small,
  box/.style={draw, rectangle, minimum width=10.6cm, minimum height=0.9cm,
              align=center, inner sep=2pt},
  lbl/.style={font=\footnotesize, align=left, text width=4.6cm},
  arr/.style={-{Stealth[length=2mm]}, thick}]
  \node[box] (l1) at (0,0) {导数级：解析 Jacobian/Hessian vs 有限差分};
  \node[box] (l2) at (0,1.25) {优化级：KKT 残差、目标值、边界、求解路径};
  \node[box] (l3) at (0,2.5) {系统级：PF 回放、结果审计、投影 round-trip};
  \node[box] (l4) at (0,3.75) {外部基准：手算解、MATPOWER、DSP/BPA、参考目标值};
  \draw[arr] (l1.north) -- (l2.south);
  \draw[arr] (l2.north) -- (l3.south);
  \draw[arr] (l3.north) -- (l4.south);
  \node[lbl, anchor=west] at (5.6,0) {微算例即可，不依赖外部工具};
  \node[lbl, anchor=west] at (5.6,1.25) {内置于默认 ctest 套件};
  \node[lbl, anchor=west] at (5.6,2.5) {跨模块一致性};
  \node[lbl, anchor=west] at (5.6,3.75) {部分为文档实测，非逐轮断言};
\end{tikzpicture}
\caption{OPF 验证的四层结构。下层失败会立即定位到导数或模型装配错误；
上层失败则指向求解器全局化或外部语义差异。}
\label{fig:ver:layers}
\end{figure}

证据按可信度分三档，本章表格与行文中严格区分：
\begin{itemize}
  \item \textbf{默认 ctest 断言}：注册于 \srcpath{tests/CMakeLists.txt}:31--34、109、115、
        185、189、282、288 的测试目标，每次构建均可复跑；
  \item \textbf{隐藏性能用例}：带 Catch2 隐藏标签 \fld{[.performance]} 的用例
        （如 \srcpath{tests/test\_opf\_solver\_backends.cpp}:914--943、945--1024，
        \srcpath{tests/test\_reactive\_power\_opt.cpp}:84--110），默认测试运行不执行，
        其结果不构成逐轮回归判据；
  \item \textbf{文档实测}：\srcpath{docs/large\_hybrid\_ipm\_diagnostics.md}、
        \srcpath{docs/reactive\_power\_optimization\_validation.md}、
        \srcpath{docs/lcc\_dat\_opf\_report/} 中记录的墙钟时间、迭代数与残差，
        是特定机器、特定构建下的一次性记录，引用时标“文档实测”。
\end{itemize}

单位制同全书：容差中 pu 为标幺值（$S_{\mathrm{base}}=100$~MVA），功率容差为
MW/MVAr 有名值，角度容差为度或 rad（随测试原文注明），成本币种未定。

\subsection{验证矩阵总表}\label{subsec:ver:matrix}

表~\ref{tab:ver:matrix} 汇总全部验证通道。各通道细节见后续小节；表中“容差/通过准则”
列均为测试断言或文档声明的原文数值。

\begin{footnotesize}
\begin{longtable}{@{}L{2.0cm}L{3.0cm}L{2.9cm}L{3.1cm}L{3.9cm}@{}}
\caption{OPF 验证矩阵总表}\label{tab:ver:matrix}\\
\toprule
\textbf{验证通道} & \textbf{算例（规模）} & \textbf{对照基准} & \textbf{容差/通过准则} & \textbf{出处（文件:行号）}\\
\midrule
\endfirsthead
\multicolumn{5}{@{}l}{\footnotesize（续表）}\\
\toprule
\textbf{验证通道} & \textbf{算例（规模）} & \textbf{对照基准} & \textbf{容差/通过准则} & \textbf{出处（文件:行号）}\\
\midrule
\endhead
\midrule
\multicolumn{5}{r@{}}{\footnotesize（接下页）}\\
\endfoot
\bottomrule
\endlastfoot
手算解析解 &
\fld{opf\_hand\_01}--\fld{04}（1--2 母线 JSON 微算例） &
README 闭式手算（merit order、支路拥塞、负荷削减、线性化成本） &
LP 精确；QP 参考目标 815.333 &
\srcpath{external\_data/opf\_hand\_cases/README.md}:21--296（未接入 ctest，GUI 手检）\\
MATPOWER 小算例 &
case5/9/14/30/39/57（$\le 57$ 节点） &
NativeAC 与 NativeQP 各自收敛 + Newton PF 回放 &
可行性 $10^{-6}$，PF 残差 $10^{-8}$ &
\srcpath{tests/test\_acopf\_dcopf\_crossval.cpp}:52--54、117--144、250--303；
\srcpath{tests/test\_matpower\_cases.cpp}:618--638\\
MATPOWER 大算例 &
case2383wp、case2869pegase（2869 节点）、case6515rte（6515 节点） &
Parity IPM 收敛/有限性；文档参考目标值 133999/386107 &
case2869：收敛且迭代 $<100$、稀疏后端；case6515rte：仅断言有限 &
\srcpath{tests/test\_opf\_solver\_backends.cpp}:1058--1071、1154--1196；
\srcpath{docs/large\_hybrid\_ipm\_diagnostics.md}:263--269（文档实测）\\
AC/DC/PF 三路交叉 &
MATPOWER 小算例；IEEE 14 节点 &
AC OPF、DC OPF、Newton PF 三者一致性 &
PF 回放收敛；DC 目标 $\le$ AC 目标 $\times 1.05+1$ &
\srcpath{tests/test\_acopf\_dcopf\_crossval.cpp}:250--303；
\srcpath{tests/test\_crossmodule\_integration.cpp}:206--246\\
Native/\allowbreak Parity/\allowbreak Ipopt 后端交叉 &
case9、case30、case300\_acdc、case2000\_acdc &
三后端同案收敛性；Ipopt 有界参数矩阵 &
Parity 目标 $\le$ ED 目标 $\times 1.001+1$；Ipopt 预算内结果有限 &
\srcpath{tests/test\_opf\_solver\_backends.cpp}:1075--1150；
\srcpath{docs/large\_hybrid\_ipm\_diagnostics.md}:118--131、195--220（文档实测）\\
夜间工况 OPF &
南沙全网 JSON、画布 \fld{power\_system.json}（无常规机组） &
外部电网按成本提升为 OPF 变量后收敛 &
\fld{allow\_fallback=false} 下收敛；能量平衡 $10^{-3}$~MW &
\srcpath{tests/test\_nighttime\_opf.cpp}:46--96、164--209\\
时序 UC$\to$OPF$\to$PF 流水线 &
画布系统 24 步/4 步（含 UC） &
逐步 OPF 与 PF 均收敛；UC 可行 &
24/24、4/4 步收敛 &
\srcpath{tests/test\_nighttime\_opf.cpp}:125--162；
\srcpath{docs/time\_series\_power\_flow\_models.md}:276--349\\
三相混合 OPF &
手组 3 母线 9 相节点 + 2 直流节点 &
Full 与 GraphReduced 一致；独立耦合 PF 复现；SOCP 松弛下界 &
目标相对差 $10^{-6}$；电压 $10^{-6}$；VUF $\le 0.05$ &
\srcpath{tests/test\_three\_phase\_hybrid\_opf.cpp}:24--143、284--317、463--510\\
RPO 专项 &
case9；case2869 + 20 OLTC + 10 并联（authored）；case300\_acdc &
独立 OPF 复算 + PF 回放 + 物理网损账 &
目标相对差 $\le\max(10^{-3},10\epsilon_d)$；$|\Delta V|\le\max(5\times10^{-4},10\epsilon_p)$ &
\srcpath{tests/test\_reactive\_power\_opt.cpp}:16--110；
\srcpath{docs/reactive\_power\_optimization\_validation.md}:103--123、159--177（文档实测）\\
大规模混合 IPM 诊断 &
case300\_acdc、case2000\_acdc（2000 节点 + 8 VSC）、2869/9241/13659 &
UMFPACK/\allowbreak KLU/\allowbreak MUMPS/\allowbreak Eigen 多后端同一收敛；文档实测残差 &
测试：违反 $<10^{-5}$、迭代上限；文档：$3.49\times10^{-9}$/$9.03\times10^{-7}$ &
\srcpath{tests/test\_opf\_solver\_backends.cpp}:86--112、880--943、1132--1196；
\srcpath{docs/large\_hybrid\_ipm\_diagnostics.md}:77--131、224--289（文档实测）\\
BPA/DSP LCC 算例 &
\srcpath{2DC.dat}、\srcpath{cigre.dat}（6 交流母线、2 直流母线、2 台 LCC） &
DSP SOL 外部参考；解析 Jacobian/Hessian 有限差分；固定 tap PF 回放 &
$1500/1410$~MW $\pm2$；线损 90~MW $\pm2$；Jacobian $2\times10^{-5}$；回放 $0.5$~MW &
\srcpath{tests/test\_opf\_solver\_backends.cpp}:310--422、424--598、600--878；
\srcpath{docs/lcc\_dat\_opf\_report/LCC\_DAT\_OPF\_technical\_report.tex}:1039--1124\\
结果审计、回放与投影 &
case9、case14、IEEE14 AC/DC、4 母线串联归约 &
\srcpath{verify\_opf\_result} 审计；归约/投影前后目标一致 &
电压违反 $<0.02$~pu；归约目标差 $10^{-3}$ &
\srcpath{tests/test\_solver\_capabilities.cpp}:194--239；
\srcpath{tests/test\_graph\_roundtrip.cpp}:223--258、671--705；
\srcpath{tests/test\_multiscale\_comprehensive.cpp}:170--229\\
\end{longtable}
\end{footnotesize}

\subsection{手算解析解基准}\label{subsec:ver:hand}

\srcpath{external\_data/opf\_hand\_cases/} 提供四个刻意做到极小的 JSON 算例，每个都在
README 中给出完整闭式推导，用于对照 GUI/API 的 DC OPF 路径（直流模型无损，手算结果
精确；README 同时声明 AC/parity 路径只作冒烟，因为交流可行性细节不在手算范围内，
\srcpath{external\_data/opf\_hand\_cases/README.md}:1--19）。

\subsubsection*{算例与校验量}
\begin{itemize}
  \item \textbf{01 单母线 merit order}：60~MW 负荷，机组 1（0--40~MW，成本 10）与
        机组 2（0--100~MW，成本 20）。手算调度 $P_1=40$、$P_2=20$~MW，目标 800，
        单节点 LMP 为边际机组成本 20
        （\srcpath{external\_data/opf\_hand\_cases/README.md}:21--64）；
  \item \textbf{02 两母线支路拥塞}：100~MW 负荷，线路 $x=0.1$~pu、限值 40~MVA。
        手算 $P_1=40$、$P_2=60$~MW，目标 3400；DC 角度按
        \begin{align*}
        f_{12}=\frac{\theta_1-\theta_2}{x_{12}}\,S_{\mathrm{base}}
        \end{align*}
        得 $\theta_2=-0.04$~rad$=-2.2918^\circ$，两侧 LMP 分别为 10 与 50
        （同文件:65--124）；
  \item \textbf{03 负荷削减}：80~MW 负荷、50~MW 机组上限。默认 \fld{voll=0} 时求解器
        取有效 VOLL $=\max(100,\,1.1\,c_1)=100$，手算削减 30~MW、目标 4000
        （同文件:126--174）；
  \item \textbf{04 二次成本的 GUI DC 口径}：README 明确当前 GUI DC 端点对二次成本按
        保守边际系数的 LP 线性化处理，
        \begin{align*}
        c_i^{LP}=c_{1i}+2c_{2i}P_i^{\max},
        \end{align*}
        手算 LP 调度 $(80,0)$~MW、目标 1120；真 QP 参考解 $(36.667,43.333)$~MW、
        目标 815.333 仅供支持真 QP 的后端对照（同文件:176--296）。
\end{itemize}

\subsubsection*{已知偏差}
该目录在当前检出中\textbf{未被任何 C++ 测试或 ctest 目标引用}（全仓检索仅
\srcpath{ngnlp/chapters/07\_validation\_and\_usage.tex}:57--60 提及），定位为
GUI/API 人工校核材料；算例 02 的 README 明确“支路对偶当前不被 DC OPF 路径认证，
应校验原始调度与潮流限值而非 \fld{branch\_mu\_*}”
（\srcpath{external\_data/opf\_hand\_cases/README.md}:124），与
\ref{subsec:ver:limitations} 节的支路对偶条目一致。

\subsection{MATPOWER 算例通道}\label{subsec:ver:matpower}

\subsubsection*{小算例（$\le 57$ 节点）}
\srcpath{tests/test\_acopf\_dcopf\_crossval.cpp}:52--54 固定 case5/9/14/30/39/57
六个算例，AC OPF 走 NativeAC（\fld{use\_parity\_ipm=false}，同文件:117--130），
DC OPF 走自研 NativeQP LCQP（同文件:132--144），两者均以可行性容差 $10^{-6}$ 调用；
随后把 AC OPF 的 $P_g/Q_g/V_g$ 写回系统副本跑 Newton 潮流（容差 $10^{-8}$，
同文件:151--183）。\srcpath{tests/test\_matpower\_cases.cpp}:618--638 另有两项
DC OPF 断言：case9 目标有限且 \fld{pg\_mw} 尺寸正确，case14 所有在役机组调度落在
$[P_{\min}-10^{-3},\,P_{\max}+10^{-3}]$~MW 内。

\subsubsection*{大算例（Parity IPM）}
\srcpath{tests/test\_opf\_solver\_backends.cpp}:1058--1071 是 case2383wp（波兰冬季
峰荷）的回归守卫：注释记录该算例在引入基于梯度的目标缩放前发散（迭代顶格时最大
平稳性约 0.45），现断言收敛。同文件:1154--1176 断言 case2869pegase 收敛、迭代
$<100$ 且 KKT 维数超 1500 时实际使用稀疏分解后端（\fld{linear\_solver\_backend}
含 \fld{sparse}）。同文件:1178--1196 对 case6515rte 只断言目标有限且
$<10^{30}$——注释明确“该难例在迭代预算内不期望收敛”，这是对此前 Eigen SparseLU
下发散到 NaN（目标约 $-3\times10^{48}$）的稳健性守卫，不是收敛性声明。
文档另记录 case9241/13659 的参考目标 315837/386107 与 KLU 后端实测
（\srcpath{docs/large\_hybrid\_ipm\_diagnostics.md}:263--269，文档实测）。
诊断目标 \srcpath{opf\_numerical\_benchmark} 另对 25,000-bus ACTIVSg25k 执行
一次预热加三次计时的 Release/KLU 协议：普通 AC-PF 启动中位数为
$183.808\,\mathrm{s}$，66 次迭代、65 次分解，primal/dual 为
$3.239\times10^{-7}/1.004\times10^{-7}$；默认 $2\,\mathrm{s}$ DC Phase I
未形成可接纳候选，精确回退后中位数为 $186.588\,\mathrm{s}$，慢 $1.51\%$。
这是可求解性和预算失败语义的性能验证，不是 25k 加速声明；该长时诊断不注册到
常规 CTest。
同一目标的实验性零分解 component dispatch-dual predictor 仅把 normalized dual
$468.10733\to468.08762$（$0.0042\%$），未通过固定 $5\%$ exact residual 审计，
以 $1.978\,\mathrm{ms}$ 成本回退；单次总时长 $184.344\,\mathrm{s}$ 且 Phase II
仍为 66 次迭代、65 次分解。因此它默认关闭，结果支持优先研究完整 continuation
state 与跨重复求解 symbolic pattern 复用，而非 partial dual 或更精确的辅助 DC-IPM。
同一 ACTIVSg25k 的 prepared-session 两轮 Release/KLU 验证进一步给出：首轮
$184.645\,\mathrm{s}$/66 次迭代/65 次分解，第二轮
$8.666\,\mathrm{s}$/4/3，且 0 次 symbolic analyze；最终 primal/dual 为
$3.389\times10^{-8}/2.092\times10^{-8}$，峰值 RSS $4709.6\,\mathrm{MiB}$。
这是完整状态与 symbolic reuse 的重复求解收益；第二轮 3 次分解均为新 numeric
factorization，不是旧数值因子的无审计复用。

\subsubsection*{已知偏差}
综合对比用例（\srcpath{tests/test\_acopf\_dcopf\_crossval.cpp}:554--694）跳过全部
$>1000$ 节点的算例（同文件:629--633）与 scenario/contab 文件（同文件:569--574），
且只统计收敛计数（\fld{acopf\_pass>0}），不对目标值与参考解逐一相等断言；
case57 的 PF(DCOPF) 回放收敛被降级为 \texttt{WARN} 而非断言（同文件:273--282，
理由：DC OPF 不提供电压与无功，难例上不保证交流可行）。

\subsection{跨求解器交叉验证}\label{subsec:ver:crossval}

\subsubsection*{AC OPF vs DC OPF vs Newton PF}
\srcpath{tests/test\_acopf\_dcopf\_crossval.cpp}:250--303 对六个小算例逐一断言：
AC OPF 收敛则 PF(AC OPF 调度） 必收敛（同文件:268--271）；DC OPF 目标必须为正
（同文件:287--291）。跨模块流水线测试进一步断言松弛性质
\begin{align*}
C_{\mathrm{DCOPF}} \le 1.05\,C_{\mathrm{ACOPF}} + 1
\end{align*}
（\srcpath{tests/test\_crossmodule\_integration.cpp}:206--209，判据 P1-A4）、
OPF 调度下 Newton PF 必收敛（P1-A6，同文件:232）以及逐母线电压一致
$|V_m^{PF}-V_m^{OPF}|<0.05$~pu（P1-A7，同文件:234--246）。

\subsubsection*{Native / Parity / Ipopt 三路}
\srcpath{tests/test\_opf\_solver\_backends.cpp}:1075--1112 在 case30 上分别断言
EconomicDispatch（merit-order + PF 捷径）、ParityIPM（全空间非线性 OPF）两后端收敛；
同文件:1114--1130 断言真 OPF 目标不超过经济调度捷径的 $1.001$ 倍加 1。Ipopt 后端的
断言是\textbf{优雅降级}：注释声明嵌入式 Ipopt/MUMPS 默认被门控关闭，选择 Ipopt 时必须
回退到 Parity IPM 且 \fld{solver\_path} 报告为 \fld{ParityIPM}（同文件:1100--1108；
编译期门控宏 \srcpath{HACDCPF\_HAVE\_IPOPT}，\srcpath{src/optimal\_power\_flow/ac\_opf.cpp}:33；
CMake 选项 \srcpath{HACDCPF\_ENABLE\_IPOPT} 默认 OFF、macOS 默认 ON，见
第~\ref{sec:overview}~节）。case300\_acdc 上另断言 Ipopt 迭代/容差选项透传生效：
1 次迭代预算必须报 \texttt{Max iterations exceeded} 且不收敛，800 次预算收敛且
违反 $<10^{-6}$（同文件:114--149）。

\subsubsection*{已知偏差}
\begin{itemize}
  \item AC 与 DC 目标值\textbf{不做直接相等断言}：测试注释说明两者成本口径不同，
        直接比较无意义（\srcpath{tests/test\_acopf\_dcopf\_crossval.cpp}:284--291）；
  \item AC OPF 结果不暴露支路潮流，$P_f$ 对比代码留空（同文件:222--227）；
  \item “DC OPF 调度质量”判据宽松：相对调度差只要求 $<1.0$（即 100\%，
        同文件:536--547；注意 546 行尾注释“Less than 50\%”与实际断言 \fld{rel\_diff<1.0}
        不一致，以断言为准）；
  \item 默认构建下 Ipopt 选择静默回退 Parity IPM，因此默认 ctest 中并不存在真正的
        三路目标值相等断言；Ipopt 侧的数值证据来自隐藏性能用例与文档实测
        （\srcpath{docs/large\_hybrid\_ipm\_diagnostics.md}:195--220）；
  \item 独立 \srcpath{power\_models} 模块的 OPF 交叉验证目前主要体现在
        三相混合 OPF 的 NativeIPM/Ipopt 适配器路径（见 \ref{subsec:ver:threephase} 节），
        AC OPF 的 AML builder 无独立对照测试；其源码等价模型见
        \srcpath{docs/modules/power\_models/power\_models\_manual.tex}。
\end{itemize}

\subsection{夜间工况 OPF 与外部电网提升}\label{subsec:ver:nighttime}

\subsubsection*{回归起源}
\srcpath{tests/test\_nighttime\_opf.cpp}:1--5 的文件头注释记录了本通道的回归起源：
修复前 \srcpath{solve\_ac\_opf} 在 \fld{sys.ac.generators} 为空（$n_g=0$）时立即返回
失败；修复后带正成本的外部电网被提升为 OPF 发电机变量，夜间（无光伏、无常规机组）
工况得以收敛。

\subsubsection*{算例与判据}
\begin{itemize}
  \item \textbf{画布 \fld{power\_system.json}}（无常规机组、1 外部电网、1 双绕组
        变压器、含 VSC/DC-DC/储能/直流光伏）：Parity IPM、禁 fallback 下断言收敛；
        储能功率与 authored 设定一致（$2\times10^{-4}$~MW），VSC 无功跟随设定
        （$5\times10^{-2}$~MVAr），外部电网无功 $0<Q<0.20$~MVAr，DC-DC 效率与
        等值电阻的直流平衡按公式逐项核对（$2\times10^{-4}$~MW），交流供给与负荷
        平衡到 $10^{-3}$~MW（\srcpath{tests/test\_nighttime\_opf.cpp}:46--96）；
  \item \textbf{南沙全网}（\fld{nansha\_full\_network.json}，无常规机组、外部电网带
        成本）：夜间用例先断言前提（机组为空、至少一个外部电网成本为正），再断言
        \fld{allow\_fallback=false} 下收敛；昼间对照用例同样收敛（同文件:164--209）；
  \item \textbf{同母线独立调度}：50~MW 负荷母线上同时挂 100 成本机组与 10 成本外部
        电网，断言机组 $P_g\approx0$、外部电网承担全部 50~MW（$10^{-2}$~MW），
        AC 与 DC OPF 各一项（同文件:211--255、312--350）；
  \item \textbf{投影正确性}：死岛剥离后发电机结果保持 authored 顺序（同文件:257--310）；
        闭合开关合并母线后外部电网/机组/负荷全部保留并正确调度（同文件:352--409）。
\end{itemize}

\subsubsection*{已知偏差}
本通道覆盖 2 个工程 JSON 算例与 3 个手工微系统，均为“收敛 + 少量物理量”断言，
没有外部工具对照；外部电网的无功调节能力、成本曲线形态在大系统上的行为未单独验证。

\subsection{时序 UC$\to$OPF$\to$PF 流水线校核}\label{subsec:ver:tspf}

时序生产模拟的三阶段流水线（UC MILP $\to$ 逐步 AC OPF $\to$ 逐步混合潮流回放）
模型细节属第~\ref{sec:overview}~节范围，此处只列与 OPF 可信性直接相关的校核机制
（\srcpath{docs/time\_series\_power\_flow\_models.md}:276--349）：

\begin{itemize}
  \item \textbf{跟踪带}：对每台在线机组，OPF 功率盒围绕 UC 基点收紧，
        \begin{align*}
        \beta_g=\max\bigl(\beta^{abs},\ \beta^{rel}\max(|p^{UC}_g|,1)\bigr),
        \end{align*}
        松弛机 $(\beta^{rel},\beta^{abs})=(20\%,5\,\mathrm{MW})$，非松弛机
        $(2\%,1\,\mathrm{MW})$；其工程语义是 OPF 只作类 AGC 再调度，吸收无损 UC
        看不见的网损与直流近似误差（\srcpath{docs/time\_series\_power\_flow\_models.md}:284--300）；
  \item \textbf{降级路径}：某步 OPF 不收敛时，该步退回对原始 UC 快照直接跑 PF，
        流水线不中断（同文件:307--308）；
  \item \textbf{交叉指标}：每步记录 $\max_b|V_{m,b}^{OPF}-V_{m,b}^{PF}|$ 与经首母线
        参考对齐的相角差、双侧物理网损，作为流水线“信任证书”（同文件:323--339）；
  \item \textbf{储能一致性}：OPF 故意不重调度储能，避免功率与 SOC 轨迹互相矛盾
        （同文件:302--305）；
  \item \textbf{线程安全}：按日并行时逐步 OPF 被强制到线程安全的 Parity IPM，
        不可重入的 Ipopt/MUMPS 路径不会并发进入（同文件:412--414）。
\end{itemize}

测试侧：\srcpath{tests/test\_nighttime\_opf.cpp}:125--162 断言画布系统 24 步动态
OPF 全部 24 步 OPF 与 PF 收敛，以及含 UC 的 4 步默认日生产路径 UC 可行且
4/4 步收敛。

\subsection{三相混合 OPF 测试}\label{subsec:ver:threephase}

模型与算法见第~\ref{sec:threephase}~节；本小节只列验证证据。全部用例共用同一个
手组算例：3 个三相母线（9 个相节点）、2 个直流节点、1 台 VSC（效率 0.98、容量
0.25~pu）、3 台相机组，电压不平衡度上限 \fld{vuf\_max}$=0.05$
（\srcpath{tests/test\_three\_phase\_hybrid\_opf.cpp}:24--102）。

\subsubsection*{验证项与容差}
\begin{itemize}
  \item \textbf{Full vs GraphReduced}：两种模型变体在同一算例上均收敛，目标相对差
        $\le10^{-6}$，全电压复数差 $<10^{-6}$，电压违反 $<10^{-8}$，VUF 不超过
        上限 $+10^{-7}$，换流器电流不平衡度一致到 $10^{-6}$
        （同文件:106--143）；
  \item \textbf{图归约序列}：两断面序列求解与独立求解一致（目标 $10^{-7}$、电压
        $10^{-6}$），归约不丢时变负荷节点（同文件:145--172）；
  \item \textbf{GFL/GFM 动态平衡}：在 OPF 运行点上初始化 GFL（PLL）与 GFM（下垂）
        动态模型并 trim 到网络平衡，状态导数 $<10^{-7}$；解析等式 Jacobian 与
        Lagrangian Hessian 误差均 $<10^{-5}$（同文件:174--282）；
  \item \textbf{独立耦合 PF 复现}：由 OPF 结果构造三相混合 PF 算例，独立 Newton
        PF 复现电压到 $2\times10^{-6}$（同文件:284--317）；
  \item \textbf{Native Full 认证归约解}：把 GraphReduced 的原始--对偶点传输到 Full
        模型热启动，初始/最终原始与对偶残差、互补均 $\le10^{-6}$，目标一致到
        $10^{-6}$（同文件:319--401）；
  \item \textbf{约束 oracle}：惰性不等式求解与全行求解目标一致到 $10^{-7}$，
        被略去行的裕度满足激活阈值（同文件:403--461）；
  \item \textbf{SOCP 松弛下界}：抬升松弛给出经对偶证书认证的目标下界，不超过
        非线性解 $+10^{-7}$，外逼近轮次下界单调不降（同文件:463--510）。
\end{itemize}

\subsubsection*{已知偏差}
算例规模极小（9 相节点），不构成大规模三相系统的性能证据；松弛路径对恒电流负荷
显式拒绝并报告 \fld{model\_limitations}（同文件:512--522）；测试使用 AML 引擎层的
\fld{SolverBackend::Ipopt}/\fld{NativeIPM}，其可用性随 MIPSolvers 后端配置变化。


\subsection{RPO 专项验证}\label{subsec:ver:rpo}

模型（离散调压设备搜索 + 连续 AC/DC OPF 补全）见第~\ref{sec:rpo}~节；验证证据分
测试断言与文档实测两部分。

\subsubsection*{测试断言（\srcpath{tests/test\_reactive\_power\_opt.cpp}）}
\begin{itemize}
  \item \textbf{诚实局部证书}：case9 网损目标 RPO 收敛，且结果明确
        \fld{globally\_certified=false}、\fld{optimality\_gap\_available=false}、
        \fld{model\_limitations} 非空，算法标识为
        \fld{discrete\_coordinate\_search\_with\_ac\_opf}；无可调离散设备时目标
        等于基线网损（同文件:16--42）；
  \item \textbf{可投切并联}：case9 加 2 步电容器组后，电压偏差目标不差于基线
        （$+10^{-10}$），档位动作在铭牌范围内；\fld{gap} 恒为 1.0（不报告虚假
        全局间隙，同文件:44--82）；
  \item \textbf{控制清单}：\srcpath{inspect\_rpo\_controls} 对可调/固定/档位越界/
        停运设备分别给出排除原因，档位窗口 \fld{max\_tap\_move} 生效（同文件:112--182）；
  \item \textbf{MATPOWER 固定变比不伪造}：MATPOWER 导入的支路变比保持固定，
        不再虚构 4001 档 OLTC；\srcpath{case300\_acdc} 仅 authored 3 台
        $[-4,4]$、每档 1.25\% 的 OLTC 可调（同文件:184--211）；
  \item \textbf{档位动作联动}：带 \fld{source\_branch\_idx} 的 OLTC 动作按相对变比
        同步更新关联 AC 支路 \fld{tap}，保证 RPO 内层、独立 OPF 与 PF 回放使用
        同一调压动作（同文件:213--243）；
  \item case300 种子求解用例带 \fld{[.performance]} 标签，默认不运行（同文件:84--110）。
\end{itemize}

\subsubsection*{全过程交叉验证（文档口径）}
每次 RPO 后固定执行“保存离散档位 $\to$ 独立 OPF 复算 $\to$ PF 回放 $\to$ 网损账”
四步（\srcpath{docs/reactive\_power\_optimization\_validation.md}:103--123）。
独立 OPF 一致性判据为
\begin{align*}
\frac{|f_1-f_2|}{\max(1,|f_1|)}\le\max(10^{-3},10\epsilon_d),\qquad
\max_b|\Delta V_b|\le\max(5\times10^{-4},10\epsilon_p),
\end{align*}
其中 $\epsilon_d$、$\epsilon_p$ 分别为平稳性与可行性容差；非凸 OPF 可能存在目标
等价但 P/Q 分配不同的局部点，故 P/Q 最大差异只展示、不触发假失败。纯 AC 的 PF
电压差判据为 $10^{-3}$~pu；富混合 AC/DC 在 authored-space 调度未完整回写时显示
“可用数值检查通过，PF 回放部分覆盖”（同文件:113--123）。

\subsubsection*{大规模实测（文档实测）}
case2869 + 20 台 OLTC + 10 台并联的 authored 系统上：单位成本经济基线 74 次迭代
/25~s 收敛到参考目标 133999（违反 $1.1\times10^{-6}$）；暖启动链 2--7 次迭代、
平均约 0.46~s/次；完整 MinActiveLoss RPO 动作 22 台设备，物理网损
$1561.8\to1558.8$~MW。IEEE24 三区域回归中 vdev（MatchRPO 路径）与旧版逐位一致；
loss 路径因改用等价经济基点，基线/最终网损均优于旧版，combined 正常收敛
（\srcpath{docs/reactive\_power\_optimization\_validation.md}:159--177）。

\subsubsection*{已知偏差}
稳健内层 \fld{LossEconomic} 按网损最优分配机组无功、不直接为电压最优：case2869 上
vdev RPO 虽收敛，但 vdev $0.1392\to0.1401$ 基本未降；3000+ 完整搜索深度受时间限
（默认 120~s 内得可行改进点），离散评估并行化未完成
（\srcpath{docs/reactive\_power\_optimization\_validation.md}:179--191）。

\subsection{大规模混合 IPM 诊断证据}\label{subsec:ver:ipm}

\subsubsection*{case300\_acdc 根因与修复}
\srcpath{docs/large\_hybrid\_ipm\_diagnostics.md}:45--75 记录：故障并非模型不可行
（Ipopt 可把原始残差压到 $2.5\times10^{-9}$），而是原生 IPM 全步更新缺全局化，
轨迹进入病态区后所有稀疏分解失败。修复是对维数 $\ge512$ 的 KKT 增加残差 filter
回溯（原始可行性/对偶平稳性/互补三轴，步长减半最多 14 次），位于
\srcpath{src/optimal\_power\_flow/parity\_ipm.cpp}。修复后 case300\_acdc 原生 IPM
42 次迭代收敛，最大约束残差约 $1.74\times10^{-7}$、平稳性约 $2.77\times10^{-7}$，
41 接受步/17 回溯拒绝；强制 UMFPACK、KLU、Eigen SparseLU、Dense LU 四后端均通过
同一算例（同文件:77--88，文档实测）。测试侧对应断言为
\srcpath{tests/test\_opf\_solver\_backends.cpp}:86--112（收敛、接受步 $>0$ 且
$\le$ 迭代数、拒绝步 $>0$）。

\subsubsection*{2000 节点富混合实测}
\srcpath{docs/large\_hybrid\_ipm\_diagnostics.md}:118--131 记录（文档实测，Release）：
原生 parity IPM + 稀疏 UMFPACK 90 次迭代/7.07~s 收敛，最大违反
$3.49\times10^{-9}$、平稳性 $9.03\times10^{-7}$；内嵌 Ipopt + MUMPS 100 次预算
48.15~s 后原始残差 $4.96\times10^{-4}$、对偶 1.49，未达停止判据。测试断言为
\srcpath{tests/test\_opf\_solver\_backends.cpp}:880--912（Parity IPM 收敛、
违反 $<10^{-5}$、迭代 $<160$）与同文件:914--943（Ipopt 有界基准，仅断言有限，
隐藏性能用例）。据此文档建议 2000 节点级富混合系统优先使用原生 IPM，Ipopt 保留为
小中型问题后端与交叉诊断路径。

\subsubsection*{后端结构分工与暖启动}
2026-07-19 更新的诊断结论（同文件:224--289，文档实测）：线性代数默认顺序
MUMPS $\to$ UMFPACK $\to$ KLU $\to$ Eigen；Newton 系统默认增广 KKT +
W\"achter--Biegler $\delta_W$ 惯性校正；纯 AC 大算例自动走 KLU（BTF 预分解契合
近块三角电网 KKT，2869/9241/13659 三案收敛到同一参考目标且更快），混合 AC/DC 因
增广 KKT 惯性信息只有 MUMPS 报告而必须走 MUMPS——文档明确 KLU/LU 不报告惯性时
混合算例失去正则而停滞。\fld{ac\_pf\_warm\_start} 选项使 case13659pegase 以 189 次
迭代/约 41~s 收敛于参考目标 386107，但在良态算例上反而更慢、个别 rte 算例不收敛，
因此是选择性启用而非默认（同文件:235--239）。

\subsubsection*{已知偏差}
arm64 本地静态 MUMPS 5.7.3 在全部六类 NLP 测试中第 1 次迭代即
\texttt{Ipopt restoration failed}，当前默认仍用 Homebrew MUMPS 动态库
（同文件:111--116、190--193）；Jacobian 系数尺度达 $10^{12}$ 时默认 gradient
scaling 会压小 scaled 对偶误差，适配器在启用未缩放分量门槛后明确返回
\texttt{search direction too small} 而非假收敛——此类模型必须在建模层无量纲化
（同文件:184--188）。

\subsection{BPA/DSP LCC 算例验证}\label{subsec:ver:lcc}

LCC 直流算例的 OPF 验证是平台唯一带\textbf{外部商业软件参考解}的 OPF 通道。
完整证据链见专项报告
\srcpath{docs/lcc\_dat\_opf\_report/LCC\_DAT\_OPF\_technical\_report.tex}，
此处汇总与本手册相关的部分。

\subsubsection*{算例与基准}
\srcpath{data/dsp/cigre.dat}（DSP Samples/CIGRE）：100~MVA 基准，6 交流母线、
4 交流支路、2 直流母线、1 条 $10~\Omega$ 直流线路、2 台 LCC；DSP 参考解为
整流/逆变 $U_d=500/470$~kV、$I_d=3$~kA、交流有功 $-1500/+1410$~MW、线损
$I_d^2R=90$~MW（同报告 :982--1009）。\srcpath{data/dsp/2DC.dat} 为补充交叉校核算例
（同报告 :1116--1124）。

\subsubsection*{验证项与容差}
\begin{itemize}
  \item \textbf{导数级}：LCC parity 平衡方程解析 Jacobian 对中心有限差分
        $<2\times10^{-5}$，覆盖阀侧母线、换相母线与直流电压三类变量
        （\srcpath{tests/test\_opf\_solver\_backends.cpp}:310--422）；隔离 LCC 的
        局部 Hessian 对称性 $<10^{-10}$，对 Jacobian 差分相对容差 $5\times10^{-3}$
        （同文件:424--598）；
  \item \textbf{raw CIGRE OPF}：原始 DAT 直接导入（BS 平衡机 $P_{\min}=-9999$~MW
        为卡片语义），Parity IPM 收敛，$P_g$ 复现 $1500/-1410$~MW $\pm2$、线损
        90~MW $\pm2$、违反 $<10^{-6}$；OPF 结果 JSON 双向往返不丢
        \fld{lcc\_transfers}（同文件:600--671）；报告记录 7 次迭代收敛、
        primal/dual 残差 $1.7\times10^{-11}$/$6.0\times10^{-7}$（同报告 :1095--1108，
        文档实测）；
  \item \textbf{不可行变体安全失败}：受端 $P_{\min}=0$ 时模型真不可行，断言返回
        有限、可序列化的失败结果，不伪造无损直流解、不从错误尺寸 primal 提取物理
        运行点（同文件:673--710）；
  \item \textbf{固定 tap PF 回放}：把 OPF 调度写回副本（显式关闭 PF 专用 R 卡
        分接头外环以保证同一 Ybus），LCC 端口 $P/Q/P_{dc}$ 与 OPF 一致到
        0.5~MW/MVAr（同文件:814--846）；
  \item \textbf{后端守卫}：EconomicDispatch 后端对混合 AC/DC/LCC 显式拒绝；
        DC OPF 显式拒绝并打 \fld{dc-opf:rejected-lcc} 范围标签；不支持的控制模式
        组合与缺失换相电抗在 readiness 检查即拒绝（同文件:848--877）。
\end{itemize}

\subsubsection*{当前验证状态（报告声明）}
报告“当前验证状态”表声明：DAT 接口与 PF、DSP 对比、LCC OPF（175 项断言/4 个
用例，含三变量 Jacobian/Hessian、raw CIGRE、固定 tap 回放与后端保护）、
$P_{\min}=0$ 不可行变体、浏览器 CIGRE E2E 均通过
（\srcpath{docs/lcc\_dat\_opf\_report/LCC\_DAT\_OPF\_technical\_report.tex}:1039--1060）。
同一节明确\textbf{完整 \srcpath{test\_opf\_solver\_backends} 不在该轮验收范围内}，
且历史全量运行存在与 LCC 无关的后端/大算例基线失败，不能写成全绿（同文件
:1058--1060）。

\subsubsection*{已知偏差}
OPF 中换流变压器 tap 是固定 formulation 输入，不复现 DSP 的调档运行点；报告接受
准则也不要求复现调档点（同文件:1036--1037、1086--1114）。2DC 卡片
$P_{\min}=P_{\max}=0$，是 PF/Jacobian 参考而非 OPF-ready（测试同文件:712--715）。
其余局限（准稳态、无换相重叠迭代、clamp 不可二阶光滑等）见
\ref{subsec:ver:limitations} 节。

\subsection{结果审计、PF 回放与投影 round-trip}\label{subsec:ver:audit}

\begin{itemize}
  \item \textbf{\srcpath{verify\_opf\_result} 审计}：case9 收敛结果经审计后
        电压限值违反 $<0.02$~pu、机组限值违反 $<1.0$~MW，报告/重算目标均有限；
        结构性不可行系统（负荷无源）审计必须给出不可行提示
        （\srcpath{tests/test\_solver\_capabilities.cpp}:194--239）；
  \item \textbf{混合系统禁 AC-only 回退}：含 DC 母线与 VSC 的系统在极端容差下
        不收敛时，状态不得含 \texttt{economic-dispatch + AC PF fallback}，且不可行
        提示中必须出现 \texttt{AC-only economic-dispatch fallback suppressed}
        （同文件:241--302）；
  \item \textbf{能力标志与实际行为一致}：\fld{supports\_ac\_opf}/
        \fld{supports\_dc\_opf} 恒真且实解收敛（同文件:135--156）；
  \item \textbf{投影 round-trip}：纯 AC 重编号后 PF 与 AC OPF 结果向量尺寸与
        authored 母线一致（\srcpath{tests/test\_graph\_roundtrip.cpp}:223--258）；
        串联归约前后 DC OPF 目标一致到 $10^{-3}$（同文件:671--705）；
  \item \textbf{AML 数据装配}：孤立 DC 母线及其关联支路不进入 OPF 模型
        （\srcpath{tests/test\_hacdcpf.cpp}:234--261）；
  \item \textbf{全约束族冒烟}：多尺度综合算例在支路限值、换流器容量/电流/调制
        约束全部启用下收敛，违反 $<100\times$ 可行性容差
        （\srcpath{tests/test\_multiscale\_comprehensive.cpp}:170--229）；
  \item \textbf{入口冒烟}：IEEE14 AC/DC 的 AC/DC OPF 入口返回非空状态串
        （\srcpath{tests/test\_hacdcpf.cpp}:372--400）；
  \item \textbf{浏览器端到端}：RPO GUI 全流程（输入、运行、结果与诚实证书）由
        Playwright E2E 覆盖，注册于 \srcpath{tests/CMakeLists.txt}:532--537。
\end{itemize}


\subsection{已知局限汇总}\label{subsec:ver:limitations}

表~\ref{tab:ver:limitations} 从测试中的 skip/WARN 声明、门控分支、结果结构中的
诚实标签与文档“局限/已知问题”段落逐条收集，未做美化。表中“表现/依据”列为代码
或文档原文语义的转述。

\begin{footnotesize}
\begin{longtable}{@{}L{4.2cm}L{1.9cm}L{5.2cm}L{4.0cm}@{}}
\caption{OPF 模块已知局限汇总}\label{tab:ver:limitations}\\
\toprule
\textbf{局限} & \textbf{影响范围} & \textbf{表现/依据} & \textbf{出处（文件:行号）}\\
\midrule
\endfirsthead
\multicolumn{4}{@{}l}{\footnotesize（续表）}\\
\toprule
\textbf{局限} & \textbf{影响范围} & \textbf{表现/依据} & \textbf{出处（文件:行号）}\\
\midrule
\endhead
\midrule
\multicolumn{4}{r@{}}{\footnotesize（接下页）}\\
\endfoot
\bottomrule
\endlastfoot
嵌入式 Ipopt 默认被门控关闭 &
AC OPF 后端 &
\srcpath{HACDCPF\_ENABLE\_IPOPT} 默认 OFF（macOS ON）；选 Ipopt 而后端不可用时静默回退 Parity IPM，\fld{solver\_path} 报告 \fld{ParityIPM} &
\srcpath{tests/test\_opf\_solver\_backends.cpp}:1100--1108；\srcpath{src/optimal\_power\_flow/ac\_opf.cpp}:33\\
Ipopt 在 2000 节点混合算例预算内不达标 &
大规模混合 &
100 次预算后原始 $4.96\times10^{-4}$、对偶 1.49，未收敛；文档建议生产路径用原生 IPM &
\srcpath{docs/large\_hybrid\_ipm\_diagnostics.md}:118--131、211--220（文档实测）\\
case6515rte 不期望收敛 &
大 AC 算例 &
测试仅断言目标有限（$<10^{30}$）与稀疏后端，注释明确不期望预算内收敛 &
\srcpath{tests/test\_opf\_solver\_backends.cpp}:1178--1196\\
Ipopt 有界基准与参数矩阵为隐藏用例 &
回归可见性 &
带 \fld{[.performance]} 标签，默认 ctest 不运行；case300 RPO 种子用例同 &
\srcpath{tests/test\_opf\_solver\_backends.cpp}:914--946；\srcpath{tests/test\_reactive\_power\_opt.cpp}:84--85\\
AC 与 DC OPF 目标值不直接可比 &
交叉验证 &
测试注释声明成本口径不同，直接比较无意义；只断言各自收敛与 DC 目标为正 &
\srcpath{tests/test\_acopf\_dcopf\_crossval.cpp}:284--291\\
AC OPF 结果不暴露支路潮流 &
交叉验证 &
$P_f$ 对比代码留空注释“skip for now” &
\srcpath{tests/test\_acopf\_dcopf\_crossval.cpp}:222--227\\
DC OPF 调度质量判据宽松 &
交叉验证 &
相对调度差断言 $<1.0$（100\%）；546 行尾注释“50\%”与断言不符，以断言为准 &
\srcpath{tests/test\_acopf\_dcopf\_crossval.cpp}:536--547\\
case57 的 PF(DCOPF) 回放不保证收敛 &
交叉验证 &
收敛性降级为 \texttt{WARN}，非断言 &
\srcpath{tests/test\_acopf\_dcopf\_crossval.cpp}:273--282\\
综合对比限于 $\le1000$ 节点 &
交叉验证 &
更大算例 SKIP；scenario/contab 文件排除 &
\srcpath{tests/test\_acopf\_dcopf\_crossval.cpp}:569--574、629--633\\
DC OPF 支路对偶未认证 &
DC OPF LMP &
含 \fld{pf} 变量界时 native simplex 不支持对偶提取，\fld{branch\_mu\_valid=false} &
\srcpath{tests/test\_acopf\_dcopf\_crossval.cpp}:794--799、868--872\\
LMP 提取依赖辅助 LP &
DC OPF LMP &
仅 \fld{include\_branch\_limits=false}（无 \fld{pf} 变量）时辅助 LP 成功 &
\srcpath{tests/test\_acopf\_dcopf\_crossval.cpp}:698--702、716--719\\
Ipopt 乘子不全时 LMP 标不可用 &
AC OPF &
适配器不返回完整约束乘子时 LMP 必须标为不可用 &
\srcpath{docs/lcc\_dat\_opf\_report/LCC\_DAT\_OPF\_technical\_report.tex}:1227--1228\\
混合系统抑制 AC-only 经济调度回退 &
AC OPF &
含 DC/LCC 时禁止丢弃混合网络；提示 \texttt{fallback suppressed} &
\srcpath{tests/test\_solver\_capabilities.cpp}:241--302\\
EconomicDispatch 后端拒绝混合/LCC &
AC OPF &
状态串含 \texttt{hybrid AC/DC/LCC}，不收敛 &
\srcpath{tests/test\_opf\_solver\_backends.cpp}:848--854\\
DC OPF 拒绝 LCC &
DC OPF &
显式失败并打 \fld{dc-opf:rejected-lcc} 范围标签 &
\srcpath{tests/test\_opf\_solver\_backends.cpp}:856--859\\
AML builder 不支持 LCC &
AML 建模层 &
入口/能力范围必须显式声明，不把 LCC 当 0 注入 &
\srcpath{docs/lcc\_dat\_opf\_report/LCC\_DAT\_OPF\_technical\_report.tex}:954\\
三相混合 OPF 不支持在役 LCC &
三相 OPF &
GUI/HTTP 入口显式拒绝并提示改用平衡 AC/DC Parity 或 Ipopt &
\srcpath{docs/lcc\_dat\_opf\_report/LCC\_DAT\_OPF\_technical\_report.tex}:955\\
LCC 换流变 tap 在 OPF 中为固定输入 &
LCC OPF &
结果须声明 \texttt{taps are fixed inputs} 与 \texttt{no tap-changer control}；联合优化需后续把 $t_k$ 纳入决策 &
\srcpath{docs/lcc\_dat\_opf\_report/LCC\_DAT\_OPF\_technical\_report.tex}:1086--1114、1210--1212\\
LCC 为快照准稳态模型 &
LCC &
不含换相重叠角显式迭代、换流失败、谐波、保护与控制动态；不能推断暂态安全 &
\srcpath{docs/lcc\_dat\_opf\_report/LCC\_DAT\_OPF\_technical\_report.tex}:1206--1207\\
LCC 额定电流 clamp 不可二阶光滑 &
LCC OPF &
活动区切换点非二阶光滑；应报告绑定状态，局部 KKT 非全局证书 &
\srcpath{docs/lcc\_dat\_opf\_report/LCC\_DAT\_OPF\_technical\_report.tex}:1213--1214\\
2DC 卡片非 OPF-ready &
LCC 测试 &
$P_{\min}=P_{\max}=0$，仅作 PF/Jacobian 参考 &
\srcpath{tests/test\_opf\_solver\_backends.cpp}:712--715\\
完整 backends 测试目标存在历史基线失败 &
测试套件 &
LCC 报告声明完整 \srcpath{test\_opf\_solver\_backends} 不在该轮验收范围，历史全量有与 LCC 无关的后端/大算例失败 &
\srcpath{docs/lcc\_dat\_opf\_report/LCC\_DAT\_OPF\_technical\_report.tex}:1058--1060\\
RPO 不提供全局最优性证书 &
RPO &
\fld{globally\_certified=false}、\fld{gap=1.0}；局部灵敏度面不用于安全剪枝 &
\srcpath{tests/test\_reactive\_power\_opt.cpp}:16--42、79--81；\srcpath{docs/reactive\_power\_optimization\_validation.md}:82--87\\
RPO 稳健内层不直接为电压最优 &
RPO 大电网 &
case2869 上 vdev RPO 收敛但 vdev 基本未降（$0.1392\to0.1401$） &
\srcpath{docs/reactive\_power\_optimization\_validation.md}:179--191\\
RPO 富混合回放部分覆盖 &
RPO &
混合换流器调度未完整回写 authored-space PF 时标“部分覆盖”；三相单体 RPO 未覆盖 &
\srcpath{docs/reactive\_power\_optimization\_validation.md}:95--99、113--123；\srcpath{docs/large\_hybrid\_ipm\_diagnostics.md}:133--138\\
三相松弛拒绝恒电流负荷 &
三相 OPF &
\fld{fixed-current} 显式拒绝并报告 \fld{model\_limitations} &
\srcpath{tests/test\_three\_phase\_hybrid\_opf.cpp}:512--522\\
手算算例未自动化 &
验证覆盖 &
\srcpath{external\_data/opf\_hand\_cases/} 无 C++ 测试/ctest 引用，仅 GUI 人工校核 &
\srcpath{external\_data/opf\_hand\_cases/README.md}:1--19\\
arm64 静态 MUMPS 数值/ABI 问题 &
Ipopt 后端 &
源码构建 MUMPS 5.7.3 全部用例第 1 次迭代 \texttt{restoration failed}；默认仍用 Homebrew 动态库 &
\srcpath{docs/large\_hybrid\_ipm\_diagnostics.md}:111--116、190--193\\
尺度 $10^{12}$ 模型的缩放假象 &
Ipopt 适配器 &
默认 gradient scaling 压小 scaled 对偶误差；需建模层无量纲化，不能仅靠收紧 \fld{tol} &
\srcpath{docs/large\_hybrid\_ipm\_diagnostics.md}:184--188\\
时序 OPF 单步失败时降级 &
时序流水线 &
该步退回 UC 快照直接 PF，结果精度由交叉指标揭示，不标记为 OPF 质量 &
\srcpath{docs/time\_series\_power\_flow\_models.md}:307--308\\
外部电网夜间通道覆盖有限 &
夜间 OPF &
仅 2 个工程 JSON + 3 个微系统，无外部工具对照 &
\srcpath{tests/test\_nighttime\_opf.cpp}:46--96、164--209\\
\end{longtable}
\end{footnotesize}

\subsection{可信边界结论}\label{subsec:ver:boundary}

综合上述证据，当前 OPF 模块的可信边界可划分为四档。

\subsubsection*{验证充分（默认 ctest 断言 + 外部/文档基准一致）}
\begin{itemize}
  \item \textbf{纯 AC 中小规模经济 OPF}（case5--case300 级）：NativeAC 与 Parity IPM
        均收敛，PF 回放、审计、LMP 符号/量级、拥塞再调度等系统性判据完整；
  \item \textbf{纯 AC 大规模 Parity IPM}：case2383wp、case2869pegase 有默认断言；
        文档实测 case2869/9241/13659 收敛到参考目标值 133999/315837/386107；
  \item \textbf{交直流混合 Parity IPM}：case300\_acdc（四种线性代数后端）与
        case2000\_acdc（2000 节点 + 多端直流 + 8 VSC）有默认断言与文档实测残差；
        暖启动、filter 回溯、稀疏 KKT 路径均有针对断言；有界 Phase I 在
        小型 Hybrid 上认证 primal/dual，在 case300/case2000 严重冷启动上
        只保证单调/预算/诚实非接受，未接受状态不改变 Phase II 起点；
  \item \textbf{DC OPF}：MATPOWER 小算例收敛、限值、孤岛削减、LMP 符号/均匀性/
        拥塞经济性断言完整。
\end{itemize}

\subsubsection*{有测试但覆盖有限（特定算例或单元级断言）}
\begin{itemize}
  \item \textbf{LCC OPF}：2 个 BPA DAT 算例 + DSP 外部参考 + 有限差分导数 + 固定
        tap 回放；多端、定电流/定 $\alpha$ 控制组合、tap 联合优化未覆盖；
  \item \textbf{三相混合 OPF}：9 相节点手组算例上的 Full/归约/松弛/oracle/动态
        平衡全链断言；另有 360 个三相母线、2175 变量的 Phase I 稀疏预算与
        primal/dual-fit 证书，但无同规模完整 Phase II 或外部工具 OPF 对照；
  \item \textbf{RPO}：case9 单元断言 + case2869 authored 文档实测；全局最优性
        明确不声称，大电网电压偏差改善幅度有限；
  \item \textbf{夜间/外部电网 OPF}：2 个工程算例 + 3 个微系统；
  \item \textbf{时序流水线}：24 步动态 OPF 与 4 步 UC 全链断言；多日并行、年度
        8760 步的 OPF 质量由降级路径与交叉指标保障而非逐断言。
\end{itemize}

\subsubsection*{仅有界性/有限性保证}
case6515rte（40 次预算内只断言有限）；case2000\_acdc 的 Ipopt 路径（100 次预算
内文档实测未达停止判据）；case300\_acdc 的 Ipopt 路径需 800 次预算与特定容差组合
方收敛，且平稳性收紧到 $10^{-4}$ 反而退化（文档实测）。这些路径可用于交叉诊断，
不应作为生产精度来源。

\subsubsection*{明确不支持并被守卫}
DC OPF 与 AML builder 的 LCC、三相混合 OPF 的在役 LCC、EconomicDispatch 后端的
混合/LCC 算例、三相松弛的恒电流负荷、混合系统的 AC-only 回退——均在入口或
readiness 检查显式拒绝并打范围标签，不会静默给出误导性结果。

\subsubsection*{总体判断}
OPF 模块在“纯 AC 经济调度”与“平衡 AC/DC 混合 Parity IPM”两条主线上具备从
导数到外部基准的完整证据链，可用于生产；LCC、三相、RPO、时序四条扩展线各有
针对性测试与诚实的能力标签，但其证据是单元/算例级的，使用时应以结果中的
\fld{model\_scope}、\fld{model\_limitations} 与审计字段为准。任何超出上述边界的
规模或工况（$>3000$ 节点富混合、含 LCC 的多端直流优化、三相不平衡单体 RPO），
当前只有工程合理性推断，没有验证证据支撑。

\section{跨模块工业级技术评价基线}

本章规定本手册所述模块进入工程算例、批量研究或生产服务前必须共同满足的评价基线。
具体模型公式、变量、约束和专用算法以正文相应章节为准；本章不扩大源码已经实现的范围。

\subsection{输入、输出、边界条件与数据结构审计}
输入审计应逐项检查有限性、单位、上下界、枚举值和引用完整性。系统对象必须先按所需层级完成
校验；AC 与 DC 母线采用域限定映射，组件稳定 \texttt{index}、向量位置和临时图索引不得混用。
输出审计必须记录单位、索引空间、模型范围、收敛或可行状态、后端、容差、迭代/节点计数和告警。
空系统、空时域、孤岛、重复 ID、零容量、退化约束与不可用可选后端均属于边界测试集。

\subsection{工程假设与数值稳定性分析}
模型使用的平衡、线性化、稳态、独立性、概率分布、时间离散或网络等值假设必须与应用场景匹配。
数值评价至少包含变量缩放、绝对/相对残差、可行性违反、条件数或枢轴质量、步长/阻尼、迭代停滞
和随机种子复现性。若采用正则化、平滑、回退、松弛或近似，应同时报告触发条件、参数值、对主指标
的敏感性以及是否改变模型口径；不得用“已返回结果”替代收敛或有效性证明。

\subsection{复杂度与资源分析}
令 $n_x$、$n_c$、$n_t$、$n_s$、$|V|$、$|E|$ 分别表示变量、约束、时段、场景、节点和边数。
报告应分别给出建模、求解、后处理的时间与内存。图遍历通常为 $O(|V|+|E|)$；逐时段/逐场景
流程至少线性增长为 $O(n_t n_s)$ 次子问题调用；稠密线性代数最坏可达 $O(n_x^3)$，稀疏方法则应
报告结构非零数、填充率和因子分解峰值内存；MILP/NLP 不给出虚假的多项式最坏界，而报告节点数、
间隙、时限和实测扩展曲线。复杂度结论必须基于固定硬件、固定后端和可复现算例族。

\subsection{异常工况处理}
非法输入应在进入求解器前返回结构化错误；不可行、无界、不收敛、时限、内存不足和后端缺失必须
相互区分。部分结果只有在对应 \fld{ValidityFlags}、\fld{model\_scope}、
\fld{model\_limitations} 或等价字段允许时才可使用。异常路径不得静默丢弃 AC/DC 资产、制造平衡
节点、把 NaN 改写为零或把近似解标为全局最优；系统原始对象应保持可恢复和可重试。

\subsection{与其他模块的接口关系}
接口链路遵循“工程语义模型 $\to$ 校验 $\to$ 投影/装配 $\to$ 求解或分析 $\to$ 结果回投影”。
跨模块传递时保留稳定 ID、单位、符号、时间基准和场景身份；消费潮流、OPF、动态或可靠性结果前，
先检查其收敛、覆盖范围和有效性标志。HTTP/JSON/Python 层只做语义保持适配，不重新解释求解结果。

\subsection{测试与验证建议}
最低验证矩阵包括：手算小例、标准算例、边界/异常输入、随机性质测试、算法或后端交叉校核、
序列化往返、稳定 ID 回投影、AC/DC 同号母线、重复运行确定性和规模扩展测试。数值算法还应加入
残差独立重算、雅可比有限差分或自动微分校核；近似模型应与高保真模型比较并给出误差分布，而非
只报告单个平均值。回归报告必须列明构建配置、算例版本、容差、通过/失败/跳过数及未验证范围。

\subsection{工业级评估指标}
\begin{itemize}
  \item \textbf{正确性}：守恒残差、约束最大违反、解析/基准误差、结果回投影一致性；
  \item \textbf{鲁棒性}：收敛率、不可行识别率、异常分类准确率、扰动敏感性与随机复现率；
  \item \textbf{性能}：端到端时延、吞吐量、峰值内存、并行效率与规模增长斜率；
  \item \textbf{可追溯性}：输入模式版本、稳定 ID 覆盖率、模型范围/限制字段完整率、告警可定位率；
  \item \textbf{工程有效性}：与标准、外部引擎或现场数据的偏差，以及误判/漏判对业务指标的影响。
\end{itemize}
指标应预先给出验收阈值和失败处置；未达到阈值时只能标记为实验性或受限适用。

\subsection{适用范围、局限性与后续改进}
本手册只评价当前源码覆盖的模型、后端和数据结构，不对未建模物理过程、未安装求解器或未执行的
外部验证作保证。后续改进应按“缺口 $\to$ 可量化指标 $\to$ 源码/测试变更 $\to$ 独立验证”闭环，
优先消除会造成错误结论的语义缺口，其次提升数值鲁棒性与性能，最后扩展模型范围。


\end{document}
